\ifx\danslelivre\undefined
\RequirePackage{luatex85}
\documentclass[10pt]{../fichiers-configuration-latex/smfart}
\usepackage{polyglossia}
%\usepackage[francais]{babel}
\usepackage{fontspec}
\usepackage[math-style=ISO,bold-style=upright,sans-style=upright]{unicode-math}
%\usepackage[mathletters]{ucs}
%\usepackage[utf8]{inputenc}
%\usepackage[T1]{fontenc}
\usepackage[unicode,colorlinks=true,linkcolor=blue,citecolor=Sepia,urlcolor=magenta]{hyperref}
\setmainlanguage{french}
\setmainfont[Ligatures=TeX]{Linux Libertine O}
% variante possible : TeXGyrePagella ?
\setmathfont{XITS Math}

% maths
\usepackage{amsmath}
%\usepackage{amssymb}
%\usepackage{mathrsfs}
%\usepackage{stmaryrd}
%\usepackage{euscript}

% couleurs
\usepackage[usenames,dvipsnames]{xcolor}

% graphiques
\usepackage{graphicx}
\usepackage[all]{xy}
% Xy-pic 3.8.9's PDF driver assumes that \fontdimen8\textfont3 is a
% legacy TeX rule width.  Under LuaTeX with unicode-math/XITS Math it is
% instead 6.62pt, which turns every arrow shaft into a thick black bar.
% Override only xypdf's private cached width; do not alter the math font.
\makeatletter
\AtBeginDocument{%
  \@ifundefined{xP@preclw}{}{%
    \edef\xP@preclw{0.4pt}%
    \edef\xP@lw{\xP@dim\xP@preclw}%
    \PackageInfo{sga4-xypdf-compat}{Xy-pic PDF line width set to \xP@preclw}%
  }%
}
\makeatother
%\usepackage{tikz}
%\usetikzlibrary{matrix,arrows,calc}
%\tikzset{column sep=3em,row sep=2em,text height=1.5ex, text depth=0.25ex,>=angle 90}
%\tikzset{description/.style={fill=white,inner sep=1pt}}

% hyperliens, 'etiquettes
\usepackage{xr-hyper}
\usepackage{../fichiers-configuration-latex/smfthm}
%\ifx\cacheretiquettes\undefined
%\usepackage[inline]{showlabels} \else \fi % pour rendre visibles les 'etiquettes
%\usepackage{../fichiers-configuration-latex/makeidx} % pour faire un index
%\usepackage[pagebackref,unicode]{hyperref}  % pour voir o`u sont cit'ees les r'ef'erences
\hypersetup{unicode,colorlinks=true,linkcolor=blue,citecolor=Sepia,urlcolor=magenta,
  hypertexnames=false}
%\hypersetup{breaklinks=true} %permet le retour `a la ligne dans les liens trop longs. Incompatible avec pdflatex
%\renewcommand*{\backref}[1]{{\small [cit\'e page(s)~#1.]}}

\usepackage{xspace}
\usepackage{pifont} % pour expose XVII

\renewcommand{\thefootnote}{\roman{footnote}} % pour mettre les notes en bas de page sous la forme (ii) et non (2), qui ressemble trop (?) a un exposant.

\newcounter{pageoriginal}
\newcommand{\pageoriginale}{\refstepcounter{pageoriginal}\marginpar{\footnotesize\textbf{\thepageoriginal}}\xspace}
\let\pageoriginaled\pageoriginale

\usepackage{ifthen}
\newboolean{original}
\def\sisi{\ifthenelse{\boolean{original}}}
\setboolean{original}{false}

\newcounter{cnote}
\def\nde#1{\ifthenelse{\boolean{original}}{}{\stepcounter{cnote}\def\thefootnote{\arabic{cnote}}\protect\footnote[\value{cnote}]{N.D.E.: #1}}}
\def\ndemark{\ifthenelse{\boolean{original}}{}{\stepcounter{cnote}\def\thefootnote{\arabic{cnote}}\protect\footnotemark[\value{cnote}]}}
\def\ndetext#1{\ifthenelse{\boolean{original}}{}{\def\thefootnote{\arabic{cnote}}\protect\footnotetext[\value{cnote}]{N.D.E.: #1}}}

\setmainlanguage{french}
\setotherlanguage{english}

% unicode-math deliberately disables the legacy \Arrowvert glyph.  The
% source uses it once as a rotated equality marker; \Vert has the intended
% appearance and keeps the historical diagram buildable with LuaLaTeX.
\renewcommand{\Arrowvert}{\Vert}


%%
%% Pour afficher le numéro de version Git (exige que LuaTeX soit appelé avec --shell-escape)
%%

\catcode`\%=12
\directlua{
function command_output(cmd)
  local output_pipe = io.popen(cmd)
  local line_iter = output_pipe and output_pipe:lines()
  local line = line_iter and line_iter()
  if output_pipe then output_pipe:close() end
  return line
end
local git_command1 = [[git log --pretty=format:"%h %ai" -1]]
local git_command2 = [[if git diff --quiet HEAD ; then echo clean ; else echo dirty ; fi]]
function git_status()
  local status = command_output(git_command1)
  if status then
    tex.tprint({luatexbase.catcodetables.CatcodeTableLaTeX, string.char(0x5c).."texttt{"}, {-2,status}, {luatexbase.catcodetables.CatcodeTableLaTeX, "}"})
    if not (command_output(git_command2) == "clean") then tex.sprint(-2, " (modifications non validées)") end
  else
    tex.sprint(-2, "?")
  end
end
}
\catcode`\%=14
\def\showgitstatus{\directlua{git_status()}}


% macros pour LaTeX
\renewcommand{\mathsf}{}

% Police différente pour U+2113 SCRIPT SMALL L
\setmathfont[range="2113]{latinmodern-math.otf}

% Give every explicit SGA label its own destination and preserve the printed
% number supplied as the optional argument.  The former definition inherited
% whichever counter anchor happened to precede it and ignored #1 entirely.
\makeatletter
\def\etiquette[#1]#2{%
  \phantomsection
  \protected@edef\@currentlabel{#1}%
  \label{#2}%
}
\makeatother
\def\eqhyperref[#1]#2{(\hyperref[#1]{#2})}
\def\pagehyperref[#1]#2{(page \hyperref[#1]{#2})}
\def\hyperrefext[#1][#2]#3{\hyperref[#2]{#1 #3}}
\def\eqhyperrefext[#1][#2]#3{\hyperref[#2]{#1 #3}}     %cf. exposes 17
\def\neweqhyperrefext[#1][#2]#3{\hyperref[#2]{#1 #3}}  %cf. exposes 17

\def\shat#1{\protect\widehat{#1}}
\def\stilde#1{\protect\widetilde{#1}}
%%%%%%%%%%%%%%%%%%


\def\etcn{etc.\xspace}
\def\etco{etc...\xspace}
\def\etc{\sisi{\etco}{\etcn}}
\def\ldotsn{\sisi{...}{\etcn}}

\def\cf{{cf.}\kern.3em}
\def\Cf{{Cf.}\kern.3em}
\def\ie{{i.e.}\kern.3em}
\def\iev{{i.e.},\xspace}
\def\loccito{loc.\kern2pt cit.\xspace}
\def\loccitn{\textit{loc.\kern2pt cit.}\xspace}
\def\loccit{\sisi{\loccito}{\loccitn}}
\def\EGA{\textup{EGA}\kern.2em}
\def\SGA{\textup{SGA}\kern.3em}

\def\quoi{}

\def\quav{\quad,\quad}
\def\quaet{\quad \text{et}\quad}
\def\fibreproduct#1#2#3{#1 ×_{#3} #2}
\let\fprod\fibreproduct
\def\ccheck#1#2{\mbox{${\check{\check{\rm #1}}}^{\lower2pt\hbox{$\scriptstyle #2$}}$}}
\def\sheaf#1{\protect\mathcal{#1}}
\def\vieuxsheaf#1{{#1}} % cf. problème avec sheaf{a\Phi_i} par exemple...
\DeclareMathOperator{\Hom}{Hom}
\DeclareMathOperator{\can}{can}
\DeclareMathOperator{\Ind}{Ind}
\DeclareMathOperator{\ind}{ind}
%\DeclareMathOperator{\Hhomtop}{\mathcal{H}\!\mathit{omtop}}
\DeclareMathOperator{\Hhomtop}{\mathbf{Homtop}}
\DeclareMathOperator{\Point}{Point}
\DeclareMathOperator{\Pro}{Pro}
\DeclareMathOperator{\Top}{Top}
\DeclareMathOperator{\Mod}{Mod}
\DeclareMathOperator{\Cart}{Cart}
%\DeclareMathOperator{\Mmorsite}{\mathcal{M}\!\mathit{orsite}}
\DeclareMathOperator{\Mmorsite}{\mathbf{Morsite}}
\DeclareMathOperator{\loc}{loc}
\DeclareMathOperator{\Loc}{Loc}
\DeclareMathOperator{\Pic}{Pic}
\DeclareMathOperator{\Et}{Et}
\DeclareMathOperator{\card}{card}
\DeclareMathOperator{\Ens}{Ens}
\DeclareMathOperator{\h}{\underline{H}}
\DeclareMathOperator{\Fonct}{Fonct}
\DeclareMathOperator{\Ker}{Ker}
\DeclareMathOperator{\but}{but}
\DeclareMathOperator{\Ob}{Ob}
\DeclareMathOperator{\ob}{ob}
\DeclareMathOperator{\Lib}{Lib}
\DeclareMathOperator{\esj}{esj}
\DeclareMathOperator{\id}{id}
\DeclareMathOperator{\Id}{Id}
\DeclareMathOperator{\ad}{ad}
\DeclareMathOperator{\Iim}{Im}
\DeclareMathOperator{\pr}{pr}
\DeclareMathOperator{\Cat}{Cat}
\DeclareMathOperator{\car}{car}
\DeclareMathOperator{\SysInd}{SysInd}
\DeclareMathOperator{\Sup}{Sup}
\DeclareMathOperator{\Coker}{Coker}
\DeclareMathOperator{\Filt}{Filt}
\DeclareMathOperator{\acc}{acc}
\DeclareMathOperator{\sur}{sur}
\DeclareMathOperator{\cat}{cat}
\DeclareMathOperator{\Ord}{Ord}
\DeclareMathOperator{\Gr}{Gr}
\DeclareMathOperator{\Cov}{Cov}
\DeclareMathOperator{\Caf}{Caf}
\DeclareMathOperator{\coker}{coker}
\DeclareMathOperator{\coim}{coim}
\DeclareMathOperator{\im}{im}
\DeclareMathOperator{\Ouv}{Ouv}
\DeclareMathOperator{\diag}{diag}
\DeclareMathOperator{\Esp}{Esp}
\DeclareMathOperator{\TOP}{TOP}
\DeclareMathOperator{\Unif}{Unif}
\DeclareMathOperator{\Iind}{\mathbf{Ind}}
\DeclareMathOperator{\Hhomcart}{\mathbf{Homcart}}
\DeclareMathOperator{\Hhomtopess}{\mathbf{Homtopess}}
\DeclareMathOperator{\Ppoints}{\mathbf{Points}}
\DeclareMathOperator{\Ppoint}{\mathbf{Point}}
\DeclareMathOperator{\Ttors}{\mathbf{Tors}}
\DeclareMathOperator{\Ssex}{\mathbf{Sex}}
\DeclareMathOperator{\Ffib}{\mathbf{Fib}}
\DeclareMathOperator{\Vvois}{\mathbf{Vois}}
\DeclareMathOperator{\Ppointess}{\mathbf{Pointess}}
\DeclareMathOperator{\scrHom}{\mathbf{Hom}}
\DeclareMathOperator{\sheafPic}{\mathbf{Pic}}
\DeclareMathOperator{\sheafAut}{\mathbf{Aut}}
\DeclareMathOperator{\Hhomtoplocan}{\mathbf{Homtoplocan}}
\DeclareMathOperator{\Hhomtopan}{\mathbf{Homtopan}}
\DeclareMathOperator{\Pptgeom}{\mathbf{Ptgéom}}
\DeclareMathOperator{\Morsite}{Morsite}
\DeclareMathOperator{\Fib}{Fib}
\DeclareMathOperator{\Fais}{Fais}
\DeclareMathOperator{\kar}{kar}
\DeclareMathOperator{\Groupes}{Groupes}
\DeclareMathOperator{\Restr}{Restr}
\DeclareMathOperator{\Isom}{Isom}
\DeclareMathOperator{\Prol}{Prol}
\DeclareMathOperator{\Res}{Res}
\DeclareMathOperator{\Tors}{Tors}
\DeclareMathOperator{\esp}{esp}
\DeclareMathOperator{\es}{es}
\DeclareMathOperator{\od}{od}
\DeclareMathOperator{\gr}{gr}
\DeclareMathOperator{\sob}{sob}
\DeclareMathOperator{\dfn}{dfn}
\DeclareMathOperator{\Sch}{Sch}
\DeclareMathOperator{\ouv}{ouv}
\DeclareMathOperator{\cons}{cons}
\DeclareMathOperator{\Nerf}{Nerf}
\DeclareMathOperator{\Points}{Points}
\DeclareMathOperator{\Espsob}{Espsob}
\DeclareMathOperator{\Homtop}{Homtop}
\DeclareMathOperator{\Pt}{Pt}
\DeclareMathOperator{\Quot}{Quot}
\DeclareMathOperator{\glob}{glob}
\DeclareMathOperator{\Homfib}{Homfib}
\DeclareMathOperator{\Qouv}{Qouv}
\DeclareMathOperator{\Quotouv}{Quotouv}
\DeclareMathOperator{\ens}{ens}
\DeclareMathOperator{\ab}{ab}
\DeclareMathOperator{\an}{an}
\DeclareMathOperator{\fs}{fs}
\DeclareMathOperator{\red}{red}
\DeclareMathOperator{\Ab}{Ab}
\DeclareMathOperator{\cosq}{cosq}
\DeclareMathOperator{\Ann}{Ann}
\DeclareMathOperator{\Zar}{Zar}


%%%%%%%%%%%%
%%%%% SGA4-2
\DeclareMathOperator{\Sex}{Sex}
\DeclareMathOperator{\vois}{vois}
\DeclareMathOperator{\Vois}{Vois}
\DeclareMathOperator{\Ext}{Ext}
\DeclareMathOperator{\ext}{ext}
\DeclareMathOperator{\Fer}{Fer}
\DeclareMathOperator{\Add}{Add}
\DeclareMathOperator{\Cocart}{Cocart}
\DeclareMathOperator{\cosk}{cosk}
\DeclareMathOperator{\plat}{plat}

\DeclareMathOperator{\faisc}{faisc}
\DeclareMathOperator{\sk}{sk}
\DeclareMathOperator{\ch}{ch}

\def\HH{\mathbf{H}}
\def\qc{\mathrm{qc}}
\def\et{\mathrm{et}}
\def\coh{\mathrm{coh}}

\def\Hvr#1{%
\begin{array}[b]{l}
{\kern2pt\scriptstyle r}\\[-8pt]
{\kern1pt\scriptstyle \vee}\\[-5pt]
H^{#1}
\end{array}}

\def\Hvv#1{%
\begin{array}[b]{l}
\kern2pt\rotatebox{90}{${\scriptscriptstyle\ll}$}\\[-6pt]
H^{#1}
\end{array}}

\def\Hvi#1{%
\begin{array}[b]{l}
{\kern0.2pt\scriptscriptstyle \infty}\\[-8pt]
{\kern1pt\scriptstyle \vee}\\[-5pt]
H^{#1}
\end{array}}

\def\Hv#1{%
\begin{array}[b]{l}
{\kern1pt\scriptstyle \vee}\\[-5pt]
H^{#1}
\end{array}}

\DeclareMathOperator{\cart}{cart}
\DeclareMathOperator{\Ffaisc}{\mathbf{Faisc}}


\def\Ev#1{%
\begin{array}[b]{l}
{\kern1pt\scriptstyle \vee}\\[-5pt]
E'_{#1}
\end{array}}

\DeclareMathOperator{\Cont}{Cont}
\DeclareMathOperator{\Ccont}{\mathscr{C}\!\mathit{ont}}
\DeclareMathOperator{\Site}{Site}
\newcommand\MM{\mathbf{M}}
\def\limtop#1{{\displaystyle{\varprojlim_{#1}\!\mbox{\rm top }}}}
\def\Limtop{\mathop{\varprojlim\!\mbox{\rm top}}}
\DeclareMathOperator{\spec}{spec}
\def\ss{\scriptstyle{\mathscr{S}}}
\DeclareMathOperator{\Gal}{Gal}
\def\smallm{\scriptstyle{\mathscr{M}}}
\DeclareMathOperator{\End}{End}
\DeclareMathOperator{\ass}{ass}
\DeclareMathOperator{\zar}{zar}
\DeclareMathOperator{\Ogg}{Ogg}
\DeclareMathOperator{\res}{res}
\DeclareMathOperator{\tr}{tr}
\DeclareMathOperator{\Supp}{Supp}
\DeclareMathOperator{\codim}{codim}
\DeclareMathOperator{\rk}{rk}
\DeclareMathOperator{\cur}{cur}
\DeclareMathOperator{\cdqc}{cdqc}
\newcommand\DD{\mathbf{D}}
\DeclareMathOperator{\cl}{cl}
\def\JJ{\mathbf{J}}
\def\FF{\mathbf{F}}
\def\BB{\mathbf{B}}
\DeclareMathOperator{\Alget}{Alget}
\DeclareMathOperator{\Ass}{Ass}
\DeclareMathOperator{\supp}{supp}
\DeclareMathOperator{\rad}{rad}
\DeclareMathOperator{\prof}{prof}
\DeclareMathOperator{\cint}{int}
\DeclareMathOperator{\sic}{sic}
\DeclareMathOperator{\Tor}{Tor}
\DeclareMathOperator{\tors}{tors}
\DeclareMathOperator{\Sym}{Sym}
\DeclareMathOperator{\parf}{parf}
\DeclareMathOperator{\tor}{tor}
\DeclareMathOperator{\defn}{defn}
\DeclareMathOperator{\Div}{Div}
\DeclareMathOperator{\Pol}{Pol}
\DeclareMathOperator{\Fex}{Fex}
\DeclareMathOperator{\Lie}{Lie}
\DeclareMathOperator{\tot}{tot}
\DeclareMathOperator{\Aut}{Aut}
\DeclareMathOperator{\HOM}{HOM}
\DeclareMathOperator{\pch}{pch}
\DeclareMathOperator{\Ch}{Ch}
\DeclareMathOperator{\EXT}{EXT}
\DeclareMathOperator{\PIC}{PIC}
\DeclareMathOperator{\Rf}{Rf}
\DeclareMathOperator{\inj}{inj}
\newcommand\VV{\mathbf{V}}

\def\Xv{%
\begin{array}[b]{l}
{\kern1pt\scriptstyle \vee}\\[-5pt]
X
\end{array}}

%%%%%%%
\newcommand\TT{\mathbf{T}}
\newcommand\XX{\mathbf{X}}
\newcommand\EE{\mathbf{E}}
%%%%%%

\def\qedsymbol{\textup{C.Q.F.D.}}
\let\CQFD\qed
\let\cqfd\qed

%%%%%\bigdot
\font\bigsymb=cmsy10 at 4pt
\def\bigdot{{\kern1.2pt\raise 1.5pt\hbox{\bigsymb\char15}}}
%%%%%

\def\Lright{\mathop{\vcenter{\hbox{\hbox{$\longrightarrow$}\lower8pt \hbox{\kern-16pt\hbox{$\longrightarrow$}}}}}}

\def\Aa{{\mathcal A}}
\def\Bb{{\mathcal B}}
\def\Cc{{\mathcal C}}
\def\Dd{{\mathcal D}}
\def\Ee{{\mathcal E}}
\def\Ff{{\mathcal F}}
\def\Gg{{\mathcal G}}
\def\Hh{{\mathcal H}}
\def\Ii{{\mathcal I}}
\def\Jj{{\mathcal J}}
\def\Kk{{\mathcal K}}
\def\Ll{{\mathcal L}}
\def\Mm{{\mathcal M}}
\def\Nn{{\mathcal N}}
\def\Oo{{\mathcal O}}
\def\Pp{{\mathcal P}}
\def\Qq{{\mathcal Q}}
\def\Rr{{\mathcal R}}
\def\Ss{{\mathcal S}}
\def\Tt{{\mathcal T}}
\def\Uu{{\mathcal U}}
\def\Vv{{\mathcal V}}
\def\Ww{{\mathcal W}}
\def\Xx{{\mathcal X}}
\def\Yy{{\mathcal Y}}
\def\Zz{{\mathcal Z}}

\newcommand{\CC}{\mathbf{C}}
\newcommand{\GG}{\mathbf{G}}
\newcommand{\KK}{\mathbf{K}}
\newcommand{\LL}{\mathbf{L}}
\newcommand{\NN}{\mathbf{N}}
\newcommand{\PP}{\mathbf{P}}
\newcommand\QQ{\mathbf{Q}}
\newcommand\RR{\mathbf{R}}
\newcommand\ZZ{\mathbf{Z}}

\def\to{\mathchoice{\longrightarrow}{\rightarrow}{\rightarrow}{\rightarrow}}
\def\from{\mathchoice{\longleftarrow}{\leftarrow}{\leftarrow}{\leftarrow}}
\def\mto{\mathchoice{\longmapsto}{\mapsto}{\mapsto}{\mapsto}}
\def\mfrom{\mathrel{\from\joinrel\mapstochar\,}}
\def\hto{\mathrel{\lhook\joinrel\to}}
\def\hfrom{\mathrel{\from\joinrel\rhook}}

\def\implies{\mathchoice{\Longrightarrow}{\Rightarrow}{\Rightarrow}{\Rightarrow}}
\def\implied{\mathchoice{\Longleftarrow}{\Leftarrow}{\Leftarrow}{\Leftarrow}}
\def\iff{\mathchoice{\Longleftrightarrow}{\Leftrightarrow}{\Leftrightarrow}{\Leftrightarrow}}

\def\To#1{\mathchoice{\xrightarrow{\labelstyle~~#1~}}{\stackrel{#1}{\longrightarrow}}{}{}}
\def\Mto#1{\mathrel{\mapstochar\To{#1}}}
\def\From#1{\mathchoice{\xleftarrow{\labelstyle~~#1~}}{\stackrel{#1}{\longleftarrow}}{}{}}

\newcommand{\isomto}{\overset{\sim}{\longrightarrow}}
\newcommand{\isomfrom}{\overset{\sim}{\longleftarrow}}

\newcommand{\IFF}{\unskip~$\iff$~\ignorespaces}
\newcommand{\THEN}{\unskip~$\implies$~\ignorespaces}

\let\leq\leqslant
\let\geq\geqslant
\let\le\leqslant
\let\ge\geqslant

\def\cf{{cf.}\kern.3em}
\def\Cf{{Cf.}\kern.3em}
\def\ie{{i.e.}\kern.3em}
\def\iev{{i.e.},\xspace}
%\def\loccito{loc.\kern2pt cit.\xspace}
%\def\loccitn{\textit{loc.\kern2pt cit.}\xspace}
%\def\loccit{\sisi{\loccito}{\loccitn}}
\def\EGA{\textup{EGA}\kern.2em}
\def\SGA{\textup{SGA}\kern.3em}

\let\oldS\S
\def\S{\oldS\kern .15em}
\def\ptbl{.\kern .2em}
\def\resp{resp.\kern.3em}
\def\Exp{Exp.\kern.3em}

\let\bar\overline
\let\hat\widehat
\let\tilde\widetilde
\let\dpl\displaystyle
\let\ldots\dots
\let\emptyset\varnothing
\let\skpt\skippointrait
\let\epsilon\varepsilon

\def\sfrac#1#2{{#1}/{#2}}

\DeclareMathOperator{\cd}{cd}
%\DeclareMathOperator{\Hhom}{\mathcal{H}\!\mathit{om}}
\DeclareMathOperator{\Hhom}{\mathbf{Hom}}
%\DeclareMathOperator{\Eext}{\mathcal{E}\!\mathit{xt}}
\DeclareMathOperator{\Eext}{\mathbf{Ext}}
\DeclareMathOperator{\Spec}{Spec}
\DeclareMathOperator{\Tr}{Tr}

\let\emptyset=\varnothing

%%%%%%%%%%%%%%%
%%% expose XVII
%%%%%%%%%%%%%%%

\makeatletter
\let\old@optexpose\optexpose
\newcommand{\empty@optexpose}[1]{}
\newcommand{\paren@optexpose}[1]{\old@optexpose{#1}(}
\newcommand{\refnoexpose}[1]{\let\optexpose\empty@optexpose\ref{#1}\let\optexpose\old@optexpose}%\refnoexpose{X.cor5.2} engendre "5.2" au lieu de "X 5.2"
\newcommand{\neweqref}[1]{\let\optexpose\paren@optexpose\textup{\ref{#1})}\let\optexpose\old@optexpose}%\neweqref{XII.eq4.1} engendre "XII (4.1)" au lieu de "(XII 4.1)"
\makeatother
\newcommand{\refvb}[2][VI${}_{\text{B}}$ ]{\sisi{#1\refnoexpose{#2}}{\ref{#2}}}

\newcommand{\leftexp}[2]{{\vphantom{#2}}^{#1}{#2}}
\newcommand{\AAA}{\mathbf{A}}
\newcommand{\Sf}{\mathfrak{S}}
\DeclareMathOperator*{\quotprojlim}{\text{``$\varprojlim$''}}
\DeclareMathOperator*{\quotinjlim}{\text{``$\varinjlim$''}}
\DeclareMathOperator{\Fl}{F\sisi{\ell}{l}}
\DeclareMathOperator{\Pl}{P\sisi{\ell}{l}}
%\DeclareMathOperator{\cl}{c\sisi{\ell}{l}}
\DeclareMathOperator{\tdf}{tdf}
\DeclareMathOperator{\rg}{rg}
\newcommand{\iInd}{\sisi{I}{i}nd}
\newcommand{\iind}{ind}
\newcommand{\ppro}{pro}
\newcommand{\tTor}{\sisi{T}{t}or}
\newcommand{\ttor}{tor}
\newcommand{\chap}{Chap.}
\newcommand{\Ffex}{\Ff\mathit{ex}}
\newcommand{\sheafk}{\underline{k}}
\newcommand{\sheafn}{\underline{n}}
\newcommand{\sheafGamma}{\underline{\Gamma}}

\sisi{
\DeclareMathOperator{\simarrow}{\rule{.15mm}{1.8mm}\kern -3.4pt\rightsquigarrow}
\newcommand{\sheafell}{\underline{\ell}}
\newcommand{\sheafZZ}{\underline{\ZZ}}
}{
\newcommand{\nderef}[1]{la N.D.E.~(\ref{#1}) page~\pageref{#1}}
}

%%%%%%%%%%%%%%%%%
%%%%% expose XVII
%%%%%%%%%%%%%%%%%



\externaldocument{../01/I.tex}
\externaldocument{../02/II.tex}
\externaldocument{../03/III.tex}
\externaldocument{../04/IV.tex}
\externaldocument{../05/V.tex}
\externaldocument{../05bis/Vbis.tex}
%\externaldocument{../06/VI.tex}
\externaldocument{../07/VII.tex}
\externaldocument{../08/VIII.tex}
\externaldocument{../09/IX.tex}
\externaldocument{../10/X.tex}
\externaldocument{../11/XI.tex}
\externaldocument{../12/XII.tex}
\externaldocument{../13/XIII.tex}
\externaldocument{../14/XIV.tex}
\externaldocument{../15/XV.tex}
\externaldocument{../16/XVI.tex}
\externaldocument{../17/XVII.tex}
\externaldocument{../18/XVIII.tex}
\externaldocument{../19/XIX.tex}

\setboolean{original}{false}

\begin{document}

\begin{center}
VI. Conditions de finitude. Topos et sites fibrés. Applications
aux questions de passage à la limite.\\
A. Grothendieck et J.-L. Verdier
\end{center}

\selectlanguage{english}version : \showgitstatus\selectlanguage{french}
\tableofcontents

\else
\chapter{Conditions de finitude. Topos et sites fibrés. Applications
aux questions de passage à la limite.}
%\addtocontents{toc}{A. Grothendieck et J.-L. Verdier}
\begin{center}
A. Grothendieck et J.-L. Verdier
\end{center}
\fi

\setcounter{pageoriginal}{164}

\setcounter{section}{-1}
\section{Introduction}\etiquette[0]{VI.sec0}

Dans\pageoriginale le présent exposé, nous étudions deux
genres de questions
intimement liées.

Tout d'abord, nous faisons une étude systématique des conditions de
finitude pour les topos, en étudiant successivement la notion de
\emph{quasi-compacité}, de \emph{quasi-séparation} et de
\emph{cohérence} (= quasi-compacité $+$ quasi-séparation),
inspirée
des notions analogues bien connues en théorie des schémas, dans le
cas d'un objet $X$ d'un topos et celui d'une flèche $u : X \to Y$
d'un topos (\S 1), puis pour le topos $E$ lui-même (\S 2), enfin
pour un morphisme de topos $f : E \to E'$ (\S 3). Le \S 4 étudie ces
notions dans le cas particulier d'un topos obtenu par le procédé de
recollement de deux topos (IV 9). Pour les paragraphes suivants, ce
qu'il suffit essentiellement de retenir de ces développements un peu
techniques, c'est la définition d'un \emph{topos cohérent} comme un
topos équivalent à un topos de la forme $C^{\sim}$, où $C$ est
un petit site dans lequel les limites projectives finies sont
représentables et où toute famille couvrante admet une
sous-famille couvrante finie; et celle d'un \emph{morphisme
cohérent} $f : E \to E'$ de topos cohérents comme étant un morphisme
qui peut se décrire (à équivalence près) comme un morphisme associé
à un morphisme de sites $f : C \to C'$, où $C$ et $C'$ sont
comme ci-dessus, et où le foncteur correspondant $f^{\ast} :
C'\to C$ est exact à gauche.

Il est sans doute utile de noter qu'avec les notions de finitude
utilisés ici, et notamment celle de topos cohérent, nous nous
éloignons résolument (et pour la première fois) des espaces
topologiques familiers aux analystes et aux « topologues »
\footnote{Les guillemets indiquent qu'il s'agit d'un
  « \ldots » qui ignore la notion de topos (\cf IV 0.4 pour le
  sens du mot « topologie »).}. Ainsi le topos associé (IV 2.1) à un espace
topologique séparé $X$ n'est cohérent que si $X$ est un ensemble fini!
C'est dire que le présent exposé est entièrement orienté vers les
types de topos provenant de la géométrie algébrique et l'algèbre.
En fait, tous les topos utilisés jusqu'à présent dans ces disciplines (sauf bien sûr
pour la géométrie algébrique par voie transcendante sur le corps des
complexes!) se trouvent être localement cohérents.

En deuxième lieu; nous développons deux \emph{théorèmes de passage à
la limite inductive} pour la cohomologie des topos. L'un (\hyperref[VI.cor5.2]{5.2}) nous
dit que les foncteurs $H^{i}(E, -)$ \pageoriginale sur un topos
cohérent commutent aux limites inductives de faisceaux abéliens. L'autre (\hyperref[VI.subsec8.7]{8.7}) dit
essentiellement que si un topos $X$ est représenté comme une limite
projective filtrante de topos cohérents $X_i$, avec des morphismes
de transition $f_{ij} : X_j \to X_i$ cohérents, alors la cohomologie
de $X$ (à coefficients dans un faisceau abélien quelconque) se
calcule comme limite inductive des cohomologies des $X_i$. Pour
donner un sens précis à cet énoncé, il convient surtout de préciser
la notion de « limite projective de topos », ce qui est fait dans
les \S\S 7.8. C'est là une notion géométrique fort utile, y compris
sans doute dans d'autres contextes que celui des limites projectives
de schémas (qui avait servi de modèle à M.\textsc{Artin} dans sa définition
initiale de cette notion). Nous montrons par exemple (\hyperref[VI.subsec8.4]{8.4}) comment
on peut unifier les diverses notions de « localisé en un point »
(d'un espace noethérien, d'un schéma pour sa topologie habituelle,
voire pour sa topologie étale), en le définissant comme la limite
projective des « voisinages » de ce point, -en parfait accord
avec l'idée intuitive de la notion de localisation en un point.

Les deux théorèmes de passage à la limite seront réexplicités dans
l'exposé suivant dans le cadre particulier de la cohomologie étale des
schémas (\hyperrefext[VII][VII.prop3.3]{3.3} et \hyperref[VI.prop5.7]{5.7}); dans ce cas, ils seront d'ailleurs
constamment utilisés dans toute la suite de ce Séminaire, et
pratiquement partout où on a travaillé jusqu'à présent avec la
cohomologie étale. Le lecteur disposé à admettre ces deux théorèmes,
dans le cas particulier où ils sont énoncés dans l'exposé VII, et
intéressé exclusivement par les applications de la cohomologie étale,
peut donc omettre la lecture du présent exposé. Il est vrai que la
notion \emph{d'objet constructible} d'un topos (\ie dont le morphisme
structural $X \to e$ dans l'objet final $e$ est cohérent) jouera
également un rôle important dans toute la suite du séminaire (alors
qu'elle ne joue qu'un rôle très secondaire dans le présent exposé, ne
serait-ce que parce que dans le cas d'un topos cohérent, elle coïncide
avec la notion d'objet cohérent, qui a ici la vedette). Mais cette
notion est développée dans l'exposé IX, dans le cas particulier des
faisceaux étales sur les schémas, d'une façon indépendante de
l'étude générale faite dans le présent exposé, et elle présente
d'ailleurs dans ce cas des phénomènes spéciaux fort utiles, qui ont servi
de base à l'exposé qui en est fait dans IX (à commencer par la
définition IX 2.3). La rédaction de IX étant antérieure de cinq années
à la présente\pageoriginale rédaction de l'exposé VI, et néanmoins
fort utilisable pour l'usager, nous nous sommes bornés à la fin de l'exposé IX de
prouver l'équivalence de la notion de constructibilité utilisée dans
IX avec celle introduite ici.

\section{Conditions de finitude pour les objets et flèches d'un
  topos}\etiquette[1]{VI.sec1}

\begin{enonce*}{Définition 1.1}\etiquette[1.1]{VI.def1.1}
Soit $C$ un site. Un objet $X$ de $C$ est dit quasi-compact si pour
toute famille couvrante $(X_i \to X)_{i\in I}$, il existe une partie
finie $J$ de l'ensemble d'indices $I$ telle que la sous-famille $(X_i
\to X)_{i \in J}$ soit encore couvrante.
\end{enonce*}

Comme un $𝒰$-topos $E$ est toujours considéré comme un site
(IV 1.1.1), la définition précédente s'applique en particulier à un
objet de $E$. On se ramène d'ailleurs à ce cas, grâce à la

\begin{enonce*}{Proposition 1.2}\etiquette[1.2]{VI.prop1.2}
Soient $C$ un $𝒰$-site, $\epsilon : C \to C^{\sim}$ le foncteur
canonique, $X$ un objet de $C$. Alors $X$ est un objet quasi-compact
du site $C$ si et seulement si $\epsilon(X)$ est un objet
quasi-compact du topos $C^{\sim}$.
\end{enonce*}

Supposons que $X$ soit quasi-compact, prouvons que $\epsilon(X)$
l'est. Soit $(F_i \to \epsilon(X))_{i \in I}$ une famille
épimorphique. Prenons pour tout $i \in I$ une famille épimorphique
$\epsilon(X_{ij}) \to F_i$, $j \in J_i$. Le composé $\epsilon(X_{ij})
\to \epsilon(X)$ s'identifie à un élément $\alpha_{ij}$ de
$\epsilon(X)(X_{ij})$; on ne peut affirmer en général qu'il provient
d'un morphisme $X_{ij} \to X$, mais par construction du faisceau
associé $\epsilon(X)$, on sait que l'on peut trouver (pour $i$, $j$
fixés) une famille couvrante $Y_k \to X_{ij}$ telle que les éléments
images de $\alpha_{ij}$ dans les $\epsilon(X)(Y_k)$ proviennent
d'éléments de $X(Y_k)$, \ie de morphismes $Y_k \to X$. Donc quitte à
raffiner la famille épimorphique $\epsilon(X_{ij}) \to F_i$, on peut
supposer que les composés $\epsilon(X_{ij}) \to \epsilon(X)$
proviennent de morphismes $X_{ij} \to X$.\pageoriginale Pour $i$, $j$
variables, la famille des $\epsilon(X_{ij}) \to \epsilon(X)$ est épimorphique, donc
(II \hyperref[VI.cor4.4]{4.4}) la famille des $X_{ij} \to X$ est couvrante. Par hypothèse sur
$X$, elle admet une sous-famille couvrante finie, et il résulte de
nouveau une famille épimorphismes correspondants $\epsilon(X_{ij}) \to
\epsilon(X)$ forment une famille épimorphique. Comme ceux-ci se
factorisent par un nombre fini des morphismes $F_i \to \epsilon(X)$,
cette famille finie est déjà épimorphique. Cela prouve que
$\epsilon(X)$ est quasi-compact.

Le fait que $\epsilon(X)$ quasi-compact implique $X$ quasi-compact est
d'autre part trivial, grâce à \loccit

Le résultat précédent montre donc que la notion de quasi-compacité
dans les sites se ramène à la notion en question dans les topos.

\begin{enonce*}{Proposition 1.3}\etiquette[1.3]{VI.prop1.3}
Soient $C$ un site, $(X_i \to X)_{i \in I}$ une famille couvrante
finie, avec les $X_i$ quasi-compacts. Alors $X$ est quasi-compact.
\end{enonce*}

Pour le voir, on peut supposer grâce à \hyperref[VI.prop1.2]{1.2} que $C$ est
un topos. Soit $(X'_j \to X)$ une famille couvrante \ie épimorphique
dans $C$, alors pour tout $i \in I$ la famille déduite par le
changement de base $X_{i} \to X$ est épimorphique, donc (puisque
$X_i$ est quasi-compact) admet une sous-famille épimorphique finie,
correspondant à une partie finie $J_i \subset J$ de l'ensemble
d'indices $J$. Posant $J'= \bigcup_{i\in I} J_i$, $J'$ est une
partie finie de $J$ puisque $I$ est fini, et la famille $(X'_j \to
X)_{j \in {J'}}$ est déjà épimorphique puisqu'elle l'est localement,
\cqfd

\begin{enonce*}{Corollaire 1.4}\etiquette[1.4]{VI.cor1.4}
Soit\pageoriginale $(X_i)_{i \in I}$ une famille d'objets d'un topos
$E$, et soit
$X$ la somme. Pour que $X$ soit quasi-compact, il faut et il suffit
que tous les $X_i$ soient quasi-compacts, et que tous les $X_i$ sauf
un nombre fini soient isomorphes à « l'objet vide » $∅_{E}$
de $E$.
\end{enonce*}

La suffisance résulte en effet aussitôt de \hyperref[VI.prop1.3]{1.3} (car la somme ne change
pas quand on laisse tomber les sommandes vides). Pour la nécessité on
note qu'une sous-famille finie $(X_i)_{i \in J}$ doit couvrir $X$, ce
qui implique que pour $i \in I-J$ on a $X_i =
\fprod{X_{i}}{\amalg_{j=J}}{X}X_{j} =
\fprod{\amalg_{j=J}X_{i}}{X_j}{X} = ∅$ (II
  \hyperref[VI.lem4.5.1]{4.5.1}); d'autre part, pour montrer que tout $X_i$ est quasi-compact,
  on note que $X \simeq X_i \amalg {X'}_i$, et que toute famille
  couvrante $Y_j \to X_i$ de $X_i$ définit une famille couvrante $Y_j
  \amalg {X'}_i \to \sisi{X_j}{X_i} \amalg {X'}_i = X$ de $X$, qui admet une
  sous-famille finie couvrante, ce qui implique évidemment que la
  famille finie correspondante des $Y_j \to X_i$ est aussi couvrante,
  ce qui montre que $X_i$ est quasi-compact.

\begin{enonce*}[remark]{Remarque 1.5.1}\etiquette[1.5.1]{VI.rem1.5.1}
Le fait pour un objet $X$ d'un topos $E$ d'être quasi-compact ne
dépend que du topos induit $\sisi{E_{/X'}}{E_{/X}}$ et signifie que l'objet final de
ce dernier est quasi-compact. De même, un objet $X$ d'un site $C$ est
quasi-compact si et seulement si l'objet final du site induit
$C_{/X}$(III 5.1) est quasi-compact.
\end{enonce*}

\setcounter{subsection}{5}
\setcounter{subsubsection}{1}
\subsubsection{}\etiquette[1.5.2]{VI.subsubsec1.5.2}
Soit $X$ un objet d'un topos $E$. Pour que $X$ soit quasi-compact, il
faut et il suffit que toute famille filtrante croissante $(X_i)_{i \in
I}$ de sous-objets de $X$ couvrant $X$ contienne un $X_i$ égal à
$X$. Cette condition ne dépend donc\pageoriginale que de l'ensemble ordonné des
sous-objets de l'objet $X$ (\ie des ouverts du topos induit $E_{/X}$).

\subsubsection{}\etiquette[1.5.3]{VI.subsubsec1.5.3}
Considérons une famille génératrice $(X_i)_{i \in I}$ de $E$. Il
résulte alors des définitions et de \hyperref[VI.prop1.3]{1.3} que : pour
qu'un objet $X$ de $E$ soit quasi-compact, il faut qu'il admette une
famille couvrante par un nombre fini des $X_i$, \ie que $X$ soit
isomorphe à un objet quotient d'une somme finie d'objets de $E$, et
cette condition est suffisante si les $X_i$ sont quasi-compacts.
Notons d'autre part qu'avec l'hypothèse faite sur $E$, on peut
évidemment, quitte à remplacer la famille génératrice envisagée par
une sous-famille, supposer que $I \in 𝒰$. Utilisant le fait
que l'ensemble des quotients d'un objet de $E$ est petit (I 7.5), on
trouve alors que la sous-catégorie pleine $E_{\qc}$ de $E$ formée
des objets quasi-compacts de $E$ est équivalente à une catégorie
$\in 𝒰$, \ie que le cardinal de l'ensemble des classes
d'isomorphie d'objets de $E_{\qc}$ est $\in 𝒰$.

\begin{enonce*}[remark]{Exemples 1.6.1}\etiquette[1.6.1]{VI.exem1.6.1}
Soit $X$ un espace topologique. Un ouvert $U$ de $X$ est un objet
quasi-compact du site $\Ouv(X)$ des ouverts de $X$ si et seulement si
$U$ est un espace quasi-compact pour la topologie induite par
$X$. Plus généralement, un faisceau $F$ sur $X$ est un objet
quasi-compact du topos $\Top(X)$ si et seulement si l'espace étalé $X'$
sur $X$ associé à $F$ est quasi-compact. (Cela résulte de l'assertion
précédente, de 1.5 et du fait que $\Top (X)_{/F}$ est équivalent à
$\Top(X')$ (IV 5.7)).
\end{enonce*}

\setcounter{subsection}{6}
\setcounter{subsubsection}{1}
\subsubsection{}\etiquette[1.6.2]{VI.subsubsec1.6.2}
Considérons\pageoriginale sur la catégorie ($\Sch$) des schémas
(éléments de $𝒰$)
une des topologies $T_i$ de \cite{SGA4-06:6} \hyperref[VI.def6.3]{6.3} (par exemple la topologie de
Zariski, ou la topologie étale, ou la topologie fppf, ou la topologie
fpqc). Pour qu'un schéma soit quasi-compact au sens de cette
topologie, il faut et il suffit que ce soit un schéma quasi-compact au
sens habituel. (Avec les notations de \loccit, on utilise seulement le
fait que toute famille $\in P'$ est finie.)

\begin{enonce*}{Définition 1.7}\etiquette[1.7]{VI.def1.7}
Soient $E$ un topos, $f : X \to Y$ une flèche de $E$. On dit que $f$
est un morphisme quasi-compact si pour toute flèche $Y'\to Y$, avec
$Y'$ quasi-compact, l'objet $X'= \fprod{X}{Y'}{Y}$ est également
quasi-compact. On dit que $f$ est quasi-séparé si le morphisme
diagonal $X \to \fprod{X}{X}{Y}$ est quasi-compact. On dit que $f$ est
cohérent si $f$ est quasi-compact et quasi-séparé.
\end{enonce*}

\setcounter{subsection}{7}
\setcounter{subsubsection}{0}
\subsubsection{}\etiquette[1.7.1]{VI.subsubsec1.7.1}
Explicitant la définition d'un morphisme $f : X \to Y$ quasi-séparé on
voit qu'elle signifie aussi que pour toute double flèche $g_1, g_2 :
X'\rightrightarrows X$ « au-dessus de $Y$ » \ie telle que $fg_1 =
fg_2$, $X'$ quasi-compact implique $\Ker(g_1, g_2)$ quasi-compact.

\begin{enonce*}{Proposition 1.8}\etiquette[1.8]{VI.prop1.8}
Soit $E$ un topos.
\begin{enumerate}
\renewcommand{\theenumi}{\roman{enumi}}
\renewcommand{\labelenumi}{(\theenumi)}
\item Tout isomorphisme dans $E$ est un morphisme cohérent, \ie est
  quasi-compact et quasi-séparé. Le composé de deux morphismes
  quasi-compacts (\resp quasi-séparés, \resp cohérents) est
  quasi-compact (\resp quasi-séparé, \resp cohérent).

\item Soient\pageoriginale $f : X \to Y$ et $g : Y' \to Y$ des
  morphismes dans $E$,   et $f': X'= \fprod{X}{Y'}{Y} \sisi{\to}{\to Y'}$ déduit de $f$ par le changement de
  base $g$. Si $f$ est quasi-compact (\resp quasi-séparé, \resp
  cohérent), il en est de même de $f'$.

\item Soient $f : X \to Y$ et $g : Y \to Z$ \sisi{et $Y \to Z$ }{}des
  morphismes dans $E$. Si $gf$ est quasi-séparé, $f$ est quasi-séparé.

\item Soient $f : X \to Y$ et $g : Y \to Z$ des morphismes dans $E$,
  avec $g$ quasi-séparé. Si $gf$ est quasi-compact (\resp cohérent),
  alors $f$ est quasi-compact (\resp cohérent).
\end{enumerate}
\end{enonce*}

Les démonstrations sont bien connues (\Cf \EGA IV 1 ou EGA I 2ème
édition). Les énoncés (i) et (ii) dans le cas « quasi-compact »
résultent trivialement des définitions; dans le cas
« quasi-séparé », la stabilité par changement de base (ii) résulte aussitôt de la
définition et du cas précédent, de même que le fait que tout
isomorphisme soit quasi-séparé. La stabilité par composition $X \to
Y \to Z$ se voit à l'aide du diagramme
\[\vcenter{
\xymatrix@R=1.5cm@C=1.5cm{
X \ar[d] \ar@/_3.5pc/[dd]^(.2){\Delta_{X/Y}}_{\Delta_{X/Z}}&\\
\fprod{X}{X}{Y} \ar[d]^{u} \ar[rr]&  &Y \ar[d]^{\Delta_{Y/Z}}\\
\fprod{X}{X}{Z}  \ar[rr]&        & \fprod{Y}{Y}{Z}
}}\leqno{(*)}
\]
à carré\pageoriginale cartésien. Par hypothèse $\Delta_{Y/Z}$ est
quasi-compact, donc aussi $u$ qui s'en déduit par changement de
base, et $\sisi{\Delta_{Y/Z}}{\Delta_{X/Y}}$ est quasi-compact, donc aussi le composé $u
\Delta_{X/Y} = \Delta_{X/Z}$. Cela établit (i) et (ii), le cas « cohérent » résultant de la conjonction des deux cas précédents.
Pour prouver le premier cas envisagé dans (iv), on considère le
diagramme suivant à carrés cartésiens
\[\vcenter{
\xymatrix@C=2.75cm@R=1.5cm{
  &     X \ar[r]^{gf}          & Z\\
  X \ar[r]^{\Gamma_f=(\id_X,f)}  \ar[d]
  \ar@/_2pc/[rr]_(.55){f\qquad\qquad \fprod{f}{\id_{Y}}{Z}}  &\fprod{X}{Y}{Z} \ar[r]^{\pr_2}\ar[d]\ar[u]^{\pr_1}& Y\ar[u]_{g}\\
Y \ar[r]^{\Delta_{Y/Z}}&  \fprod{Y}{Y}{Z}& \quad .
}}\leqno{(**)}
\]
Par l'hypothèse $g$ est quasi-séparé, $\Delta_{Y/Z}$ est
quasi-compact, donc $\sisi{\Gamma_F}{\Gamma_f}$ est quasi-compact par changement de
base; de même par hypothèse $gf$ est quasi-compact, donc $\pr_2$ l'est
aussi par changement de base ; par suite $f = \pr_2 \Gamma_f$ est
quasi-compact comme composé de deux morphismes quasi-compacts. Le cas
respé de (iv) résulte aussitôt du cas précédent et de (iii), qu'il
reste à prouver maintenant. Pour ceci on reprend le diagramme
$(\ast)$, en notant que l'hypothèse $gf$ quasi-séparé signifie que
$\Delta_{X/Z} = u \Delta_{X/Y}$ est quasi-compact, et que la
conclusion $f$ quasi-séparé signifie que $\Delta_{X/Y}$ est
quasi-compact\sisi{, et que la conclusion $f$ quasi-séparé signifie que
$\Delta_{X/Y}$ est quasi-compact}{}, ce qui résulte de (iv) une fois
qu'on a prouvé que $u$ est\pageoriginale quasi-séparé. Or $u$ est
manifestement un
monomorphisme, et on a en effet le résultat trivial:

\begin{enonce*}{Corollaire 1.8.1}\etiquette[1.8.1]{VI.cor1.8.1}
Tout monomorphisme est quasi-séparé.
\end{enonce*}

En effet, si $u : X \to Y$ est un monomorphisme, $\Delta_{X/Y}$ est
un isomorphisme donc est quasi-compact, \cqfd

\begin{enonce*}{Corollaire 1.8.2}\etiquette[1.8.2]{VI.cor1.8.2}
Soient $f:X \to Y$, $f':X' \to Y'$ deux $S$-morphismes dans $E$. Si
$f$, $f'$ sont quasi-compacts (\resp quasi-séparés, \resp
cohérents), il en est de même de $\fprod{f}{f'}{S} : \fprod{X}{X'}{S} \to \fprod{Y}{Y'}{S}$.
\end{enonce*}

Cela résulte de façon bien connue de la conjonction de (i) et (ii).

\begin{enonce*}{Remarques 1.9.1}\etiquette[1.9.1]{VI.rem1.9.1}
Soient $E$ une catégorie quelconque, $Q$ une partie de ob stable par
isomorphisme (la donnée de $Q$ revient donc à celle d'une
sous-catégorie strictement pleine de $E$), jouant le rôle d'objets
« quasi-compacts ». Supposons que dans $E$ les produits fibrés
soient représentables. Alors on peut paraphraser la définition
\hyperref[VI.def1.7]{1.7} pour définir la notion de morphismes
$Q$-quasi-compacts, $Q$-quasi-séparés, et $Q$-cohérents. Alors
\hyperref[VI.prop1.8]{1.8} et \hyperref[VI.cor1.8.1]{1.8.1}. sont valables, et
\hyperref[VI.cor1.8.2]{1.8.2} également lorsque dans $E$ les produits de deux
objets sont représentables.
\end{enonce*}

\setcounter{subsection}{9}
\setcounter{subsubsection}{1}
\subsubsection{}\etiquette[1.9.2]{VI.subsubsec1.9.2}
Soient $S$ un objet de $E$, et $f : X \to Y$ un morphisme de
$E_{/S}$. Il résulte aussitôt de la remarque 1.5 que pour que $f$
soit quasi-compact (\resp quasi-séparé, \resp cohérent) il faut et il
suffit que le morphisme correspondant\pageoriginale dans $E$ (déduit
de $f$ par le foncteur « oubli de $S$ ») le soit. En
particulier, si $f : X \to Y$ est un morphisme dans $E$, prenant
ci-dessus $S=Y$, on voit que $f$ est quasi-compact (\resp
quasi-séparé, \resp cohérent) si et seulement si
 le morphisme structural de l'objet $(X, f)$ du topos induit
$E_{/Y}$, de but l'objet final dudit topos, est quasi-compact (\resp
quasi-séparé, \resp cohérent). On dit aussi que (moyennant le
morphisme structural $f$) $X$ \emph{est quasi-compact (\emph{\resp}
  quasi-séparé, \emph{\resp} cohérent) au-dessus de} $Y$, locution qu'on peut
interpréter indifféremment comme se rapportant au topos $E$, ou au
topos induit $E_{/Y}$ (dont $Y$ est considéré comme un objet final).

\subsubsection{}\etiquette[1.9.3]{VI.subsubsec1.9.3}
Un objet $X$ d'un  topos ob $E$ sera dit \emph{constructible} s'il est
cohérent au-dessus de l'objet final. Tout morphisme d'un objet
constructible dans un autre est cohérent (\hyperref[VI.prop1.8]{1.8} (iv)). La sous-catégorie
pleine $E_{\cons}$ de $E$ formée des objets constructibles de $E$ est
stable par $\varprojlim$ finies, comme il résulte aussitôt de \hyperref[VI.cor1.8.2]{1.8.2} et
\hyperref[VI.prop1.8]{1.8} (ii).

\subsubsection{}\etiquette[1.9.4]{VI.subsubsec1.9.4}
Les notions \hyperref[VI.def1.7]{1.7} n'ont guère d'intérêt que dans
les topos $E$ qui
admettent une famille génératrice formée d'objets
quasi-compacts. Dans
ce cas, ces notions sont « de nature locale en bas » (IV 8.5.4),
comme on va voir maintenant.

\begin{enonce*}{Proposition 1.10}\etiquette[1.10]{VI.prop1.10}
Supposons\pageoriginale que le topos $E$ admette une sous-catégorie
génératrice $C$ formée d'objets quasi-compacts, et soit $f : X \to Y$ un morphisme
  dans $E$.
\begin{enumerate}
\renewcommand{\theenumi}{\roman{enumi}}
\renewcommand{\labelenumi}{(\theenumi)}
\item Pour que $f$ soit quasi-compact, il faut et il suffit que pour
  tout morphisme $S \to Y$, avec $S \in \Ob C$, $\fprod{X}{S}{Y}$ soit
  quasi-compact.

\item Soit $Y_i \to Y$, $i \in I$, une famille couvrante. Pour que $f$
  soit quasi-compact (\emph{\resp} quasi-séparé, \emph{\resp}
  cohérent) il faut et il suffit que pour tout $i \in I$, $f_i : X_i = \fprod{X}{Y_i}{Y}
  \to Y_i$ le soit.
\end{enumerate}
\end{enonce*}

Prouvons (i). La nécessité est triviale par définition.
Pour la suffisance, il faut prouver que pour tout objet quasi-compact
$Y'$ et tout morphisme $Y'\to Y$, $X'= \fprod{X}{Y'}{Y}$ est
quasi-compact. Or il existe une famille couvrante $S_i \to Y'$, les
$S_i$ dans $\Ob C$, et comme $Y'$ est quasi-compact, on peut prendre la
famille couvrante en question \emph{finie}. Alors par hypothèse les
${X'}_i = \fprod{X}{S_i}{Y} \simeq \fprod{X'}{S_i}{Y'}$ sont
quasi-compacts, et comme ils couvrent $X'$ (les $S_i$ couvrant $Y'$),
$X'$ est quasi-compact en vertu de \hyperref[VI.prop1.3]{1.3}.

Prouvons (ii). Il suffit de traiter le cas « quasi-compact », car
le cas « quasi-séparé » s'en déduit par la définition \hyperref[VI.def1.7]{1.7}, et le
cas « cohérent » également, par conjonction des deux cas
précédents. Le « il faut » a déjà été vu dans \hyperref[VI.prop1.8]{1.8} (ii), il reste à
voir que si les $f_i$ sont quasi-compacts, alors $f$ l'est, \ie que
pour tout $Y'$ quasi-compact et tout morphisme $Y'\to Y$, l'objet $X'=
\fprod{X}{Y'}{Y}$ est quasi-compact. Or comme $(Y_i \to Y)$ est
couvrante,\pageoriginale il existe des $S_{\alpha}$ au-dessus de $Y$
couvrant $Y'$, tels que les morphismes composés $S_{\alpha}\to Y'\to
Y$ se factorisent chacun par un $Y_{i}$. Comme $C$ est génératrice, on
peut prendre les $S_{\alpha}$ dans $\Ob C$, et
comme $Y'$ est
quasi-compact, on peut prendre la famille couvrante finie. Alors les
$\fprod{X'}{S_{\alpha}}{Y'}$ forment une famille couvrante finie de
$X'$, et on est réduit par \hyperref[VI.prop1.3]{1.3} à prouver que chaque
$\fprod{X'}{S_{\alpha}}{Y'} = \fprod{X}{S_{\alpha}}{Y}$ est quasi
compact. Or, choisissant une factorisation de $S_{\alpha} \to Y$ par
un $Y_i$, l'objet envisagé est aussi isomorphe à
$\fprod{X_i}{\sisi{S}{S_{\alpha}}}{Y_i}$, qui est bien quasi-compact puisque $X_i$ est
quasi-compact sur $Y_i$ et $S_{\alpha}$ est quasi-compact.

\begin{enonce*}{Corollaire 1.11}\etiquette[1.11]{VI.cor1.11}
Soient $g : Y \to Z$ un morphisme dans le topos $E$, $(f_i : X_i \to
Y)_{i \in I}$ une famille couvrante de morphismes de but $Y$. Alors:
\begin{enumerate}
\renewcommand{\theenumi}{\roman{enumi}}
\renewcommand{\labelenumi}{(\theenumi)}
\item Supposons $I$ finie. Si les $gf_i$ sont quasi-compacts, il en
  est de même de $g$.

\item Supposons les $f_i$ quasi-compacts. Si les $gf_i$ sont
  quasi-séparés, il en est de même de $g$. Si $I$ est fini, et si les
  $gf_i$ sont cohérents, $g$ est cohérent.
\end{enumerate}
\end{enonce*}

Le cas (i) résulte aussitôt des définitions et de \hyperref[VI.prop1.3]{1.3},
(sans hypothèse sur $E$). Pour prouver (ii), il suffit en vertu de
(i) de traiter le cas non respé. Considérons le diagramme
\hyperref[VI.prop1.8]{1.8} $(\ast)$ avec $X$ remplacé par un $X_i$. Comme la
famille des $f_i$ est couvrante, il en est de même de la famille des
$\fprod{f_i}{\id_Y}{Z}$ qui s'en déduit par changement de base,
donc\pageoriginale en vertu de \hyperref[VI.prop1.10]{1.10} (ii), pour prouver
que $\Delta_{Y /Z}$ est quasi-compacte, il suffit de prouver que les
$\Gamma_{f_i}$ le sont. Or $gf_i$ étant quasi-séparé par hypothèse,
il en est de même de $\pr_2 : \fprod{X_i}{Y}{Z} \to Y$ qui s'en
déduit par changement de base, d'autre part $\pr_2 \Gamma_{f_i} =
f_i$ est quasi-compact par hypothèse. Il en est donc de même de
$\Gamma_{f_i}$ en vertu de \hyperref[VI.prop1.8]{1.8} (iv) appliqué au
morphisme composé de $\pr_2$ et de $\Gamma_{f_i}$, \cqfd

\begin{enonce*}{Corollaire 1.12}\etiquette[1.12]{VI.cor1.12}
Soient $(X_i)_{i \in I}$ une famille finie d'objets de $E$, $X$ la
somme des $X_i$, $f_i : X_i \to Y$ des morphismes dans $E$, et $f : X
\to Y$ le morphisme correspondant. Pour que $f$ soit quasi-compact
(\resp quasi-séparé, \resp cohérent) il faut et il suffit que les
$f_i$ le soient ; dans le cas « quasi-séparé », l'hypothèse « I
finie » peut-être omise.
\end{enonce*}

La suffisance de la condition est un cas particulier de \hyperref[VI.cor1.11]{1.11}, compte
tenu que les morphismes canoniques $f_i : X_i \to X$ forment une
famille couvrante finie, et que ce sont des morphismes quasi-compacts,
comme il résulte aussitôt de la définition et de \hyperref[VI.cor1.4]{1.4}. La nécessité
résulte de la transitivité \hyperref[VI.prop1.8]{1.8} (i), compte tenu que les $f_i$ sont
même cohérents (car ils sont quasi-séparés en vertu de \hyperref[VI.cor1.8.1]{1.8.1}).

\begin{enonce*}{Définition 1.13}\etiquette[1.13]{VI.def1.13}
Soit $X$ un objet d'un topos $E$. On dit que $X$ est un \emph{objet
quasi-séparé du topos} $E$ si pour tout objet $S$ quasi-compact
de $E$, tout morphisme de $S$ dans $X$ est quasi-compact, \ie (\hyperref[VI.def1.7]{1.7}) si
pour deux objets quasicompacts $S$, $T$ de $E$ et des morphismes $S
\to X$, $T \to X$, le produit fibré $\fprod{S}{T}{X}$ est toujours
quasi-compact. On dit que $X$ est un objet cohérent du topos $E$ si
$X$ est quasi-compact et quasi-séparé.
\end{enonce*}

\begin{enonce*}{Proposition 1.14}\etiquette[1.14]{VI.prop1.14}
Soit\pageoriginale $f : X \to Y$ un morphisme dans un topos $E$.
\begin{enumerate}
\renewcommand{\theenumi}{\roman{enumi}}
\renewcommand{\labelenumi}{(\theenumi)}
\item Si $Y$ est quasi-compact, alors $f$ quasi-compact implique $X$
  quasi-compact. Si $Y$ est quasi-séparé, alors $X$ quasi-compact
  implique $f$ quasi-compact. Si $Y$ est cohérent, $f$ quasi-compact
  équivaut à $X$ quasi-compact.

\item Si $Y$ est quasi-compact (\emph{\resp} quasi-séparé,
  \emph{\resp} cohérent) et
  si $f$ est quasi-compact (\emph{\resp} quasi-séparé,
\emph{\resp} cohérent) alors
  $X$ est quasi-compact (\emph{\resp} quasi-séparé, \emph{\resp} cohérent).
\end{enumerate}
\end{enonce*}

Les premières deux assertions de (i) sont contenues trivialement dans
les définitions de morphisme quasi-compact (\hyperref[VI.def1.7]{1.7}) \resp d'objet
quasi-séparé (\hyperref[VI.def1.13]{1.13}), la troisième résulte de la conjonction des deux
premières. Le premier cas de (ii) n'est qu'une redite, le troisième
résulte encore de la conjonction des deux premiers, reste à traiter le
cas quasi-séparé. Soit donc $S$  un objet quasi-compact de $E$ et $g:
S\to X$ un morphisme, prouvons que $\sisi{f}{g}$ est quasi-compact. Comme $Y$
est quasi-séparé, $fg$ est quasi-compact, et comme $f$ est
quasi-séparé, il en résulte (\hyperref[VI.prop1.8]{1.8} (iii)) que $g$ est quasi-compact,
\cqfd

\begin{enonce*}{Corollaire 1.15}\etiquette[1.15]{VI.cor1.15}
Tout sous-objet d'un objet quasi-séparé de $E$ est quasi-séparé. Soit
$(X_i)_{i \in I}$ une famille d'objets de $E$, avec $I \in \sisi{U}{𝒰}$; alors
$X = \amalg_i X_i$ est quasi-séparé (\emph{\resp}
cohérent) si et seulement si les $X_i$ le sont.
\end{enonce*}

La\pageoriginale première assertion résulte de \hyperref[VI.prop1.14]{1.14} (ii) et de \hyperref[VI.cor1.8.1]{1.8.1}, mais est
également triviale sur la définition. Pour la deuxième non respée,
comme les $X_i \to X$ sont des monomorphismes, il reste à prouver le
« il suffit », qui résulte facilement des définitions et de \hyperref[VI.prop1.3]{1.3}; le
cas respé résulte du cas traité et de \hyperref[VI.cor1.4]{1.4}.

\begin{enonce*}{Corollaire 1.16}\etiquette[1.16]{VI.cor1.16}
Sous les conditions de \hyperref[VI.prop1.10]{1.10} (ii), si les $Y_i$ sont cohérents, alors
$f$ est quasi-compact si et seulement si les $X_i = \fprod{X}{Y_i}{Y}$
sont des objets quasi-compacts de $E$.
\end{enonce*}

En effet, en vertu du critère \hyperref[VI.prop1.10]{1.10} (ii), il faut exprimer que les $f_i
: X_i \to Y_i$ sont quasi-compacts, ce qui équivaut à $X_i$
quasi-compacts en vertu de \hyperref[VI.prop1.14]{1.14} (i).

\begin{enonce*}{Corollaire 1.17}\etiquette[1.17]{VI.cor1.17}
Soient $E$ un topos admettant une famille génératrice formée
d'objets quasi-compacts, et $(g_i : Y_i \to Y)_{i \in I}$ une
famille couvrante (\resp couvrante finie) dans $E$, avec les $Y_i$
cohérents. Pour que $Y$ soit quasi-séparé (\resp cohérent), il faut
et il suffit que pour tout $i \in I$, $g_i$ soit quasi-compact, ou
encore que pour tout couple d'indices $i$, $j \in I$,
$\fprod{Y_i}{Y_j}{Y}$ soit un objet quasi-compact de $E$.
\end{enonce*}

Le cas respé résulte aussitôt du cas non respé et de \hyperref[VI.prop1.3]{1.3}. La nécessité
de la première condition est triviale par définition, prouvons qu'elle
est suffisante. Soit donc $f : X \to Y$ un morphisme, avec\pageoriginale $X$
quasi-compact, prouvons que $f$ est quasi-compact. Or comme les $g_i$
sont quasi-compacts, il s'ensuit que les $X_i = \fprod{X}{Y_i}{Y}$
sont quasi-compacts, ce qui implique que $f$ l'est en vertu de
\hyperref[VI.cor1.16]{1.16}. Exprimant enfin la quasi-compacité des $g_i$ à l'aide du même
critère \hyperref[VI.cor1.16]{1.16}, on trouve le deuxième critère de \hyperref[VI.cor1.17]{1.17}. En particulier:

\begin{enonce*}{Corollaire 1.17.1}\etiquette[1.17.1]{VI.cor1.17.1}
Soient $E$ comme dans \hyperref[VI.cor1.17]{1.17}, $X$ un objet cohérent de
$E$, $\Rr \rightrightarrows X$ une relation d'équivalence dans $X$,
$Y = X / \Rr$. Alors les conditions suivantes sont équivalentes:
\begin{enumerate}
\renewcommand{\theenumi}{\roman{enumi}}
\renewcommand{\labelenumi}{(\theenumi)}
\item $Y$ est cohérent,
\item $f : X \to Y$ est quasi-compact,
\item $\Rr$ est quasi-compact.
\end{enumerate}
\end{enonce*}

\begin{enonce*}{Corollaire 1.18}\etiquette[1.18]{VI.cor1.18}
Soient $E$ un topos admettant une sous-catégorie génératrice $C$
formée d'objets quasi-compacts, $Y$ un objet de $E$. Pour que $Y$ soit
quasi-séparé, il faut et il suffit que tout morphisme $X \to Y$, avec
$X \in \ob C$, soit quasi-compact.
\end{enonce*}

La nécessité est triviale par définition, les $\sisi{Y}{X} \in \ob C$ étant
quasi-compacts par hypothèse. Pour prouver la suffisance, soit $Z$
un objet quasi-compact de $E$, prouvons que tout morphisme $\sisi{Z'}{Z}\to Y$
est quasi-compact. Pour ceci, il suffit en vertu de
\hyperref[VI.prop1.10]{1.10} (i) de vérifier que pour tout $X \to Y$, avec $X
\in \ob C$, $\fprod{Z}{X}{Y}$ est quasi-compact, ce qui résulte en
effet du fait que les $X \to Y$ sont quasi-compacts par hypothèse,
et que $Z$ est quasi-compact.

\begin{enonce*}{Corollaire 1.19}\etiquette[1.19]{VI.cor1.19}
Soient \pageoriginale $E$ un topos admettant une sous-catégorie
génératrice formée d'objets quasi-compacts, $(g_i : Y_i \to Y)_{i \in
  I}$ une famille couvrante dans $E$, avec les $Y_i$ quasi-séparés et
les $g_i$ quasi-compacts. Alors $Y$ est quasi-séparé.
\end{enonce*}

En effet, soit $X$ un objet quasi-compact de $E$, montrons que tout
morphisme $f : X \to Y$ est quasi-compact. En vertu de \hyperref[VI.prop1.10]{1.10}(ii) il
suffit de prouver que les morphismes induits $f_i : X_i =
\fprod{X}{Y_i}{Y} \to Y_i$ sont quasi-compacts. Or comme $g_i$ est
quasi-compact, $X_i$ est quasi-compact ($X$ l'étant), et comme $Y_i$
est quasi-séparé, il s'ensuit bien que $f_i$ est quasi-compact.

\begin{enonce*}[remark]{Remarque 1.20.1}\etiquette[1.20.1]{VI.rem1.20.1}
Le lecteur habitué à l'usage des termes « quasi-compact » et « quasi-séparé » en géométrie algébrique s'attendra à certaines
autres relations entre les notions « absolues »(concernant les
objets du topos $E$) et les notions « relatives » (concernant
des flèches de $E$) envisagées, par exemple : si $f: X \to Y$ est un
morphisme dans $E$ tel que $X$ soit quasi-compact-séparé, alors $f$
est quasi-séparé. De telles propriétés ne sont valables que
moyennant des conditions de finitude restrictive sur le topos $E$,
plus fortes que celles envisagées dans \hyperref[VI.prop1.10]{1.10}, conditions
que nous étudierons au numéro suivant.
\end{enonce*}

\setcounter{subsection}{20}
\setcounter{subsubsection}{1}
\subsubsection{}\etiquette[1.20.2]{VI.subsubsec1.20.2}
Il résulte encore de 1.5 que le fait que pour un objet $X$ du topos
$E$ d'être quasi-séparé ne dépend que du topos induit $E_{/X}$, et
signifie que l'objet final dudit topos est quasi-séparé, \ie que
dans $E_{/X}$, le produit de deux objets quasi-compacts est
quasi-compact.

\subsection{Cas des sites}\etiquette[1.21]{VI.subsec1.21}\pageoriginale

Soient  $C$ un $𝒰$-site, $E$ le topos des $𝒰$-faisceaux
sur $C$, $f$ une flèche de $C$, $X$ un objet de $C$. On dit que $f$ est
un \emph{morphisme quasi-compact} (\resp quasi-séparé) \emph{du site}
$C$ si $\epsilon(f)$ est un morphisme quasi-compact (\resp
quasi-séparé) du topos $E$; et de même que $X$ est un \emph{objet
  quasi-séparé du site} $C$ si $\epsilon(X)$ est un
objet quasi-séparé du topos $E$ (comparer 1.2). Lorsque $C$ est un
$𝒰$-topos muni de sa topologie canonique, alors le foncteur
canonique $\epsilon : C \to E$ est une équivalence de catégories, et
par suite la définition précédente est compatible avec les
définitions antérieures 1.7 et \hyperref[VI.def1.13]{1.13}. Notons également
que la définition envisagée ne change pas quand on remplace
l'univers $𝒰$ par un autre univers $𝒱$ tel que $C$
soit aussi un $𝒱$-site, du moins lorsque $C$ admet une
famille topologiquement génératrice formée d'objets quasi-compacts.
Pour le voir, on est ramené aussitôt au cas où $𝒰
\subset 𝒱$, donc $E$ est un $𝒱$-site, et à prouver
que pour toute flèche $f$ du $𝒰$-topos $E$, $f$ est
quasi-compacte (\resp quasi-séparé) si et seulement si elle définit
une flèche quasi-compacte dans le topos $E'$ des
$𝒱$-faisceaux sur $E$, et de même que pour tout objet $X$ de
$E$, $X$ est quasi-séparé si et seulement si l'objet correspondant
de $E'$ est quasi-séparé. Notons que nous connaissons déjà (sans
hypothèse sur $E$) l'énoncé analogue pour le cas d'un objet
quasi-compact (\hyperref[VI.prop1.2]{1.2}). Le cas « $f$ quasi-compact »
s'en déduit, grâce au critère \hyperref[VI.prop1.10]{1.10} (i) appliqué
successivement dans $E$ et\pageoriginale dans $E'$, pour la même
famille génératrice $(X_i)_{i \in I}$ de $E$ formée d'objets
quasi-compacts; d'où également le cas « $f$ quasi-séparé »,
qui signifie que le morphisme diagonal correspondant est
quasi-compact. Enfin pour le cas « $X$ quasi-séparé », on est
encore ramené au cas « $f$ quasi-compact », grâce au critère
\hyperref[VI.cor1.18]{1.18}.

\subsection{}\etiquette[1.22]{VI.subsec1.22}
{\textit{\textbf{Exemples.}}} Reprenons l'exemple
\hyperref[VI.subsubsec1.6.2]{1.6.2} de la catégorie $(\Sch)$ avec une topologie
$T_i$, qui est un $𝒱$-site pour $𝒱$ convenable,
admettant la sous-catégorie des schémas affines comme catégorie
génératrice topologique formée d'objets quasi-compacts. On laisse au
lecteur le soin de vérifier qu'une flèche $f: X \to Y$ de ce site
est quasi-compacte (\resp quasi-séparé) si et seulement si c'est un
morphisme quasi-compact (\resp quasi-séparé) de schémas au sens
habituel (\cite{SGA4-06:3} ou \cite{SGA4-06:4}); de même un objet du site est
quasi-compact (\resp quasi-séparé) si et seulement si c'est un
schéma quasi-compact (\resp quasi-séparé) au sens habituel de
loc.cit. Il y a lieu de dire d'ailleurs, comme ici, \emph{morphisme
cohérent de schémas, schéma
  cohérent}, au lieu de « morphisme quasi-compact et quasi-séparé de
schémas », « schéma quasi-compact et quasi-séparé », comme on le
fait d'ailleurs dans \EGA I $2^{\text{ème}}$ édition.

\setcounter{subsection}{22}
\subsubsection{}\etiquette[1.22.1]{VI.subsubsec1.22.1}

Pour un schéma donné $X \in 𝒰$, on peut aussi regarder le topos
$\Top(X)$ associé à l'espace topologique sous-jacent, ou le topos
$\Top (X_{\et})$ (\resp $\Top(X_{\text{fppf}})$), associé au site des
$X$-schémas étales (\resp localement\pageoriginale de présentation
finie) éléments de $𝒰$, avec la topologie étale (\resp fppf). C'est un
$𝒰$-topos qui admet une famille génératrice formée d'objets cohérents,
correspondants aux objets affines du site de définition;

Ceci dit, l'objet final de ce topos est quasi-compact (\resp
quasi-séparé, \resp cohérent) si et seulement si le schéma $X$ est
quasi-compact (\resp quasi-séparé, \resp cohérent). De même, un
morphisme du site des espaces étalés sur $X$ (\resp du site
$X_{\text{et}}$ des schémas étales sur $X$, \resp du site
$X_{\text{fppf}}$ des schémas localement de présentation finie sur
$X$) est quasi-compact (\resp quasi-séparé, \resp cohérent) si et
seulement si c'est un morphisme de schémas qui est (au sens habituel
de \loccit) quasi-compact (\resp quasi-séparé, \resp cohérent).

\begin{enonce*}{Théorème 1.23}\etiquette[1.23]{VI.thm1.23}
Soient $E$ un $𝒰$-topos, $X$ un objet de $E$.
\begin{enumerate}
\renewcommand{\theenumi}{\roman{enumi}}
\renewcommand{\labelenumi}{(\theenumi)}
\item Pour que $X$ soit quasi-compact, il faut et il suffit que pour
  tout système inductif filtrant $(Y_i)_{i \in I}$ de $E$, avec $I \in
  𝒰$, l'application naturelle
\begin{equation*}
\varinjlim  \Hom(X, Y_i) \to \Hom(X, \varinjlim Y_i)\tag{1.23.1}
\end{equation*}\etiquette[1.23.1]{VI.eq1.23.1}
soit injective.

\item Supposons que $E$ admette une sous-catégorie génératrice $C$
  formée d'objets quasi-compacts. Pour que $X$ soit cohérent, il faut
  que pour toute donnée comme dans \emph{(i)}, l'application
  \emph{(1.23.1)} soit bijective.
\end{enumerate}
\end{enonce*}

\begin{enumerate}
\renewcommand{\theenumi}{\roman{enumi}}
\renewcommand{\labelenumi}{(\theenumi)}
\item \underline{Nécessité}.
Supposons\pageoriginale $X$ quasi-compact, il faut prouver que
  si $i\in I$ et $f_i$, $g_i : X \rightrightarrows Y_i$ sont tels que
  les composés $f$, $g$ de $f_i$, $g_i$ avec $Y_i \to Y = \varinjlim
  Y_j$ sont égaux, alors il existe $j \geq i$ tel que les composés
  $f_j$, $g_j$ avec $Y_i \to Y_j$ sont déjà égaux. Or comme dans $E$
  les limites inductives filtrantes commutent aux limites projectives
  finies, et en particulier aux noyaux, on a
$$
\Ker (f, g) = \mathop{\lim}_{j\geq i} \Ker (f_j, g_j){\quoi}.
$$
Par hypothèse $\Ker(f, g) = X$, donc $X$ est la limite de la famille
filtrante croissante de ses sous-objets $\Ker(f_j, g_j)$, donc comme
$X$ est quasi-compact, il est égal à un de ces sous-objets, donc il
existe $j \geq i$ tel que $f_j = g_{j}$.

\underline{Suffisance}. Pour prouver que $X$ est quasi-compact, il revient au
même (\hyperref[VI.subsubsec1.5.2]{1.5.2}) de prouver que pour toute famille filtrante croissante
$(X_i)_{i \in I}$ de sous-objets de $X$ dont la $\varinjlim$ est $X$,
il existe $i \in I$ tel que $X_i = X$. Or comme dans $E$ tout
monomorphisme est effectif (II 4.0), il s'ensuit que
$$
X_i = \Ker (f_i, g_i : X \rightrightarrows X \amalg_{X_{i}} X){\quoi}.
$$
Soit $Y_i = X \amalg_{X_i} X$, et considérons le système inductif formé
par les $Y_i$. Si $Y$ en est la limite inductive, on a $Y \simeq X
\amalg_{\varinjlim X_i} X = X \amalg_X X = X$.
Par suite, pour un $i \in I$ fixé, $f_i$, $g_i$ donnent le
même élément de $\Hom(X, Y)$, donc par hypothèse il existe $j \geq i$
tel que $f_j = g_j$ \ie $X_j = X$, \cqfd

\item Il reste\pageoriginale à prouver que sous les conditions indiquées, tout
  morphisme $X \to Y$ provient d'un morphisme $X \to Y_i$ pour un $i
  \in I$ convenable. Comme $C$ engendre $E$ et que les $Y_i$ couvrent
  $Y$, nous pouvons couvrir $X$ par des objets $X_{\alpha}$ de $C$, de
  telle façon que chacun des composés $X_{\alpha} \to X \to Y$ se
  factorise par un des $Y_i$. Comme $X$ est quasi-compact, on peut
  supposer la famille des $X_{\alpha}$ finie, donc que les morphismes
  $X_{\alpha} \to Y$ se factorisent par un $Y_i$ fixe en des
  $g_{\alpha i} : X_{\alpha} \to Y_i$. Comme $X$ est quasi-séparé, les
  $\fprod{X_{\alpha}}{X_{\beta}}{X}$ sont quasi-compacts. D'ailleurs
  pour tout $(\alpha, \beta)$, les composés de $\pr_1$ et $\pr_2$ sur
  $\fprod{X_{\alpha}}{X_{\beta}}{X}$ avec $X_{\alpha}$, $X_{\beta} \to
  Y_i$ sont
\[
\xymatrix{
\fprod{X_{\alpha}}{X_{\beta}}{X} \ar[rd] \ar[dd]& &\\
                            &X_{\alpha} \ar[ld] \ar[rrrr]&   & & & Y_i
\ar[lld]\\
X    \ar[rrr]                       &          & & Y  &
}
\]
tels que leurs composés avec $Y_i \to Y$ sont égaux. En
  vertu de(i), il existe donc un $j \geq i$ tel que les composés avec
  $Y_i \to Y_j$ soient déjà égaux. Comme l'ensemble des couples
  $(\alpha, \beta)$ est fini, on peut prendre $j$ indépendant de ce
  couple. Par suite les morphismes $g_{\alpha j}$ proviennent par
  composition d'un morphisme $g_j : X \to Y_j$, dont le composé avec
  $Y_j \to Y$ est d'ailleurs le morphisme $X \to Y$ donné, puisqu'il
  en est ainsi après composition avec les $Y_j \to X$. Cela achève la
  démonstration.
\end{enumerate}

\begin{enonce*}[remark]{Remarques 1.23.2}\etiquette[1.23.2]{VI.rem1.23.2}
L'hypothèse sur $E$ dans (ii) peut être remplacée par une hypothèse
supplémentaire \emph{sur} $X$, savoir que $X$ admet un système fondamental de
familles couvrantes par des objets quasi-compacts.
\end{enonce*}

\begin{enonce*}{Corollaire 1.24}\etiquette[1.24]{VI.cor1.24}
Soient\pageoriginale $E$ un topos, et soit $C$ une sous-catégorie
strictement pleine
de $E$ formée d'objets cohérents. Utilisant le fait que dans $E$ les
$𝒰$-limites inductives filtrantes sont représentables, on trouve \emph{(I
  8.6.1)} un foncteur canonique.
\begin{equation*}
\Ind (C) \to E{\quoi}.\tag{1.24.1}
\end{equation*}\etiquette[1.24.1]{VI.eq1.24.1}
Ce foncteur est pleinement fidèle. Si $E$ admet une sous-catégorie
génératrice formée d'objets cohérents, et si $C$ est stable par
facteurs directs, alors les conditions suivantes sont équivalentes:
\begin{enumerate}
\renewcommand{\theenumi}{\roman{enumi}}
\renewcommand{\labelenumi}{(\theenumi)}
\item Le foncteur \emph{(1.24.1)} est essentiellement surjectif, \ie c'est
  une équivalence de catégories.

\item $C$ est une sous-catégorie génératrice de $E$, stable par
  limites inductives finies.

\item La catégorie $C$ est égale à la sous-catégorie pleine $E_{\coh}$
  de $E$ formée de tous les objets cohérents de $E$, et la condition
  nécessaire \emph{\hyperref[VI.thm1.23]{1.23} (ii)} de cohérence pour un objet de $E$ est
  aussi suffisante.
\end{enumerate}
\end{enonce*}

Soit $E_{PF}$ la sous-catégorie pleine de $E$ définie dans I 8.7.6. En
vertu de I 8.7.5 a), l'énoncé \hyperref[VI.thm1.23]{1.23} (ii) équivaut à la deuxième
inclusion de
$$
C \subset E_{\coh} \subset E_{PF}{\quoi},
$$
et l'inclusion composée $C \supset E_{PF}$ signifie que (1.24.1) est
pleinement fidèle. Comme\pageoriginale par hypothèse la sous-catégorie
$E_{\coh}$ de
$E$ est génératrice, il en est a fortiori de même de $E_{PF}$; d'autre
part dans $E$ les limites projectives finies sont représentables, de
sorte que nous pouvons appliquer I 8.7.7. La condition (iii) de \hyperref[VI.cor1.24]{1.24}
s'écrit aussi $C = E_{\coh} = E_{PF}$ et équivaut à la condition $C =
E_{PF}$, qui n'est autre que la condition (ii bis) de \loccit (compte
tenu que par hypothèse,  $E_{\coh}$ donc $E_{PF}$ est une
sous-catégorie génératrice de $E$). De même la condition (ii) de \hyperref[VI.cor1.24]{1.24}
n'est autre que la condition (iii bis) de \loccit. Donc l'équivalence
des conditions (i), (ii bis) et (iii bis) dans \loccit donne
l'équivalence des conditions (i), (ii) et (iii) de \hyperref[VI.cor1.24]{1.24}, \cqfd

La démonstration (ou plus précisément, l'invocation de I 8.7.7 (ii
bis)) fournit également la partie a) du

\begin{enonce*}{Corollaire 1.24.2}\etiquette[1.24.2]{VI.cor1.24.2}
Soit $E$ un topos admettant une sous-famille génératrice formée
d'objets cohérents, et soit $E_{PF}$ la sous-catégorie strictement
pleine des objets $X$ de $E$ tels que le foncteur covariant $\Hom(X,
-)$ qu'il représente commute aux limites inductives filtrantes.
Alors:
\begin{enumerate}
\renewcommand{\theenumi}{\alph{enumi}}
\renewcommand{\labelenumi}{\theenumi)}
\item Le foncteur naturel
$$
\Ind (E_{PF}) \to E
$$
est une équivalence de catégories.

\item Pour qu'un objet $X$ de $E$ appartienne à $E_{PF}$, il faut et
  il suffit qu'il soit isomorphe à un facteur direct ($=$ image d'un
  projecteur (I, 10.6)) du conoyau\pageoriginale d'une double flèche $X_1
  \rightrightarrows X_0$, avec $X_1$ et $X_0$ cohérents.
\item La sous-catégorie $E_{PF}$ de $E$ est stable par limites
  inductives finies, et est équivalente à une petite catégorie.
\end{enumerate}
\end{enonce*}

Il reste à prouver b) et c). La première assertion de c) est triviale
grâce à I 2.8, et implique le « il suffit » de b) grâce à \hyperref[VI.thm1.23]{1.23}
(ii). Pour le « il faut », on utilise \hyperref[VI.thm1.23]{1.23} (i) qui montre que $X$
est quasi-compact, d'où l'existence d'un épimorphisme $X_{0} \to X$, avec
$X_0$ cohérent, compte tenu du fait que $E_{\coh}$ engendre $E$ et est
stable par sommes finies. Soit $R = \fprod{X_0}{X_0}{X}$, de sorte que
$X$ s'identifie au conoyau de $R \rightrightarrows X_0$. En vertu de
a), on a $R = \varinjlim R_i$, limite inductive filtrante des coker
$(R_i \rightrightarrows X_0) = X_i$. D'après l'hypothèse $X \in \Ob
E_{PF}$, pour $i$ assez grand le morphisme $u_i : X'_i \to X$ admet
une section $v : X \to X'_i$. Soit $w = vu_i : X'_i \to
X'_i$, de sorte que $w^2 = w$ et $X \simeq \Coker (w, \id_{X'_i}
: X'_i \rightrightarrows X'_i)$. Comme $R_i$ est à nouveau
quasi-compact, on le couvre par $X_{1} \to R_i$ avec $X_i$
cohérent, d'où $X_i \simeq \Coker (X_1 \rightrightarrows X_0)$, ce qui
prouve b).

Enfin, comme $E_{PF}\subset E_{\text{qu.cpct}}$, la dernière assertion de c)
résulte de \hyperref[VI.subsubsec1.5.3]{1.5.3}.

\begin{enonce*}{Corollaire 1.25}\etiquette[1.25]{VI.cor1.25}
Soit $E$ un topos tel que la sous-catégorie pleine $C$ de $E$ formée
des objets cohérents soit génératrice. Considérons le foncteur
canonique
$$
\varphi : \Ind(C) \to E{\quoi}.
$$
Les conditions\pageoriginale suivantes sont équivalentes:
\begin{itemize}
\item[(i)] Le foncteur $\varphi$ est essentiellement surjectif, \ie tout
  objet de $E$ est limite inductive filtrante d'objets cohérents.

\item[(i bis)] Le foncteur $\varphi$ est une équivalence de catégories.

\item[(ii)] Toute limite inductive finie dans $E$ d'objets cohérents est
  un objet cohérent.

\item[(ii bis)] Le conoyau dans $E$ d'une double flèche d'objets
  cohérents est un objet cohérent.

\item[(iii)] La réciproque de \emph{\hyperref[VI.thm1.23]{1.23} (ii)} est valable.
\end{itemize}
\end{enonce*}

C'est un cas particulier de \hyperref[VI.cor1.24]{1.24}, compte tenu que $C$ est stable par
sommes finies (\hyperref[VI.cor1.4]{1.4}), ce qui implique l'équivalence de (ii) et de (ii
bis), et compte tenu du

\begin{enonce*}{Lemme 1.25.1}\etiquette[1.25.1]{VI.lem1.25.1}
Soit $E$ un topos. Tout facteur direct $X'$ d'un objet cohérent $X$ de
$E$ est cohérent.
\end{enonce*}

En effet, $X'$ est quasi-compact comme quotient de $X$ (\hyperref[VI.prop1.3]{1.3}), et
quasi-séparé comme sous-objet de $X$ (\hyperref[VI.cor1.15]{1.15}).

\begin{enonce*}{Corollaire 1.26}\etiquette[1.26]{VI.cor1.26}
Soit $E$ un topos satisfaisant aux conditions équivalentes
\emph{\hyperref[VI.cor1.25]{1.25}.} Alors, il en est de même de tout topos induit.
\end{enonce*}

Cela résulte aussitôt du critère (ii), compte tenu de \hyperref[VI.subsubsec1.20.2]{1.20.2} et du
fait que l'existence d'une sous-catégorie génératrice de $E$ formée
d'objets cohérents implique manifestement la même propriété pour un
topos induit.

\begin{enonce*}[remark]{Remarque 1.27}\etiquette[1.27]{VI.rem1.27}
Parmi\pageoriginale les exemples importants de topos satisfaisant les
conditions de \hyperref[VI.cor1.25]{1.25}, signalons les topos noethériens
(\hyperref[VI.subsec2.14]{2.14}), et le topos zariskien ou étale
(\hyperrefext[VII][VII.sec1]{1}) d'un schéma $X$ (\cf IX, note p.42). Mais
même lorsque
le topos $E$ est \emph{cohérent}(\hyperref[VI.def2.3]{2.3}) (\ie la sous-catégorie pleine
$C$ des objets cohérents est génératrice et stable par limites
projectives finies), il n'est pas toujours vrai que $C$ soit
\emph{parfait} (\hyperref[VI.def2.9.1]{2.9.1}) \ie satisfasse aux conditions équivalentes de
\hyperref[VI.cor1.25]{1.25}; \cf \hyperref[VI.exer1.28]{1.28} ci-dessous.
\end{enonce*}

\begin{enonce*}[remark]{Exercice 1.28}\etiquette[1.28]{VI.exer1.28}
Soient $E$ un topos, $p_1$, $p_2 : X_1 \rightrightarrows X_0$ une double
flèche dans $E$. Définir une suite croissante de sous-objets $R_n$ $(n
\geq 0)$ de $X_0\times X_0$ par la condition $R_0 = \Sup(\text{Im}(p_1,
p_2), \text{Im}(p_2, p_1))$, et $R_n = \pr_{13} (\pr_{12}^{-1} (R_{n-1}),
\pr_{23}^{-1}(R_{n-1}))$ pour tout $n \geq 1$, où $\pr_{ij} (1 \leq i
< j \leq 3)$ désignent les projections qu'on devine du produit triple
de $X_0$ vers le produit double.
\begin{enumerate}
\renewcommand{\theenumi}{\alph{enumi}}
\renewcommand{\labelenumi}{\theenumi)}
\item Montrer qu'on obtient ainsi une suite croissante de sous-objets
  de $X_0\times X_0$, que $R = \varinjlim R_n$ est un graphe
  d'équivalence (I 10.4) dans $X_0$, et que le morphisme canonique
$$
X = \Coker (p_1, p_2) \to X_0 /R
$$
est un isomorphisme.

\item En conclure que si $X_0$ est quasi-compact et $X = \Coker(p_1,
  p_2)$ est quasi-séparé (donc cohérent), alors la suite des $R_n$ est
  stationnaire. Inversement,\pageoriginale si cette suite est
  stationnaire, et si on suppose $X_0$
cohérent et $X_1$ quasi-compact alors $X$ est cohérent. (Pour ce
dernier point, écrire $R_n \simeq \fprod{(R_{n-1}, \pr_2)}{(R_{n-1},
  \pr_1)}{X_0}$ et en conclure que $R_n$ est quasi-compact sur $X_0$
pour tout $n$, puis utiliser 1.19.)

\item Construire un exemple de deux applications d'ensembles $p_1$, $p_2
  : X_1 \rightrightarrows X_0$, telles que la suite des $(R_n)$ ne
  soit pas stationnaire.

\item Soit $C$ une petite catégorie où les limites inductives finies
  et les limites projectives finies sont représentables, et où tout
  morphisme se factorise en un épimorphisme effectif suivi d'un
  monomorphisme. Répéter pour une double flèche $(p_1, p_2)$ de $C$ la
  construction des $R_n$. Munissant $C$ de la topologie canonique,
  montrer que le foncteur canonique $\epsilon : C \to C^{\sim}$ est exact, et
  commute par suite à la formation des $R_n$.

\item Prendre pour $C$ une petite sous-catégorie pleine de $(\Ens)$,
  stable à isomorphisme près par limites projectives finies et limites
  inductives finies, et contenant les ensembles $X_0$, $X_1$ de c). En
  conclure que le topos $C^{\sim}$ (qui est un topos
  \emph{cohérent} (\hyperref[VI.def2.3]{2.3}), \ie engendré par la sous-catégorie de ses
  objets cohérents, laquelle est stable par limites projectives finies)
  ne satisfait pas aux conditions équivalentes de \hyperref[VI.cor1.25]{1.25}.

\item Soient $k$ un corps, $C$ le site fppf des schémas localement de
  présentation finie sur $k$ (\SGA 3 IV 6.3). Montrer qu'il existe un
  schéma affine réduit $X$ sur $k$ qui admet un automorphisme $f$ qui
  n'est pas d'ordre fini (prendre par exemple l'espace affine $E^2$
  muni de l'automorphisme $f(x) = x$, $f(y) = x + y$).\pageoriginale
  Montrer que si
   $p_1$, $p_2 : X_1 \rightrightarrows X_0$ est une double flèche dans
  $C$, telle que la double flèche correspondante dans $(\Ens) p_1
  (\bar{k}):p_2(\bar{k}): X_1(\bar{k}) \rightrightarrows X_0
  (\bar{k})$ donne une suite $(R_n)$ non stationnaire, alors il en
  est de même de la double flèche de $C^{\sim}$ définie par $p_1$,
  $p_2$, donc, si $X_0$ est de type fini sur $k$, alors le conoyau
  dans $C^{\sim}$ de $(p_1, p_2)$ n'est pas cohérent. En conclure que
  le conoyau dans $C^{\sim}$ du couple $(f, \id_X)$ n'est pas
  cohérent, donc que le topos $C^{\sim}$ ne satisfait pas aux
  conditions équivalentes de \hyperref[VI.cor1.25]{1.25}. En conclure plus généralement que
  si $S$ est un schéma non vide, $C$ le site fppf des schémas
  localement de présentation finie sur $S$, alors le topos $C^{\sim}$
  ne satisfait pas aux conditions de \hyperref[VI.cor1.25]{1.25}. (On fera attention que,
  contrairement à l'apparence, $C^{\sim}$ n'est donc noethérien (\hyperref[VI.def2.11]{2.11})
  que si $S$ est vide, comme il résulte de ce qui précède et de \hyperref[VI.subsec2.14]{2.14}.
\end{enumerate}
\end{enonce*}

\begin{enonce*}{Définition 1.30}\etiquette[1.30]{VI.def1.30}
Soit $E$ un topos. Un objet $X$ de $E$ est dit un \emph{objet
prénoethérien}\footnote{Par la notion plus forte d'objet
  \emph{noethérien}, \cf \hyperref[VI.def2.11]{2.11} ci-dessous.} du topos $E$ s'il satisfait
aux deux conditions équivalentes suivantes:
\begin{enumerate}
\renewcommand{\theenumi}{\roman{enumi}}
\renewcommand{\labelenumi}{(\theenumi)}

\item Tout sous-objet de $X$ est quasi-compact.

\item Toute suite croissante de sous-objets de $X$ est stationnaire.
\end{enumerate}
\end{enonce*}

\setcounter{subsection}{30}
\setcounter{subsubsection}{0}
\subsubsection{}\etiquette[1.30.1]{VI.subsubsec1.30.1}
L'équivalence des conditions (i) et (ii) est claire. On notera que ces
conditions ne dépendent encore que du topos induit $E_{/X}$, et même
seulement de l'ensemble ordonné des ouverts de $E_{/X}$, isomorphe à
l'ensemble ordonné des sous-objets de $X$. Elles sont stables par
passage à un sous-objet de $X$.


\subsection{}\etiquette[1.31]{VI.subsec1.31}
On prouve\pageoriginale comme pour \hyperref[VI.prop1.3]{1.3} et
\hyperref[VI.cor1.4]{1.4} les faits suivants : si $(X_i \to X)$ est une
famille couvrante finie, avec les $X_i$ prénoethériens, alors $X$
est prénoethérien; en particulier un quotient d'un objet
prénoethérien de $E$ est prénoethérien. Si $(X_i)_{i \in I}$ est une
famille d'objets de $E$, alors leur somme $X$ est un objet
prénoethérien de $E$ si et seulement si tous les $X_i$ sauf un
nombre fini sont isomorphes à $∅_{E}$, et tous les $X_i$
sont prénoethériens.

\subsection{}\etiquette[1.32]{VI.subsec1.32}
Évidemment un objet prénoethérien d'un topos $E$ est quasi-compact;
lorsque $E$ admet une famille génératrice formée d'objets
prénoethériens, la réciproque est vraie (comme il résulte aussitôt
de \hyperref[VI.subsec1.31]{1.31}): les objets prénoethériens de $E$ sont alors
ses objets quasi-compacts.

Supposons que $E$ admette une famille génératrice $(X_i)_{i\in I}$
formée d'objets quasi-compacts, alors les trois conditions suivantes
sont équivalentes:
\begin{enumerate}
\renewcommand{\theenumi}{\roman{enumi}}
\renewcommand{\labelenumi}{(\theenumi)}
\item Les $X_i$ sont prénoethériens.

\item Tout objet quasi-compact de $E$ est prénoethérien.

\item Tout monomorphisme dans $E$ est quasi-compact.
\end{enumerate}
\noindent
(On vient de voir l'équivalence de (i) et (ii), (ii) $\Rightarrow$
(iii) est trivial à partir des définitions, et (iii) $\Rightarrow$
(ii) sur la définition \hyperref[VI.def1.30]{1.30} (i).)

On notera\pageoriginale que (iii) a comme conséquence que \emph{tout
  morphisme de
  $E$ est quasi-séparé}; donc (\hyperref[VI.prop1.14]{1.14} (ii)) tout objet $X$ de $E$
  au-dessus d'un objet quasi-séparé $X$ de $E$ est lui-même
  quasi-séparé.

\begin{enonce*}[remark]{Exemple 1.33}\etiquette[1.33]{VI.exem1.33}
(\emph{Topos classifiant d'un groupe}). Soient $G$ un groupe $\in
𝒰$, $B_G$ son topos classifiant (IV 2.4) \ie la catégorie
des ensembles $\in 𝒰$ où $G$ opère à gauche. Les objets
connexes-non vides de $B_G$ (IV 4.3.5) sont les ensembles à
opérateurs qui sont non vides et sur lesquels $G$ opère
transitivement, et tout objet $E$ de $B_G$ est somme de ses
composantes connexes, qui sont les sous-objets non vides minimaux de
$E$ (savoir les orbites de $G$ dans $E$). Par suite $E$ est un objet
quasi-compact de $B_G$ si et seulement si l'ensemble $\pi_0(E)$ des
orbites de $G$ dans $E$ est fini, et alors $E$ est même un objet
prénoethérien de $B_G$. Ces objets (et même le seul objet $G_s$, qui
est connexe non-vide) forment une sous-catégorie génératrice de
$B_G$. Pour que l'objet final $e$ de $B_G$ soit quasi-séparé, \ie
pour que le produit de deux objets quasi-compacts de $B_G$ soit
quasi-compact, il faut et il suffit que le produit $G_s\times G_s$
soit quasi-compact, ce qui équivaut manifestement au fait que $G$
est fini. On en conclut, plus généralement, qu'un objet $E$ de $B_G$
est quasi-séparé si et seulement si les stabilisateurs de ses points
sont des sous-groupes \emph{finis} de $G$; \ie s'il est isomorphe à
une somme d'espaces homogènes $G/H$, avec $H$ \emph{fini}.) En
effet, en vertu de \hyperref[VI.cor1.15]{1.15} on peut supposer dans ce critère
$E$ connexe, mais si $x \in E$ a comme groupe de stabilité $G_x$, on
sait (IV 5.8)\pageoriginale que $B_G / E$ est équivalent au topos
$B_{G_x}$, donc son objet final est quasi-séparé si et seulement si
$G_x$ est fini. On conclut de ce critère que tout objet de $B_G$ qui
se trouve au-dessus d'un objet quasi-séparé est quasi-séparé, a
fortiori, le produit de deux objets quasi-séparés de $B_G$ est
quasi-séparé. On conclut aussi que les objets cohérents de $B_G$
sont les objets isomorphes à des sommes finies d'espaces homogènes
$G/H$, avec $H$ \emph{fini}. Comme $G_s$ est cohérent, on voit que
les objets cohérents de $B_G$ forment déjà une sous-catégorie
génératrice de $B_G$. De plus, cette sous-catégorie est stable par
produits fibrés (comme il résulte du fait que tout objet au-dessus
d'un objet quasi-séparé est quasi-séparé). (Dans la terminologie
\hyperref[VI.def2.11]{2.11} plus bas $B_G$ est un topos \emph{localement
noethérien}, mais il n'est quasi-séparé que si $G$ est fini, et
alors $B_G$ est même noethérien.)

On voit tout de suite qu'un morphisme $f : E'\to E$ est
quasi-compact si et seulement si l'application induite $\pi_0(f) :
\pi_0 (E')\to \pi_{0}(E)$ pour les ensembles d'orbites est propre
\ie à fibres finies; on retrouve qu'un monomorphisme est
quasi-compact (\hyperref[VI.subsec1.32]{1.32}), donc tout morphisme est
quasi-séparé. Par suite les morphismes cohérents dans $B_G$ sont
simplement les morphismes quasi-compacts, qu'on vient de
caractériser.

Soit $E$ un objet du topos $B_G$. On vérifie facilement que $E$ est
une limite inductive filtrante d'objets cohérents de $B_G$ si et
seulement si les stabilisateurs des points de $E$ sont des groupes
ind-finis (\ie limites\pageoriginale inductives filtrantes de leurs
sous-groupes finis). En particulier, si $G$ n'est pas ind-fini (par
exemple si $G =  \ZZ$), l'objet final du topos $B_G$ n'est pas
limite inductive d'objets cohérents; plus précisément, pour que le topos $B_G$
satisfasse aux conditions équivalentes de \hyperref[VI.cor1.25]{1.25}, il faut et il suffit
que son objet final soit limite inductive filtrante d'objets
cohérents, ou encore que le groupe $G$ soit ind-fini.
\end{enonce*}


\section{Conditions de finitude pour un topos}\etiquette[2]{VI.sec2}

\begin{enonce*}{Proposition 2.1}\etiquette[2.1]{VI.prop2.1}
Soit $E$ un $𝒰$-topos. Les conditions suivantes sont équivalentes
\begin{enumerate}
\renewcommand{\theenumi}{\roman{enumi}}
\renewcommand{\labelenumi}{(\theenumi)}
\item Il existe une sous-catégorie pleine génératrice $C$ de $E$,
  formée d'objets quasi-compacts, et stable par produits fibrés.

\item Il existe un $𝒰$-site $C$, tel que tout objet de $C$ soit
  quasi-compact, que dans $C$ les produits fibrés soient
  représentables, et que $E$ soit équivalent à la catégorie des
  faisceaux $C^{\sim}$.
\end{enumerate}

Supposons que ces conditions sont satisfaites. Alors on peut même
choisir dans \emph{(i) (\resp (ii))} la catégorie $C$
$𝒰$-petite. D'autre part, pour tout $X \in \ob C$, l'objet
$X$ \emph{(\resp $\epsilon (X)$)} du topos $E$ est \emph{cohérent}.
\end{enonce*}

L'implication (i) $\Rightarrow$ (ii) résulte de IV 1.2.1., et le
fait que dans (i) (\resp (ii)) on peut prendre $C$ $𝒰$-petite se voit
aussitôt par l'argument de IV 1.2.3 appliqué à une petite
sous-catégorie pleine génératrice au sens topologique (II 3.0.1) de
$C$. Il reste à prouver (ii) $\Rightarrow$ (i) et la dernière
assertion de \hyperref[VI.prop2.1]{2.1}, ce qui résulte aussitôt du

\begin{enonce*}{Corollaire 2.1.1}\etiquette[2.1.1]{VI.cor2.1.1}
Soit\pageoriginale $C$ un $𝒰$-site où les produits fibrés soient
représentables et dont tout objet soit quasi-compact, et soient $E : C^{\sim}$,
$\varepsilon : C \to E$ le foncteur canonique. Alors pour tout $X \in
\ob C$, $\varepsilon(X)$ est un objet cohérent de $E$. La
sous-catégorie pleine de $E$ formée des objets cohérents de $E$ qui se
trouvent au-dessus de quelque $\varepsilon(X)$ est stable par produits
fibrés (et est évidemment génératrice dans $E$).
\end{enonce*}

On sait déjà que $\epsilon(X)$ est quasi-compact (\hyperref[VI.prop1.2]{1.2}),
prouvons qu'il est quasi-séparé. Soient donc $F$, $G$ deux objets
quasi-compacts de $E$ et $F \to \varepsilon(X)$, $G \to \epsilon(X)$
deux morphismes, prouvons que le produit fibré
$\fprod{F}{G}{\epsilon(X)}$ est quasi-compact. Recouvrant $F$ et $G$
par un nombre fini d'objets de la forme $\epsilon(Y_i)$ \resp
$\epsilon(Z_j)$, et procédant comme dans \hyperref[VI.prop1.2]{1.2} pour nous
ramener au cas où les composés $\epsilon(Y_i) \to \epsilon(X)$,
$\epsilon(Z_j) \to \epsilon(X)$ proviennent de morphismes $Y_i \to
X$ \resp $Z_j \to Y$, on est ramené grâce à \hyperref[VI.prop1.3]{1.3} au cas
où les morphismes $F \to \epsilon(X)$ et $G \to \epsilon(X)$
sont les transformés par $\epsilon$ de morphismes $Y \to X$, $Z \to
X$. Mais comme les produits fibrés sont représentables dans $C$ et
que $\epsilon$ y commute, on est réduit à prouver que
$\epsilon(\fprod{Y}{Z}{X})$ est quasi-compact, ce qui a été déjà
noté plus haut pour tout objet de la forme $\epsilon(U)$.

Soient maintenant $G \to F$, $H \to F$ des morphismes d'objets
cohérents de $E$, tels qu'il existe un morphisme $F \to
\epsilon(X)$, pour quelque $X \in \ob C$, prouvons que
$\fprod{G}{H}{F}$ est cohérent. Comme $G$ et $H$ sont quasi-compacts
et $F$ quasi-séparé, il résulte déjà des définitions que le produit
fibré est quasi-compact.\pageoriginale Reste à voir qu'il est
quasi-séparé, ou ce qui revient au même (\hyperref[VI.cor1.15]{1.15}) que $\fprod{G}{H}{\epsilon(X)}$ est
quasi-séparé. Cela résulte du fait plus général:

\begin{enonce*}{Corollaire 2.1.2}\etiquette[2.1.2]{VI.cor2.1.2}
Sous les conditions de \hyperref[VI.cor2.1.1]{2.1.1}, pour tout objet quasi-séparé $G$ de $E$,
tout morphisme $G \to \epsilon(X)$ est quasi-séparé. Si $H \to
\epsilon(X)$ est un deuxième morphisme, avec $H$ quasi-séparé,
alors $\fprod{G}{H}{\epsilon(X)}$ est quasi-séparé.
\end{enonce*}

La première assertion implique la deuxième, car elle implique par
changement de base que $\fprod{G}{H}{\epsilon(X)} \to H$ est quasi-séparé,
donc, puisque $H$ est quasi-séparé, $\fprod{G}{H}{\epsilon(X)}$
l'est aussi (\hyperref[VI.prop1.14]{1.14} (ii)). Pour prouver que $G \to \epsilon(X)$ est quasi-séparé,
nous allons utiliser le

\begin{enonce*}{Lemme 2.1.3}\etiquette[2.1.3]{VI.lem2.1.3}
Soient $E$ un topos, $f : X \to Y$ un morphisme dans $E$. Supposons
que $E$ admette une famille génératrice d'objets quasi-compacts, et
qu'il existe une famille couvrante $(g_i : X_i \to X)_{i \in I}$ de
$X$ telle que les $X_i$ soient quasi-compacts, et que pour tout couple
d'indices $i$, $j \in I$, $\fprod{X_i}{X_j}{Y}$ soit
quasi-séparé. Alors, si $X$ est un objet quasi-séparé de $E$, $f$ est
un morphisme quasi-séparé.
\end{enonce*}

En effet, la famille $(\fprod{g_i}{g_j}{Y} : \fprod{X_i}{X_j}{Y} \to
\fprod{X}{X}{Y})_{(i, j) \in I \times I}$ est couvrante, donc pour
voir que le morphisme $\Delta_{X/Y} : X \to \fprod{X}{X}{Y}$ est
quasi-compact, il suffit de voir que pour tout $(i, j) \in I \times
I$, le morphisme qui s'en\pageoriginale déduit par changement de base
$\fprod{g_i}{g_j}{Y}$ l'est (\hyperref[VI.prop1.10]{1.10} (ii)). Or ce dernier n'est autre que
le morphisme canonique $\fprod{X_i}{X_j}{X} \to \fprod{X_i}{X_j}{Y}$;
comme par hypothèse $X$ est quasi-séparé et les $X_i$ sont
quasi-compacts, $\fprod{X_{i}}{X_{j}}{X}$ est quasi-compact, donc
comme $\fprod{X_{i}}{X_{j}}{Y}$ est quasi-séparé par hypothèse, le
morphisme $\fprod{X_{i}}{X_{j}}{X}\to \fprod{X_{i}}{X_{j}}{Y}$ est
bien quasi-compact,\cqfd

Nous pouvons maintenant terminer la démonstration de \hyperref[VI.cor2.1.2]{2.1.2}, en
prouvant que tout morphisme $f : G \to \epsilon(X)$, avec 
$G$ quasi-séparé, est quasi-séparé. Pour ceci, couvrons $G$ par des
objets $\epsilon(X_i)$, tels que les morphismes composés
$\epsilon(X_i) \to \epsilon(X)$ proviennent de morphismes $X_i
\to X$. En vertu de \hyperref[VI.lem2.1.3]{2.1.3}, comme les $\epsilon(X_i)$ sont
quasi-compacts, il suffit de vérifier que les produits fibrés
$\fprod{\epsilon(X_i)}{\epsilon(X_j)}{\epsilon(X)} =
\epsilon(\fprod{X_i}{X_j}{X})$ sont quasi-séparés. Or on a déjà vu
que tout objet de $E$ de la forme $\epsilon(U)$ est cohérent. Cela
prouve \hyperref[VI.cor2.1.2]{2.1.2} et achève la démonstration de \hyperref[VI.prop2.1]{2.1}.

\begin{enonce*}{Proposition 2.2}\etiquette[2.2]{VI.prop2.2}
Soit $E$ un topos contenant une sous-catégorie pleine génératrice $C$
formée d'objets cohérents. Les conditions suivantes sur $E$ sont
équivalentes:
\begin{itemize}
\item[(i)] Tout objet quasi-séparé de $E$ est quasi-séparé sur l'objet
  final.

\item[(i bis)] Pour tout morphisme $f : X \to Y$ dans $E$, $X$
  quasi-séparé implique $f$ quasi-séparé.

\item[(i ter)] Tout objet $X$ de $C$ est quasi-séparé sur l'objet final de $E$,
  \ie le morphisme diagonal $X \to X \times X$ est quasi-compact.

\item[(ii)] Le produit\pageoriginale de deux objets quasi-séparés de $E$ est
  quasi-séparé.

\item[(ii bis)] Le produit dans $E$ de deux objets de $C$ est quasi-séparé.
\end{itemize}

Ces conditions impliquent que la sous-catégorie pleine $E_{\coh}$ de
$E$ formée des objets cohérents de $E$ est stable par produits fibrés
(et à fortiori, que $E$ satisfait aux conditions équivalentes de \hyperref[VI.prop2.1]{2.1}).
\end{enonce*}

Évidemment (i bis) $\Rightarrow$ (i), et l'implication inverse
résulte de \hyperref[VI.prop1.8]{1.8} (iii). On a (i) $\Rightarrow$ (ii) par
un argument déjà fait : si $X$ et $Y$ sont quasi-séparés, comme par
l'hypothèse (i) $X$ est quasi-séparé sur $e$, $X \times Y$ est
quasi-séparé sur $Y$ par changement de base, donc $X$ est
quasi-séparé en vertu de \hyperref[VI.prop1.14]{1.14} (ii). Le même argument
montre que (i ter) $\Rightarrow$ (ii bis), et d'autre part (i)
$\Rightarrow$ (i ter) et (ii) $\Rightarrow$ (ii bis) sont triviaux,
de sorte qu'il reste à établir (ii bis) $\Rightarrow$ (i). Mais si
$X$ est un objet quasi-séparé de $E$, recouvrant $X$ par des objets
$X_i$ de $C$ (ce qui est possible, $C$ étant génératrice), comme les
$X_i$ sont quasi-compacts par hypothèse sur $C$, on peut appliquer
\hyperref[VI.lem2.1.3]{2.1.3} au morphisme $X \to e$, de sorte que pour prouver
que ce dernier est quasi-séparé, on est ramené précisément à voir
que les $X_i \times X_j$ sont quasi-séparés, ce qui est l'hypothèse
(ii bis).

\begin{enonce*}[remark]{Remarque 2.2.1}\etiquette[2.2.1]{VI.rem2.2.1}
La condition (i ter) équivaut aussi à la condition (i quater). Pour
toute double flèche $g_1$, $g_2 : Y \rightrightarrows X$  dans $C$, le
noyau dans $E$ est quasi-compact (ou encore (\hyperref[VI.cor1.15]{1.15}) cohérent).
\end{enonce*}

Pour le\pageoriginale voir, il suffit d'appliquer le critère \hyperref[VI.prop1.10]{1.10} (i)
au morphisme diagonal $X \to X \times X$ envisagé dans (i ter).

\begin{enonce*}{Définition 2.3}\etiquette[2.3]{VI.def2.3}
Un topos $E$ satisfaisant les conditions équivalentes de \emph{\hyperref[VI.prop2.2]{2.2} (\resp de
\hyperref[VI.prop2.1]{2.1})} est appelé un \emph{topos algébrique (\resp} un \emph{topos localement
algébrique}, ou \emph{localement cohérent)}. Un objet $X$ d'un topos $E$ est
appelé \emph{objet algébrique du topos} $E$ si le topos induit $E_{/X}$ est
algébrique. Un topos $E$ est dit \emph{quasi-séparé (\resp cohérent)} s'il
est algébrique et si son objet final est quasi-séparé (\resp
cohérent).
\end{enonce*}

\setcounter{subsection}{3}
\subsubsection{}\etiquette[2.3.1]{VI.subsubsec2.3.1}
C'est dans les topos algébriques que les notions de finitude
développées au ${\rm n^{\circ}}$~1 ont des propriétés pleinement
satisfaisantes, analogues aux propriétés des notions de même nom dans
la catégorie des schémas. Tous les topos localement algébriques (=
localement cohérents) rencontrés en pratique sont en fait algébriques,
donc l'intérêt pratique de cette notion semble pour l'instant assez
réduit. On peut cependant construire des topos localement algébriques
et non algébriques (\hyperref[VI.exer2.17]{2.17} f).

\subsection{}\etiquette[2.4]{VI.subsec2.4}
Voici quelques remarques générales concernant les notions introduites
dans \hyperref[VI.def2.3]{2.3}, qui convaincront le lecteur (du moins on
l'espère) que la terminologie introduite est cohérente.

\subsubsection{}\etiquette[2.4.1]{VI.subsubsec2.4.1}
Sous les conditions de \hyperref[VI.prop2.1]{2.1} (i) \resp (ii), pour tout $X \in C$, $X$
(\resp $\epsilon(X)$) est un objet (cohérent) algébrique de $E$
(\hyperref[VI.cor2.1.2]{2.1.2} et 1.19.2),
en d'autres\pageoriginale termes le topos induit
$E_{/X}$ (\resp $E_{/\epsilon(X)} \approxeq (C_{/X})^{\sim}$) est
algébrique, donc
\emph{cohérent}. Il revient au même de dire que \emph{lorsque
$C$ admet un objet
final, alors $E$ lui-même est algébrique, donc cohérent.}

\subsubsection{}\etiquette[2.4.2]{VI.subsubsec2.4.2}
Si $E$ est un topos algébrique, alors tout topos induit $E_{/X}$ est
également algébrique (\hyperref[VI.prop2.2]{2.2} (i bis)); pour qu'un topos $E$ soit
localement algébrique = localement cohérent, il faut et il suffit
qu'il existe une famille $(X_i)_{i \in I}$ d'objets de $X$, couvrant
l'objet final, telle que les $X_i$ soient algébriques (\resp
cohérents) \ie telle que les topos induits $E_{/X_{i}}$ soient algébriques
(\resp cohérents). -La suffisance résulte aussitôt des définitions, la
nécessité des observations \hyperref[VI.subsubsec2.4.1]{2.4.1}. Bien entendu, un topos algébrique
est localement algébrique, et si $E$ est un topos localement
algébrique, tout topos induit $E_{/X}$ est localement algébrique.

\subsubsection{}\etiquette[2.4.3]{VI.subsubsec2.4.3}
Si $X$ est un objet algébrique d'un topos $E$, alors tout objet $X'$
de $E$ qui se trouve au-dessus de $X$ est encore algébrique (\cf
\hyperref[VI.subsubsec2.4.2]{2.4.2}). En particulier, tout objet d'un topos algébrique est
algébrique, et inversement bien sûr.

\subsubsection{}\etiquette[2.4.4]{VI.subsubsec2.4.4}
La sous-catégorie pleine de $E$ formée des objets
\emph{quasi-séparés} qui sont \emph{algébriques} est stable par
produits fibrés et par sous-objets (\hyperref[VI.cor1.15]{1.15}), donc aussi
par noyaux. Si $E$ est algébrique, cette catégorie (qui n'est alors
autre que la catégorie des objets quasi-séparés de $E$
sans\pageoriginale plus) est également stable par produit de deux
objets. Si $E$ est quasi-séparé (et alors seulement) cette catégorie
est stable par limites projectives finies quelconques, ou encore
contient l'objet final de $E$.

La sous-catégorie pleine $C$ de $E$ formée des objets \emph{cohérents}
qui sont \emph{algébriques} est stable par produits fibrés. Si $E$ est
localement algébrique \ie localement cohérent, dire que $E$ est
algébrique revient à dire que $C$ est également stable par noyaux de
doubles flèches (\hyperref[VI.rem2.2.1]{2.2.1}); et dire que $E$ est quasi-séparé revient à
dire que de plus $C$ est stable par produit de deux objets (\cf
\hyperref[VI.cor1.17]{1.17}). Enfin dire que $E$ est cohérent revient à dire que $C$ est
stable par limites projectives finies quelconques, ou encore qu'il
contient l'objet final de $E$.

\subsubsection{}\etiquette[2.4.5]{VI.subsubsec2.4.5}
Soit $E$ un topos. Pour que $E$ soit localement cohérent (\resp
algébrique, \resp quasi-séparé, \resp cohérent) il faut et il suffit
qu'il admette une sous-catégorie pleine génératrice $C$ formée
d'objets quasi-compacts (qui seront automatiquement cohérents
(\hyperref[VI.prop2.1]{2.1})), et qui soit \emph{stable} par produits fibrés
(\resp par produits fibrés et par noyaux, \resp par produits fibrés
et par produits de deux objets, \resp par limites projectives finies
quelconques); ou encore que $E$ soit équivalent à un topos
$C^{\sim}$, où $C$ est un $𝒰$-site dont tout objet est
quasi-compact, et où les produits fibrés et les produits de deux
objets (\resp toutes les limites projectives finies) sont
représentables. (Pour la nécessité, on prend \pageoriginale pour $C$
la catégorie des objets de $E$ qui sont cohérents et
\emph{algébriques}, et on applique \hyperref[VI.subsec2.4]{2.4}; pour la
suffisance, on applique \hyperref[VI.prop2.2]{2.2} (i ter)). Dans cet énoncé,
on peut évidemment supposer encore $C$ $𝒰$-petit.

\subsubsection{}\etiquette[2.4.6]{VI.subsubsec2.4.6}
Soit $E$ un topos, Si $E$ est algébrique, un objet $X$ de $E$ est
quasi-séparé (\resp cohérent) si et seulement si le topos induit
$E_{/X}$ l'est. Lorsque $E$ n'est plus supposé algébrique, on peut
dire seulement que $E_{/X}$ est quasi-séparé (\resp cohérent) si et
seulement si $X$ est quasi-séparé (\resp cohérent) et
\emph{algébrique}. Ce n'est pas grave, vu qu'on n'aura sans doute pas
à utiliser ces notions en dehors du cas $E$ algébrique.

\subsubsection{}\etiquette[2.4.7]{VI.subsubsec2.4.7}
Pour que la notion d'objet cohérent (\resp quasi-séparé d'un
topos $E$ ait des propriétés satisfaisantes, il faut se borner aux
objets qui sont de plus algébriques, condition automatiquement
satisfaite si $E$ est algébrique. Comme nous n'aurons sans doute
jamais à travailler avec des $E$ localement algébriques qui ne soient
algébriques, la question de réviser la terminologie introduite dans
\hyperref[VI.def1.13]{1.13}, en la réservant éventuellement aux seuls objets algébriques, ne
se pose donc pas en termes bien aigus!

D'autre part, nous nous sommes abstenus de définir la notion de topos
quasi-compact. Dans l'esprit du présent numéro, il s'imposerait
d'appeler ainsi les topos \emph{algébriques} sont l'objet final est
quasi-compact.

Nous \pageoriginale avons hésité à introduire un tel usage, car si
$X$ est un espace topologique, il ne serait pas vrai que $X$ est
quasi-compact si et seulement si le topos associé $\Top(X)$ l'est.
On notera à ce propos que $\Top(X)$ est localement cohérent si et
seulement si $X$ admet une base d'ouverts formée d'ouverts
quasi-compacts $U$ tels que pour $U_j$, $U_k \subset U_i$, $U_j \cap
U_k$ soit encore quasi-compact (et alors $\Top(X)$ est même
algébrique). Cette condition est vérifiée pour les espaces
topologiques sous-jacents aux schémas, mais rarement pour les
espaces rencontrés par les analystes, fussent-ils (les espaces)
compacts. Explicitons encore, pour la commodité des références:

\begin{enonce*}{Proposition 2.5}\etiquette[2.5]{VI.prop2.5}
Soit $E$ un topos algébrique, $f : X \to Y$ un morphisme dans $E$. Si
$Y$ est un objet quasi-séparé (\resp cohérent) de $E$, alors, pour que
$X$ soit un objet quasi-séparé (\resp cohérent) de $E$, il faut et il
suffit que le morphisme $f$ soit quasi-séparé (\resp cohérent).
\end{enonce*}

La suffisance a déjà été vue dans \hyperref[VI.prop1.14]{1.14}(ii), et est vraie sans
hypothèse sur $E$. La nécessité dans le cas non respé est vraie sans
supposer même $Y$ quasi-séparé, et n'est  autre que la définition \hyperref[VI.def2.3]{2.3}
via \hyperref[VI.prop2.2]{2.2} i bis). Il reste à prouver que si $Y$ est cohérent et $X$
cohérent, alors $f$ est cohérent. Comme on sait déjà que $f$ est
quasi-séparé, il suffit de voir que $f$ est quasi-compact, ce qui
résulte du fait que $X$ est quasi-compact et $Y$ quasi-séparé (\hyperref[VI.def1.13]{1.13}).

\begin{enonce*}{Corollaire 2.6}\etiquette[2.6]{VI.cor2.6}
Soient\pageoriginale $E$ un topos algébrique, $f: X \to Y$ un
morphisme dans $E$, $(Y_i \to Y)_{i \in I}$ une famille couvrante, avec les $Y_i$
quasi-séparés (\resp  cohérents). Pour que $f$ soit quasi-séparé
(\resp quasi-compact, \resp cohérent) il faut et il suffit que pour
tout $i \in I$, $X_i = \fprod{X}{Y_i}{Y}$ soit un objet quasi-séparé
(\resp quasi-compact, \resp cohérent) de $E$.
\end{enonce*}

Il suffit de conjuguer \hyperref[VI.prop1.10]{1.10} (ii) et \hyperref[VI.prop2.5]{2.5}
dans le cas « quasi séparé » ou « cohérent »; le cas « quasi-compact » a déjà été vu dans \hyperref[VI.cor1.16]{1.16}. On notera que
la nécessité de la condition énoncée dans \hyperref[VI.cor2.6]{2.6} est
également valable sans hypothèse sur le topos $E$.

On trouve comme cas particulier de \hyperref[VI.cor2.6]{2.6}:

\begin{enonce*}{Corollaire 2.7}\etiquette[2.7]{VI.cor2.7}
Soient $E$ un topos cohérent, $X$ un objet de $E$. Pour que $X$ soit
un objet quasi-compact (\resp quasi-séparé, \resp cohérent) de $E$, il
faut et il suffit qu'il soit quasi-compact (\resp quasi-séparé, \resp
cohérent) au-dessus de l'objet final (\hyperref[VI.subsubsec1.9.2]{1.9.2}).
\end{enonce*}

En particulier, $X$ est constructible (\hyperref[VI.subsubsec1.9.2]{1.9.2}) si et seulement s'il
est cohérent.

\begin{enonce*}{Corollaire 2.8}\etiquette[2.8]{VI.cor2.8}
Soient $E$ un topos localement cohérent, $f:X \to Y$ un morphisme dans
$E$. Pour que $f$ soit quasi-séparé, il faut et il suffit que pour
tout objet quasi-séparé $Y'$ de $E_{/Y}$, $f^{\ast} (Y') =
\fprod{X}{Y'}{Y}$ soit un objet quasi-séparé de $E_{/X}$ (ou de $E$,
cela revient au même (\hyperref[VI.subsubsec1.20.2]{1.20.2})).
\end{enonce*}

Le \pageoriginale « il faut » est valable sans condition sur $E$,
car si $Y'$ est quasi-séparé, comme $f':X'= \fprod{X}{Y'}{Y} \to Y'$ est quasi-séparé par
changement de base, $X'$ est quasi-séparé (\hyperref[VI.prop1.14]{1.14} (ii)). Pour la
réciproque, notons qu'il existe une famille couvrante $Y_i \to Y$ de
$Y$ par des objets quasi-séparés et algébriques. En vertu de \hyperref[VI.prop1.10]{1.10}
(ii), pour vérifier que $f$ est quasi-séparé, il suffit de vérifier
que les $f_i : X_i = \fprod{X}{Y_i}{Y} \to Y_i$ le sont. Or
l'hypothèse sur $f^{\ast}$ implique que les $X_i$ sont quasi-séparés,
donc les $f_i$ le sont puisque $E_{/Y_i}$ est un topos algébrique.

\begin{enonce*}[remark]{Remarque 2.8.1}\etiquette[2.8.1]{VI.rem2.8.1}
Les énoncés \hyperref[VI.prop2.5]{2.5} et \hyperref[VI.cor2.6]{2.6} restent valables si on y remplace l'hypothèse
« $E$ algébrique » par l'hypothèse plus générale « $Y$ est
algébrique », comme on voit trivialement en appliquant l'énoncé
primitif au topos induit $E_{/Y}$, qui est algébrique. De même, dans
\hyperref[VI.cor2.8]{2.8} il suffit de supposer que $E_{/Y}$ (au lieu de $E$) soit
localement cohérent.
\end{enonce*}

\begin{enonce*}{Proposition 2.9}\etiquette[2.9]{VI.prop2.9}
Soient $E$ un topos, $C$ une sous-catégorie de $E$. Les conditions
suivantes sont équivalentes:
\begin{enumerate}
\renewcommand{\theenumi}{\roman{enumi}}
\renewcommand{\labelenumi}{(\theenumi)}
\item $E$ est localement cohérent (\resp cohérent), satisfait aux
  conditions équivalentes de \hyperref[VI.cor1.25]{1.25}, et $C$ est la sous-catégorie
  pleine $E_{\coh}$ de $E$ formée des objets cohérents.

\item $C$ est une sous-catégorie strictement pleine génératrice de
  $E$, stable par limites inductives finies et par produits fibrés
  (\resp et par limites projectives finies), et est formée d'objets
  quasi-compacts.
\end{enumerate}
\end{enonce*}

L'implication\pageoriginale (i) $\Rightarrow$ (ii) résulte trivialement des
définitions et de la forme \hyperref[VI.cor1.25]{1.25} (ii) des conditions envisagées dans
\hyperref[VI.cor1.25]{1.25}; prouvons l'implication inverse. Le fait que $C$ soit une
sous-catégorie pleine génératrice, formée d'objet quasi-compacts, et
stable par produits fibrés (\resp par limites projectives finies)
implique, par définition, que $E$ est localement cohérent (\resp
cohérent). Le fait que de plus $C$ soit strictement pleine et stable
par limites inductives finies implique alors que $C = E_{\coh}$ (en
vertu de \hyperref[VI.cor1.24]{1.24} (ii) $\Rightarrow$ (iii)), et que la condition \hyperref[VI.cor1.25]{1.25}
(ii) est satisfaite, \cqfd

\begin{enonce*}{Définition 2.9.1}\etiquette[2.9.1]{VI.def2.9.1}
Un $𝒰$-topos $E$ est dit parfait s'il est cohérent (\hyperref[VI.def2.3]{2.3}) et s'il
satisfait aux conditions équivalentes de \hyperref[VI.cor1.25]{1.25}. \ie si la
sous-catégorie $E_{\coh}$ de $E$ formée des objets cohérents de $E$ est
stable par limites inductives finies. On dit que $E$ est localement
parfait s'il existe une famille $(X_i)_{i \in I}$ d'objets de $E$
couvrant l'objet final, telle que pour tout $i \in I$, le topos induit
$E_{/X_i}$ soit parfait.
\end{enonce*}

\setcounter{subsection}{9}
\setcounter{subsubsection}{1}
\subsubsection{}\etiquette[2.9.2]{VI.subsubsec2.9.2}

Il résulte aussitôt de cette définition qu'un topos parfait (\resp
localement parfait) est cohérent (\resp localement cohérent). Comme
exemples de topos parfaits (\resp localement parfaits), signalons
les topos noethériens (\resp localement noethériens) introduits dans
\hyperref[VI.def2.11]{2.11} ci-dessous (\cf \hyperref[VI.subsec2.14]{2.14}), le topos
zariskien et le topos étale d'un schéma cohérent (\resp d'un schéma
quelconque) (IX, note page 42). Comme contre-exemple,\pageoriginale
signalons le topos fppf d'un schéma $S$, qui est localement cohérent
(et même cohérent si $S$ est un schéma cohérent, \ie quasi-compact
et quasi-séparé) (\hyperrefext[VII][VII.lem5.6]{5.6}), mais qui n'est pas localement
parfait si $S \neq ∅$ (\hyperref[VI.exer1.28]{1.28} f)).

Dans un topos localement parfait, la notion d'objet constructible
est particulièrement stable (\hyperref[VI.prop2.9.3]{2.9.3} ci-dessous), ce qui
est une raison pourquoi la notion semble intéressante. Une autre
raison est dans le fait que comme celle de topos cohérent ou
localement cohérent, la notion de topos parfait ou localement
parfait est stable par rapport à la formation du sous-topos fermé
complémentaire d'un ouvert (\hyperref[VI.prop4.6]{4.6}), et qu'elle se
comporte de façon particulièrement simple pour l'opération de
recollement de topos (\hyperref[VI.cor4.11]{4.11}, \hyperref[VI.cor4.14]{4.14}).

\begin{enonce*}{Proposition 2.9.3}\etiquette[2.9.3]{VI.prop2.9.3}
Soit $E$ un topos localement parfait (\hyperref[VI.def2.9.1]{2.9.1}). Alors la
sous-catégorie
pleine $E_{\cons}$ de $E$ formée des objets constructibles de $E$
(\hyperref[VI.subsubsec1.9.3]{1.9.3}) est stable par limites inductives finies (et
aussi par limites
projectives finies bien sûr, en vertu de (\hyperref[VI.subsubsec1.9.3]{1.9.3})).
\end{enonce*}

Comme la propriété pour un objet de $E$ d'être constructible est
locale sur $E$ (\hyperref[VI.prop1.10]{1.10} (ii)), on est ramené au cas
où $E$ est un topos
parfait. Mais alors $E_{\cons} = E_{\coh}$ (\hyperref[VI.cor2.7]{2.7}), et la conclusion
résulte de la définition \hyperref[VI.def2.9.1]{2.9.1}.

\begin{enonce*}{Corollaire 2.9.4}\etiquette[2.9.4]{VI.cor2.9.4}
Soit\pageoriginale $E$ un topos localement parfait. Alors $E$ est parfait si et
seulement si $E$ est cohérent. Pour tout objet $X$ de $E$, le topos
induit $E_{/X}$ est localement parfait.
\end{enonce*}

\begin{enonce*}[remark]{Problème 2.9.5}\etiquette[2.9.5]{VI.prob2.9.5}
Il est concevable que les topos parfaits soient assez particuliers
pour se prêter à une \emph{théorie de structure} assez explicite
(suivant un modèle proposé par M.\textsc{Artin} à l'époque du
séminaire oral). Appelons \emph{topos fini} un topos équivalent à un
topos de la forme $\widehat{C}$, où $C$  est une catégorie finie
(\cf exercice \hyperref[VI.exer3.11]{3.11}), et \emph{topos profini} un topos
qui est une limite projective filtrante de topos finis (au sens de
6. plus bas, qui s'applique grâce à \hyperref[VI.exer3.12]{3.12} c)). Comme un
topos fini est noethérien (\hyperref[VI.exer2.17]{2.17} g)) donc parfait, et
qu'une limite projective filtrante de topos parfaits est parfait
(6.), on voit qu'un topos profini est parfait. La question qui se
pose serait de savoir si réciproquement tout topos parfait est
profini. Cela équivaut à la question si toute sous-catégorie finie
de $E_{\coh}$ est contenue dans une sous-catégorie pleine $F_{0}$ de
$E_{\coh}$ qui est équivalente à une catégorie $F_{\coh}$, où
$F$ est un topos fini (\cf \hyperref[VI.exer3.11]{3.11}). (Si la réponse était
négative, il y aurait lieu de trouver des conditions intrinsèques
supplémentaires maniables pour un topos parfait qui assurent qu'il
est profini.) Il resterait enfin à étudier la structure des topos
profinis en termes d'une notion convenable de « catégorie profinie
karoubienne », inspirée de IV 7.6 h). Cette\pageoriginale étude
n'a pas été faite encore même dans le cas particulier
où $E$ est
le topos zariskien, ou étale, d'un brave schéma cohérent $X$, et
devrait donner alors des invariants profinis plus fins que l'espace compact
$X_{\cons}$ (\EGA IV 1.) associé à $X$.

Dans le même ordre d'idées, signalons la question suivante: Soit
$E$ un topos parfait, $E'$ le topos induit sur un ouvert de $E$, et
$E''$ le sous-topos fermé correspondant (IV 9.9). Si $E'$ et $E''$
sont des topos finis, en est-il de même de $E$?
\end{enonce*}

\begin{enonce*}{Proposition 2.10}\etiquette[2.10]{VI.prop2.10}
Soit $E$ un topos. Les conditions suivantes sont équivalentes:
\begin{itemize}
\item[(i)] $E$ admet une sous-catégorie pleine génératrice
  stable par
  produits fibrés, formée d'objets prénoethériens de $E$.

\item[(ii)] $E$ admet une sous-catégorie pleine
  génératrice formée d'objets
  prénoethériens quasi-séparés (\ie prénoethériens cohérents).

\item[(iii)] $E$ est localement cohérent (\hyperref[VI.def2.3]{2.3}) et tout objet quasi-compact de
  $E$ est prénoethérien.

\item[(iii bis)] $E$ est localement cohérent et admet une famille
  génératrice formée d'objets prénoethériens de $E$.
\end{itemize}
\end{enonce*}

L'équivalence de (iii) et (iii bis) résulte de \hyperref[VI.subsec1.32]{1.32},
et il est trivial que (i) $\Rightarrow$ (iii bis) et (iii)
$\Rightarrow$ (i), donc (i) équivaut à (iii) et (iii bis). Enfin (i)
$\Rightarrow$ (ii) résulte de la dernière assertion de
\hyperref[VI.prop2.1]{2.1},\pageoriginale et il reste à prouver (ii)
$\Rightarrow$ (i). Cette implication résulte du fait que la
sous-catégorie pleine de $E$ formée des objets noethériens
quasi-séparés, si elle est génératrice, est stable par produits
fibrés, car un tel produit fibré est quasi-compact en vertu des
définitions, donc noethérien par la première assertion de
\hyperref[VI.subsec1.32]{1.32}; et il est quasi-séparé par la dernière
assertion de \hyperref[VI.subsec1.32]{1.32}.

\begin{enonce*}{Définition 2.11}\etiquette[2.11]{VI.def2.11}
Un topos $E$ est dit topos localement noethérien s'il satisfait aux
conditions équivalentes de \hyperref[VI.prop2.10]{2.10}, noethérien si de
plus son objet final
est cohérent (ou ce qui revient au même (\hyperref[VI.subsec1.32]{1.32})
prénoethérien et
quasi-séparé). Un objet $X$ d'un topos $E$ est dit objet
noethérien de
$E$ si le topos induit $E_{/X}$ est un topos noethérien.
\end{enonce*}

\setcounter{subsection}{11}
\subsection{}\etiquette[2.12]{VI.subsec2.12}
Si $E$ est un topos localement noethérien, alors il en est de même de
tout topos induit $E_{/X}$; inversement, si on peut trouver une
famille $(X_i)$ couvrant l'objet final telle que les $E_{/X_i}$ soient
localement noethériens, il en est de même de $E$. Si $C$ est une
sous-catégorie pleine génératrice de $E$ formée d'objets
prénoethériens et stable par produits fibrés (\hyperref[VI.def2.11]{2.11} (i)), alors pour
tout $X \in C$, $E_{/X}$ est un topos noethérien, \ie les objets de
$C$ sont même \emph{noethériens}, (car si $E$ est localement
noethérien et $X \in \Ob E$, alors $E_{/X}$ est noethérien si et
seulement si $X$ est cohérent). Il s'ensuit qu'un topos $E$ est
localement noethérien si et seulement si on peut recouvrir l'objet
final de $E$ par des objets noethériens $X_i$.

\subsection{}\etiquette[2.13]{VI.subsec2.13}
Dans un\pageoriginale topos localement noethérien $E$, tout monomorphisme est
quasi-compact, et tout morphisme est quasi-séparé (\hyperref[VI.subsec1.32]{1.32}). A fortiori,
$E$ est un topos \emph{algébrique} (\hyperref[VI.def2.3]{2.3}) et non seulement localement
cohérent \ie localement algébrique. Donc un topos $E$ est noethérien
si et seulement si il est localement noethérien et \emph{cohérent}
(\hyperref[VI.def2.3]{2.3}). En d'autres termes, si $E$ est un topos localement noethérien,
alors les objets noethériens de $E$ ne sont autres que les objets
\emph{cohérents} de $E$.

\subsection{}\etiquette[2.14]{VI.subsec2.14}
Soit $E$ un topos localement noethérien, Si $E$ est quasi-séparé,
\ie si son objet final est quasi-séparé, alors \emph{tout} objet de
$E$ est quasi-séparé (\hyperref[VI.subsec1.32]{1.32}), donc les objets
prénoethériens de $E$ sont identiques aux objets noethériens (\ie
cohérents) de $E$. Comme un objet quotient d'un objet prénoethérien
est prénoethérien, il s'ensuit que $E$ satisfait aux conditions
équivalentes de \hyperref[VI.cor1.25]{1.25} (puisqu'il satisfait la condition
\hyperref[VI.cor1.25]{1.25} (ii bis)). En particulier, si $C$ est la
sous-catégorie pleine de $E$ formée des objets noethériens
(cohérents) de $E$, alors le foncteur naturel
$$
\Ind(C) \to E
$$
est une équivalence de catégories. On conclut de ceci qu'\emph{un topos
noethérien est parfait} (\hyperref[VI.def2.9.1]{2.9.1}), donc qu'un topos localement
noethérien est localement parfait.

Si on suppose seulement $E$ localement noethérien, il sera encore
vrai que tout objet de $E$ est limite inductive de ses sous-objets
prénoethériens (car dans un topos admettant une sous-catégorie
génératrice formée\pageoriginale d'objets quasi-compacts, tout objet
est limite inductive filtrante de ses sous-objets quasi-compacts).
Mais comme un objet prénoethérien de $E$ n'est plus nécessairement
cohérent, cet énoncé n'a alors qu'une utilité très limitée, et on
sait (\hyperref[VI.exem1.33]{1.33}) que les conditions de \hyperref[VI.cor1.25]{1.25} ne
sont plus nécessairement vérifiées.

\begin{enonce*}[remark]{Exemple 2.15}\etiquette[2.15]{VI.exem2.15}
(\emph{Espaces topologiques}.) Soit $X$ un espace topologique. Alors
les conditions suivantes sont équivalentes : (i) l'espace
topologique $X$ est noethérien (\ie tout ouvert de $X$ est
quasi-compact), \ie toute suite croissante d'ouverts de $X$ est
stationnaire); (ii) le topos $\Top(X)$ est noethérien; (iii) l'objet
final $X$ de $\Top(X)$ est noethérien; (iv) l'objet final $X$ de
$\Top(X)$ est prénoethérien.

En effet, on a trivialement (ii) $\Leftrightarrow$ (iii) $\Rightarrow$
(iv) $\Leftrightarrow$ (i), d'autre part (i) implique que la
sous-catégorie génératrice Ouv(X) de $\Top(X)$, qui est stable par
limites projectives finies, est formée d'objets prénoethériens, d'où (ii).

On voit de même que l'espace topologique $X$ est localement
noethérien (\ie est réunion d'ouverts noethériens) si et seulement
si le topos $\Top(X)$ est localement noethérien.
\end{enonce*}

\setcounter{subsection}{15}
\setcounter{subsubsection}{0}
\subsubsection{}\etiquette[2.15.1]{VI.subsubsec2.15.1}
Ainsi, nos définitions \hyperref[VI.def1.30]{1.30}, \hyperref[VI.def2.11]{2.11} sont compatibles avec la
terminologie reçue pour les espaces topologiques. D'autre part, nous
avons tenu dans le cas général à donner au terme « objet
noethérien » un sens\pageoriginale plus fort que celui de la notion
plus naïve de \hyperref[VI.def1.30]{1.30}, pour que l'ensemble des propriétés qui s'attachent à cette
notion (plutôt que la seule structure grammaticale de la définition)
soit bien en accord avec l'intuition qui s'attache aux espaces
topologiques et aux schémas \emph{noethériens}.

Le lecteur trouvera d'autres exemples dans les exercices suivants.

\begin{enonce*}[remark]{Exercice 2.16}\etiquette[2.16]{VI.exer2.16}
(\emph{Espaces à opérateurs et topos classifiants}).

Soit $X$ un espace topologique sur lequel opère un groupe discret
$G$, d'où (IV 2.3) un topos $E = \Top(X, G)$. On interprète les
objets $X'$ de ce topos comme des espaces étalés sur $X$ à groupe
d'opérateurs $G$, et on note que le topos induit $E_{/X}$, est
canoniquement équivalent à $\Top(X', G)$.
\begin{enumerate}
\renewcommand{\theenumi}{\alph{enumi}}
\renewcommand{\labelenumi}{\theenumi)}
\item Pour que l'objet final du topos $E$ soit quasi-compact, il faut
  et il suffit que l'espace topologique quotient $X/G$ soit
  quasi-compact (Utiliser IV 8.4.1).

\item Soit $P$ l'objet de $E$ définit par l'espace étalé trivial $X
  \times G$ avec opération diagonale de $G$. Montrer que le topos
  induit $E_{/P}$ est équivalent au topos $\Top(X)$ (comparer IV
  5.8.3). En conclure que $E$ est localement cohérent (\resp
  localement noethérien) si et seulement si $\Top(X)$ est localement
  cohérent (\cf \hyperref[VI.subsubsec2.4.7]{2.4.7}) (\resp localement noethérien (\cf
  \hyperref[VI.exem2.15]{2.15})). Montrer que si $E$ est localement cohérent, \ie localement
  algébrique, il est même algébrique, en utilisant la famille
  génératrice formée\pageoriginale des objets $T_U$, où $U$ est un
  ouvert cohérent   de $X$, $T_U = G \times U$ avec opération $g \cdot (u, g') =
  (u, gg')$ de $G$, et le morphisme structural $p_U : T_U \to X$ défini
  par $P_U(u, g) = g\cdot u$.

\item Supposons $E$ algébrique, \ie $\Top(X)$ algébrique
  (\hyperref[VI.subsubsec2.4.7]{2.4.7}). Montrer que le morphisme structural $P \to e$ est
  quasi-séparé. Montrer que $E$ est quasi-séparé si et seulement si
  $\Top(X)$ est quasi-séparé (ou encore l'espace topologique $X$ est
  quasi-séparé, \ie l'intersection de deux ouverts quasi-compacts de
  $X$ est un ouvert quasi-compact), \emph{et} $G$ est fini ou $X$
  vide. (Comparer \hyperref[VI.exem1.33]{1.33}.)

\item Montrer que le morphisme structural $P \to e$ est quasi-compact
  si et seulement si $G$ est fini. Montrer que $E$ est cohérent (\resp
  noethérien) si et seulement si $\Top(X)$ l'est, \emph{et} $G$ est
  fini ou $X$ vide.

\item Soient $E$ un topos, $G$ un Groupe de $E$, d'où un topos
  classifiant $B_G$ (IV 2.4). Faire l'étude des conditions de finitude
  dans $B_G$, en s'inspirant de ce qui précède. Même question
  lorsqu'on part d'un pro-groupe strict $𝒢 = (G_i)_{i \in I}$ de
  $E$. En particulier, on verra que si $E$ est cohérent (\resp
  noethérien) et les $G_i$ sont cohérents, alors $B_{𝒢}$ est cohérent
  (\resp noethérien).
\end{enumerate}
\end{enonce*}

\begin{enonce*}[remark]{Exercice 2.17}\etiquette[2.17]{VI.exer2.17}
(\emph{Topos de la forme $\widehat{C}$}.)
Soient $C$ une catégorie équivalente à une catégorie $\in 𝒰$, et $E =
\widehat{C}$ la catégorie des préfaisceaux sur $C$. Par le foncteur
canonique $\varepsilon : C \to \widehat{C}$, on identifie $C$ à une
sous-catégorie pleine de $E$.
\begin{enumerate}
\renewcommand{\theenumi}{\alph {enumi}}
\renewcommand{\labelenumi}{(\theenumi)}
\item $C$ est\pageoriginale une sous-catégorie génératrice de $E$
  formée d'objets  quasi-compacts. Pour que ceux-ci soient prénoethériens (\cf \hyperref[VI.subsec1.32]{1.32}),
  il faut et il suffit que pour tout objet $X$ de $C$, tout crible de
  $X$ soit engendré par une famille finie de morphismes $X_i \to X$ de
  $C$. Pour que l'objet final de $E$ soit quasi-compact (\resp
  prénoethérien), il faut et il suffit qu'il existe une sous-catégorie
  finale de $C$ dont l'ensemble d'objets soit fini (\resp que tout
  crible de $C$ soit engendré par une sous-catégorie de $C$ dont
  l'ensemble d'objets soit fini).

\item Pour que $E$ admette une sous-catégorie génératrice formée
  d'objets cohérents (ou encore, quasi-séparés), il faut et il suffit
  que les objets de $C$ soient cohérents dans $E$, ou encore que pour
  deux flèches $Y \xrightarrow{f} X$,  $Z \xrightarrow{g} X$ dans $C$,
  l'objet $\fprod{Y}{Z}{X}$ de $E$ soit quasi-compact \ie on peut
  trouver une famille finie de carrés commutatifs dans $C$
\[
\xymatrix@=1cm{
 & T_i \ar[ld] \ar[rd]& \\
Y\ar[rd]_{f}&    & Z\ar[ld]^{g}\\
 & X  &,
}
\]
telle que  tout autre carré commutatif dans $C$
\[
\xymatrix@=1cm{
 & T \ar[ld] \ar[rd]& \\
Y\ar[rd]_{f}&    & Z\ar[ld]^{g}\\
 & X  &,
}
\]
provienne\pageoriginale d'un morphisme $T \to T_i$. (Noter que si
$(F_i)_{i \in I}$ est une famille génératrice de $E$, tout $X \in \ob C$ est isomorphe à
un facteur direct d'un des $F_i$.)

\item Pour que $E$ soit localement cohérent, \ie soit engendré par une
  famille d'objets cohérents et \emph{algébriques}, il faut et il
  suffit que les objets $X$ de $C$ soient des objets cohérents et
  algébriques de $E$, \ie que la condition de b) soit vérifiée, ainsi
  que la condition suivante: pour toute double flèche $f$, $g : Y
  \rightrightarrows Z$ de $C$ \emph{au-dessus d'un objet $X$ de} $C$,
  le noyau de cette double flèche dans $E$ est un objet quasi-compact
  de $E$, \ie il existe une famille finie de diagramme commutatifs
  dans $C$
\[
\xymatrix@C=1.5cm{
T_i \ar[r]& Y \ar@<2pt>[r]^{f,g}\ar@<-2pt>[r]&Z
}
\]
telle que tout autre diagramme commutatif dans $C$
\[
\xymatrix@C=1.5cm{
T \ar[r]& Y \ar@<2pt>[r]^{f,g}\ar@<-2pt>[r]&Z
}
\]
provienne d'un morphisme $T \to T_i$. Pour que $E$ soit algébrique,
il faut et il suffit qu'il satisfasse à la condition de b) et à la
condition précédente, mais où on prend une double flèche
quelconque $(f, g)$ de $C$ (pas nécessairement au-dessus d'un objet
$X$ de $C$). Pour que $E$ soit quasi-séparé, il faut et il suffit
qu'il satisfasse les conditions de b), et que de plus pour deux
objets $X$, $Y\in\ob C$, le produit $X\times Y$ dans $E$ soit
quasi-compact. Pour que $E$ soit cohérent, il faut et il suffit
qu'il satisfasse aux deux conditions précédentes, et qu'il existe une
sous-catégorie finale de $C$ dont l'ensemble sous-jacent soit fini.

\item Donner\pageoriginale un exemple où les objets de la
  sous-catégorie génératrice $C$ de $E$ sont prénoethériens, mais où $E$ n'admet pas de famille
  génératrice formée d'objets cohérents (\ie (b)) où les objets de $C$
  ne sont pas tous cohérents). Prendre pour ceci pour $C$ la catégorie
  ayant des objets $X$, $T_i$ $(i = 0, 1, \ldots)$ et $e$ (l'objet
  final), les seuls morphismes entre ces objets en plus des morphismes
  structuraux dans $e$ et des identités, étant des morphismes $u_i :
  T_i \to X$, $v_i : X \to T_i$ soumis aux conditions $u_i v_i =
  \id_X$, et enfin $p_i = v_i u_i$ (satisfaisant nécessairement $p_i^2
  = \id_{T_i}$). On vérifie que les seuls cribles de $X$ ou d'un $T_i$
  sont les deux cribles triviaux, et que $e$ admet exactement un
  crible non trivial, donc en vertu de a), les objets de $C$ sont des
  objets prénoethériens de $E$. Cependant, le produit $X \times X$
  dans $E$ n'est pas quasi-compact, car il ne satisfait pas au critère
  de b).

\item Donner un exemple où la sous-catégorie génératrice $C$ de $E$
  est formée d'objets cohérents de $E$, mais où $E$ n'est pas
  localement cohérent. Prendre pour ceci pour $C$ la catégorie dont
  l'ensemble des objets est formée d'objets distincts $e$ (l'objet
  final), $Y$, $Z$ et $T_i$ $(i=0, 1, \ldots)$, avec comme seuls
  morphismes entre ces objets, en plus des morphismes dans l'objet
  final $e$ et des morphismes identiques, des morphismes $f$, $g : Y
  \rightrightarrows Y$, soumis aux conditions $fu_i = gu_i$. On
  vérifiera que $C$ satisfait à la condition de b) (avec un peu de
  patience; on trouve que tous les produits fibrés d'objets de $C$
  sont isomorphes à $∅_E$ ou sont dans $C$,\pageoriginale à
  l'exception de $\fprod{Z}{Z}{e}$ et $\fprod{Y}{Z}{e}$ qui sont recouverts par deux
  éléments de $C$), mais évidemment $\Ker(f, g)$ ne satisfait pas à la
  condition de c).

\item Donner un exemple où $C$ est stable par produits fibrés (a
  fortiori, $E$ est un topos localement cohérent \ie localement
  algébrique) mais où $E$ n'est pas un topos algébrique. (Prendre
  l'exemple donné dans d), et la sous-catégorie pleine $C'$ de $E$
  formée des objets $Y$, $Z$, $T_{i}$ $(i > 0)$ et $∅_{E}$).

\item Supposons la catégorie $C$ finie. Prouver que le topos $E$ est
  noethérien, que ses objets noethériens (\ie quasi-compacts) sont les
  contrafoncteurs $F : C^{\circ} \to (\Ens)$ tels que pour tout objet
  $X$ de $C$, $F(X)$ soit un ensemble \emph{fini}, et que les
  foncteurs fibres sur $E$ transforment objets noethériens en
  ensembles finis (\cf IV 7.6. h)).

\item Supposons que tout morphisme $f: X \to X$ de $C$ se factorise en
  un composé $ip$, où $i$ est un monomorphisme et où $p$ admet un
  inverse à droite. Soit $X$ un objet de $C$, $I(X)$ l'ensemble de ses
  sous-objets au sens de $C$, considéré comme un ensemble ordonné,
  $S(I(X))$ l'ensemble des parties $U$ de $I(X)$ telles que pour deux
  éléments $X'$, $X'' \in I(X)$, $X'\in U$ et $X''\leq X'$ implique
  $X''\in U$. Montrer que l'ensemble ordonné des sous-objets dans $E$
  de $X$ est isomorphe à l'ensemble $S(I(X))$ (ordonné par
  inclusion). En conclure que si $I(X)$ est fini, alors $X$ est un
  objet prénoethérien de $E$. En particulier, si (en plus de la
  condition de factorisation ci-dessus) $C$ est stable par produits
  fibrés (\resp par limites projectives finies) et\pageoriginale si
  pour tout $X \in
  \ob C$, l'ensemble $I(X)$ des sous-objets de $X$ dans $C$ est fini,
  alors le topos $E$ est localement noethérien (\resp noethérien).

\item Soit $C$ la catégorie des ensembles finis (ou, au choix, des
  ensembles finis non vides) $\in 𝒰$, de sorte que $E = \widehat{C}$
  est la catégorie des ensembles simpliciaux augmentés (\resp des
  ensembles simpliciaux tout court). Montrer que $E$ est un topos
  noethérien. (Utiliser h).) Prenant un pro-objet $(X_i)_{i \in I}$
  non essentiellement constant de $C$, montrer que $E$ admet des
  foncteurs fibres qui ne transforment pas objets noethériens en
  objets noethériens (\ie en ensembles finis).

\item On considère le diagramme d'implications suivant de propriétés
  pour un topos $E$:
\[
\xymatrix@R=.75cm{
**[r]\txt{E est noethérien} \ar@{=>}[d] \ar@{=>}[r]& **[r]\txt{E est
  localement  noethérien} \ar@{=>}[d] & \ar@{=>}[rdd]\\
**[r]\txt{E est cohérent} \ar@{=>}[r] & **[r]\txt{E
est algébrique} \ar@{=>}[d]\\
 & **[r]\txt<9.5pc>{E est loc. cohérent, i.e.\\E \raisebox{-1ex}{cohalg}
  engendre E} \ar@{=>}[d]& & \txt{E \raisebox{-1ex}{prénoeth} engendre
  E} \ar@{=>}[ldd] \\
 & **[r]\txt{E \raisebox{-1ex}{coh} engendre E} \ar@{=>}[d]&\\
 & **[r]\txt{E \raisebox{-1ex}{qucpct} engendre E}  &&.\\
}
\]
Montrer que toutes les implications de ce diagramme sont strictes,
et qu'il n'y a pas entre les notions envisagées d'autres implications
que les implications composées du diagramme précédent. Ici $E_{\coh},
E_{\text{qucpct}}$, $E_{\text{prénoeth}}$,
$E_{\text{cohalg}}$\pageoriginale
désignent respectivement les sous-catégories pleines de $E$ formées des objets cohérents, \resp
quasi-compacts, \resp prénoethériens, \resp cohérents et algébriques de
$E$. (On se bornera à des topos de la forme $\widehat{C}$, en
utilisant les résultats énoncés dans d), e), f).)

\item Résoudre la question suivante (dont le rédacteur de ces lignes
  ignore la réponse) : $E = \widehat{C}$ peut-il être localement
  noethérien sans que les $\Hom(X, Y)$ (pour $X$, $Y \in \ob C$) soient
  finis ?
\end{enumerate}
\end{enonce*}

\begin{enonce*}[remark]{Exercice 2.18}\etiquette[2.18]{VI.exer2.18}
Soit $E$ un topos.
\begin{enumerate}
\renewcommand{\theenumi}{\alph{enumi}}
\renewcommand{\labelenumi}{\theenumi)}
\item Soit $X$ un objet prénoethérien de $E$. Montrer que $X$ est
  isomorphe à une somme finie d'objets connexes (IV 4.3.5) de
  $E$. (Montrer par l'absurde qu'une suite croissante de partitions de
  $X$ (IV 8.7) est stationnaire.)

\item Soient $X$ un objet de $E$, $(f_i : X_i \to X)_{i \in I}$ une
  famille couvrante de $X$. Montrer que si chacun des $X_i$ est
  isomorphe à une somme d'objets connexes de $E$, il en est de même de
  $X$. (Se ramener au cas où les $X_i$ sont connexes, puis au cas où
  ce sont des sous-objets de $X$, et considérer alors sur l'ensemble
  d'indices $I$ la relation d'équivalence $R$ engendrée par la
  relation $X_i \cap X_j \neq ∅_E$, et montrer que si $J = I/R$,
  $X$ est somme des $X(i) = \displaystyle{\mathop{\Sup}_{i \in I}} X_i
  .)$

\item Conclure de b) que si on désigne par $(X_i)_{i \in I}$ une
  famille génératrice de $E$, alors $E$ est localement connexe (IV
  8.7. $\ell$)) si\pageoriginale et seulement si chacun des $X_i$ $(i \in I)$
  est isomorphe à une somme d'objets connexes de $E$.

\item Conclure de a) et c) que si $E$ admet une sous-catégorie
  génératrice formée d'objets prénoethériens, en particulier si $E$
  est localement noethérien, alors $E$ est localement connexe, et a
  fortiori est isomorphe au topos somme (IV 8.7 b)) d'une famille de
  topos connexes (IV 8.7 e)) \ie dont l'objet final est connexe. (En
  particulier, on peut associer à $E$, pour tout morphisme $f : P \to
  E$, où $P$ est un topos « connexe non vide et simplement
  connexe » (IV 2.7.5) un pro-groupe fondamental $\pi_1(E, f)$; on
  notera que si $E$ est localement noethérien « non vide » on peut
  toujours trouver un tel $f$, avec $P$ le topos ponctuel, grâce à
  \textsc{Deligne} ([9]).
\end{enumerate}
\end{enonce*}

\section{Conditions de finitude pour un morphisme de topos}\etiquette[3]{VI.sec3}

\begin{enonce*}{Définition 3.1}\etiquette[3.1]{VI.def3.1}
Soit $f : E'\to E$ un morphisme de topos. On dit que $f$ est
\emph{quasi-compact (\resp quasi-séparé)} si pour tout objet quasi-compact
(\resp quasi-séparé) $X$ de $E$, $f^{\ast} (X)$ est quasi-compact
(\resp quasi-séparé). On dit que $f$ est cohérent si $f$ est
quasi-compact et quasi-séparé.
\end{enonce*}

\setcounter{subsection}{1}
\setcounter{subsubsection}{0}
\subsubsection{}\etiquette[3.1.1]{VI.subsubsec3.1.1}
On notera que si $f$ possède une des propriétés précédentes, alors
tout morphisme de topos isomorphe à $f$ possède la même propriété. Il
est trivial également que le composé de deux morphismes de topos
quasi-compacts (\resp quasi-séparés, \resp cohérents) est encore
quasi-compact (\resp quasi-séparé, \resp cohérent).

\begin{enonce*}{Proposition 3.2}\etiquette[3.2]{VI.prop3.2}
Avec les\pageoriginale notations de \hyperref[VI.def3.1]{3.1}, soit $(X_i)_{i \in I}$ une famille
génératrice d'objets quasi-compacts (\resp cohérents) de
$E$. Supposons, dans le cas respé, que $E'$ admette une famille
génératrice formée d'objets quasi-compacts. Pour que $f$ soit
quasi-compact (\resp cohérent), il faut et il suffit que pour tout $i
\in I$, $f^{\ast}(X_i)$ soit quasi-compact (\resp cohérent).
\end{enonce*}

La nécessité de la condition est triviale. Pour la suffisance dans
le premier cas, on note que pour tout objet quasi-compact $X$ de
$E$, il y a une famille épimorphique finie de morphismes $X_i \to
X$, donc il y a une famille épimorphique finie $f^{\ast} (X_i) \to
f^{\ast}(X)$, avec les $f^{\ast}(X_i)$ quasi-compacts par hypothèse,
donc $f^{\ast}(X)$ est quasi-compact (\hyperref[VI.prop1.3]{1.3}). Dans le
deuxième cas, il reste à voir que si $X$ est un objet quasi-séparé
de $E$, alors $f^{\ast}(X)$ est quasi-séparé. L'hypothèse sur $X$
peut s'exprimer (\hyperref[VI.cor1.17]{1.17}) par l'existence d'une famille
épimorphique de morphismes $X_i \to X$ telle que les
$\fprod{X_i}{X_j}{X}$ soient quasi-compacts. Ceci dit, on aura une
famille couvrante $f^{\ast}(X_i) \to f^{\ast}(X)$, telle que les
produits fibrés $\fprod{f^{\ast}(X_i)}{f^{\ast}(X_j)}{f^{\ast}(X)}$
sont quasi-compacts (puisque en vertu de a), $f^{\ast}$ transforme
objets quasi-compacts en objets quasi-compacts). Comme les
$f^{\ast}(X_i)$ sont cohérents, on conclut encore à l'aide de
\hyperref[VI.cor1.17]{1.17}.

\begin{enonce*}{Corollaire 3.3}\etiquette[3.3]{VI.cor3.3}
Soient $C$ et $C'$ deux $𝒰$-sites où les produits fibrés soient
représentables et où toute famille couvrante admet une sous-famille
couvrante finie, $g : C \to C'$ un morphisme de sites \emph{(IV
4.9.1)}. Alors le morphisme de topos $f: C^{\sim} \to {C'}^{\sim}$
défini par $g$ est cohérent.
\end{enonce*}

En effet,\pageoriginale $\epsilon_C(C)$ (\resp $\epsilon_{C'}(C')$) est une
famille génératrice de $C$ (\resp de $C'$) satisfaisant les conditions
respées de \hyperref[VI.prop3.2]{3.2} (\hyperref[VI.cor2.1.1]{2.1.1}).

\begin{enonce*}{Proposition 3.4}\etiquette[3.4]{VI.prop3.4}
Soient $E$ un topos localement cohérent (\hyperref[VI.def2.3]{2.3}), $f:X \to Y$ une
flèche de $E$, d'où un morphisme de topos induits \emph{(IV (5.5.2))} $E_{/X}
\to E_{/Y}$. Pour que ce dernier soit un morphisme de topos
quasi-compact (\resp quasi-séparé, \resp cohérent), il faut et il
suffit que $f$ soit un morphisme quasi-compact (\resp quasi-séparé,
\resp cohérent).
\end{enonce*}

Le cas « quasi-compact » est trivial en vertu des définitions (sans
condition sur $E$), le cas « quasi-séparé » n'est autre que \hyperref[VI.cor2.8]{2.8},
enfin le cas « cohérent » résulte de la conjonction des deux cas
précédents.

\begin{enonce*}{Proposition 3.5}\etiquette[3.5]{VI.prop3.5}
Soit $f : E'\to E$ un morphisme de topos localement cohérents. Soit
$(X_i)_{i \in I}$ une famille génératrice dans $E$ formée d'objets
cohérents algébriques. Les conditions suivantes sont équivalentes:
\begin{itemize}
\item[(i)] Pour toute flèche $u$ quasi-compacte de $E$, $f^{\ast}(u)$
est quasi-compacte.

\item[(i bis)] Pour toute flèche $u : X_i \to X_j$, $f^{\ast}(u)$ est
  quasi-compacte.

\item[(ii)] Pour tout objet cohérent $Y'$ de $E'$, tout objet cohérent
  algébrique $Y$ de $E$, et toute flèche $v : Y'\to f^{\ast}(Y)$, le
  morphisme de topos correspondant
$$
f_v : {E'}_{/Y'} \to E_{/Y}
$$
composé du\pageoriginale morphisme de localisation ${E'}_{/Y'} \to
{E'}_{/f^{\ast}(Y)}$ déduit de $v$ \emph{(IV (5.5.2))} et du morphisme
${E'}_{/f^{\ast}(Y)} \to E_{/Y}$ induit par $f$ \emph{(IV (5.10.1))} est
cohérent.

\item[(ii bis)] Même condition que dans \emph{(ii)}, mais en se bornant à un
  ensemble de données $v_{\alpha} : {Y'}_{\alpha} \to
  f^{\ast}(Y_{\alpha})$ tel que les ${Y'}_{\alpha}$ soient algébriques
  et recouvrent l'objet final de $E'$.

\item[(ii ter)] (Lorsqu'on se donne une famille génératrice
  $(X'_{i'})_{i'\in I'}$ dans $E'$ formée d'objets cohérents, et
  pour tout $i'\in I'$, un $i \in I$ et une flèche $v_{i'} :
  X'_{i'} \to f(X_i)$.) Pour tout $i' \in I'$, le morphisme de
  topos $E'_{/{X'}_{i'}} \to E_{/X}$ défini par $v_i$ est cohérent.
\end{itemize}
\end{enonce*}

Comme toute flèche $X_i \to X_j$ est quasi-compacte, il est évident
que (i) $\Rightarrow$ (i bis). Inversement, supposons (i bis) vérifié
et prouvons (i). Soit $u : X \to Y$ une flèche quasi-compacte de
$E$. En termes d'une famille épimorphique $X_i \to Y$, l'hypothèse sur
$u$ s'exprime donc par le fait que les $\fprod{X}{X_i}{Y}$ sont
quasi-compacts (\hyperref[VI.cor1.16]{1.16}), \ie qu'il existe pour chaque $i$ une famille
finie épimorphique de morphismes $X_j \to \fprod{X}{X_i}{Y}$
(\hyperref[VI.prop1.3]{1.3}). Alors on a une famille épimorphique correspondante
$f^{\ast}(X_i) \to f^{\ast}(Y)$, telle que pour tout $i$ on ait une
famille épimorphique finie $f^{\ast}(X_j) \to
f^{\ast}(\fprod{X}{X_i}{Y}) \simeq
\fprod{f^{\ast}(X)}{f^{\ast}(X_i)}{f^{\ast}(Y)}$, dont le composé avec
la projection $\pr_2$ dans $f^{\ast}(X_i)$ est un morphisme
quasi-compact $f^{\ast}(X_j) \to f^{\ast}(X_i)$, grâce à l'hypothèse
(i bis). Il s'ensuit (\hyperref[VI.cor1.11]{1.11} (i)) que pour tout $i$, $\pr_2 :
\fprod{f^{\ast}(X)}{f^{\ast}(X_i)}{f^{\ast}(Y)}$ est quasi-compact,
donc (\hyperref[VI.prop1.10]{1.10} (ii)) $f^{\ast}(u) : f^{\ast}(X) \to f^{\ast}(Y)$ est quasi-compact.

Il reste\pageoriginale à prouver les implications (i) $\Rightarrow$
(ii) et  (ii bis) $\Rightarrow$ (i), l'implication  (ii) $\Rightarrow$ (ii bis)
étant triviale, et (ii ter) étant un cas particulier de (ii
bis). L'implication  (i) $\Rightarrow$ (ii) est immédiate : en effet,
un objet $X$ de $E_{/Y}$ est quasi-compact (\resp quasi-séparé) si et
seulement si le morphisme structural $X \to Y$ (\resp le morphisme
diagonal $X \to \fprod{X}{X}{Y}$) est quasi-compact (\hyperref[VI.cor2.7]{2.7}), et on a le
même critère de quasi-compacité et de quasi-séparation pour
$f^{\ast}_v (X')$. Prouvons enfin  (ii bis) $\Rightarrow$ (i). Soit
donc $u : X \to Y$ une flèche quasi-compacte dans $E$, prouvons que
$f^{\ast}(u)$ est quasi-compacte. Comme les $Y'_{\alpha}$
recouvrent l'objet final de $E'$, il suffit de prouver que les
$f^{\ast}(u) | Y'_{\alpha}$ sont quasi-compacts (\hyperref[VI.prop1.10]{1.10} (ii)). Or,
ces morphismes ne sont autres que les $f^{\ast}_{v_\alpha}(u |
Y_{\alpha})$, et on est donc réduit à prouver ceci:

\begin{enonce*}{Corollaire 3.6}\etiquette[3.6]{VI.cor3.6}
Soit $f : E'\to E$ un morphisme cohérent de topos localement
cohérents. Alors pour tout flèche quasi-compacte $u$ de $E$,
$f^{\ast}(u)$ est quasi-compacte.
\end{enonce*}

Grâce à l'implication (i bis) $\Rightarrow$ (i) de \hyperref[VI.prop3.4]{3.4} déjà prouvée,
il suffit de prouver que si $u : X \to Y$ est une flèche de $E$, avec
$X$, $Y$ cohérents, alors $f^{\ast}(u)$ est un morphisme
quasi-compact, ce qui résulte
aussitôt de l'hypothèse sur $f$, impliquant que $f^{\ast}(X)$
et $f^{\ast}(Y)$ sont également cohérents.

\begin{enonce*}{Définition 3.7}\etiquette[3.7]{VI.def3.7}
Soit $f : E'\to E$\pageoriginale un morphisme de topos localement
cohérents. On dit que $f$ est \emph{localement cohérent} s'il satisfait aux conditions
équivalentes de \hyperref[VI.prop3.5]{3.5}.
\end{enonce*}

Remarquons que « localement » signifie local en haut.

\setcounter{subsection}{7}
\setcounter{subsubsection}{0}
\subsubsection{}\etiquette[3.7.1]{VI.subsubsec3.7.1}
Cette notion ne dépend encore que de la classe d'isomorphie du
morphisme de topos $f$, et elle est manifestement stable par
composition de morphismes de topos. De plus, en vertu de \hyperref[VI.prop3.5]{3.5}, si $f$
est cohérent il est localement cohérent, la réciproque étant vraie si
$E'$ et $E$ sont cohérents (en vertu de \hyperref[VI.prop3.5]{3.5} (ii bis)). Du critère \hyperref[VI.prop3.5]{3.5}
(ii bis) résulte aussitôt le critère suivant:

\begin{enonce*}{Corollaire 3.7}\etiquette[3.7]{VI.cor3.7}
Soient $f : E'\to E$ un morphisme de topos localement cohérents,
$(Y_i)_{i \in I}$, $({Y'}_i)_{i \in I}$ deux familles d'objets de
$E$ et de $E'$, et pour tout $i \in I$ $v_i : Y'_i \to
f^{\ast}(Y_i)$ un morphisme, d'où un morphisme de topos
$$
f_{v_i} : E'_{/ Y'_i} \to E_{/Y_i}{\quoi}.
$$
Supposons que les ${Y'}_i$ recouvrent l'objet final de $E'$. Alors $f$
est localement cohérent si et seulement si pour tout $i \in I$,
$f_{v_i}$ l'est.
\end{enonce*}

En particulier, appliquant ceci au morphisme identique d'un topos
induit, on trouve:

\begin{enonce*}{Corollaire 3.8}\etiquette[3.8]{VI.cor3.8}
Soient $E$ un topos localement cohérent, $v : X \to Y$ une flèche de
$E$. Alors le morphisme des topos induits $E_{/X} \to E_{/Y}$ défini
par $v$ est localement cohérent.
\end{enonce*}

\begin{enonce*}{Remarque 3.9}\etiquette[3.9]{VI.rem3.9}
Il semble\pageoriginale que tous les morphismes de topos algébriques
qu'on ait rencontrés en pratique soient localement cohérents (c'est
pourquoi le terme « localement cohérent » n'est pas appelé sans
doute à un très grand usage). Il revient au même d'affirmer que les
morphismes de topos cohérents qu'on a rencontrés en pratique sont
cohérents. On peut cependant construire des morphismes non cohérents
de topos cohérents, et plus particulièrement, un morphisme non
cohérent du topos ponctuel \emph{(IV 2.2)} dans un topos noethérien
$E$, \emph{\cf \hyperref[VI.exer2.17]{2.17} i)}.
\end{enonce*}

\begin{enonce*}[remark]{Exemples 3.10}\etiquette[3.10]{VI.exem3.10}
(\emph{Topos associés aux schémas}.)
Reprenons l'exemple \hyperref[VI.subsec1.22]{1.22} de la catégorie ($\Sch$) avec des topologies
$T_i$ de \SGA IV 6.3. Soit, pour tout schéma $X$, $X_{T_i}$ le
$𝒱$-topos défini par le site induit $(\Sch)_{/X}$. Alors le morphisme
de schémas $f : X \to Y$ définit un morphisme de topos $f_{T_i} :
X_{T_i} \to Y_{T_i}$, et il résulte de \hyperref[VI.prop3.4]{3.4} et de \hyperref[VI.subsec1.22]{1.22} que ce dernier
morphisme de topos est quasi-compact (\resp quasi-séparé, \resp
cohérent) si et seulement si le morphisme de schémas $f$ est
quasi-compact (\resp quasi-séparé, \resp cohérent) au sens habituel de
\EGA IV 1. En tout état de cause, $f_{T_i}$ est localement cohérent en
vertu de \hyperref[VI.cor3.8]{3.8}.

On peut aussi associer à un schéma $X$ les $𝒰$-topos $\Top(X)$ et
$\Top(X_{\text{ét}})$ comme dans \hyperref[VI.subsubsec1.22.1]{1.22.1}, et à tout morphisme de schémas $f :
X \to Y$ sont alors associés des morphismes $\Top(f)$ et
$\Top(f_{\text{ét}})$ des topos correspondants. Soit $T(f)$ l'un de ces deux
morphismes\pageoriginale de topos. On vérifie immédiatement via \hyperref[VI.prop3.2]{3.2}
que ce morphisme de topos est toujours localement cohérent, et qu'il est quasi-compact
(\resp quasi-séparé, \resp cohérent) si et seulement si le morphisme
de schémas $f$ est quasi-compact (\resp quasi-séparé, \resp cohérent)
au sens habituel.

Ces observations montrent donc encore que la terminologie introduite
dans le présent numéro est compatible avec la terminologie reçue en
théorie des schémas, et mérite donc d'être acceptée par le lecteur
le plus récalcitrant.
\end{enonce*}

\begin{enonce*}[remark]{Exercice 3.11}\etiquette[3.11]{VI.exer3.11}
(\emph{Topos cohérents et prétopos}.)
\begin{enumerate}
\renewcommand{\theenumi}{\alph {enumi}}
\renewcommand{\labelenumi}{\theenumi)}
\item On appelle $𝒰$-prétopos (ou simplement prétopos), une catégorie
  $C$ satisfaisant aux conditions suivantes:
\begin{enumerate}
\renewcommand{\theenumii}{\arabic{enumii}}
\renewcommand{\labelenumii}{\theenumii)}
\item Les limites projectives finies dans $C$ sont représentables.

\item Les sommes finies dans $C$ sont représentables, elles sont
  disjointes et universelles.

\item Les relations d'équivalence dans $C$ sont effectives, et tout
  épimorphisme dans $C$ est effectif universel.

\item $C$ est équivalente à une catégorie $\in 𝒰$.
\end{enumerate}
Montrer que si $E$ est un $𝒰$-topos cohérent, alors la
sous-catégorie pleine $E_{\coh}$ de $E$ formée des objets cohérents
de $E$ est un $𝒰$-prétopos, et que le foncteur d'inclusion
$E_{\coh} \to E$ est exact à gauche, commute aux\pageoriginale
sommes finies et au passage au quotient par des relations
d'équivalence. (Utiliser \hyperref[VI.subsubsec1.5.3]{1.5.3}, \hyperref[VI.cor1.15]{1.15}
et \hyperref[VI.cor1.17.1]{1.17.1}). Montrer que la topologie induite par $E$ sur
$C$ est la topologie « précanonique », \ie la topologie dont les
familles couvrantes $X_i \to X$ sont celles qui admettent une
sous-famille \emph{finie} couvrante pour la topologie canonique de
$C$. Par suite $E$ se reconstitue à équivalence près par la
connaissance du prétopos $C = E_{\coh}$, comme le topos $C^{\sim}$
($C$ étant munie de sa topologie précanonique).

\item Soit $C$ un $𝒰$-prétopos. Munissons $C$ de la topologie
  précanonique. Montrons que tout morphisme $f$ de $C$ se factorise en
  $f'' f'$ , avec $f'$ un épimorphisme et $f''$ un
  monomorphisme. Montrer que le topos $E = C^{\sim}$ est cohérent, et
  que le foncteur canonique $\epsilon : C \to E$ induit une
  équivalence de $C$ avec la sous-catégorie $E_{\coh}$ de $E$. (Montrer
  d'abord que $\varepsilon$ est un foncteur pleinement fidèle exact à
  gauche commutant aux sommes finies et au passage au quotient par une
  relation d'équivalence, puis qu'un sous-objet cohérent dans $E$ d'un
  objet de $C$ est dans l'image essentielle de $C$.) Par suite, le
  $𝒰$-prétopos $C$ se reconstitue à équivalence de catégories près
  quand on connaît le $𝒰$-topos associé $E = C^{\sim}$.

\item Soient $C$ et $C'$ deux $𝒰$-prétopos. Montrer que pour qu'un
  foncteur $\varphi : C'\to C$ soit un morphisme de sites de $C$ dans
  $C'$ (pour les topologies précanoniques), il faut et il suffit que
  $\varphi$ soit exact à gauche, et commute aux sommes finies et au
  passage au quotient par une relation d'équivalence.

\item Avec les\pageoriginale notations de c), supposons que $C$ et $C'$ soient
  associés à deux $𝒰$-topos cohérents $E$, $E'$ comme dans a). Montrer
  que le foncteur canonique (IV 4.9.3)
$$
\Mmorsite (C, C') \to \Hhomtop(E, E')
$$
induit une équivalence du premier membre (explicité dans c)) avec la
sous-catégorie pleine du deuxième formé des morphismes de topos $E \to
E'$ qui sont \emph{cohérents}, un foncteur quasi-inverse étant obtenu
en associant à tout morphisme cohérent de topos $f : E \to E'$ le
foncteur $E'_{\coh} = C' \to E_{\coh} = C$ induit par $f^{\ast}$.

\item Soient $C$ un $𝒰$-prétopos et $E = C^{\sim}$. Exprimer
  directement en termes de propriétés d'exactitude de $C$ les
  conditions équivalentes de \hyperref[VI.cor1.25]{1.25} (Consulter \hyperref[VI.exer1.28]{1.28} b.) Montrer qu'un
  objet $X$ de $C$ est noethérien dans $E$ si et seulement si toute
  suite croissante de sous-objets de $X$ dans $C$ est stationnaire,
  donc que $E$ est noethérien si et seulement si tout objet de $C$
  satisfait à la condition précédente.
\end{enumerate}
\end{enonce*}

\begin{enonce*}[remark]{Exercice 3.12}\etiquette[3.12]{VI.exer3.12}
Un topos $E$ est appelé un topos \emph{fini} s'il existe une catégorie
finie $C$ telle que $E$ soit équivalent à $\widehat{C}$.
\begin{enumerate}
\renewcommand{\theenumi}{\alph{enumi}}
\renewcommand{\labelenumi}{\theenumi)}
\item Pour que $E$ soit un topos fini, il faut et il suffit que $E$
  admette suffisamment de points essentiels (IV 7.6 b)), et que la
  catégorie $\Ppointess(E)$ des points essentiels de $E$ soit
  équivalente à une catégorie finie. (Utiliser IV 7.6 d) et h).)

\item Supposons\pageoriginale $E$ fini. Prouver que tout point de $E$ est
  essentiel. Prouver que les 2-foncteurs $E \mapsto \Ppoint(E)$ et $C
  \mapsto \widehat{C}$ établissent des 2-équivalences entre la 2-catégorie
  des topos finis $E$, et la 2-catégorie des catégories karoubiennes
  (IV 7.5 a)) $C$ qui sont équivalentes à des catégories finies, et
  que tout morphisme de topos finis est essentiel (IV 7.6
  a)). (Utiliser a) et IV 7.6 h).)

\item Supposons $E = \widehat{C}$, avec $C$ équivalente à une
  catégorie finie, et soit $F$ un topos cohérent. Rappelons (IV
  (4.6.3.1.)) que les morphismes de topos $f : E \to F$ correspondent
  aux foncteurs $u : C \to \Ppoint(F)$. Montrer que pour que $f$ soit
  cohérent, il faut et il suffit que $f$ transforme point cohérent de
  $E$ (\hyperref[VI.def3.1]{3.1}) en point cohérent de $F$, ou encore que pour tout $X \in
  \ob C$, $u(X)$ soit un point \emph{cohérent} de $F$. (Utiliser \hyperref[VI.prop3.2]{3.2} et
  2,17 g).) En conclure que tout morphisme $f : E \to F$ d'un topos
  fini dans un topos fini est cohérent.

\item Soit $X$ un espace topologique sobre (IV 4.2.1). Pour que le
  topos $\Top(X)$ (IV 2.1) soit fini, il faut et il suffit que $X$
  soit un ensemble fini. (Utiliser IV 7.1.6)
\end{enumerate}
\end{enonce*}

\section{Conditions de finitude dans un topos obtenu par
  recollement}\etiquette[4]{VI.sec4}

Le présent paragraphe ne sera plus utilisé dans la suite du Séminaire.

\subsection{}\etiquette[4.1]{VI.subsec4.1}
Soient $E$ un topos, $U$ un ouvert de $E$ (\ie un sous-objet de
l'objet final de $E$), et considérons les sous-topos ouverts et fermés
correspondants de\pageoriginale $E$ (IV 9)
$$
E'= E_{/U}, E'' = E_{\complement U}{\quoi}.
$$
Nous désignerons par
$$
j : E'\to E, i : E''\to E
$$
les morphismes de topos canoniques. Rappelons (\hyperref[VI.prop3.4]{3.4}) que pour que $j$
soit un morphisme quasi-compact (\resp cohérent) de topos, il faut et
il suffit que l'inclusion
$$
j : U \to e
$$
de $U$ dans l'objet final de $E$ (inclusion que nous noterons
également $j$) soit un morphisme quasi-compact (\resp cohérent) dans
$E$; d'autre part, $j : E'\to E$ est toujours quasi-séparé (car $j : U
\to e$ l'est (\hyperref[VI.cor1.8.1]{1.8.1})). Nous nous proposons de donner des critères pour
que le topos $E''$ soit cohérent, et le morphisme de topos $i : E''\to
E$ soit cohérent. Notons d'abord:

\begin{enonce*}{Proposition 4.2}\etiquette[4.2]{VI.prop4.2}
Les notations étant celles de \hyperref[VI.subsec4.1]{4.1}, le morphisme d'inclusion $i :
E''\to E$ est quasi-compact, \ie pour tout objet quasi-compact $X$ de
$E$, l'objet $i^{\ast} (X)$ de $E''$ est quasi-compact; de plus, si
$X$ est prénoethérien (\hyperref[VI.def1.30]{1.30}), $i^{\ast}(X)$ est prénoethérien.
\end{enonce*}

Cela résulte de la définition 1.1 et du

\begin{enonce*}{Lemme 4.2.1}\etiquette[4.2.1]{VI.lem4.2.1}
L'application\pageoriginale $Y \mapsto i^{\ast}(Y)$ induit un isomorphisme
d'ensembles ordonnés entre l'ensemble des sous-objets de $X$ qui
contiennent le sous-objet $X_U = X \times U$ de $X$, et l'ensemble des
sous-objets de $i^{\ast}(X)$.
\end{enonce*}

Notons que, $i_{\ast}$ étant conservatif (car pleinement fidèle) et
exact à gauche (en particulier, commutant aux produits fibrés), il
s'ensuit aussitôt qu'un morphisme $u : Y \to Z$ dans $E''$ est un
monomorphisme si et seulement si $i_{\ast}(u) : i_{\ast}(Y) \to
i_{\ast}(Z)$ l'est . Il s'ensuit aussitôt, identifiant (par
$i_{\ast}$) $E'$ à la sous-catégorie pleine de $E$ formée des objets
$Z$ satisfaisant aux conditions équivalentes de IV 9.7 1), que les
sous-objets dans $E'$ de l'objet $Z$ de $E''$ s'identifient aux
sous-objets $Y$ de $Z$ dans $E$ qui veulent bien appartenir à $E''$,
ou, ce qui revient au même, qui contiennent $Z_U \xrightarrow{\sim}
U$. Or pour $Z$ de la forme $i_{\ast} i^{\ast} (X) = X_{\complement
U}$ donné par la somme amalgamée IV (9.5.1),
$$
i_{\ast} i^{\ast} (X) = X_{\complement U} = X \amalg_{X_U} U {\quoi},
$$
la donnée d'un sous-objet $Y$ de ce dernier dans $E$ équivaut à la
donnée d'un couple formé d'un sous-objet $Y_1$ de $X$ et d'un sous-objet
$Y_2$ de $U$, avec la condition que les images inverses de ces
sous-objets dans $X_U$ coïncident. La condition que $Y$ contienne $U$
signifie alors que $Y_2 = U$, et la condition qui reste sur $Y_1$ est
que $Y_1$ contienne $X_U$. Cela établit donc une bijection entre
l'ensemble des sous-objets de $i^{\ast}(X)$ dans $E''$, et l'ensemble
des sous-objets de $X$ qui contiennent $X_U$. Il est
clair\pageoriginale que c'est un isomorphisme d'ensembles ordonnés,
et que l'application inverse est bien celle annoncée
dans 3.14.1.

\begin{enonce*}{Corollaire 4.3}\etiquette[4.3]{VI.cor4.3}
\begin{enumerate}
\renewcommand{\theenumi}{\alph{enumi}}
\renewcommand{\labelenumi}{\theenumi)}
\item Soit $X$ un objet de $E$. Pour que $X$ soit quasi-compact, il
  suffit que $i^{\ast}(X)$ et $j^{\ast}(X)$ le soient, et cette
  condition est également nécessaire si $j : U \to e$ est
  quasi-compact.

\item Soit $Y$ un objet de $E''$. Pour que $Y$ soit quasi-compact, il
  suffit que $i_{\ast} (Y)$ le soit, et cette condition est également
  nécessaire si $U$ est quasi-compact.
\end{enumerate}
\end{enonce*}

\noindent
\underline{Démonstration}
\begin{enumerate}
\renewcommand{\theenumi}{\alph{enumi}}
\renewcommand{\labelenumi}{\theenumi)}
\item La suffisance résulte de la définition 1.1 et du fait que le
  couple $(i^{\ast}, j^{\ast})$ est conservatif (IV 9.11 3)). La
  nécessité résulte du fait que $i$ et $j$ sont quasi-compacts (en
  vertu de 3.14 et de l'hypothèse que $j$ est quasi-compact).

\item Comme $Y \simeq i^{\ast} i_{\ast}(Y)$, la suffisance résulte de
  \hyperref[VI.prop4.2]{4.2}. La nécessité résulte de la suffisance dans a), compte tenu que
  $i^{\ast}i_{\ast}(Y) \simeq Y$ et $j^{\ast}i_{\ast}(Y)\simeq U$.
\end{enumerate}

\begin{enonce*}{Corollaire 4.4}\etiquette[4.4]{VI.cor4.4}
Soit $Y$ un objet de $E''$. Si $U$ est quasi-compact ou si $E$ admet
une famille génératrice formée d'objets quasi-compacts, alors pour que
$Y$ soit quasi-séparé (\resp cohérent), il \emph{suffit} qu'il en soit ainsi
de $i_{\ast} (Y)$. Si $U$ est quasi-séparé (\resp cohérent) et si $j :
U \to e$  est un morphisme quasi-compact, alors pour que $Y$ soit
quasi-séparé (\resp cohérent), il faut que $i_{\ast}(Y)$ le soit.
\end{enonce*}

Le cas\pageoriginale respé résulte du cas non respé, compte tenu de
\hyperref[VI.cor4.3]{4.3} b). Supposons $i_{\ast}(Y)$ quasi-séparé, et
prouvons qu'il en est de même de $Y$, sous l'une des deux hypothèses
faites. Il faut donc prouver que pour deux objets quasi-compacts
$Y'$  et $Y''$ au-dessus de $Y$, le produit fibré
$\fprod{Y'}{Y''}{Y}$ est quasi-compact. Lorsqu'on suppose $Y$
quasi-compact, on note qu'il suffit de prouver que
$i_{\ast}(\fprod{Y'}{Y''}{Y})$ est quasi-compact (\hyperref[VI.cor4.3]{4.3}
b)), ce qui résulte du fait que $i_{\ast}(Y')$ et $i_{\ast}(Y'')$ le
sont  (\hyperref[VI.cor4.3]{4.3} b)), utilisant ici la quasi-compacité de
$U)$, que $i_{\ast}$ commute aux produits fibrés, et l'hypothèse
$i_{\ast}(Y)$ quasi-séparé. Lorsqu'on suppose que $E$ admet une
famille génératrice formée d'objets quasi-compacts, alors
$i_{\ast}(Y')$ est limite inductive filtrante de ses sous-objets
quasi-compacts ${X'}_{\alpha}$, donc $Y' \simeq i^{\ast}
i_{\ast}(Y')$ est limite inductive filtrante de sous-objets
$i^{\ast}({X'}_{\alpha})$, et comme $Y'$ est quasi-compact, il est
égal à un des $i^{\ast}({X'}_{\alpha})$. De même $Y''$ est de la
forme $i^{\ast}({X''}_{\beta})$, où ${X''}_{\beta}$ est un
sous-objet quasi-compact de $i_{\ast}(Y'')$. Mais alors
$\fprod{Y'}{Y''}{Y} =
i^{\ast}(\fprod{X'_{\alpha}}{X''_{\beta}}{i_{\ast}Y})$, et comme
$\fprod{X'_{\alpha}}{X''_{\beta}}{i_{\ast}(Y)}$ est quasi-compact en
vertu de l'hypothèse $i_{\ast}(Y)$ quasi-séparé, il s'ensuit que
$\fprod{Y'}{Y''}{Y}$ l'est aussi en vertu de \hyperref[VI.prop4.2]{4.2}.

Inversement, supposant $U$ quasi-séparé et $j : U \to e$
quasi-compact, montrons que si $Y$ est quasi-séparé, il en est de
même de $i_{\ast}(Y)$, \ie que pour deux objets $X'$, $X''$
quasi-compacts au-dessus de $i_{\ast}(Y)$, le produit fibré
$\fprod{X'}{X''}{i_{\ast}(Y)}$ est quasi-compact. Pour ceci,
appliquant \hyperref[VI.cor4.3]{4.3} a), il suffit de prouver que son image
par $i^{\ast}$ est quasi-compact\pageoriginale (ce qui résulte de
\hyperref[VI.prop4.2]{4.2} et de l'hypothèse que $Y$ est quasi-séparé)
\emph{et} que son image par $j^{\ast}$ est quasi-compact; or cette
dernière est $j^{\ast}(X')\times j^{\ast}(X'')$, et est bien
quasi-compacte car $j^{\ast}(X')$ et $j^{\ast}(X'')$ le sont ($j$
étant quasi-compact) et $U$ est quasi-séparé.

\begin{enonce*}{Corollaire 4.5}\etiquette[4.5]{VI.cor4.5}
Supposons que $E$ admette une sous-catégorie génératrice formée
d'objets quasi-compacts (\emph{\resp} prénoethériens), alors il en est de
même de $E''$.
\end{enonce*}

Cela résulte de \hyperref[VI.prop4.2]{4.2} et du

\begin{enonce*}{Lemme 4.5.1}\etiquette[4.5.1]{VI.lem4.5.1}
Si $(X_{\alpha})_{\alpha}$ est une famille génératrice dans $E$, alors
$(i^{\ast}(X_{\alpha}))_{\alpha}$ est une famille génératrice dans $E''$.
\end{enonce*}

Cela signifie en effet que la famille des foncteurs
$$
Y \to \Hom (i^{\ast}(X_{\alpha}), Y)
$$
sur $E''$ est conservative, or on a $\Hom(i^{\ast}(X_{\alpha}), Y)
\simeq \Hom(X_{\alpha}, i_{\ast} (Y))$, et il suffit d'utiliser le
fait que le foncteur $i_{\ast}$ est conservatif.

\begin{enonce*}{Proposition 4.6}\etiquette[4.6]{VI.prop4.6}
Les notations sont celles de \hyperref[VI.subsec4.1]{4.1}.
\begin{enumerate}
\renewcommand{\theenumi}{\alph{enumi}}
\renewcommand{\labelenumi}{\theenumi)}
\item Si $E$ admet une famille génératrice formée d'objets cohérents,
  il en est de même pour le sous-topos fermé $E''$, et le morphisme
  d'inclusion $i : E''\to E$ est cohérent.

\item Supposons que le morphisme $j : U \to e$ dans $E$ soit quasi-compact.
Si\pageoriginale $E$ est localement cohérent (\emph{\resp} cohérent, \emph{\resp}
quasi-séparé, \emph{\resp} localement noethérien, \emph{\resp} noethérien,
\emph{\resp} localement parfait, \emph{\resp} parfait) il en est de même de $E''$, et le morphisme
d'inclusion $i : E''\to E$ est cohérent.
\end{enumerate}
\end{enonce*}

\begin{enumerate}
\renewcommand{\theenumi}{\alph{enumi}}
\renewcommand{\labelenumi}{\theenumi)}
\item La première assertion résultera de la seconde et de \hyperref[VI.lem4.5.1]{4.5.1}. En
  tous cas, \hyperref[VI.cor4.5]{4.5} implique que $E''$ admet une famille génératrice
  formée d'objets quasi-compacts, donc \hyperref[VI.prop3.2]{3.2} s'applique et nous montre
  qu'il suffit de vérifier que pour tout objet cohérent $X$ de $E$,
  l'objet $i^{\ast}(X)$ de $E''$ est cohérent. Utilisant la
  compatibilité de la formation du topos complémentaire d'un ouvert
  avec la localisation (IV 9.13 b)), on est ramené au cas où $X$ est
  l'objet final de $E$, donc à prouver que l'objet final de $E''$ est
  cohérent. Or, \hyperref[VI.cor4.4]{4.4} s'applique, donc il suffit de prouver que
  $i_{\ast}(e_{E''}) = e_E$ est cohérent, ce qui est bien le cas.

\item Comme sous chacune des hypothèses faites dans b), $E$ admet une
  famille génératrice formée d'objets cohérents, il résulte déjà de a)
  que $i : E'' \to E$ est cohérent. De plus, \hyperref[VI.lem4.5.1]{4.5.1} implique alors que
  $E''$ admet la sous-catégorie pleine ${E''}_{\coh}$ formée de ses
  objets cohérents comme sous-catégorie génératrice. Supposons $E$
  cohérent, et montrons que $E''$ l'est aussi, \ie (\hyperref[VI.subsubsec2.4.5]{2.4.5}) montrons
  que la sous-catégorie ${E''}_{\coh}$ de $E''$ est stable par limites
  projectives finies, sachant qu'il en est ainsi pour la
  sous-catégorie $E_{\coh}$ de $E$; or cela résulte aussitôt du critère
  \hyperref[VI.cor4.4]{4.4} respé (qui s'applique, car $U$ est maintenant cohérent, $e$
  l'étant\pageoriginale et $j : U \to e$ étant quasi-compact), compte tenu que
  $i_{\ast}$ est exact à gauche. Par localisation, utilisant encore IV
  9.13 b), on en conclut que si $E$ est localement cohérent, il en est
  de même de $E''$. Si $E$ est quasi-séparé, alors $E''$ l'est aussi:
  en effet son objet final est quasi-séparé en vertu de \hyperref[VI.cor4.4]{4.4}, celui de
  $E$ l'étant, et il reste à prouver que le produit de deux objets
  quasi-séparés de $E''$ est quasi-séparé (sachant qu'il en est ainsi
  dans $E''$), ce qui résulte encore du critère \hyperref[VI.cor4.4]{4.4} et du fait que
  $i_{\ast}$ commute aux produits (compte tenu que le sous-objet $U$
  de $e$ est quasi-séparé, $e$ l'étant, de sorte que \hyperref[VI.cor4.4]{4.4}. s'applique).
\end{enumerate}

Comme un topos est localement noethérien (\resp noethérien) si et
seulement si il est localement cohérent (\resp cohérent) et admet une
famille génératrice formée d'objets prénoethériens (\hyperref[VI.prop2.10]{2.10} (iii
bis)), il résulte de ce qui précède et de \hyperref[VI.cor4.5]{4.5} que si $E$ est
localement noethérien (\resp noethérien), il en est de même de
$E''$. Supposons maintenant $E$ parfait, et prouvons que $E''$
l'est. Comme on sait déjà qu'il est cohérent, il reste à vérifier
qu'il satisfait au critère \hyperref[VI.cor1.25]{1.25} (i), \ie que tout objet $Y$ de $E''$
est limite inductive filtrante d'objets  cohérents, sachant que
l'énoncé analogue est vrai dans $E$; or $i_{\ast}(Y)$ étant limite
inductive filtrante d'objets cohérents $X_{\alpha}$, $Y \simeq
i^{\ast}i_{\ast}(Y)$ est limite inductive filtrante des objets
cohérents $i^{\ast}(X_{\alpha})$. Par localisation (IV 9.13 b)) on en
conclut que si $E$ est localement parfait, il en est de même de
$E''$. Cela achève la démonstration de \hyperref[VI.prop4.6]{4.6}.

\begin{enonce*}[remark]{Remarque 4.6.1}\etiquette[4.6.1]{VI.rem4.6.1}
En fait,\pageoriginale dans \hyperref[VI.prop4.6]{4.6} b) l'hypothèse que $j :
U \to e$ soit quasi-compact est inutile, sauf peut-être dans le cas
« $E$ quasi-séparé ». Il suffit en effet, en vertu de la
démonstration qui précède, de voir que $E$ cohérent implique $E''$
cohérent. Or $U$ est limite inductive filtrante de ses sous-objets
quasi-compacts $U_{\alpha}$, et on vérifiera alors dans
\hyperref[VI.sec7]{7}. qu'alors $E''$ s'identifie au topos limite projective
(au sens de \hyperref[VI.sec7]{7}.) des sous-topos fermés $E''_i$
complémentaire des $E_{/U_i}$, qui en vertu de \hyperref[VI.prop4.6]{4.6} b)
sont des topos cohérents à morphismes de transition cohérents. Il
s'ensuit alors que $E''$ est un topos cohérent (\hyperref[VI.sec7]{7}).
\end{enonce*}

\begin{enonce*}{Corollaire 4.7}\etiquette[4.7]{VI.cor4.7}
Supposons que la sous-catégorie $E_{\coh}$ de $E$ formée des objets
cohérents soit génératrice, et soit $X$ un objet de $E$.
\begin{enumerate}
\renewcommand{\theenumi}{\alph{enumi}}
\renewcommand{\labelenumi}{\theenumi)}
\item Pour que $X$ soit quasi-séparé, il faut que $i^{\ast}(X)$ et
  $j^{\ast}(X)$ le soient, et cette condition est aussi suffisante si $j
  :U \to e$ est quasi-compact.

\item Supposons que $j : U \to e$ soit quasi-compact. Alors $X$ est
  cohérent si et seulement si $i^{\ast}(X)$ et $j^{\ast}(X)$ le sont.

\item Supposons $E$ cohérent et $j : U \to e$ quasi-compact. Alors $E$
  est parfait (\hyperref[VI.def2.9.1]{2.9.1}) si et seulement si $E'$ et $E''$ le sont.
\end{enumerate}
\end{enonce*}

\begin{enumerate}
\renewcommand{\theenumi}{\alph{enumi}}
\renewcommand{\labelenumi}{\theenumi)}
\item On a signalé dans \hyperref[VI.subsec4.1]{4.1} que $j : E'\to E$ est quasi-séparé, et
  d'autre part on sait par \hyperref[VI.prop4.6]{4.6} a) qu'il en est de même de $i$, d'où la
  nécessité. Pour la suffisance, soient $X'$ et $X''$ des objets
  quasi-compacts au-dessus de $X$, il faut prouver que
  $\fprod{X'}{X''}{X}$ est quasi-compact, et pour ceci
  il\pageoriginale suffit de
  prouver que ses images par $i^{\ast}$ et $j^{\ast}$ le sont (\hyperref[VI.cor4.3]{4.3}
  a)). Comme $i$ et $j$ sont quasi-compacts, en vertu de \hyperref[VI.prop4.2]{4.2} et de
  l'hypothèse $j : U \to e$ quasi-compact, la conclusion résulte alors
  du fait que $i^{\ast}$ et $j^{\ast}$ commutent aux produits fibrés,
  et de l'hypothèse que $i^{\ast}(X)$ et $j^{\ast}(X)$ sont
  quasi-séparés.

\item Résulte de la conjonction de a) et de \hyperref[VI.cor4.3]{4.3} a).

\item La nécessité a été déjà vue (\hyperref[VI.prop4.6]{4.6} b)). La suffisance résulte du
  fait que, $E$ étant cohérent, $E'$ et $E''$ le sont (\hyperref[VI.prop4.6]{4.6} b)), de
  sorte que pour un des topos envisagés, le fait qu'il soit parfait
  signifie que la catégorie de ses objets cohérents est stable par
  $\varinjlim$ finies. On conclut donc par b).
\end{enumerate}

\begin{enonce*}[remark]{Exercice 4.8}\etiquette[4.8]{VI.exer4.8}
\begin{enumerate}
\renewcommand{\theenumi}{\alph{enumi}}
\renewcommand{\labelenumi}{\theenumi)}
\item Résoudre la question suivante (dont le rédacteur avoue à sa
  confusion ignorer la réponse) : avec les notations de
  \hyperref[VI.subsec4.1]{4.1}, si $E$
  est quasi-séparé \\ (\resp algébrique) en est-il de
  même de $E''$?

\item Supposons que $E$ soit équivalent à un topos de la forme
  $\widehat{C}$, où $C$ est une catégorie équivalente à une catégorie
  $\in 𝒰$. Prouver directement, en utilisant \hyperref[VI.exer2.17]{2.17} c) et IV 9.24, que
  si $E$ est localement cohérent (\resp cohérent, \resp algébrique,
  \resp quasi-séparé, \resp localement noethérien, \resp noethérien)
  il en est de même de $E''$.
\end{enumerate}
\end{enonce*}

\setcounter{subsection}{8}
\subsection{}\etiquette[4.9]{VI.subsec4.9}
Gardons les notations de \hyperref[VI.subsec4.1]{4.1}, et rappelons (IV 9.16) que la donnée
d'une situation $(E, U)$, formée par un topos $E$ et un ouvert $U$ de
$E$, équivaut essentiellement à celle d'un triple $(E', E'', f)$, où
$E'$ et $E''$ sont des\pageoriginale topos et où
\begin{equation*}
f : E'\to E''\tag{4.9.1}
\end{equation*}\etiquette[4.9.1]{VI.eq4.9.1}
est un « foncteur de recollement », \ie un foncteur \emph{exact à
  gauche et accessible}; partant de $(E, U)$ comme dans \hyperref[VI.subsec4.1]{4.1}, le
  foncteur de recollement associé est donné par
\begin{equation*}
  f = i^{\ast} j_{\ast}{\quoi}.\tag{4.9.2}
\end{equation*}\etiquette[4.9.2]{VI.eq4.9.2}

Nous nous proposons d'exprimer, en termes des données $E'$, $E''$,
$f$ le fait que $E$ soit un topos cohérent (\resp parfait) et que
$U$ soit un objet quasi-compact, ou ce qui revient au même, que $E$
et $E'$ soient cohérents. Nous savons déjà que ceci entraîne
que $E''$ est également cohérent (\hyperref[VI.prop4.6]{4.6} b)), de sorte que
la question revient à la suivante: étant donnés deux topos
\emph{cohérents} $E'$ et $E''$ et un foncteur de recollement $f$
\eqhyperref[VI.eq4.9.1]{4.9.1}, à quelles conditions sur $f$ le topos recollé $E$
est-il cohérent (\resp parfait)? Nous savons d'ailleurs
(\hyperref[VI.cor4.7]{4.7} c)) que $E$ est parfait si et seulement si $E$ est
cohérent, et $E'$ et $E''$ sont parfaits; donc le cas respé du
problème posé se ramène au cas non respé. Signalons
cependant que la
solution (\hyperref[VI.prop4.10]{4.10}) de ce dernier est plus jolie si on
suppose déjà $E'$ parfait.

\setcounter{subsection}{9}
\setcounter{subsubsection}{2}
\subsubsection{}\etiquette[4.9.3]{VI.subsubsec4.9.3}

Notons d'abord que si $E$ est cohérent, alors on reconstruit la
sous-catégorie pleine $E_{\coh}$ de $E$ formée des objets $X = (X',
X'', u : X''\to f(X'))$ de $E$ qui sont cohérents, comme étant la
sous-catégorie $E_{0}$ de $E$ formée des $X$ pour\pageoriginale
lesquels $X'= j^{\ast} (X)$ et $X''= i^{\ast}(X)$ sont cohérents
(\hyperref[VI.cor4.7]{4.7}
b)). Une condition nécessaire pour que $E$ soit cohérent est
donc que la
sous-catégorie $E_{0}$ précédente soit
génératrice. Cette condition
est également suffisante, car en vertu de \hyperref[VI.cor4.3]{4.3} a) $E_0$
est formée
d'objets quasi-compacts, d'autre part il est clair que $E_{0}$ est
stable par limites projectives finies, et on conclut
par 3.4.5.

Rappelons maintenant (IV 9.18) que la donnée d'un foncteur de
recollement \eqhyperref[VI.eq4.9.1]{4.9.1} équivaut à celle d'un \emph{faisceau} sur $E''$, à
valeurs dans $\Pro(E')^{0}$:
\begin{equation*}
G \in \Ffaisc (E'', \Pro(E')^{0}), G : {E''}^{0} \to \Pro(E')^{0}
\tag{4.9.4}{\quoi},
\end{equation*}\etiquette[4.9.4]{VI.eq4.9.4}
ou ce qui revient au même, à celle d'un foncteur $g = G^{0}$,
\begin{equation*}
g:E''\to \Pro (E')\tag{4.9.5}
\end{equation*}\etiquette[4.9.5]{VI.eq4.9.5}
qui \emph{commute aux limites inductives}. Si $C$ est une
sous-catégorie génératrice de $E''$, qu'on munit de la topologie
induite, on sait (II \hyperref[VI.subsec6.10]{6.10}) que la donnée de $G$ équivaut aussi (à
isomorphisme unique près) à celle d'un faisceau sur $C$ à valeurs dans
$\Pro(E')^{0}$
\begin{equation*}
G_C \in \Ffaisc (C, \Pro(E')^0), G_C : C^{\circ} \to \Pro (E')^{0}{\quoi},\tag{4.9.6}
\end{equation*}\etiquette[4.9.6]{VI.eq4.9.6}
ou, ce qui revient au même, à celle d'un foncteur $g_C = G_C^{0}$,
\begin{equation*}
g_C : C \to \Pro (E'){\quoi},\tag{4.9.7}
\end{equation*}\etiquette[4.9.7]{VI.eq4.9.7}
satisfaisant\pageoriginale aux conditions d'exactitude à gauche qu'on sait (II
6.2.1)). Bien entendu, $g_C$ n'est autre que la \emph{restriction de
$g$ $(4.9.5)$ à $C$}. Le cas le plus intéressant pour nous est celui où
on prend $C = E'_{\coh}$, d'où des objets
\begin{align*}
& G_0 \in \Ffaisc (E''_{\coh}, \Pro (E')^{0}), G_0 : {E''}_{\coh}^0 \to
 \Pro (E')^0 {\quoi},\tag{4.9.8}\\
& g_0 = G^0_0 : E''_{\coh} \to \Pro(E'){\quoi},\tag{4.9.9}
\end{align*}
dont chacun revient encore à la donnée de $f$.

\setcounter{subsubsection}{9}
\subsubsection{}\etiquette[4.9.10]{VI.subsubsec4.9.10}

Supposons la sous-catégorie pleine génératrice $C$ de $E''$ formée
d'objets quasi-compacts et qu'elle soit stable dans $E''$ par sommes
finies, par passage au quotient par des relations d'équivalence, et
par produits fibrés : c'est le cas par exemple pour $C = E''_{\coh}$
(\hyperref[VI.cor1.15]{1.15} et \hyperref[VI.cor1.17.1]{1.17.1}). Il est alors immédiat, si $P$ est une catégorie où
les limites projectives finies sont représentables (par exemple $P =
\Pro(E')^{0}$), qu'un foncteur $G_C : G^{0} \to P$ est un faisceau à
valeurs dans $P$ si et seulement si le foncteur $g_C = G_C^0 : C \to
P^0$ commute aux sommes finies et au passage au quotient par une
relation d'équivalence. Ceci précise en particulier quels sont les
faisceaux \eqhyperref[VI.eq4.9.6]{4.9.6} (exprimant donc les foncteurs de recollement $f :
E'\to E''$).

Ces rappels étant posés, nous pouvons donner la solution au problème
posé dans \hyperref[VI.subsec4.9]{4.9}:

\begin{enonce*}{Proposition 4.10}\etiquette[4.10]{VI.prop4.10}
Soient $E'$, $E''$\pageoriginale deux topos cohérents, et $f : E'\to
E''$ un foncteur de recollement \emph{(IV 9.10)}, \ie un foncteur
exact à gauche et accessible. Soit $E$ le topos qu'on en déduit par
recollement \emph{(IV 9.16)}, et désignons par $E'_{\coh}$ (\resp
$E'_{PF}$) la sous-catégorie strictement pleine de $E'$ formée des
objets cohérents (\resp des objets $X$ tels que le foncteur
covariant $\Hom(X, -)$ représenté par $X$ commute aux petites
limites inductives filtrantes). Considérons les conditions
suivantes:
\begin{itemize}
\item[(i)] $E$ est cohérent.

\item[(ii)] Pour tout objet $X'$ de $E'$, il existe une famille de
  morphismes
$$
v_{\alpha} : Y'_{\alpha} \to X'
$$
de but $X'$, à sources des objets cohérents de $E'$, telle que la
famille des
$$
f(v_{\alpha}) : f(Y'_{\alpha}) \to f(X')
$$
dans $E''$ soit couvrante.

\item[(ii bis)] $f$ commute aux (petites) limites inductives
  filtrantes, et \emph{(ii)} est vrai pour tout $X'\in \Ob E'_{PF}$.

\item[(ii ter)] Le foncteur canonique \emph{(IV 9.20.1)}
\begin{equation*}
\pi : \Ppoint(E'') \to \Pro(E')\tag{4.10.1}
\end{equation*}\etiquette[4.10.1]{VI.eq4.10.1}
défini par le foncteur de recollement $f$ se factorise (à isomorphisme
près) par $\Pro(E'_{\coh})$ (via le foncteur pleinement fidèle
$\Pro(E'_{\coh}) \to \Pro(E')$ provenant de l'inclusion $E'_{\coh}
\to E'$):
\begin{equation*}
\pi_0 : \Ppoint (E'') \to \Pro(E'_{\coh}){\quoi}.\tag{4.10.2}
\end{equation*}\etiquette[4.10.2]{VI.eq4.10.2}

\item[(iii)] Le\pageoriginale foncteur $f$ commute aux (petites)
  limites inductives
  filtrantes.

\item[(iii bis)] Le foncteur $g_0$ (4.9.9) se factorise (à
  isomorphisme près) par un foncteur
\begin{equation*}
E''_{\coh} \to \Pro(E'_{PF}){\quoi},\tag{4.10.3}
\end{equation*}\etiquette[4.10.3]{VI.eq4.10.3}
via le foncteur pleinement fidèle $\Pro(E'_{PF}) \to \Pro(E')$ déduit
de l'inclusion $E'_{PF} \to E'$.

\item[(iii ter)] Le foncteur canonique $\pi$ \eqhyperref[VI.eq4.10.1]{4.10.1} se factorise (à
  isomorphisme près) en
\begin{equation*}
\pi_1 : \Ppoint(E'') \to \Pro(E'_{PF}){\quoi}.\tag{4.10.4}
\end{equation*}\etiquette[4.10.4]{VI.eq4.10.4}

\item[(iv)] Le foncteur $g_0$ (4.9.9) se factorise (à isomorphisme près)
  par un foncteur
\begin{equation*}
E''_{\coh} \to \Pro(E'_{\coh}){\quoi}.\tag{4.10.5}
\end{equation*}\etiquette[4.10.5]{VI.eq4.10.5}
On a alors le diagramme d'implications:
\begin{equation*}
\text{\rm (iv) $\Rightarrow$ (i) $\Leftrightarrow$ (ii) $\Leftrightarrow$ (ii bis)
  $\Leftrightarrow$ (ii ter) $\Rightarrow$ (iii) $\Leftrightarrow$ (iii bis)
  $\Leftrightarrow$ (iii ter)}{\quoi}.\tag{4.10.6}
\end{equation*}\etiquette[4.10.6]{VI.eq4.10.6}
\end{itemize}
\end{enonce*}

Signalons tout de suite le

\begin{enonce*}{Corollaire 4.11}\etiquette[4.11]{VI.cor4.11}
\begin{enumerate}
\renewcommand{\theenumi}{\alph{enumi}}
\renewcommand{\labelenumi}{\theenumi)}
\item Supposons que $E'$ soit parfait. Alors toutes les conditions
  envisagées dans \hyperref[VI.prop4.10]{4.10} sont équivalentes, en particulier $E$ est
  cohérent si et seulement si $f$ commute aux limites inductives
  filtrantes, ou encore si et seulement si le foncteur $g_0$ (4.9.9)
  se factorise (à isomorphisme près)\pageoriginale par $\Pro(E'_{\coh})$.

\item Pour que $E$ soit parfait, il faut et il suffit que $E'$ et
  $E''$ le soient, et que $f$ (\resp $g_0$) satisfasse à la condition
  énoncée dans {\rm a)}.
\end{enumerate}
\end{enonce*}

En effet, a) résulte du fait que $E'$ parfait signifie $E'_{\coh} =
E'_{PF}$, de sorte que (iii bis) implique (iv). L'assertion b)
s'ensuit, comme il résulte des remarques préliminaires de
4.9.


\setcounter{subsection}{11}
\subsection{Démonstration de 4.10}\etiquette[4.12]{VI.subsec4.12}

\begin{enumerate}
\renewcommand{\theenumi}{\alph{enumi}}
\renewcommand{\labelenumi}{\theenumi)}
\item Explicitons la condition équivalente à (i) obtenue dans \hyperref[VI.subsubsec4.9.3]{4.9.3},
  savoir que la sous-catégorie pleine $E_0$ de $E$ est génératrice \ie
  que pour tout objet $X = (X', X'', u : X''\to f(X'))$ de $E$, la
  famille de tous les morphismes
$$
v = (v', v'') : Y = (Y', Y'', v : Y''\to f(Y')) \to X {\quoi},
$$
avec $Y \in \Ob E_0$ \ie $Y'$, $Y''$ cohérents, est épimorphique. Comme
le couple de foncteurs $(i^{\ast}, j^{\ast})$ est conservatif, cela
signifie aussi que la famille des morphismes correspondants $Y' \to
X'$ est épimorphique dans $E'$, et que la famille des morphismes
correspondants $Y''\to X''$ est épimorphique dans $E''$. Or c'est
clair pour la première, comme on voit en prenant des $Y \in \Ob E_0$
\emph{au-dessus de} $U$ \ie tels que $Y'' = ∅_{E''}$ : on
trouve la famille de toutes les flèches dans $E'$ de but $X'$, à
source cohérente, qui est bien épimorphique puisque $E'$ est cohérent
donc la sous-catégorie ${E'}_{\coh}$ est génératrice. Il reste donc à
exprimer que la famille des $Y''\to X''$ est épimorphique, quel que
soit l'objet donné $X$ de $E$.

\item Or supposons\pageoriginale satisfaite la condition (ii) pour l'objet $X'$
  envisagé ici. Il en résulte que l'on peut trouver une famille
  épimorphique de morphismes $Y''_{\beta} \to X''$, dont les
  composées avec $u : X''\to f(X')$ se relèvent chacun en un morphisme
  $Y''_{\beta} \to f({Y'_{\alpha}})$ pour $\alpha = \phi(\beta)$
  convenable; comme $E''_{\coh}$ est une sous-catégorie génératrice
  de $E''$, on peut même supposer les $Y''_{\beta}$ cohérents. Mais
  alors les diagrammes commutatifs
\[
\xymatrix@C=2.25cm@R=1.25cm{
Y''_{\beta}  \ar[r] \ar[d]& f(Y'_{\varphi(\beta)}) \ar[d]\\
X'' \ar[r] & f(X')
}
\]
fournissent la famille cherchée de morphismes dans $E$, donnant une
famille épimorphique $Y''_{\beta} \to X''$. Donc (ii) $\Rightarrow$ (i).

Réciproquement, supposons (i) et prouvons (ii). On applique la
condition explicitée dans a) au cas de l'objet
$$
X = j_{\ast}(X') = (X', f(X'), \id = f(X') \xrightarrow{\sim} f(X')){\quoi};
$$
comme les $Y''\to X''\simeq f(X')$ forment une famille épimorphique,
il en est de même a fortiori de la famille des $f(v') : f(Y') \to
f(X')$, d'où la conclusion, puisque les $Y'$ sont cohérents par
hypothèse. Donc
$$
{\rm(i)} \Leftrightarrow {\rm(ii)}{\quoi}.
$$

\item Notons\pageoriginale maintenant que la condition (iv) signifie
  aussi que pour   tout objet $X''$ de $F''$, le foncteur pro-représentable
$$
g(X'') : X'\to \Hom (X'', f(X'))
$$
sur $E'$ est pro-représentable par un pro-objet dont les composants
sont dans ${E'}_{\coh}$, ou ce qui revient au même, qu'il existe
dans la catégorie $E'_{/g(X'')}$, formée des couples d'un objet $X'$
de $E'$ et d'un élément $u \in g(X'') (X') =\Hom(X'', f(X'))$, $u :
X''\to f(X')$, un ensemble cofinal d'objets $(X'_{\alpha},
u_{\alpha})$, $u_{\alpha} : X'' \to f(X'_{\alpha})$, avec
$X'_{\alpha} \in \Ob E'_{\coh}$. La condition de cofinalité, comme
la catégorie $E'_{/g(X'')}$ est filtrante, signifie simplement que
pour tout objet $(X', u)$, $u : X''\to f(X')$ de $E'_{/g(X'')}$ il
existe un morphisme d'un $({X'}_i, u_i)$ dans l'objet précédent, \ie
un morphisme $v'_{\alpha} : X'_{\alpha} \to X'$ rendant commutatif
le triangle
\[\vcenter{
\xymatrix@C=2cm@R=1.5cm{
                      &f(X'_{\alpha}) \ar[d]^{f'(v'_\alpha)}\\
X''\ar@{-}[ru]^{u_{\alpha}} \ar[r]& f(X').
}}\leqno{(*)}
\]
Donc la condition (iv) équivaut à ceci:

\begin{itemize}
\item[(iv bis)] Pour tout flèche $u : X''\to f(X')$, avec $X'\in \Ob
  E'$ et $X''\in \Ob E'_{\coh}$, il existe une flèche $v : Y'\to X'$
  dans $E'$ telle que $u$ se factorise à travers $f(v') : f(X') \to
  f(X')$.

Or il est clair que la condition (iv bis) implique (ii),
  puisque pour $X'$ fixé, la famille des $u : X''\to
  f(X')$\pageoriginale de source un objet cohérent est épimorphique, de sorte qu'il en est a fortiori
  ainsi de la famille de tous les $f(v'_{\alpha})$ dans les
  diagrammes correspondants $(\ast)$. Donc
$$
{\rm(iv)} \Rightarrow {\rm(i)}{\quoi}.
$$
\end{itemize}

\item Explicitons la condition (ii ter). Soit $p$ un point de $E''$,
  alors par définition $\pi(p)$ est le pro-objet de $E'$ qui
  pro-représente le foncteur $h : X'\to f(X')_{p}$, ou encore la
  pro-objet défini par le foncteur canonique $E'_{/h} \to E'$, où
  $E'_{/h}$  désigne la catégorie des couples $(X', u)$, avec $X'\in
  \Ob E'$ et $u \in h(X')$ \ie $u \in f(X')_p$. Dire que ce pro-objet
  est isomorphe à un pro-objet provenant de $\Pro(E'_{\coh})$ signifie
  que la sous-catégorie de $E'_{/h}$ formée des $(X', u)$ avec
  $X'\in \Ob E'_{\coh}$  y est cofinale, \ie que pour tout objet $(X',
  u)$, $u \in f(X')_p$, de $E'_{/h}$, il existe un objet $(X'_0,
  u_0)$ qui le majore, avec $X'_0$ cohérent, \ie qu'il existe un
  morphisme $X'_0 \to X'$ ($X'_0$ cohérent) tel que $u$ soit
  dans $\text{Im}(f(X'_0) \to f(X'))$. Dire que ceci est vrai pour tout
  point $p$ signifie que pour tout objet $X'$ de $E'$, la famille des
$f(X'_{0})\to f(X')$, avec $X'_0 \in \Ob E'_{\coh/X'}$, est telle que pour tout point $p$ de $X''$, la famille des $f(X'_0)_p \to f(X')_p$ soit surjective. Comme on verra dans
  l'appendice (10. ) qu'un topos cohérent $E''$ a assez de points, on
  voit que la condition obtenue signifie aussi que pour tout $X'\in \Ob
  E$, la famille des $f(X'_0) \to f(X')$ est épimorphique. Mais cela
  n'est autre que la condition (ii), donc
$$
{\rm(ii)} \Leftrightarrow \text{(ii ter)}{\quoi}.
$$

\item La condition\pageoriginale que $f$ commute aux limites
  inductives filtrantes   (condition (iii)) équivaut, comme la famille des foncteurs fibres
  de $E''$ est conservative, à celle que les foncteurs composés $h :
  X'\to f(X')_p$ le sont. Prouvons que cela signifie que le pro-objet
  $\pi(p)$ qui pro-représente ce foncteur est dans l'image essentielle
  de $\Pro(E'_{PF})$: cela prouvera alors les implications
$$
{\rm(iii)} \Leftrightarrow \text{(iii ter)} \Leftarrow \text{(ii ter)}
\Rightarrow \text{(iii bis)}{\quoi}.
$$
Nous sommes donc amenés à prouver le

\begin{enonce*}{Lemme 4.12.1}\etiquette[4.12.1]{VI.lem4.12.1}
Soient $E'$ un topos tel que la sous-catégorie $E'_{\coh}$ y soit
génératrice, et $X' \in \Ob \Pro(E')$. Pour que $X'$ appartienne à
l'image essentielle de $\Pro(E'_{PF})$, il faut et il suffit que le
foncteur $h : E'\to (\Ens)$ qu'il pro-représente commute aux limites
inductives filtrantes.
\end{enonce*}

Comme une limite inductive filtrante de foncteurs commutant aux
limites inductives filtrantes possède la même propriété, la suffisance
résulte de la définition de $E'_{PF}$. Pour la nécessité, il faut
prouver que si le foncteur $h$ commute aux limites inductives
filtrantes, alors dans la catégorie filtrante $E'_{/h}$ des couples
$(X', u)$ avec $u \in h(X')$, la sous-catégorie formée des couples
$(X', u)$ avec $X'\in \Ob E'_{PF}$ est cofinale, \ie que pour tout
$(X', u)$ dans $E'_{/h}$, il existe un morphisme $X'_{0} \to
X'$, avec $X'_{0} \in \Ob E'_{PF}$, tel que $u \in \text{Im}
(h(X'_0) \to h(X'))$.\pageoriginale Or on sait que $X'$ est limite inductive
filtrante d'objets $X'_i$ de $E'_{PF}$ (\hyperref[VI.cor1.24.2]{1.24.2} a)), donc $h(X')$
est limite inductive filtrante des $h(X'_i)$, d'où la conclusion.

\item Il reste seulement à prouver les implications
$$
\text{(ii bis)} \Rightarrow \text{(ii)\quad et\quad  (iii)} \Leftrightarrow \text{(iii bis)}{\quoi}.
$$
La première implication résulte trivialement du fait que tout objet
$X'$ de $E'$ est limite inductive filtrante d'objets $X'_i$ de
$E'_{PF}$. Pour la deuxième, il suffit de noter que la famille des
foncteurs
$$
Y'' \longmapsto \Hom(X'', Y''){\quoi},
$$
pour $X'' \in \Ob {E''}_{\coh}$, est conservative et formée de
foncteurs commutant aux limites inductives filtrantes (\hyperref[VI.thm1.23]{1.23} (ii)),
donc le foncteur $f : E'\to E''$ commute aux limites inductives
filtrantes si et seulement s'il en est ainsi des foncteurs composés
$X'\to \Hom(X', f(X')) = \Hom_{\Pro(E')}(g_0 (X''), X')$, ce qui
équivaut au fait que les $g_0(X'')$ sont dans l'image essentielle
de $\Pro(E'_{PF})$, en vertu de \hyperref[VI.lem4.12.1]{4.12.1}. \cqfd
\end{enumerate}

\begin{enonce*}[remark]{Remarques 4.13}\etiquette[4.13]{VI.rem4.13}
Il convient de préciser, en langage faisceautique, la signification de
la condition \hyperref[VI.prop4.10]{4.10} (iv). Considérons le diagramme commutatif de
foncteurs
\[\vcenter{
\xymatrix@C=2.5cm@R=1.5cm{
\Pro(E'_{\coh})^{\circ}  \ar@{^{(}->}[r]^{\alpha}  \ar[d]^{\beta_0}&
\Pro(E')^{\circ} \ar[d]^{\beta}\\
\text{$\Ev{\coh}$} & \text{$\Ev{}$}\ar[l]^{\gamma} ,
}}\leqno{(4.13.1)}
\]\etiquette[4.13.1]{VI.eq4.13.1}
où\pageoriginale les signes $\vee$ désignent les catégories de
foncteurs à valeurs dans $(\Ens)$, et la deuxième flèche horizontale
$\gamma$ est la flèche de restriction des foncteurs, toutes les
autres flèches désignant les foncteurs pleinement fidèles bien
connus. Notons que la flèche d'inclusion $\alpha$ \emph{ne commute
pas en général aux limites
  projectives}, mêmes finies, donc en général un faisceau $F$ sur un
site $C$ (tel que $E''$) définit un préfaisceau $\alpha \circ F$
sur ce même site, à valeurs dans $\Pro(E'')^{0}$, qui n'est pas
nécessairement un faisceau. Néanmoins, si un \emph{préfaisceau} $F$
sur $C$ à valeurs dans $\Pro(E'_{\coh})^{0}$ est tel que le
préfaisceau $\alpha F$ correspondant à valeurs dans $\Pro(E')^{0}$ est
un faisceau, alors $F$ est lui-même un faisceau. Cela provient en
effet du fait que la condition d'être un faisceau est une propriété
d'exactitude à gauche (II 6.2  1)), et que les foncteurs $\beta$,
$\beta_0$ et $\gamma$ commutent aux petites limites projectives.
\end{enonce*}

D'autre part, il résulte de la caractérisation \hyperref[VI.subsubsec4.9.10]{4.9.10} des faisceaux
sur $E''_{coh}$ à valeurs dans une catégorie $P$ quelconque, que
pour tout foncteur $\alpha : P \to Q$ tel que $\alpha^{0} : P^0 \to
Q^0$ commute aux \emph{sommes finies} et au \pageoriginale
\emph{passage au quotient par des relations d'équivalence}, le foncteur $g \mapsto \alpha
\circ g'$ de $\Hom({E''}^0_{\coh}, P)$ dans $\Hom({E''}^0_{\coh},
Q)$ transforme faisceaux en faisceaux. D'autre part, nous savons que
dans $E'_{\coh}$ et dans $E'$ les sommes finies et les quotients
par les relations d'équivalence sont représentables, et que le
foncteur d'inclusion $E'_{\coh} \to E'$ y commute (\hyperref[VI.cor1.15]{1.15} et
\hyperref[VI.cor1.17.1]{1.17.1}). Il en est donc de même pour les catégories
$\Pro(E'_{\coh})$ et $\Pro(E')$ et le foncteur d'inclusion entre ces
catégories (I 8.9.5 b) et 8.9.7), lequel n'est autre que $\alpha^{0}$,
où $\alpha$ est le foncteur intervenant dans
\eqhyperref[VI.eq4.13.1]{4.13.1}. De ceci on
conclut que si
\begin{equation*}
\varphi : {E'}^0_{\coh} \to \Pro (E'_{\coh})^0\tag{4.13.2}
\end{equation*}\etiquette[4.13.2]{VI.eq4.13.2}
est un foncteur, alors $\varphi$ \emph{est un faisceau sur
$E''_{\coh}$ à valeurs dans $\Pro(E'_{\coh})^0$ si et seulement si
$\alpha \circ \varphi$ est un faisceau sur $E''_{\coh}$ à valeurs
dans $\Pro(E')^{0}$. Donc la catégorie des foncteurs de recollement
$f : E'\to E''$} qui satisfont à la condition \hyperref[VI.prop4.10]{4.10} (iv)
est canoniquement équivalente à la catégorie des faisceaux $\varphi$
sur $E''_{\coh}$ à valeurs dans $\Pro(E'_{\coh})^0$. On en conclut,
compte tenu de \hyperref[VI.prop4.10]{4.10}, qu'un tel faisceau $\varphi$
définit canoniquement un foncteur de recollement (à isomorphisme
unique près), et  que le topos $E$ déduit de ce dernier est
\emph{cohérent}; de plus (\hyperref[VI.cor4.11]{4.11}) \emph{si $E'$ est
parfait}, on obtient ainsi \emph{tous} les topos \emph{cohérents}
qu'on peut déduire des topos $E'$, $E''$ par recollement.

Notons\pageoriginale d'autre part que $E'_{\coh}$ est équivalente à
une petite catégorie et est stable par limites projectives finies, de sorte que
(I 8.10.14) $\Pro(E'_{\coh})^0$ est équivalente, par le foncteur
pleinement fidèle $\beta_0$ de \eqhyperref[VI.eq4.13.1]{4.13.1}, à la
sous-catégorie pleine de
$E'_{\coh} = \Hom(E'_{\coh}, (\Ens))$ formée des foncteurs qui
sont exacts à gauche. Donc la donnée d'un faisceau
\eqhyperref[VI.eq4.13.2]{4.13.2} équivaut à
celle d'un foncteur
\begin{equation*}
F : {E''}_{\coh}^0 \times E'_{\coh} \to (\Ens)\tag{4.13.3}
\end{equation*}\etiquette[4.13.3]{VI.eq4.13.3}
qui soit un faisceau par rapport au premier argument et qui soit exact
à gauche en le second argument; ou encore, à celle d'un foncteur
\begin{equation*}
f_0 : E'_{\coh} \to E''\tag{4.13.4}
\end{equation*}\etiquette[4.13.4]{VI.eq4.13.4}
qui soit exact à gauche. Si $f$ est le foncteur de recollement
satisfaisant à la condition \hyperref[VI.prop4.10]{4.10} (iv) associé à $\varphi$, on vérifie
immédiatement que $f_0$ est canoniquement isomorphe à la restriction
du foncteur $f : E'\to E''$. On voit donc que \emph{sous les
  conditions envisagées $f_0$ détermine $f$} (à isomorphisme canonique
près). Signalons enfin, pour résumer le contenu de \hyperref[VI.cor4.11]{4.11}:


\begin{enonce*}{Corollaire 4.14}\etiquette[4.14]{VI.cor4.14}
Lorsque $E'$ est un topos parfait, alors la catégorie des foncteurs de
recollement $f : E'\to E''$ qui donnent par recollement un topos
cohérent $E$ est équivalente par $f \mapsto f_0 = f / E'_{\coh}$ à
la catégorie des foncteurs $f_0 : E'_{\coh} \to E''$ qui sont
exacts à gauche, ou encore à la\pageoriginale catégorie des faisceaux
$\varphi$ sur le site $E''_{\coh}$ (ou, ce qui revient au même, sur le topos
$E''$) à valeurs dans la catégorie $\Pro(E'_{\coh})^0$ \emph{(\cf \hyperref[VI.subsubsec4.9.10]{4.9.10})}.
\end{enonce*}

On observera que la première assertion résulte également
\hyperref[VI.prop4.10]{4.10} en notant que $E' \approxeq \Ind (E'_{\coh})$
(\hyperref[VI.cor1.25]{1.25}), donc que la catégorie des foncteurs $f : E'\to
P$ commutant aux limites inductives filtrantes (où $P$ est une
catégorie où les petites limites inductives filtrantes sont
représentables) est équivalente, par le foncteur restriction, à la
catégorie des foncteurs quelconque $f : E'_{\coh} \to P$; il est
immédiat que sous ces conditions, si dans $P$ les $\varinjlim$
finies sont représentables, que $f$ est exact à gauche si et
seulement si $f_0$ l'est.

\begin{enonce*}[remark]{Remarques 4.15}\etiquette[4.15]{VI.rem4.15}
\begin{enumerate}
\renewcommand{\theenumi}{\alph{enumi}}
\renewcommand{\labelenumi}{\theenumi)}
\item La démonstration donnée de \hyperref[VI.prop4.10]{4.10} montre qu'on obtient des
  conditions équivalentes à (ii ter) \resp (iii ter) en remplaçant
  la catégorie $\Ppoint(E'')$ par une sous-catégorie qui définisse une
  famille \emph{conservative} de foncteurs fibres.

\item Prenant pour $E''$ le topos ponctuel (IV 2.2), on voit aussitôt
  que les conditions (ii ter) et (iii ter) ne sont équivalentes
  \emph{que si} $E'_{\coh} = E'_{PF}$, \ie si $E'$ est
  parfait. Donc la condition « $E$ cohérent » n'équivaut à la
  condition « $f$ commute aux limites inductives filtrantes » que
  si $E'$ est parfait. Par contre, prenant plus généralement pour
  $E''$ un topos cohérent de la forme $\widehat{C}$, on voit que
  \emph{dans ce cas la condition} (i) \emph{équivaut à} (iv) \emph{pour tout
  $E'$}. En fait, le rédacteur n'est pas arrivé à construire dans le
  cas général un exemple où $E$ soit cohérent, sans que la condition
  (iv) soit vérifiée. C'est honteux!
\end{enumerate}
\end{enonce*}

\section{Commutation des foncteurs $H^i(X, -)$ aux limites inductives
  filtrantes}\etiquette[5]{VI.sec5}\pageoriginale

\begin{enonce*}{Théorème 5.1}\etiquette[5.1]{VI.thm5.1}
Soient $E'$ un topos algébrique (\hyperref[VI.def2.3]{2.3}), $E$ un topos localement
cohérent (\hyperref[VI.def2.3]{2.3}), $f : E' \to E$ un morphisme cohérent de topos
(\hyperref[VI.def3.1]{3.1}). Alors, pour tout entier $q$, les foncteurs $R^q f_{\ast}$ \emph{(V
5.0)} commutent aux limites inductives filtrantes de faisceaux abéliens.
\end{enonce*}

\begin{enonce*}{Corollaire 5.2}\etiquette[5.2]{VI.cor5.2}
Soit $E'$ un topos cohérent (\hyperref[VI.def2.3]{2.3}). Pour tout entier $q$, le foncteur
$H^q(E', -)$ commute aux limites inductives filtrantes de faisceaux abéliens.
\end{enonce*}

\begin{enonce*}{Corollaire 5.3}\etiquette[5.3]{VI.cor5.3}
Soient $E$ un topos (\resp un topos algébrique) (\hyperref[VI.def2.3]{2.3}), $X$ un objet de
$E$ algébrique et cohérent (\resp cohérent) (\hyperref[VI.def2.3]{2.3}). Pour tout entier
$q$, le foncteur $H^q (X, -)$ commute aux limites inductives
filtrantes de faisceaux abéliens.
\end{enonce*}

Le corollaire \hyperref[VI.cor5.2]{5.2}, se déduit de \hyperref[VI.thm5.1]{5.1} en prenant pour $E$ le topos
ponctuel et pour $f : E'\to E$ l'unique morphisme (IV 2.2). Montrons
que \hyperref[VI.cor5.2]{5.2} entraîne \hyperref[VI.cor5.3]{5.3}. Soit $j_X : E_{/X} \to E$ le morphisme de
localisation. On a un isomorphisme canonique, fonctoriel en le
faisceau abélien (V 2.2.1) $H^q (X, F) \simeq H^q(E_{/X}, j^{\ast}_X
F)$. Le foncteur $j^{\ast}_X$ commute aux limites inductives et le
topos $E_{/X}$ est cohérent (\hyperref[VI.subsubsec2.4.6]{2.4.6}) d'où \hyperref[VI.cor5.3]{5.3}. Montrons que \hyperref[VI.cor5.3]{5.3}
entraîne \hyperref[VI.thm5.1]{5.1}. Soit $E_{\text{cohalg}}$ la sous-catégorie pleine de $E$
définie par les objets algébriques et cohérents de $E$. C'est une
catégorie génératrice (\hyperref[VI.subsubsec2.4.5]{2.4.5}). Pour tout faisceau abélien $F$, $R^q
f_{\ast} F$ est le faisceau associé au préfaisceau $X \mapsto
H^q(f^{\ast} X, F)$ $(X \in \ob C)$ (\hyperrefext[V][V.prop5.1]{5.1}). Comme $f$ est cohérent,
$f^{\ast}X$ est cohérent pour tout objet $X$ de $C$ (\hyperref[VI.def3.1]{3.1}). D'après \hyperref[VI.cor5.3]{5.3}
le préfaisceau $X \mapsto H^q (f^{\ast} X, F)$ commute aux
limites inductives filtrantes de l'argument $F$ d'où \hyperref[VI.thm5.1]{5.1}. en utilisant
le fait que le foncteur faisceau associé commute aux limites
inductives. (On remarquera qu'on n'utilise que la propriété suivante
du morphisme $f$: Il existe une famille génératrice $S$ de $E$ telle
que pour tout $X$ de $S$, $f^{\ast}(X)$ soit algébrique et
cohérent). Il reste donc à démontrer \hyperref[VI.cor5.2]{5.2}. On sait déjà que le foncteur
$H^{\circ} (E', -)$ commute aux limites inductives filtrantes de
faisceaux abéliens (1.2.3). D'après un résultat classique de \cite{SGA4-06:7}, il
suffit pour démontrer \hyperref[VI.cor5.2]{5.2}, de montrer qu'une limite inductive
filtrante de faisceaux acycliques pour le foncteur $H^{\circ}(E', -)$
est acyclique pour $H^{\circ} (E', -)$. Ceci résulte du lemme suivant
appliqué au cas où $C = E'_{\coh}$:

\begin{enonce*}{Lemme 5.4}\etiquette[5.4]{VI.lem5.4}
Soient $E$\pageoriginale un topos localement cohérent (\hyperref[VI.def2.3]{2.3}), $C$ une
sous-catégorie pleine de $E$, génératrice et stable par produit fibré telle que tout
objet de $C$ soit quasi-compact (\hyperref[VI.def1.1]{1.1}). Une limite inductive filtrante
de faisceaux $C$-acycliques \emph{(\hyperrefext[V][V.subsec4.2]{4.2})} est un faisceau $C$-acyclique.
\end{enonce*}

Soit $(F_i)_{i \in I}$, un système inductif filtrant de faisceaux
$C$-acycliques et notons $F$ sa limite inductive. Tout objet $Y$ de
$C$ est cohérent (\hyperref[VI.prop2.1]{2.1}) et par suite $F(Y) =
\displaystyle\mathop{\varinjlim}_i F_i (Y)$
(1.2.3). Pour montrer que $F$ est $C$-acyclique, il suffit de montrer
que $\check{H}^{q} (Y, F) = 0$ pour tout $q > 0$ et tout $Y$ de $C$
(\hyperrefext[V][V.prop4.3]{4.3}). Soient $X$ un objet de $C$, $\Xx = (X_{\alpha} \to
X)_{\alpha \in \Lambda}$ une famille couvrante finie par des objets de
$C^{\bigdot}$. On a $H^{q} (\Xx, F) = H^q (C^{\bigdot}(\Xx, F))$
(\hyperrefext[V][V.subsubsec2.4.3]{2.4.3}). D'après ce qui précède, et en utilisant le fait que
$C^{\bigdot}(\Xx, F)$ ne fait intervenir que des produits finis de
groupes (\hyperrefext[V][V.subsubsec2.3.3]{2.3.3}), on a $C^{\bigdot}(\Xx, F) =
\varinjlim_i C^{\bigdot}(\Xx, F_i)$. Donc $H^q
(\Xx, F) = H^q (C^{\bigdot}(\mathfrak{X}, F)) = H^q
(\varinjlim_i C^{\bigdot}(\mathfrak{X}, F_i)) =
\varinjlim_i C^{\bigdot}(\mathfrak{X}, F_i)) = 0$ pour $q
> 0$ (\hyperrefext[V][V.prop4.3]{4.3}). Par suite $\check{H}^q (X, F) =
\varinjlim_{\mathfrak{X}}  H^q (\mathfrak{X}, F) =
0$.

\begin{enonce*}{Corollaire 5.5}\etiquette[5.5]{VI.cor5.5}
Soient $E$ un topos cohérent (\hyperref[VI.def2.3]{2.3}), ${(X_i)_{i \in I}}$
une famille filtrante décroissante de sous-objets cohérents de
l'objet final. Notons $\Phi$ la famille des fermés de $E$ \emph{(IV
9)} contenue dans le fermé complémentaire de l'un des $X_i$. La
famille $\Phi$ est une famille de supports de $E$
\emph{(\hyperrefext[V][V.def6.12]{6.12})} et pour tout entier $q$, les foncteurs
$H^q_{\Phi} (E, -)$ et $H^q_{\Phi}$ \emph{(\hyperrefext[V][V.subsec6.13]{6.13})}
commutent aux limites inductives filtrantes.
\end{enonce*}

Soit $Z_i$ le fermé complémentaire de $X_i$. Pour tout faisceau
abélien $F$ on a une suite exacte (\hyperrefext[V][V.prop6.5]{6.5}):
$$
0 \to H^{0}_{Z_i} (E, F) \to H^{0} (E, F) \to H^{0} (X_i, F) \to
H^{1}_{Z_i} (E, F) \to \ldots\quad {\quoi}.
$$
Les foncteurs $H^q(E, F)$ et $H^q(X_i, F)$ commutent aux limites
inductives de l'argument $F$ (\hyperref[VI.cor5.2]{5.2}, \hyperref[VI.cor5.3]{5.3}). Par
suite $H^q_{Z_i} (E, F)$ commute aux limites inductives filtrantes
de l'argument $F$, d'où la propriété analogue pour les foncteurs
$H^q_{\Phi}(E, F)$, en passant à la limite inductive sur les $Z_i$
(\hyperrefext[V][V.subsec6.13]{6.13}). L'assertion concernant les foncteurs
$ℋ^q_{\Phi}$ s'en déduit en localisant et en passant au
faisceau associé.

\begin{enonce*}[remark]{Définition 5.6}\etiquette[5.6]{VI.def5.6}
Soient\pageoriginale $E$ un topos annelé d'anneau $A$ et $q$ un entier. Un
$A$-Module $G$ est dit de $q$-présentation finie s'il existe une
famille $X_i$ d'objets de $E$ couvrant l'objet final de $E$, telle que
pour tout $i$ on ait une suite exacte de Modules sur $E_{/X_i}$:
\begin{equation*}
L_q \to L_{q-1} \to \cdots \to L_1 \to L_0 \to G_{/X_i} \to 0\tag{5.5.1}
\end{equation*}\etiquette[5.5.1]{VI.eq5.5.1}
où pour tout $p$, $L_p \simeq A^{n_p}_{/X_i}$, $n_p$ entier.
\end{enonce*}

\begin{enonce*}{Proposition 5.7}\etiquette[5.7]{VI.prop5.7}
Soit $G$ un faisceau de $q$-présentation finie. Pour tout $i \leq
q-1$, le foncteur $\Eext^i_A (G, -)$ commute aux limites inductives
filtrantes de $A$-Modules.
\end{enonce*}

Le problème est local sur $E$ (\hyperrefext[V][V.prop6.1]{6.1}). On peut donc supposer, quitte à
se localiser aux $E_{/X_i}$, qu'on a une suite exacte du type
\eqhyperref[VI.eq5.5.1]{5.5.1}. Posons $L. = L_q \to \cdots L_1 \to L_0$. On a $\Eext^{i}_A(G, F)
: \Hh^{i} (\Hhom_A(L., F))$ et le foncteur $F \mapsto
\Hhom_A (L., F)$ commute aux limites inductives, d'où la proposition.

\begin{enonce*}{Corollaire 5.8}\etiquette[5.8]{VI.cor5.8}
Soient $(E, A)$ un topos annelé, $G$ un faisceau de $A$-Modules de
$q$-présentation finie, $X$ un objet algébrique et cohérent de
$E$. Les foncteurs $F \mapsto \Ext^{i}_A (X; G, F)$, tels que
$\frac{i(i-1)}{2} \leq q-1$, commutent aux limites inductives
filtrantes de l'argument $F$.
\end{enonce*}

Se déduit de \hyperref[VI.prop5.7]{5.7}, et de \hyperref[VI.cor5.3]{5.3} par la suite spectrale (V 6.1.3).

\setcounter{subsection}{8}
\subsection{}\etiquette[5.9]{VI.subsec5.9}

Soient $(E, A)$ un topos annelé, $\Phi$ une famille de supports de $E$
(\hyperrefext[V][V.def6.12]{6.12}), $G$ un $A$-module de $r$-présentation finie (\hyperref[VI.def5.6]{5.6}), $X$ un
objet de $E$ algébrique et cohérent (\hyperref[VI.def2.3]{2.3}). En passant à la limite
inductive sur les fermés de $\Phi$ dans les dernières suites
spectrales de V 6.9.1 et V 6.9.2 respectivement, on obtient deux
suites spectrales (\hyperrefext[V][V.subsec6.13]{6.13})
\begin{equation*}
\begin{cases}
'E^{pq}_2 = \varinjlim_{Z \in \Phi} \Ext^p_A (X; G, H^q_Z F)
\Rightarrow \Ext^{p+q}_{A, \Phi} (X; G, F){\quoi},\\
''E^{pq}_2 = \varinjlim_{Z \in \Phi} \Eext^p_A (G, H^q_Z F)
\Rightarrow \Eext^{p+q}_{A, \Phi} (G, F){\quoi}.
\end{cases}\tag{5.9.1}
\end{equation*}\etiquette[5.9.1]{VI.eq5.9.1}

Il résulte de \hyperref[VI.prop5.7]{5.7} et de \hyperref[VI.cor5.8]{5.8} qu'on a des
isomorphismes canoniques
\begin{equation*}
\begin{cases}
'E^{pq}_2 \simeq \Ext^p_A (X; G, H^q_\Phi F)&\quav \frac{p(p-1)}{2} \leq r - 1{\quoi},\\
''E^{pq}_2 \simeq \Eext^p_A (G, H^q_\Phi F) &\quav p \leq r-1{\quoi}.
\end{cases}\tag{5.9.2}
\end{equation*}\etiquette[5.9.2]{VI.eq5.9.2}

\begin{enonce*}{Corollaire 5.10}\etiquette[5.10]{VI.cor5.10}
On utilise\pageoriginale les hypothèses et les notations de \hyperref[VI.cor5.5]{5.5}.
Soient $A$ un Anneau de $E$ et $G$ un $A$-Module de $q$-présentation finie. Pour
tout entier $i$ tel que $\dfrac{i(i - 1)}{2} \leq r - 1$, les
foncteurs $F \mapsto \Ext^i_{A, \Phi} (E; G, F)$ et $F \mapsto
\Eext^i_{A, \Phi} (G, F)$ \emph{(\hyperrefext[V][V.subsec6.13]{6.13})} commutent aux limites
inductives filtrantes.
\end{enonce*}

Se déduit de \hyperref[VI.cor5.5]{5.5} et \hyperref[VI.prop5.7]{5.7} par les premières suites spectrales de V 6.9.1
et V 6.9.2 respectivement.

\section{Limites inductive et projective d'une catégorie fibrée}\etiquette[6]{VI.sec6}

\setcounter{subsection}{-1}
\subsection{}\etiquette[6.0]{VI.subsec6.0}

Nous supposons que le lecteur est familier avec la théorie des
catégories fibrées \cite{SGA4-06:5}. Ce numéro a essentiellement pour but de fixer
la terminologie et les notations concernant les catégories
fibrées. Toutes les catégories considérées dans ce numéro
appartiendront, sauf mention contraire, à un univers fixé
$\Uu$.

Soient $\Ff$, $\Gg$, $\Ee$ trois catégories,
$\pi : \Ff \to \Ee$ et ${\pi}' : \Gg \to
\Ee$ deux foncteurs. On désigne par
$\Hhom_{\Ee}(\Ff, \Gg)$ la catégorie dont les
  objets sont les foncteurs $u : \Ff \to \Gg$ tels que
  le diagramme
\[
\xymatrix@=1.75cm{
\Ff \ar@{-}[rr]^{u} \ar[rd]_{\pi}&             &\Gg \ar[ld]^{\pi'}\\
           & \Ee &
}\]
soit commutatif, et dont les morphismes sont les
$\Ee$-morphismes de foncteurs, \ie les morphismes de foncteurs
transformés par $\pi'$ en morphismes identiques.

Soit $\xi$ un objet de $\Ee$. Les objets $X$ de
$\Ff$ tels que $\pi (X) = \xi$ sont appelés \emph{les objets
  de $\Ff$ au-dessus de} $\xi$. Soit $f : \xi \to \eta$, un
morphisme de $\Ee$. Les morphismes $m$ de $\Ff$ tels
que $\pi(m) = f$ sont appelés les morphismes de $\Ff$
\emph{au-dessus de}\pageoriginale $f$. Soient $f : \xi \to \eta$ un
morphisme de $\Ee$ et $X$ un objet de $\Ff$ au-dessus de $\xi$, $Y$
un objet de $\Ff$ au-dessus de $\eta$; on désigne par
$\Hom_f(X, Y)$ l'ensemble des morphismes de $X$ dans $Y$ au-dessus de $f$.

\begin{enonce*}{Définition 6.1}\etiquette[6.1]{VI.def6.1}
\begin{enumerate}
\renewcommand{\labelenumi}{\theenumi)}
\item Un morphisme $m : X \to Y$ de $\Ff$ est dit \emph{cartésien} si
  pour tout morphisme $p : Z \to Y$ au-dessus de $\pi (m)$, il existe
  un unique morphisme $q : Z \to X$ au-dessus de l'identité de $\pi
  (X)$ tel que $m q = p$.

\item La catégorie $\Ff$ est dite \emph{préfibrée} par $\pi$
  au-dessus de $\Ee$ si pour tout morphisme $f : \xi \to \eta$
  de $\Ee$  tout objet de $\Ff$ au-dessus de $\eta$
  est but d'un morphisme cartésien au-dessus de $f$.

\item La catégorie $\Ff$ est dite \emph{fibrée} par $\pi$ au-dessus
  de $\Ee$ si elle est préfibrée et si le composé de deux
  morphismes cartésiens composables de $\Ff$ est un morphisme
  cartésien.
\end{enumerate}
\end{enonce*}

\setcounter{subsection}{1}
\subsubsection{}\etiquette[6.1.1]{VI.subsubsec6.1.1}

Soit $\xi$ un objet de $\Ee$. On appelle \emph{catégorie
  fibre} de $\Ff$ en $\xi$, et on désigne par
  $\Ff_{\xi}$, la sous-catégorie de $\Ff$ dont les
  objets sont les objets de $\Ff$ au-dessus de $\xi$ et les
  morphismes sont les morphismes de $\Ff$ au-dessus de
  $\id_{\xi}$. Soient deux objets $X$ et $Y$ de $F_{\xi}$, on désigne
  par $\Hom_{\xi} (X, Y)$ l'ensemble des morphismes de $X$ dans $Y$
  dans $\Ff_{\xi}$.

\subsubsection{}\etiquette[6.1.2]{VI.subsubsec6.1.2}

Soit $f : \xi \to \eta$ un morphisme de $\Ee$. Un objet de $Y$
au-dessus de $\eta$ est but d'un morphisme cartésien au-dessus de $f$,
si et seulement si le foncteur défini sur la fibre $\Ff_{\xi}$
à valeur dans les ensembles
$$
Z \mapsto \Hom_f(Z, Y)\quad Z \in \ob (\Ff_{\xi}){\quoi},
$$
est représentable. Lorsque la catégorie $\Ff$  est préfibrée,
le foncteur
$$
Y \mapsto \Hom_f(., Y)\quad Y \in \ob (\Ff_{\eta}){\quoi},
$$
à valeurs dans la catégorie des préfaisceaux d'ensembles sur
$\Ff_{\xi}$, est en fait à valeurs dans la catégorie des
foncteurs représentables sur $\Ff_{\xi}$, et définit donc, à
isomorphisme unique près, un foncteur $f^{\ast} : \Ff_{\eta} \to
\Ff_{\xi}$, qui est appelé le foncteur \emph{image\pageoriginale réciproque
   pour} $f$. Supposons que la catégorie $\Ff$
soit préfibrée et choisissons pour tout morphisme $f$ de $\Ee$ un foncteur
changement de base $f^{\ast}$. La catégorie $\Ff$ est alors
fibrée au-dessus de $\Ee$ si et seulement si pour tout couple
$f$ et $g$ de morphismes composables de $\Ee$ le foncteur
composé $f^{\ast}g^{\ast}$ est un foncteur image réciproque. Soit
alors
\begin{equation*}
C_{f, g} : f^{\ast} g^{\ast} \to (gf)^{\ast}\tag{6.1.2.1}
\end{equation*}\etiquette[6.1.2.1]{VI.eq6.1.2.1}
l'isomorphisme canonique. Les isomorphismes $C_{f,g}$ vérifient une
condition de cocycles, provenant de l'associativité de la composition
des morphismes dans $\Ee$:
\begin{equation*}
C_{f, hg}(f^{\ast} \circ C_{g, h}) = C_{gh, f}(C_{f, g} \circ h^{\ast})\tag{6.1.2.2}
\end{equation*}\etiquette[6.1.2.2]{VI.eq6.1.2.2}

\subsubsection{}\etiquette[6.1.3]{VI.subsubsec6.1.3}

Réciproquement lorsqu'on se donne pour tout objet $\xi$ de
$\Ee$ une catégorie $\Ff_{\xi}$, pour tout morphisme
$f : \xi \to \eta$ un foncteur $f^{\ast} : \Ff_\eta \to
\Ff_{\xi}$, et pour tout couple $(f, g)$ de morphismes de
$\Ee$ un isomorphisme de foncteurs $C_{f, g} : f^{\ast}
g^{\ast} \to (fg)^{\ast}$, tels que les $C_{f, g}$ vérifient la
condition \eqhyperref[VI.eq6.1.2.2]{6.1.2.2}, on peut construire de manière essentiellement
unique une catégorie $\Ff$ fibrée au-dessus de $\Ee$
dont les fibres « sont » les catégories $\Ff_{\xi}$ et dont
les foncteurs changement de base peuvent être choisis « égaux » aux
foncteurs donnés à l'avance \cite{SGA4-06:5}.

\subsubsection{}\etiquette[6.1.4]{VI.subsubsec6.1.4}
Soient $\Ff$ et $\Gg$ deux catégories au-dessus de
$\Ee$. On désigne par
$$
\Hhom_{\cart/\Ee} (\Ff, \Gg)
$$
la sous-catégorie pleine de $\Hhom_{\Ee} (\Ff,
\Gg)$ définie par les foncteurs qui transforment les
morphismes cartésiens de $\Ff$ en morphismes cartésiens de
$\Gg$. Les objets de \\ $\Hhom_{\cart/\Ee} (\Ff,
\Gg)$ sont appelés \emph{foncteurs cartésiens}.

Soient $C$ et $C'$ deux catégories et $S$ un ensemble de morphismes de
la catégorie $C$. On désigne par $\Hom_{S^{-1}} (C, C')$ l'ensemble des
foncteurs de $C$ dans $C'$ qui transforment les morphismes de $S$ en
isomorphismes.

Désignons par (Cat) la catégorie dont les objets sont les catégories
appartenant à l'univers et dont les morphismes sont les foncteurs
entre ces catégories (la catégorie (Cat) n'appartient pas à l'univers).

Soient $\Ff$  une catégorie au-dessus de $\Ee$ et $S$
l'ensemble des morphismes cartésiens de $\Ff$. La catégorie
$\Ff$ définit deux foncteurs sur (Cat) à valeur dans la
catégorie\pageoriginale des ensembles appartenant à l'univers:
$$
C \mapsto \Hom_{S^{-1}} (\Ff, C)\quad C \in \ob (\Cat)
$$

$C \mapsto \Hom_{\Cart /\Ee} (\Ff, C  \times
\Ee) =$ \{Foncteurs cartésiens de $\Ff$ dans $C \times
  \Ee$\}

(la catégorie $C\times \Ee$ est considérée comme une catégorie
au-dessus de $\Ee$ par le foncteur deuxième projection).

\begin{enonce*}{Proposition 6.2}\etiquette[6.2]{VI.prop6.2}
Les foncteurs $(\Cat) \to (\Ens) :$
\begin{align*}
C &\mapsto \Hom_{S^{-1}} (\Ff, C)\\
C &\mapsto \Hom_{\Cart/\Ee} (\Ff, C \times \Ee)
\end{align*}
sont canoniquement isomorphes. Ils sont représentables.
\end{enonce*}

\begin{enonce*}[remark]{Preuve}
Pour prouver la première assertion, il
suffit de remarquer que les morphismes cartésiens de $C \times
\Ee$ sont les morphismes de la forme $m \times f$, où $m$ est
un isomorphisme de $C$. Pour prouver la seconde assertion, il suffit
de prouver que le foncteur $\Hom_{S^{-1}} (\Ff, .)$ est
représentable. Ceci résulte de \cite{SGA4-06:1}. Indiquons simplement l'idée de la
démonstration. On adjoint formellement aux morphismes de $\Ff$
les inverses des morphismes de $S$. On considère la catégorie libre
engendrée par les objets de $\Ff$, les morphismes de
$\Ff$ et les inverses formels des morphismes de $S$ (catégorie des chemins). On passe au quotient
par les relations provenant des relations entre morphismes de
$\Ff$ et les relations du type:
$$
s^{-1} s = \id\quav \text{ss}^{-1} = \id\quad  s \in S{\quoi}.
$$
La catégorie ainsi obtenue est notée $\Ff(S^{-1})$. La catégorie $\Ff(S^{-1})$ munie du foncteur
canonique $Q : \Ff \to \Ff(S^{-1})$ représente le foncteur
$\Hom_{S^{-1}} (\Ff, .)$.
\end{enonce*}

\begin{enonce*}{Définition 6.3}\etiquette[6.3]{VI.def6.3}
Lorsque $\Ff$ est fibrée au-dessus de $\Ee$ le
foncteur $\Hom_{S^{-1}} (\Ff, .)$ est noté
$\varinjlim_{\mathcal{E}^{\circ}} \Ff$
et est appelé le foncteur limite inductive de $\Ff$ au-dessus
de $\mathcal{E}^{\circ}$. La catégorie qui le représente est encore
notée $\varinjlim_{\mathcal{E}^{\circ}}
\Ff$ et est appelée la catégorie limite inductive de
$\Ff$ au-dessus de $\Ee^{\circ}$.
\end{enonce*}

\setcounter{subsection}{4}
\setcounter{subsubsection}{-1}
\subsubsection{}\etiquette[6.4.0]{VI.subsubsec6.4.0}

Supposons $\Ff$ fibrée au-dessus de $\Ee$; choisissons
pour tout morphisme $f$ de $\Ee$ un foncteur changement de base
au-dessus de $f$ et soit
$$
Q :\Ff \to \varinjlim_{\mathcal{E}^{\circ}} \Ff
$$
le\pageoriginale foncteur canonique. Les foncteurs d'inclusions des fibres de
$\Ff$ dans $\Ff$ composés avec le foncteur $Q$
fournissent, pour tout objet $\xi$ de $\Ee$, un foncteur
$$
u_{\xi} : \Ff_{\xi} \to \varinjlim_{\mathcal{E}} \Ff
$$
et pour tout morphisme $f : \xi \to \eta$ de $\Ee$ un
diagramme commutatif à isomorphisme canonique près:
\[
\xymatrix@R=1cm@C=3cm{
\Ff_{\eta} \ar[rd]^{u_{\eta}} \ar[dd]_{f^\ast}&\\
                   &\displaystyle\mathop{\varinjlim}_{\Ee^{0}}\Ff\\
\Ff_{\xi} \ar[ru]_{u_{\xi}} &
}
\]
La catégorie $\varinjlim_{\mathcal{E}^{\circ}} \Ff$ apparaît alors
comme la limite inductive \emph{au sens des pseudo-foncteurs}
\cite{SGA4-06:5} du pseudo-foncteur $\Ee^{\circ} \to \Cat$ qui associe à
tout $\xi \in \ob$, $\Ff_{\xi}$ et à tout $f : \xi \to \eta$, le
foncteur $f^{\ast} : \Ff_{\eta} \to \Ff_{\xi}$. On notera toute-fois
que même lorsque $\xi \mapsto \Ff_{\xi}$ est un véritable foncteur,
la catégorie $\varinjlim_{\mathcal{E}^{\circ}} \Ff$ n'est pas en
général la limite inductive au sens de I 2 du foncteur $\xi \mapsto
\Ff_{\xi}$ (\cf \hyperref[VI.exe6.8]{6.8}).

\begin{enonce*}{Proposition 6.4}\etiquette[6.4]{VI.prop6.4}
Soit $\Ff$ une catégorie fibrée au-dessus de $\Ee$. On
suppose que la catégorie $\Ee$ possède les propriétés
suivantes:
\begin{itemize}
\item[L1)] Tout diagramme
\[
\xymatrix@C=1.5cm@R=.65cm{
\bigdot \ar[rd]& \\
        & \bigdot\\
\bigdot \ar[ru]&
}
\]
s'insère dans un diagramme commutatif:
\[
\xymatrix@C=1.25cm@R=1cm{
        & \bigdot \ar[rd]&\\
\bigdot \ar[ru]\ar[rd]&         & \bigdot\\
        & \bigdot \ar[ru]&
}
\]

\item [L2)] Pour\pageoriginale tout couple de morphisme $u, v : \cdot
  \rightrightarrows \cdot$ tel qu'il existe un morphisme $t$ vérifiant
  la relation $tu = tv$, il existe un morphisme $w$ vérifiant la
  relation $uw = vw$.

(Notons que si la catégorie $\Ee^{\circ}$ est pseudo-filtrante
\emph{(I 2.7)}, elle possède les propriétés \emph{L1)} et
\emph{L2)}). L'ensemble $S$ des
morphismes cartésiens de $\Ff$ possède alors les propriétés:

\item[Fr1)] L'ensemble $S$ est stable par composition. Les
  isomorphismes appartiennent à~$S$.

\item[Fr2)] Tout diagramme
\[
\xymatrix@C=1.5cm@R=.65cm{
\bigdot \ar[rd]& \\
        & \bigdot\\
\bigdot \ar[ru]^{s}&,
}
\]
où $s$ appartient à $S$, peut se compléter en un diagramme commutatif
\[
\xymatrix@C=1.25cm@R=1cm{
        & \bigdot \ar[rd]&\\
\bigdot \ar[ru]^{t}\ar[rd]&         & \bigdot\\
        & \bigdot \ar[ru]^{s}&
}
\]
où $t$ appartient à $S$.

\item[Fr3)] Pour tout couple de morphismes $u, v : \cdot
  \rightrightarrows \cdot$ tel qu'il existe un morphisme $s \in S$
  vérifiant la relation $su = sv$, il existe un morphisme $t \in S$
  tel que $ut = vt$.
\end{itemize}
\end{enonce*}

\begin{enonce*}[remark]{Preuve}
Laissée au lecteur à titre d'exercice.
\end{enonce*}

\begin{enonce*}{Proposition 6.5}\etiquette[6.5]{VI.prop6.5}
Soient $\Ff$\pageoriginale une catégorie, $S$ un ensemble de
morphismes de $\Ff$ possédant les propriétés \emph{Fr1), Fr2)} et \emph{Fr3)
  (\hyperref[VI.prop6.4]{6.4})}. Pour tout objet $X$ de $\Ff$, désignons par $S(X)$ la catégorie des
morphismes de $S$ de but $X$. La catégorie $S(X)$ est cofiltrante \emph{(I
2.7)}. La catégorie $\Ff(S^{-1})$ (qui représente le foncteur
$\Hom_{S^{-1}} (\Ff, \cdot)$) peut alors se décrire comme suit:
\begin{enumerate}
\renewcommand{\theenumi}{\alph{enumi}}
\renewcommand{\labelenumi}{\theenumi)}
\item Les objets de $\Ff (S^{-1})$ sont les objets de
  $\Ff$.

\item Soient $X$ et $Y$ deux objets de $\Ff(S^{-1})$. On a
$$
\Hom_{\Ff(S^{-1})} (X, Y)= \varinjlim_{S(X)}
\Hom_{\Ff} (., Y){\quoi}.
$$

\item Soient $f : X \to Y$ et $g : Y \to Z$ deux morphismes de
  $\Ff(S^{-1})$.
Soit
\[
\xymatrix{
\bigdot  \ar[d]_{t} \ar[rrd]^{f'}&\\
X        & &Y & & t\in S
}
\]
(resp.
\[
\xymatrix{
\bigdot \ar[d]_{t'} \ar[rrd]^{g'}&\\
Y                          & & Z & & t'\in S)
}
\]
un diagramme dans $\Ff$ dont l'image est $f$ (\resp $g$). Soit
\[
\xymatrix@C=2.5cm@R=1cm{
\bigdot \ar[rd]^{f''} \ar[d]_{t''}& \\
\bigdot \ar[rd]^{f'}        & \bigdot \ar[d]^{t'}\\
                            & Y      & t''\in S
}
\]
un diagramme commutatif dont l'existence est assurée par
Fr2). L'image dans $\Hom_{\Ff(S^{-1})} (X, Z)$ du diagramme
\[
\xymatrix@C=2cm{
\bigdot \ar[d]_{tt''} \ar[rd]^{g'f''}& \\
X       & Z
}
\]
est le morphisme composé $gf$.

Soit\pageoriginale
$$
Q : \Ff \to \Ff (S^{-1})
$$
le foncteur canonique.

\item Le foncteur $Q^{\ast} : \Ff(S^{-1})^{\widehat{\phantom{,}}} \to
  \Ff^{\widehat{\phantom{,}}}$ ($F \mapsto F \circ Q$ \emph{I
  5.0}) est pleinement
  fidèle et injectif sur les objets.

\item Le foncteur $Q_{!} : \Ff^{\widehat{\phantom{,}}} \to
  \Ff(S^{-1})^{\widehat{\phantom{,}}}$ (adjoint à gauche au foncteur
  $Q^{\ast}$ \emph{(I 5.1)} est exact à gauche. (Il en est donc de même du
  foncteur $Q$ lorsque dans $\Ff$ les limites projectives
  finies sont représentables).
\end{enumerate}
\end{enonce*}

\begin{enonce*}[remark]{Preuve}
Nous renvoyons pour la preuve à [1].
\end{enonce*}

\begin{enonce*}[remark]{Exercice 6.6}\etiquette[6.6]{VI.exe6.6}
Soient $\Ff$ une catégorie, $S$ un ensemble de morphismes de
$\Ff$ possédant les propriétés Fr1), Fr2) et Fr3). Pour tout
objet $X$ de $\Ff$, on note $S(X)$ la catégorie des morphismes
de but $X$.
\begin{enumerate}
\renewcommand{\theenumi}{\alph{enumi}}
\renewcommand{\labelenumi}{\theenumi)}
\item La catégorie $S(X)$ est cofiltrante (I 2)

\item Pour tout objet $X$ de $\Ff$, désignons par $J_S(X)$
  l'ensemble des cribles de $X$ qui contiennent un crible engendré par
  un morphisme de $S$. Les $J_S(X)$ définissent sur $\Ff$ une
  topologie $T$.

\item Pour la topologie $T$, les morphismes de $S$ sont bicouvrants
  (II \hyperref[VI.cor5.2]{5.2}).

\item Pour tout objet $X$ de $\Ff$, désignons par $𝒥^2(X)$ la
  catégorie des morphismes bicouvrants de but $X$. La catégorie $S(X)$
  est cofinale dans $𝒥^2(X)$ (I 8).

\item Soit $G$ un préfaisceau d'ensemble sur $\Ff$. Soit
  $\sheaf{a}$ le   foncteur associé pour la topologie $T$. Pour tout objet $X$ de
  $\Ff$, on a un isomorphisme canonique
$$
\sheaf{a} G(X) \xrightarrow{\sim} \frac{\lim}{S(X)} G( \cdot ){\quoi}.
$$
Un préfaisceau est un faisceau si et seulement s'il transforme tout
morphisme de $S$ en isomorphisme.

\item On choisit le foncteur faisceau associé de façon que sa
  restriction à $\Ff$ soit injective sur les
  objets. Désignons par $\underrightarrow{\Ff}$ la sous-catégorie
  pleine de la catégorie des faisceaux d'ensembles sur $\Ff$
  définie par les faisceaux associés aux préfaisceaux
  représentés. \pageoriginale La   catégorie $\underrightarrow{\Ff}$
   est isomorphe à la catégorie $\Ff(S^{-1})$ décrite dans \hyperref[VI.prop6.5]{6.5}. Les objets de
  $\underrightarrow{\Ff}$ forment une famille de générateurs du
  topos $\Ff^{\sim}$ des faisceaux d'ensembles sur
  $\Ff$. La topologie induite sur $\underrightarrow{\Ff}$
  par la topologie canonique de $\Ff^{\sim}$ est la topologie
  grossière (II 1.1.4).

\item Soit $Q : \Ff \to \underrightarrow{\Ff}$ le foncteur
  canonique. Le foncteur $Q$ est un morphisme du site
  $\underrightarrow{\Ff}$ (topologie grossière) dans le site
  $\Ff$ (topologie $T$ de b)) (IV 5.9) et induit un
  isomorphisme sur les topos correspondants.

\item Soit $u : \Ff \to C$ un foncteur. Le foncteur $u$ est
  continu, si et seulement s'il transforme les morphismes de $S$ en
  isomorphismes (III 1.1)

\item Soit $u : \Ff \to C$ un foncteur transformant les
  morphismes de $S$ en isomorphismes. Le foncteur $u$ se factorise
  d'une manière unique en
\[
\xymatrix@=1.5cm{
  \Ff\ar[rr]^{u} \ar[rd]_{Q}&  & C\\
            & \underrightarrow{\Ff}{} \ar[ru]_{v}&
}
\]
(On utilisera III 1.4).

\item Soit $\Pro_S \Ff$ la sous-catégorie pleine de
  $\Pro \Ff$ (I 8) définie par les proobjets dont les
  morphismes de transition sont dans $S$. Montrer que le foncteur $Q$
  se factorise en
\[
\xymatrix@M=.25cm@C=2cm{
\Ff\ar@{^{(}->}[r]^{i}& \Pro_S\Ff\ar[r]^{j}&
\Ff_{S^{-1}}
}
\]
où $j$ est la restriction à $\Pro_S \Ff$ de $\Pro Q$. Montrer que le
foncteur $j$ admet un  adjoint à gauche pleinement fidèle
$$
A : \Ff_{S^{-1}} \to \Pro_S \Ff{\quoi}.
$$
Montrer que pour tout objet $X$ de $\Ff$, $AQ(X)$ est le
pro-objet
$$
Y \to X \in S(X) \mapsto Y{\quoi}.
$$
\end{enumerate}
\end{enonce*}

\begin{enonce*}{Proposition 6.7}\etiquette[6.7]{VI.prop6.7}
Soient\pageoriginale $\Ee$ une catégorie cofiltrante (I 8.7)
et $\Ff \to \Ee$ une catégorie fibrée \emph{(\hyperref[VI.subsubsec6.1.3]{6.1.3})}.
\begin{enumerate}
\renewcommand{\labelenumi}{\theenumi)}
\item Soient $\xi$ un objet de $\Ee$ et $\Ff/\xi$ la
  catégorie fibrée sur $\Ee/\xi$ obtenue par le changement de
  base $\Ee/\xi \to \Ee$. Le foncteur canonique
  $\varinjlim_{\Ee/\xi} \Ff /
  \xi \to \varinjlim_{\Ee}
  \Ff $ est une équivalence de catégorie.

\item Soit plus généralement $\Ee' \to \Ee$ un
  foncteur cofinal \emph{(I 8)}. Soit $\Ff'\to \Ee'$ la
  catégorie fibrée déduite de $\Ff \to \Ee$ par le
  changement de base $\Ee'\hookrightarrow \Ee$. Le foncteur
  canonique:
$$
\varinjlim_{\Ee'} \Ff' \to \varinjlim_{\Ee} \Ff
$$
est une équivalence de catégories.
\end{enumerate}
\end{enonce*}

\begin{enonce*}[remark]{Preuve}
Exercice.
\end{enonce*}

\begin{enonce*}[remark]{Exercice 6.8}\etiquette[6.8]{VI.exe6.8}
\begin{enumerate}
\renewcommand{\labelenumi}{\theenumi)}
\item Soient $g : {\Ee}^{\circ} \to \Cat$ un foncteur et
  $\Gg \to \Ee$ la catégorie fibrée correspondante
  \cite{SGA4-06:5}. Montrer qu'il existe un foncteur canonique
  $\varinjlim_{\Ee^{\circ}} \Gg \to
  \varinjlim_{\Ee^{\circ}} g$. Montrer que ce foncteur
  n'est pas nécessairement une équivalence de catégories. (On pourra
  prendre pour $g$ le foncteur constant dont la valeur est l'objet
  final de Cat et pour $\Ee$ la catégorie associée à un groupe).

\item On suppose que la catégorie $\Ee^{\circ}$ est
  pseudo-filtrante. Montrer qu'alors le foncteur canonique
  $\varinjlim_{\Ee^{\circ}} \Gg \to
  \varinjlim_{\Ee^{\circ}} g$ est une équivalence de
  catégories.
\end{enumerate}
\end{enonce*}

\begin{enonce*}{Proposition 6.9}\etiquette[6.9]{VI.prop6.9}
Soit $\Ff \to \Ee$ une catégorie fibrée (\hyperref[VI.def6.1]{6.1}). Le
foncteur $(\Cat) \to (\Ens)$:
$$
C \mapsto \Hom_{\Cart/\Ee} (C \times \Ee, \Ff)
$$
est représentable par la catégorie $\Hhom_{\Cart/\Ee}
(\Ee, \Ff)$.
\end{enonce*}

Soient $C$ une catégorie, $F : C \times \Ee \to \Ff$
un $\Ee$-foncteur cartésien (\hyperref[VI.subsubsec6.1.4]{6.1.4}), $X$ un objet de $C$. Le
foncteur $\xi \mapsto F((X, \xi))$ est un foncteur cartésien noté
$F'(X)$ de $\Ee$ dans $\Ff$ qui dépend
fonctoriellement de $X$. D'où une application $F \mapsto F'$ de
$\Hom_{\Cart/\Ee} (C \Ee, \Ff)$ dans $\Hom(C,
\Hhom_{\Cart/\Ee} (\Ee, \Ff))$ fonctorielle en
$C$ dont on vérifie immédiatement que c'est une bijection.

\setcounter{subsection}{9}
\subsection{}\etiquette[6.10]{VI.subsec6.10}

La catégorie\pageoriginale $\Hhom_{\Cart/} (\Ee, \Ff)$
est appelée la catégorie des \emph{sections cartésiennes} de la catégorie fibrée
$\Ff$. Elle est aussi appelée parfois la catégorie
\emph{limite projective} de $\Ff$ suivant
$\Ee^{\circ}$. Elle est alors notée
$\varprojlim_{\Ee^{\circ}} \Ff$ \cite{SGA4-06:2}.

\subsection{}\etiquette[6.11]{VI.subsec6.11}

Choisissons un scindage de la catégorie fibrée $\Ff \to
\Ee$, \ie choisissons pour toute flèche $f : \xi \to \eta$ de
$\Ee$ un foncteur changement de base $f^{\ast} :
\Ff_{\eta} \to \Ff_{\xi}$. Pour tout objet $\xi$ de
$\Ee$, notons
$$
v_{\xi} : \varprojlim_{\Ee^{\circ}} \Ff \to \Ff_{\xi}
$$
le foncteur d'évaluation en $\xi$. Pour tout morphisme $f : \xi \to
\eta$, on a un diagramme de foncteurs commutatif à isomorphisme
canonique près:
\[
\xymatrix@C=2.25cm@R=1.25cm{
\displaystyle\mathop{\varprojlim}_{\Ee^{\circ}} \Ff
\ar@{-}[r]^{v_{\xi}} \ar[rd]^{v_\eta}& \Ff_{\xi} \ar[d]^{f^\ast}\\
                                 & \Ff_\eta.
}
\]
La catégorie $\varprojlim_{\Ee^{\circ}}
\Ff$ apparaît ainsi comme la limite projective au \emph{sens
  des pseudo foncteurs} de $\xi \mapsto \Ff_{\xi}$. Même lorsque
$\xi \mapsto \Ff_{\xi}$ est un véritable foncteur (\ie même
lorsque le scindage choisi est un clivage \cite{SGA4-06:5}) les catégories
$\varprojlim_{\Ee^{\circ}}\Ff$
(\hyperref[VI.subsec6.10]{6.10}) et
$\varprojlim_{\Ee^{\circ}}\Ff_{\xi}$
(I 2) ne sont pas, en général, équivalentes.

\section{Topos et sites fibrés}\etiquette[7]{VI.sec7}

\subsection{Topos fibrés}\etiquette[7.1]{VI.subsec7.1}

\begin{enonce*}{Définition 7.1.1}\etiquette[7.1.1]{VI.def7.1.1}
Soit $𝒰$ un univers. Un $𝒰$-topos fibré sur une
catégorie $I$ est une catégorie fibrée sur $I$ (\hyperref[VI.def6.1]{6.1}):
$$
p : F \to I
$$
dont les fibres sont des $𝒰$-topos, et dont les foncteurs images
inverses\footnote{pour la structure fibrée (\hyperref[VI.subsubsec6.1.2]{6.1.2}).} relatifs aux
morphismes $f : i \to j$ de $I$ sont des foncteurs $f^{\ast} : F_j \to
F_i$ images inverses de morphismes de topos $f\bigdot : F_i \to F_j$
\emph{(IV 3.1.2)}.
\end{enonce*}

\setcounter{subsubsection}{1}
\subsubsection{}\etiquette[7.1.2]{VI.subsubsec7.1.2}

Il revient au même de dire qu'un $𝒰$-topos fibré sur une
catégorie $I$ est une catégorie fibrée $p : F \to I$ dont les fibres
sont des $𝒰$-topos et dont les foncteurs\pageoriginale
images inverses sont exacts et commutent aux limites inductives (IV
1.6).


\subsubsection{}\etiquette[7.1.3]{VI.subsubsec7.1.3}

Choisissons pour tout morphisme $f : i \to j$ de $I$ un foncteur
image inverse $f^{\ast} : F_j \to F_i$. On a alors, pour tout couple
$i \xrightarrow{f} j \xrightarrow{g} k$ de morphismes composables de
$I$, un isomorphisme canonique $c_{f, g} : f^{\ast} g^{\ast} \to
(gf)^{\ast}$ \eqhyperref[VI.eq6.1.2.1]{6.1.2.1}  et les $c_{f, g}$ possèdent la
propriété de cocycle \eqhyperref[VI.eq6.1.2.2]{6.1.2.2}. Choisissons de plus, pour tout
$f : i \to j$ un foncteur adjoint à droite $f_{\ast} : F_{i} \to
F_j$ au foncteur $f^{\ast} : F_j \to F_i$. Le foncteur $f_{\ast} :
F_i \to F_j$ est défini à isomorphisme canonique près par sa
propriété d'être adjoint à droite à $f^{\ast}$. On en déduit, par
les résultats généraux sur les foncteurs adjoints, des isomorphismes
canoniques;
\begin{equation*}
{c'}_{g, f} : g_{\ast} f_{\ast} \to (gf)_{\ast}\tag{7.1.3.1}{\quoi}.
\end{equation*}\etiquette[7.1.3.1]{VI.eq7.1.3.1}
qui possèdent une propriété de cocycle analogue à la propriété
6.2.2.2
et qu'on laisse au lecteur le soin d'expliciter. En appliquant alors
les résultats de \hyperref[VI.subsubsec6.1.3]{6.1.3}, on obtient une catégorie fibrée sur
$I^{\circ}$:
\begin{equation*}
p': F'\to I^{\circ}{\quoi}.\tag{7.1.3.2}
\end{equation*}\etiquette[7.1.3.2]{VI.eq7.1.3.2}
On vérifie facilement que la catégorie fibrée $p': F'\to I^{\circ}$ ne
dépend pas à $I^{\circ}$-isomorphisme unique près des différents choix
utilisés pour la construire\footnote{\etiquette[112]{page112}Le rédacteur présente ses excuses
  pour le caractère non intrinsèque de cette construction.}. La fibre
en tout objet $i$ de $I^{\circ}$ de la catégorie fibrée $p': F'\to
I^{\circ}$ est canoniquement isomorphe au $𝒰$-topos $F_i$ et est
identifiée à ce dernier. Les foncteurs images inverses de la
catégorie fibrée $p': F'\to I^{\circ}$ sont des foncteurs
\emph{images directes} par des morphismes de topos.

\subsubsection{}\etiquette[7.1.4]{VI.subsubsec7.1.4}

Soit $p : F \to I$ un $𝒰$-topos fibré. Choisissons pour
chaque $f : i \to j$ un morphisme de topos (IV 3.1.2) $f\bigdot :
F_i \to F_j$ tel que le foncteur image inverse $f^{\ast}$ (par
$f\bigdot$) soit l'image inverse pour la structure fibrée
(\hyperref[VI.def6.1]{6.1}). On obtient alors des cocycles $c_{f, g}$ reliant
ces morphismes de  topos \eqhyperref[VI.eq7.1.3.1]{7.1.3.1} et \eqhyperref[VI.eq7.1.3.2]{7.1.3.2}. La
collection des morphismes $f.$, $f \in \text{Fl}(I)$, et des
cocycles $c_{f, g}$ est appelée un \emph{biscindage} du topos fibré
$F \to I$. Le choix d'un biscindage permet d'associer au topos fibré
$F \to I$ un pseudo-foncteur $i \mapsto F_i$ de la catégorie $I$
dans la 2-catégorie des\pageoriginale $𝒰$-topos. On peut
alors lui associer de manière essentiellement unique (\cf \cite{SGA4-06:5})
un $𝒰$-topos fibré $F \to I$ muni d'un biscindage tel que le
pseudo-foncteur correspondant soit canoniquement isomorphe à $i
\mapsto F_i$. Dans la pratique on notera le plus souvent un
$𝒰$-topos fibré $F \to I$ par la notation $(F_i)_{i \in I}$
en supposant implicitement qu'on a choisi un biscindage. On notera
cependant que les constructions qu'on effectue sur les topos fibrés
(telles que le topos total associé (\hyperref[VI.subsec7.4]{7.4}), la limite
projective (\hyperref[VI.subsec8.1]{8.1})) ne dépendent que des topos fibrés
et non pas des biscindages éventuellement choisis.

\begin{enonce*}{Définition 7.1.5}\etiquette[7.1.5]{VI.def7.1.5}
Soient $p : F \to I$ et $q : G \to I$ deux $𝒰$-topos fibrés sur
$I$. Un morphisme $m$ de topos fibrés de $(F, p)$ dans $(G, q)$
consiste en la donnée d'un $I$-foncteur $m^{\ast} : G \to F$ \emph{(\hyperref[VI.subsubsec6.1.4]{6.1.4})}
et en la donnée pour tout $i$ objet de $I$ d'un adjoint à droite
$m_{i\ast} : F_i \to G_i$ au foncteur $m^{\ast}_i : G_i \to F_i$
obtenu en restreignant aux fibres en $i$ le foncteur $m^{\ast}$. Cette
donnée doit être de plus soumise à la condition suivante: Pour tout
objet $i$ de $I$ le couple de foncteurs adjoints $(m^{\ast}_i,
m_{i\ast})$ est un morphisme du topos $F_i$ dans le topos $G_i$.
\end{enonce*}

\setcounter{subsubsection}{5}
\subsubsection{}\etiquette[7.1.6]{VI.subsubsec7.1.6}

Soient $m : (F, p) \to (G, q)$ un morphisme de $𝒰$-topos
fibrés sur $I$ et $p': F' \to I^{\circ}$, $q' : G'\to I^{\circ}$ les
catégories fibrées associées par le procédé \hyperref[VI.subsubsec7.1.3]{7.1.3}
aux topos fibrés $(F, p)$ et $(G, q)$ respectivement. Choisissons
des biscindages de $F$ et $G$ (\hyperref[VI.subsubsec7.1.4]{7.1.4}), notés dans
les deux cas $(f : i \to j) \mapsto f. : (f^{\ast}, f_{\ast})$, de
$F_i$ dans $F_j$ et de $G_i$ dans $G_j$ respectivement. Comme
$f^{\ast} : G \to F$ est un $I$-foncteur, on a, pour tout morphisme
$f : i \to j$, un morphisme de foncteurs (morphisme de transition)
\begin{equation*}
b_f : f^{\ast}m^{\ast}_j \to m^{\ast}_i f^{\ast}\tag{7.1.6.1}{\quoi},
\end{equation*}\etiquette[7.1.6.1]{VI.eq7.1.6.1}
et pour tout couple $i \xrightarrow{f} j \xrightarrow{g} k$ de
morphismes composables, on a
\begin{equation*}
m^{\ast}_i (c_{f, g}) b^{\ast}_f f^{\ast} (b_g) = b_{gf} c_{f, g}\tag{7.1.6.2}{\quoi}.
\end{equation*}\etiquette[7.1.6.2]{VI.eq7.1.6.2}
En passant aux foncteurs adjoints, on obtient des morphismes
\begin{equation*}
b'_f : f_{\ast}m_{i \ast} \to m_{j\ast} f_{\ast}\tag{7.1.6.3}
\end{equation*}\etiquette[7.1.6.3]{VI.eq7.1.6.3}
qui satisfont à des relations analogues aux relations \eqhyperref[VI.eq7.1.6.2]{7.1.6.2}. Comme
les foncteurs $f_{\ast}$ sont\pageoriginale des foncteurs images inverses des
catégories fibrées $(F', p')$ et $(G', q')$ (\hyperref[VI.subsubsec7.1.3]{7.1.3}), on peut
construire un $I^{\circ}$-foncteur:
\begin{equation*}
m_{\ast} : F'\to G'\tag{7.1.6.4}{\quoi},
\end{equation*}\etiquette[7.1.6.4]{VI.eq7.1.6.4}
dont les restrictions aux fibres sont les foncteurs $m_{i\ast}$ et
dont les isomorphismes de transition sont les $b'_f$. On vérifie
avec un peu de patience que le foncteur $m_{\ast}$ ainsi construit ne
dépend pas du choix des morphismes de topos $f$.\footnote{voir note du
  bas de la page \pagehyperref[page112]{112}}.

\subsubsection{}\etiquette[7.1.7]{VI.subsubsec7.1.7}

On note $\Homtop_I (F, G)$ l'ensemble des morphismes de topos fibrés
entre le topos fibré $(F, p)$ et le topos fibré $(G, q)$. L'ensemble
$\Homtop_I(F, G)$ est l'ensemble des objets d'une catégorie notée
$\Hhomtop_I(F, G)$: Si $m$ et $n$ sont deux morphismes de topos
fibrés, un morphisme de $m$ dans $n$ est un $I$-morphisme du
foncteur $n^{\ast}$ dans le foncteur $m^{\ast}$. On en déduit
d'ailleurs par adjonction un $I^{\circ}$-morphisme du foncteur
$m_{\ast}$ dans le foncteur $n_{\ast}$ (\hyperref[VI.subsubsec7.1.6]{7.1.6}). On
a donc en définitive deux foncteurs
\begin{align*}
\Hhomtop_I (F, G) & \to \Hhom_I (G, F)^{\circ}\tag{7.1.7.1}\\
\Hhomtop_I (F, G) &\to \Hhom_{I^{\circ}} (F', G')\tag{7.1.7.2}
\end{align*}
qui associent à tout morphisme $(m^{\ast}, (m_{i\ast})_{i \in \Ob I})$,
d'une part le foncteur $m^{\ast} \in \Hom_I(G, F)$, et d'autre part le
foncteur $m_{\ast} \in \Hom_{I^\circ} (F', G')$ obtenu en recollant les
$m_{i\ast}$ comme en (\hyperref[VI.subsubsec7.1.6]{7.1.6}). Le foncteur
$m^{\ast} : G \to F$ (7.1.7.1) est appelé le \emph{foncteur image
  inverse} par le morphisme de topos fibré $m : F \to G$; le  foncteur
$m_{\ast} : F'\to G'$ (7.1.7.2) est appelé le \emph{foncteur image
  directe} par le morphisme de topos fibré $m : F \to G$. Il résulte
immédiatement des définitions que les foncteurs $m \mapsto m^{\ast}$
et $m \mapsto m_{\ast}$ sont \emph{pleinement fidèles} et que leurs
images essentielles sont, d'une part, l'ensemble des $I$-foncteurs de
$G$ dans $F$ qui, fibre par fibre, sont exacts et commutent aux
limites inductives et, d'autre part, l'ensemble des $I^0$-foncteurs de
$F'$ dans $G'$ qui, fibre par fibre, possèdent un foncteur adjoint à
gauche exact. Ceci justifie les abus de langage consistant à définir
un morphisme de topos fibrés par son image directe ou son image
inverse. En fait un tel morphisme de topos\pageoriginale fibrés n'est
alors défini qu'à isomorphisme unique près.

On note $\Hhomtop_{\Cart /I}(F, G)$ la sous-catégorie pleine de
$\Hhomtop_I(F, G)$ engendrée par les morphismes $m$ tels que
$m_{\ast}$ (ou ce qui est équivalent $m^{\ast}$) soit un foncteur
cartésien (\hyperref[VI.subsubsec6.1.4]{6.1.4}). De tels morphismes sont appelés des
\emph{morphismes cartésiens}. On a donc des foncteurs pleinement
fidèles, induits par (7.1.7.1) et (7.1.7.2)
\begin{align*}
\Hhomtop_{\Cart/I} (F, G) & \to \Hhom_{\cart /I} (G,
  F)^{\circ}{\quoi},\tag{7.1.7.3}\\
\Hhomtop_{\Cart/I} (F, G) &\to \Hhom_{\cart /I^{\circ}} (F',
  G'){\quoi}.\tag{7.1.7.4}
\end{align*}

\subsubsection{}\etiquette[7.1.8]{VI.subsubsec7.1.8}

Soient $p : F \to I$, $q : G \to I$, $r : H \to I$ trois $𝒰$-topos
fibrés sur $I$ et $m : F \to G$, $n : G \to H$ deux morphismes de
$𝒰$-topos fibrés. On définit comme en IV 3.3 le composé des deux
morphismes $m$ et $n$, qu'on note $nm : F \to H$, et si $𝒱$ est un
univers tel que $𝒰 \in 𝒱$,
on définit une catégorie ($𝒱$-$𝒰$-$\Top_{/I}$)
dont les objets sont les $𝒰$-topos fibrés sur $I$ à catégories fibres
éléments de $𝒱$, et dont les morphismes sont les morphismes de
$𝒰$-topos fibrés sur $I$. De plus, comme en IV 3.3.2, l'application de
composition
$$
\Homtop_{/I} (F, G) \times \Homtop_{/I} (G, H) \to \Homtop_{/I} (F, H)
$$
est l'application induite sur les objets par un « foncteur de
composition de morphismes »:
$$
\Hhomtop_{/I} (F, G) \Hhomtop_{/I} (G, H) \to \Hhomtop_{/I} (F, H)
$$
et les considérations de IV 3.3 et IV 3.4 s'étendent sans changement
au cas des $U$-topos fibrés. Le lecteur aura d'ailleurs remarqué que
lorsque $I$ est une catégorie ponctuelle (un objet, un morphisme
identique) un $𝒰$-topos fibré sur $I$ « n'est autre » qu'un
$𝒰$-topos, et un morphisme de $𝒰$-topos fibrés
« n'est autre » qu'un morphisme de $𝒰$-topos et les constructions
précédentes se réduisent à celles de IV 3.

\subsubsection{}\etiquette[7.1.9]{VI.subsubsec7.1.9}

Soient $p : F \to I$ un $𝒰$-topos fibré et $\phi : J \to I$
un foncteur. La catégorie fibrée $p_J : F_J \to J$ déduite de $(F,
p)$ par le changement de base $\phi : J \to I$ est un
$𝒰$-topos fibré sur $J$. Soient $p': F'\to I^{\circ}$ et
$(p_J)': (F_J)'\to J^{\circ}$ les catégories fibrées déduites
respectivement de $(F, p)$ et $(F_J, p_J)$ par la construction de
\hyperref[VI.subsubsec7.1.3]{7.1.3}. On vérifie immédiatement que la catégorie
$((F_J)', (p_J)')$ est canoniquement isomorphe à la catégorie
$(p')_J : (F')_J \to J^{\circ}$ déduite\pageoriginale de $p': F'\to
I^{\circ}$ par le changement de base $\phi^{\circ} : J^{\circ} \to
I^{\circ}$. Ces deux catégories fibrées sur $J^{\circ}$ seront par
la suite identifiées et notées $p'_J : F'_J \to J^{\circ}$.
L'opération de changement de base est fonctorielle, et même
2-fonctorielle, par rapport à son argument: pour tout couple $p : F
\to I$ et $G \to I$ de $𝒰$-topos fibrés sur $I$, on a deux
diagrammes commutatifs de catégories et foncteurs:
\[\vcenter{
\xymatrix{
\Hhomtop_I(F,G) \ar[r] \ar[d]& \Hhom_I(G,F)^{\circ} \ar[d]\\
\Hhomtop_J(F_J,G_J) \ar[r]& \Hhom_J(G_J,F_J)^{\circ},
}}\leqno(7.1.9.1)
\]\etiquette[7.1.9.1]{VI.eq7.1.9.1}
\[\vcenter{
\xymatrix{
\Hhomtop_I(F,G) \ar[r] \ar[d]& \Hhom_{I^{\circ}}(F',G') \ar[d]\\
\Hhomtop_J(F_J,G_J) \ar[r]&
\Hhom_{J^{\circ}}(F'_{J^{\circ}},G_{J^{\circ}}),
}}\leqno{(7.1.9.2)}
\]\etiquette[7.1.9.2]{VI.eq7.1.9.2}
où les foncteurs verticaux sont des foncteurs changement de base
et où les
foncteurs horizontaux sont respectivement dans \eqhyperref[VI.eq7.1.9.1]{7.1.9.1}
les foncteurs
« morphisme $\mapsto$ image inverse » (7.1.7.1), et
dans \eqhyperref[VI.eq7.1.9.2]{7.1.9.2}
les foncteurs « morphisme $\mapsto$ image directe » (7.1.7.2).

\subsection{Sites fibrés}\etiquette[7.2]{VI.subsec7.2}

\subsubsection{}\etiquette[7.2.1]{VI.subsubsec7.2.1}

Un \emph{$𝒰$-site fibré} $C$ sur une catégorie $I$ est une catégorie
fibrée $p : C \to I$ dont les fibres sont munies de topologies faisant
de celles-ci des $𝒰$-sites (IV 3.0.2) telles que pour tout morphisme
$f : i \to j$ de $I$, le foncteur image inverse $f^{\ast} : C_j \to
C_i$ soit un morphisme du site $C_i$ dans le site $C_j$ (IV 4.9.1).

\subsubsection{}\etiquette[7.2.2]{VI.subsubsec7.2.2}

Soient $p : C \to I$ et $q : D \to I$ deux $𝒰$-sites fibrés sur
$I$. Un \emph{morphisme de $𝒰$-sites fibrés de $C$ dans} $D$ est un
$I$-foncteur (\hyperref[VI.subsec6.0]{6.0}) $m : D \to C$ tel que pour tout objet $i$ de $I$,
le foncteur $m_i : D_i \to C_i$ soit un morphisme du site $C_i$ dans
le site $D_i$. On désigne par $\Morsite_I(C, D)$ l'ensemble des
morphismes de sites fibrés de $C$ dans $D$. \emph{Un foncteur continu
  du site fibré $D$ dans le site fibré} $C$ est un $I$-foncteur de
$D$ dans $C$ qui induit sur les fibres des foncteurs continus (III
1). On note $\Cont_{I}(D, C)$ l'ensemble des foncteurs continus du
site fibré $D$\pageoriginale dans le site fibré $C$. On désigne par
$\Mmorsite_I (C, D)$ la sous-catégorie pleine de $\Hhom_I (D,
C)^{\circ}$ définie par l'ensemble d'objets $\Morsite_I (C, D)
\subset \Hom_I (D, C)$. On désigne de même par $\Ccont_I(D, C)$, la
sous-catégorie pleine de $\Hhom_I (D, C)$ définie par le
sous-ensemble d'objets $\Ccont_I (D, C) \subset \Hom_I (D, C)$. On a
donc deux foncteurs pleinement fidèles et injectifs sur les objets
\begin{equation*}
\Mmorsite_I (C, D) \hookrightarrow \Ccont_I (D, C)^{\circ} \hookrightarrow
\Hhom_I (D, C)^{\circ}{\quoi}.\tag{7.2.2.1}
\end{equation*}\etiquette[7.2.2.1]{VI.eq7.2.2.1}
On définit, de même qu'en \hyperref[VI.subsubsec7.1.7]{7.1.7}, les morphismes \emph{cartésiens} du
site fibré $C$ dans le site fibré $D$. On a un diagramme de
sous-catégories pleines
\[\vcenter{
\xymatrix@M=.2cm{
\Mmorsite_{\Cart/I}(C,D)\ar@{^{(}->}[r] \ar@{_{(}->}[d]&
\Hhom_{\Cart/I}(D,C)^{\circ} \ar@{_{(}->}[d]\\
\Mmorsite_I(C,D)\ar@{^{(}->}[r]& \Hhom_I(D,C)^{\circ}.
}}\leqno{(7.2.2.2)}
\]\etiquette[7.2.2.2]{VI.eq7.2.2.2}

\subsubsection{}\etiquette[7.2.3]{VI.subsubsec7.2.3}

Les morphismes de $𝒰$-sites fibrés se composent. Ceci permet de
définir la catégorie ($𝒱$-$𝒰$-Site/$I$) dont les
objets sont les $𝒰$-sites
dont les fibres appartiennent à un univers $𝒱$ (dont $𝒰$ est un
élément). La catégorie $𝒱$-$𝒰$-site est d'ailleurs munie d'une
structure de 2-catégorie et en particulier la composition des
morphismes de $𝒰$-sites fibrés est fonctorielle par rapport aux
morphismes entre morphismes. Pour tout foncteur $\phi : J \to I$ et
tout site fibré $p : C \to I$, la catégorie $p_J : C_J \to J$ déduite
de $(C, p)$ par le changement de base $\phi : J \to I$, est munie
canoniquement d'une structure de $𝒰$-site fibré sur $J$. L'opération
de changement de base est 2-fonctorielle.

\subsubsection{}\etiquette[7.2.4]{VI.subsubsec7.2.4}

Soit $p : C \to I$ une catégorie fibrée dont les fibres sont munies de
topologies faisant de celles-ci des $𝒰$-sites. Supposons que les
limites projectives soient représentables dans les catégories
fibrées. Pour que $p : C \to I$ soit un $𝒰$-site fibré il suffit que
les foncteurs images inverses soient exacts à gauche et transforment
familles couvrantes en familles couvrantes, et cette condition est
aussi suffisante lorsque les topologies des fibres sont moins fines
que la topologie canonique (IV 4.9.2). Tous les sites fibrés utilisés
dans ce séminaire seront du type décrit ci-dessus. De même, soient $p
: C \to I$ et $q : D \to I$ deux sites fibrés tels que les limites
projectives finies dans les catégories fibres soient
représentables.\pageoriginale Un foncteur cartésien $m : D \to C$ qui
induit sur les catégories fibres des foncteurs exacts à gauche et qui transforme les
familles couvrantes des catégories fibres en familles couvrantes est
un morphisme de sites fibrés de $C$ dans $D$, la réciproque étant
vraie si les topologies fibres de $C$ sont moins fines que la
topologie canonique (IV 4.9.2). Tous les morphismes de sites fibrés
utilisés dans ce séminaire seront du type décrit ci-dessus.

\subsubsection{}\etiquette[7.2.5]{VI.subsubsec7.2.5}

Un $𝒰$-topos fibré $p : F \to I$ est un $𝒰$-site fibré
lorsqu'on munit les fibres de la topologie canonique, ce que nous ferons toujours par
la suite. Soit $m:  F \to G$ un morphisme de $𝒰$-topos fibrés. Le
foncteur image inverse $m^{\ast} : G \to F$ (\hyperref[VI.subsubsec7.1.7]{7.1.7}) est un morphisme
du $𝒰$-site fibré $F$ dans le $𝒰$-site fibré $G$. On a donc un
foncteur $(𝒱\text{-}𝒰\text{-}\Top_{/I}) \to (𝒱
\text{-}U\text{-}\Site_{/I})$, foncteur qui n'est pas
fidèle en général. Mais ce foncteur se prolonge en fait naturellement
en un 2-foncteur qui induit, pour deux $𝒰$-topos fibrés $F$ et $G$ sur $I$, un foncteur:
\begin{equation*}
\Hhomtop_{I}(F, G) \to \Mmorsite_I (F, G)\tag{7.2.5.1}
\end{equation*}\etiquette[7.2.5.1]{VI.eq7.2.5.1}
qui est une \emph{équivalence de catégories}, ainsi qu'il résulte
immédiatement des définitions. D'ailleurs le foncteur \eqhyperref[VI.eq7.2.5.1]{7.2.5.1} composé
avec l'inclusion canonique $\Mmorsite_I (F, G) \hookrightarrow
\Hhom_{I}(G, F)^{\circ}$ n'est autre que le foncteur (7.1.7.1) qui, à
un morphisme de topos fibrés, associe son image inverse.

\subsubsection{}\etiquette[7.2.6]{VI.subsubsec7.2.6}

Soit $p : C \to I$ un $𝒰$-site fibré. Choisissons pour tout morphisme
$f : i \to j$ de $I$ un foncteur changement de base $f^{\ast} : C_j
\to C_{i}$. En passant aux catégories de $\sheaf{\Uu}$-faisceaux, le
foncteur $f^{\ast}$ se prolonge en un foncteur $\Top(f)^{\ast} :
C^{\sim}_{j} \to C^{\sim}_{i}$, où $\Top(f) : C^{\sim}_{i} \to
C^{\sim}_{j}$ est un morphisme de topos (IV 4.9.1.1), rendant
commutatif à isomorphisme près le diagramme:
\[\vcenter{
\xymatrix@C=2.5cm@R=1cm{
C_j \ar[r]^{f^{\ast}} \ar[d]_{\varepsilon_j}& C_i\ar[d]^{\varepsilon_i}\\
C^{\widetilde{\phantom{a}}}_j \ar[r]^{\Top(f)^{\ast}}&
C^{\widetilde{\phantom{a}}}_i
}}\leqno{(7.2.6.1)}
\]\etiquette[7.2.6.1]{VI.eq7.2.6.1}
Soient
\begin{equation}
c_{f, g} : f^{\ast} g^{\ast} \to (gf)^{\ast}\tag{7.2.6.2}
\end{equation}\etiquette[7.2.6.2]{VI.eq7.2.6.2}
les\pageoriginale isomorphismes canoniques \eqhyperref[VI.eq6.1.2.1]{6.1.2.1} et
\begin{equation*}
b_f : \Top(f)^{\ast} \epsilon_j \to \epsilon_i f^{\ast}\tag{7.2.6.3}
\end{equation*}\etiquette[7.2.6.3]{VI.eq7.2.6.3}
les isomorphismes qui « rendent commutatifs » les diagrammes
\eqhyperref[VI.eq7.2.6.1]{7.2.6.1}. En utilisant le fait que les foncteurs $\Top(f)^{\ast}$
commutent aux limites inductives, on vérifie immédiatement que pour
tout couple de morphismes composables : $i \xrightarrow{f} j
\xrightarrow{g} k$, il existe un et un seul isomorphisme:
\begin{equation*}
\Top(c_{f, g}) : \Top(f)^{\ast} \Top(g)^{\ast} \to \Top(gf)^{\ast}\tag{7.2.6.4}
\end{equation*}\etiquette[7.2.6.4]{VI.eq7.2.6.4}
tel qu'on ait
\begin{equation*}
\epsilon_i (c_{f, g})b_f \Top(f)^{\ast}(bg) = b_{gf} \Top(c_{f,g})\tag{7.2.6.5}{\quoi}.
\end{equation*}\etiquette[7.2.6.5]{VI.eq7.2.6.5}
De plus, les isomorphismes \eqhyperref[VI.eq7.2.6.4]{7.2.6.4} satisfont automatiquement la
condition de cocycles de 6.2.2.2. On peut donc construire une
catégorie fibrée
\begin{equation*}
p^{\sim} : C^{\sim / I} \to I\tag{7.2.6.6}
\end{equation*}\etiquette[7.2.6.6]{VI.eq7.2.6.6}
dont les fibres sont canoniquement isomorphes aux $𝒰$-topos
$C^{\sim}_i$ (\hyperref[VI.subsubsec6.1.3]{6.1.3}), et un foncteur cartésien
\begin{equation*}
\epsilon_{c/I} : C \to C^{\sim /I}\tag{7.2.6.7}
\end{equation*}\etiquette[7.2.6.7]{VI.eq7.2.6.7}
qui induit sur les fibres les foncteurs $\epsilon_i : C_{i} \to
C^{\sim}_i$. Il résulte immédiatement des définitions que $p^{\sim} :
C^{\sim/I} \to I$ est un $𝒰$-topos fibré sur $I$, que $\epsilon_{c/I}$
est un morphisme cartésien de $𝒰$-sites fibrés de $C^{\sim/I}$ dans
$C$, et que le $𝒰$-topos fibré $C^{\sim /I} \to I$ et le foncteur
cartésien $\epsilon_{C/I}$ ne dépendent pas, à isomorphisme canonique
près, des différents choix utilisés pour les construire (choix des
foncteurs changement de base $f^{\ast}$ et choix des prolongements
$\Top(f)^{\ast}$). Le $𝒰$-topos fibré $C^{\sim/I}$ est appelé le
\emph{$𝒰$-topos fibré associé} au $𝒰$-site fibré.

\setcounter{paragraph}{7}
\paragraph{}\etiquette[7.2.6.8]{VI.par7.2.6.8}

On vérifie que la formation du $𝒰$-topos fibré associé à un site fibré
« commute » aux opérations de changement de catégorie base.

\begin{enonce*}{Proposition 7.2.7}\etiquette[7.2.7]{VI.prop7.2.7}
Soient\pageoriginale $E$ un $𝒰$-topos fibré sur $I$ et $C$ un
$𝒰$-site fibré sur $I$. Le foncteur $m \mapsto m^{\ast} \epsilon_{C/I}$, associant à tout
morphisme de topos fibrés $m : E \to C^{\sim /I}$ le composé avec
$\epsilon_{C/I}$ du foncteur image inverse associé $m^{\ast} : C^{\sim
/ I} \to E$, induit une équivalence de catégories préservant les
objets cartésiens
$$
\Hhomtop_I (E, C^{\sim / I}) \xrightarrow{\approx} \Mmorsite_I(E,
C){\quoi}.
$$
Lorsque dans les catégories fibres de $C$ les limites projectives
finies sont représentables, le foncteur pleinement fidèle
correspondant
$$
\Hhomtop_{I} (E, C^{\sim/ I}) \to \Hhom_{I} (C, E)^{\circ}
$$
a comme image essentielle l'ensemble des $I$-foncteurs $g : C \to E$
qui sont exacts à gauche fibre par fibre et qui transforment les
familles couvrantes des catégories fibres en familles épimorphiques
des catégories fibres correspondantes.
\end{enonce*}

Cette proposition ne fait que généraliser au contexte fibré la
proposition IV 4.9.4, et la démonstration donnée en IV 4.9.4 se
généralise immédiatement.

\subsection{Exemple}\etiquette[7.3]{VI.subsec7.3}

\begin{enonce*}[remark]{Exemple 7.3.1}\etiquette[7.3.1]{VI.exem7.3.1}
(\emph{le topos fibré des morphismes d'un topos})

Soient $E$ un $𝒰$-topos, $\text{Fl}(E)$ la catégorie des
morphismes de $E$ et $p : \text{Fl} (E) \to E$ le foncteur qui associe à un morphisme son but. La
catégorie $(\text{Fl} (E), p)$ est un $𝒰$-topos fibré sur
$E$. (\hyperref[VI.def7.1.1]{7.1.1} et IV \hyperref[VI.thm5.1]{5.1}). La fibre en tout objet $X$ de $E$ est le topos $E_{/X}$. Le
foncteur image inverse $f^{\ast} : E_{/Y} \to E_{/X}$ par un morphisme
$f : X \to Y$ est le foncteur produit fibré. À tout foncteur $\phi : I
\to E$, on peut donc associer le $𝒰$-topos fibré sur $I$
\begin{equation*}
p_I : \text{Fl}(E)_I \to I{\quoi},\tag{7.3.1.1}
\end{equation*}\etiquette[7.3.1.1]{VI.eq7.3.1.1}
obtenu par le changement de base $\phi : I \to E$, dont la fibre en
tout objet $i$ de $I$ est $E_{/\phi (i)}$ (\hyperref[VI.subsubsec7.1.9]{7.1.9}). Si $𝒱$ est un
univers dont $E$ est un élément, on a, avec les notations de \hyperref[VI.subsubsec7.1.8]{7.1.8},
une application
\begin{equation*}
\Hom(I, E) \to \ob (𝒱\text{-}𝒰\text{-}\text{\rm top}_{/I}){\quoi},\tag{7.3.1.2}
\end{equation*}\etiquette[7.3.1.2]{VI.eq7.3.1.2}
qui associe à tout foncteur $\phi \in \Hom (I, E)$ le $𝒰$-topos fibré
sur $I$ déduit de $(\text{Fl}(E), p)$ par le\pageoriginale changement
de base $\phi : I \to E$. Soient $\phi_1$ et $\phi_2$ deux foncteurs de $I$ dans $E$ et $m :
\phi_1 \to \phi_2$ un morphisme de foncteurs. Notons $p_1 : \text{Fl}(E)_1
\to I$ et $p_2 : \text{Fl} (E)_2 \to I$ les $𝒰$-topos fibrés
sur $I$ déduits de $p : \text{Fl}(E) \to E$ par les changements de base $\phi_1 : I \to E$ et
$\phi_2 : I \to E$ respectivement. Choisissons, pour tout objet $i$
de $I$, un morphisme de topos (IV 5.5.2)
\begin{equation*}
\loc(m(i)) : E_{/{\phi}_{1}(i)} \to E_{/{\phi}_{2}(i)}{\quoi}.\tag{7.3.1.3}
\end{equation*}\etiquette[7.3.1.3]{VI.eq7.3.1.3}
On vérifie alors facilement qu'il existe un et un
seul\footnote{l'unicité provient de ce qu'on exige, avec les
notations
  de (IV 5.5.3), la formule : $(gf)_! = g_! f_!$.} morphisme cartésien
de topos fibrés sur $I$:
\begin{equation*}
\loc(m) : F1(E)_1 \to F1 (E)_2{\quoi},\tag{7.3.1.4}
\end{equation*}\etiquette[7.3.1.4]{VI.eq7.3.1.4}
qui induise sur les topos fibrés les morphismes $\loc(m(i))$.  Le
morphisme $\loc(m) : \text{Fl}(E)_1 \to \text{Fl} (E)_2$ est déterminé à
isomorphisme canonique près par le morphisme de foncteurs $m : \phi_1
\to \phi_2$. Si $n : \phi_2 \to \phi_3$ est un morphisme de foncteurs,
on a un isomorphisme
\begin{equation*}
c(m, n) : \loc (n) \loc(m) \simeq \loc(nm){\quoi},\tag{7.3.1.5}
\end{equation*}\etiquette[7.3.1.5]{VI.eq7.3.1.5}
et les isomorphismes $c(m, n)$ possèdent une propriété de cocycle
analogue à \eqhyperref[VI.eq6.1.2.2]{6.1.2.2}.
\end{enonce*}

\begin{enonce*}[remark]{Exemple 7.3.2}\etiquette[7.3.2]{VI.exem7.3.2}
(Le site fibré des morphismes d'un site).

Soient $C$ un $𝒰$-site où \emph{les produits fibrés sont
  représentables}, $\text{Fl}(C)$ la catégorie des morphismes de $C$ et $p :
  \text{Fl}(C) \to C$ le foncteur qui associe à un morphisme son but. La
  catégorie fibrée $(\text{Fl}(C), p)$ est un $𝒰$-site fibré sur $C$
  lorsqu'on munit les catégories fibres des topologies induites (II
  \hyperref[VI.cor5.2]{5.2}). Soient $\epsilon_C  : C \to C^{\sim}$ le foncteur canonique
  de $C$ dans la catégorie des $𝒰$-faisceaux sur $C$. En notant
  $\text{Fl}(\epsilon_C)$ l'extension naturelle de ce dernier foncteur aux
  catégories de morphismes, on a un diagramme commutatif de
  catégories et foncteurs:
\[\vcenter{
\xymatrix@C=2cm{
\text{Fl}(C) \ar[r]^{\text{Fl}(\varepsilon_C)}\ar[d]_{p}& \text{Fl}(C^{\widetilde{\phantom{a}}}) \ar[d]^{p}\\
C       \ar[r]^{\varepsilon_C} & C^{\widetilde{\phantom{a}}}
}}\leqno{(7.3.2.1)}
\]\etiquette[7.3.2.1]{VI.eq7.3.2.1}
d'où un\pageoriginale foncteur cartésien entre catégories
fibrées sur
$C$:
\[\vcenter{
\xymatrix{
\text{Fl}(C) \ar[rr] \ar[rd]_{p}&     &
\text{Fl}(C^{\widetilde{\phantom{a}}})_C \ar[ld]^{p_{/C}}\\
                   & C     &\\
}}\leqno{(7.3.2.2)}
\]\etiquette[7.3.2.2]{VI.eq7.3.2.2}
où $p_{/C} : \text{Fl}(C^{\sim})_{C} \to C$ est déduite de
$\text{Fl}(C^{\sim}) \to
C^{\sim}$ par le changement de base $\epsilon_{C} : C \to
C^{\sim}$. Le foncteur $\text{Fl}(C) \to \text{Fl}(C^{\sim})_{C}$ de
\eqhyperref[VI.eq7.3.2.2]{7.3.2.2} est un morphisme du site fibré
$\text{Fl}(C^{\sim})_{C}$ dans le site fibré
$\text{Fl}(C)$
(\hyperref[VI.subsubsec7.2.2]{7.2.2}), d'où par \hyperref[VI.prop7.2.7]{7.2.7} un morphisme de topos fibrés sur $C$:
\begin{equation*}
\can : \text{Fl}(C^{\sim})_{C} \to \text{Fl} (C)^{\sim/C}{\quoi},\tag{7.3.2.3}
\end{equation*}\etiquette[7.3.2.3]{VI.eq7.3.2.3}
où $\text{Fl}(C)^{\sim/C}$ est le topos fibré sur $C$ associé au site fibré
$\text{Fl}(C) \to C$ (\hyperref[VI.subsubsec7.2.6]{7.2.6}). Il résulte de la définition du topos fibré
associé (\hyperref[VI.subsubsec7.2.6]{7.2.6}) et de II \hyperref[VI.cor5.5]{5.5} que le morphisme can de \eqhyperref[VI.eq7.3.2.3]{7.3.2.3} est une
$C$-équivalence de $𝒰$-topos fibrés sur $C$, \ie il existe un
morphisme $\Theta : \text{Fl}(C)^{\sim/C} \to \text{Fl} (C^{\sim})_C$
de $𝒰$-topos fibrés sur $C$ tel que les composés $\Theta \circ \can$ et $\can \circ
\Theta$ soient $C$-isomorphes à l'identité.

Soit $\phi : I \to C$ un foncteur. On en déduit par changement de base
un site fibré sur $I$:
\begin{equation*}
p_I : \text{Fl} (C)_I \to I\tag{7.3.2.4}
\end{equation*}\etiquette[7.3.2.4]{VI.eq7.3.2.4}
et comme la formation du topos fibré associé commute au changement de
base (\hyperref[VI.subsubsec7.2.6]{7.2.6}), on a une $I$-équivalence de $𝒰$-topos fibrés
\begin{equation*}
\can_I : \text{Fl} (C^{\sim})_{I} \to \text{Fl}(C)^{\sim/I}_{I}\tag{7.3.2.5}
\end{equation*}\etiquette[7.3.2.5]{VI.eq7.3.2.5}
où $\text{Fl}(C^{\sim})_{I} \to I$ est le $𝒰$-topos fibré obtenu par le
changement de base $\epsilon_C \circ \phi : I \to C^{\sim}$.

Soient enfin $\phi_1 : I \to C$ et $\phi_2 : I \to C$ deux foncteurs
et $m : \phi_1 \to \phi_2$ un morphisme de foncteurs. Notons $\text{Fl}(C)_1$
et $\text{Fl}(C)_2$ les sites fibrés sur $I$ déduits de $p :
\text{Fl}(C) \to C$ par les changements de base $\phi_1$ et $\phi_2$
respectivement. On
construit comme en \hyperref[VI.exem7.3.1]{7.3.1} un morphisme de sites fibrés sur $I$:
\begin{equation*}
\loc(m) : \text{Fl} (C)_1 \to \text{Fl}(C)_2{\quoi}.\tag{7.3.2.6}
\end{equation*}\etiquette[7.3.2.6]{VI.eq7.3.2.6}
Notons\pageoriginale $\text{Fl}(C^{\sim})_1 \to I$ et
$\text{Fl}(C^{\sim})_2 \to I$ les $𝒰$-topos fibrés sur $I$
obtenus par les changements de base $\epsilon_C \circ \phi_1$ et
$\epsilon_C \circ \phi_2$ respectivement. On a un diagramme
commutatif à isomorphisme près de morphismes sur $I$:
\[\vcenter{
\xymatrix@C=2.5cm{
\text{Fl}(C^{\sim})_1 \ar[d]\ar[r]^{\loc(\varepsilon_C\ast m)}&
\text{Fl}(C^{\sim})_2 \ar[d]_{\bigdot}\\
\text{Fl}(C)_1       \ar[r]^{\loc(m)} & \text{Fl}(C)_2
}}\leqno{(7.3.2.7)}
\]\etiquette[7.3.2.7]{VI.eq7.3.2.7}
où les morphismes verticaux se déduisent par changement de base des
morphismes de \eqhyperref[VI.eq7.3.2.2]{7.3.2.2}.
\end{enonce*}

\begin{enonce*}[remark]{Exercice 7.3.3}\etiquette[7.3.3]{VI.exe7.3.3}
(\emph{Petit site et gros site fibrés}).

Cet exercice est une suite à l'exercice IV 4.10.6 dont on utilise les
hypothèses et les notations. On note $\MM$ la sous-catégorie
pleine de $\text{Fl}(S)$ définie par les morphismes de $ℳ$,
et $p : \MM \to S$ le foncteur qui associe à un morphisme son but.

$6^{\circ}$) Montrer que $(\MM, p)$ est un site fibré sur $S$ et
que, lorsque dans $S$ les produits fibrés sont représentables, le
foncteur d'inclusion $\MM \to \text{Fl}(S)$ est un morphisme cartésien du
site fibré $\text{Fl} (S) \to S$ dans le site fibré $\MM \to S$.

$7^{\circ}$) Soit $\MM^{\sim/S} \to S$ le topos fibré associé à
$\MM \to S$. Montrer qu'on a avec les notations de \hyperref[VI.exem7.3.2]{7.3.2}, un
morphisme canonique $g : \text{Fl} (S^{\sim})_{/S} \to \MM^{\sim/S}$ de
topos fibrés sur $S$ qui induit sur chaque fibre en $X \in \ob (S)$ le
morphisme $(\Prol_{X}, \Res_X)$ (IV 4.10.6 $4^{\circ})$). Montrer que,
bien que pour tout $X \in \ob (S)$ le morphisme $g_X : S^{\sim}_{/X}
\to S(X)^{\sim}$ admette une section (\loccit), le morphisme $g$
n'admet pas en général de section.
\end{enonce*}

\subsection{La topologie totale d'un topos ou d'un site fibré}\etiquette[7.4]{VI.subsec7.4}

\begin{enonce*}{Définition 7.4.1}\etiquette[7.4.1]{VI.def7.4.1}
Soit $p  : C \to I$ un $𝒰$-site fibré et pour tout $i \in \ob (I)$,
notons $\alpha_{i!} : C_i \to C$ le foncteur d'inclusion de la
catégorie fibre $C_i$ dans $C$. On appelle \emph{topologie totale} sur $C$ la
topologie la moins fine sur $C$ qui rende continu les foncteurs
$\alpha_{i!}$ pour $i \in \ob$ \emph{(I) (III 3.6)}. On appelle
\emph{site total}  la
catégorie $C$ munie de la topologie totale.
\end{enonce*}

\begin{enonce*}{Proposition 7.4.2}\etiquette[7.4.2]{VI.prop7.4.2}
Soient\pageoriginale $p : C \to I$ un $𝒰$-site fibré, $T$ la
topologie totale sur $C$ et pour tout $i \in \ob(I)$, $T_i$ la topologie de la fibre $C_i$.
\begin{enumerate}
\renewcommand{\labelenumi}{\theenumi)}
\item Soit $X$ un objet de $C$ au-dessus d'un objet $i$ de $I$. Une
  famille $(X_\beta \to X)_{\beta \in B}$ est couvrante pour $T$ si et
  seulement si elle est raffinée par une famille couvrante pour $T_i
  (Y_{\gamma} \to X)_{\gamma \in \Gamma}$ de $C_i$.

\item La topologie $T$ induit sur les catégories $C_i$, les
  topologies $T_i$.

\item Pour tout $i \in \ob (I)$, le foncteur $\alpha_{i!} : C_i \to C$
  est cocontinu.
\end{enumerate}
\end{enonce*}

\setcounter{subsubsection}{2}
\setcounter{paragraph}{0}
\paragraph{}\etiquette[7.4.2.1]{VI.par7.4.2.1}

La propriété 3) résulte de 1) et de III 2.1. Démontrons 1). Soit, pour
tout objet $X$ de $C$, $J(X)$ l'ensemble des cribles de $C/X$ décrits
dans 1). On vérifie que les $J(X)$ possèdent les propriétés $T$ 1),
$T$ 2) et $T$ 3) de II 1.1, et qu'ils définissent par suite une
topologie $T'$ sur $C$. La topologie $T'$ est la moins fine des
topologies sur $C$ pour lesquelles les familles couvrantes pour les
topologies $T_i$ sont couvrantes. La topologie $T'$ est donc moins
fine que $T$ (III 1.6). Pour montrer que $T'= T$, il suffit donc de
montrer que pour tout objet $i$ de $I$, le foncteur $\alpha_{i!} : C_i
\to C$ est continu lorsqu'on munit $C$ de la topologie $T'$; ou
encore, il suffit de montrer que la topologie induite par $T'$ sur la
catégorie $C_i$ est la topologie $T_i$, ce qui démontrera en même
temps la propriété 2). Soit $T'_i$ la topologie sur $C_i$ induite
par $T'$, et $u : R \hookrightarrow X$ un crible couvrant pour
$T'_i$. D'après III 3.2 le morphisme $\alpha_{i!} (u) : \alpha_{i!}
(R)\to \alpha_{i!} (X)$ est bicouvrant et en particulier couvrant, ce qui
revient à dire, d'après la définition de $T'$, que le crible $R
\hookrightarrow X$ est couvrant pour $T_i$. La topologie $T'_i$ est
donc moins fine que $T_i$, et pour démontrer que $T'_i$ est plus
fine que $T_i$, il suffit, par définition de la topologie induite (III
3.1), de montrer que le foncteur $\alpha_{i!} : C_i \to C$ est continu
lorsqu'on munit les catégories $C_i$ et $C$ des topologies $T_i$ et
$T'$ respectivement. En résumé, pour démontrer les propriétés 1) et
2), il suffit, via III 1.2 et  1.5, de démontrer le lemme suivant:

\begin{enonce*}{Lemme 7.4.2.2}\etiquette[7.4.2.2]{VI.lem7.4.2.2}
Avec les hypothèses et notations de \hyperref[VI.prop7.4.2]{7.4.2} et \hyperref[VI.par7.4.2.1]{7.4.2.1}, soient $X$ un
objet de $C$ au-dessus de $i$ et $R \xrightarrow{u} X$ un morphisme de
$C^{\widehat{\phantom{,}}}_i$. Les propriétés suivantes sont
équivalentes:
\begin{enumerate}
\renewcommand{\theenumi}{\roman{enumi}}
\renewcommand{\labelenumi}{\theenumi)}
\item Le morphisme\pageoriginale $R \to X$ est bicouvrant pour $T_i$.

\item Le morphisme $\alpha_{i!}(u) : \alpha_{i!}(R) \to
  \alpha_{i!}(X)$ est bicouvrant pour $T'$.
\end{enumerate}
\end{enonce*}

i) $\Rightarrow$ ii) : Il est clair que si $R \xrightarrow{u} X$ est
bicouvrant pour $T_i$, le morphisme $\alpha_{i!} (u)$ est couvrant
pour $T'$. Il reste à montrer qu'il est bicouvrant (II \hyperref[VI.cor5.2]{5.2}). Soit $Y$
un objet de $C$ au-dessus d'un objet $j$ de $I$, $m$ et $n$ $n : Y
\rightrightarrows \alpha_{i!} (R)$ deux morphismes de
$C^{\widehat{\phantom{,}}}$ tels que $\alpha_{i!} (u) m = \alpha_{i!}
(u) n$, et posons $f = p(\alpha_{i!} (u)m)$. Comme le foncteur
$\alpha_{i!}$ commute aux limites inductives (I 5.4.3)), on a (I 3.4):
$$
\Hom_{C^{\widehat{\phantom{,}}}} (Y, \alpha_{i!} (R)) \simeq
\varinjlim_{C_{i}/R} \Hom_C (Y, \alpha_{i!} (.)){\quoi}.
$$
Comme pour tout objet $Z$ de $C_i$, on a avec les notations de \hyperref[VI.subsubsec6.1.1]{6.1.1}:
$$
\Hom_{C} (Y, Z) \simeq \coprod_{w \in \Hom(j, i)} \Hom_w (Y,
Z){\quoi},
$$
on en déduit
$$
\varinjlim_{C_{i}/R} \Hom (Y, \alpha_{i!} (.)) \simeq
\coprod_{w \in \Hom(j, i)} \varinjlim_{C_{i}/R}
\Hom_{w} (Y, \alpha_{i!} (.)){\quoi}.
$$
Mais pour tout $w \in \Hom(j, i)$, on a un foncteur image inverse
$w^{\ast} : C_i \to C_j$, et on a $\Hom_w (Y, \alpha_{i!} (.)) \simeq
\Hom_{C_{j}}(Y, w^{\ast} (.))$. D'où en notant encore $w^{\ast} :
C^{\widehat{\phantom{,}}}_i \to C^{\widehat{\phantom{,}}}_j$ le
prolongement de $w^{\ast}$ aux préfaisceaux (qu'on devrait noter
d'après I 5.4 $(w^{\ast})_{!}$) et en remarquant que le foncteur
$w^{\ast} : C^{\widehat{\phantom{,}}}_i \to
C^{\widehat{\phantom{,}}}_j$ commute aux limites inductives (I
5.4.3)), on obtient un isomorphisme:
$$
\Hom_{C^{\widehat{\phantom{,}}}} (Y, \alpha_{i!} (R)) =
\coprod_{w \in \Hom(j, i)}
\Hom_{C^{\widehat{\phantom{,}}}_{j}}(Y,w^{\ast}(R)){\quoi}.
$$
Dans cet isomorphisme, les morphismes $m$ et $n : Y \rightrightarrows
\alpha_{i!} (R)$ correspondent à deux morphismes $m'$ et $n': Y
\rightrightarrows f^{\ast} R$ tels que $f^{\ast} (u)m' = f^{\ast}
(u)n'$. Comme le foncteur image inverse est un morphisme de topos, il
est en particulier continu et par suite $f^{\ast}(u)$ est bicouvrant
(III 3.2). Il existe donc une famille couvrante de $C_j$ $(Y_{\delta}
\xrightarrow{V_{\delta}} Y)_{\delta}$ telle que pour tout $\delta$ on
ait $m' v_{\delta} = n' v_{\delta}$. On en déduit que $m v_{\delta} =
n v_{\delta}$ pour tout $\delta$, et par suite que $\alpha_{i!}(u)$
est bicouvrant pour $T'$.

ii) $\Rightarrow$ i) :\pageoriginale Supposons que $\alpha_{i!} (R)
\xrightarrow{\alpha_{i!}(u)} \alpha_{i!} (X)$ soit bicouvrant.
Le foncteur $\alpha_{i!} : C_i \to C$ est cocontinu (III 2.1). Par
suite $\alpha^{\ast}_{i} \alpha_{i!}(u)$ est bicouvrant (III 2.3.2) et
(II \hyperref[VI.cor5.3]{5.3} ii)). Le foncteur $\alpha_{i!} : C_i \to C$ possède la
propriété (PPF) de I 5.14. Par suite on a un diagramme cartésien des
$C^{\widehat{\phantom{,}}}_i$ (I 5.14 3)):
\[
\xymatrix@C=2.25cm@R=1.25cm{
R  \ar[r] \ar[d]_{u}& \alpha_i^{\ast}\alpha_{i!}(R) \ar[d]^{\alpha_i^{\ast}\alpha_{i!}(u)}\\
X     \ar[r]          & \alpha_i^{\ast}\alpha_{i!}(X).
}
\]
Par suite, $u$ est bicouvrant (II \hyperref[VI.cor5.2]{5.2}).

\begin{enonce*}[remark]{Remarque 7.4.3}\etiquette[7.4.3]{VI.rem7.4.3}
\begin{enumerate}
\renewcommand{\labelenumi}{\theenumi)}
\item La proposition \hyperref[VI.prop7.4.2]{7.4.2} peut s'énoncer et se démontrer dans un
  cadre un peu plus général que celui des sites fibrés, mais
  apparemment sans intérêt: au lieu d'un site fibré $p : C \to I$, on
  considère une catégorie fibrée $p : C \to I$ dont les fibres sont
  munies de topologies et dont les foncteurs images inverses sont
  continus pour ces topologies (alors que dans le cas des sites fibrés
  on exige de plus que ces foncteurs images inverses soient des
  morphismes de sites (IV 4.9)).

\item On peut démontrer que la topologie totale d'un site fibré est la
  topologie la plus fine rendant cocontinus les foncteurs $\alpha_{i!}
  : C_i \to C$ (III 2).

\item Soit $p : C \to I$ un $𝒰$-site fibré au-dessus d'une catégorie
  équivalent à une catégorie $𝒰$-petite. Alors le site total $C$ est un
  $𝒰$-site (II 3.0.2). En effet, si pour tout objet $i$ de $I$
  $(X_{\beta})_{\beta \in B_i}$ est une petite famille topologiquement
  génératrice de $C_i$, la famille $(X_{\beta})_{\beta \in
  \amalg_{i}B_i}$ est topologiquement génératrice pour $C$
  et est $𝒰$-petite. On appelle \emph{Topos total} et on note
  $\Top(C)$ le topos des faisceaux sur le site total.

\item Soit $p : C \to I$ un $𝒰$-site fibré sur une petite
  catégorie. Comme le foncteur $\alpha_{i!} : C_i \to C$ est
  cocontinu, il donne naissance à un morphisme de topos $\alpha_i =
  (\alpha^{\ast}_i, \alpha_{i\ast})$ de $C^{\sim}_i$ dans $\Top(C)$ (IV
  4.7). Comme de plus le foncteur $\alpha_{i!}$: $C_i \to C$ est
  continu, le foncteur $\alpha^{\ast}_i : \Top(C) \to C^{\sim}_i$ admet
  un adjoint à gauche noté\pageoriginale encore $\alpha_{i!} : C^{\sim}_i \to
  \Top(C)$ qui prolonge le foncteur $\alpha_{i!} : C_i \to C$ (III 1.2
  iv)).
\end{enumerate}
\end{enonce*}

\begin{enonce*}{Proposition 7.4.4}\etiquette[7.4.4]{VI.prop7.4.4}
Soit $p : C \to I$ un $𝒰$-site fibré sur une petite catégorie $I$.
\begin{enumerate}
\renewcommand{\labelenumi}{\theenumi)}
\item Un préfaisceau $F$ sur $C$ est un faisceau pour la topologie
  totale si et seulement si pour tout objet $i$ de $I$, le préfaisceau
  $F \circ \alpha_{i!} = F | C_i$ est un faisceau sur $C_i$.

\item Un foncteur $m$ de $C$ dans un site $D$ est un foncteur continu
  du site total $C$ dans le site $D$ si et seulement si pour tout
  objet $i$ de $I$ le foncteur $m | C_i = m \circ \alpha_{i!} : C_i
  \to D$ est continu.
\end{enumerate}
\end{enonce*}

Si $F$ est un faisceau pour la topologie totale, le préfaisceau $F
\circ \alpha_{i!}$ est un faisceau sur $C_i$ car $\alpha_{i!} : C_i
\to C$ est continu. Supposons que pour tout $i$, $F \circ
\alpha_{i!}$ soit un faisceau sur $C_i$. Soit $X$ un objet de $C$
au-dessus d'un objet $i$ de $I$. Pour tout crible couvrant $R$ de
$C_i /X$, on a $F(X) \simeq F \circ \alpha_{i!} (R)$ et par suite
$F(X) \simeq F(\alpha_{i!} (R))$. Soit $R' \hookrightarrow X$ le
crible de $C/X$ image du morphisme canonique de préfaisceaux
$\alpha_{i!} (R) \to X$. Le morphisme $\alpha_{i!} (R) \to R'$ est
un épimorphisme de préfaisceaux et par suite $F(R') \to
F(\alpha_{i!} (R))$ est injectif. Comme $F(R') \to F(\alpha_{i!}
(R))$ est aussi surjective, l'application $F(R') \to F(\alpha_{i!}
(R))$ est bijective, et par suite l'application $F(X) \to F(R')$ est
bijective. Comme les cribles couvrants $R' \hookrightarrow X$ qu'on obtient ainsi sont cofinaux dans l'ensemble des cribles couvrants $X$ \eqhyperref[VI.par7.4.2.1]{7.4.2.1}),
le préfaisceau $F$ est un faisceau, car en revenant à la
construction du faisceau associé (II 3) on constate que $F$ est
isomorphe à son faisceau associé (II 3.3). La deuxième assertion se
déduit immédiatement de la première et de la définition de la
continuité (III 1.1).

\setcounter{subsubsection}{4}
\subsubsection{}\etiquette[7.4.5]{VI.subsubsec7.4.5}

Soient $p : C \to I$ un $𝒰$-site fibré sur une catégorie $I$
équivalente à une petite catégorie et
$$
f : i \to j
$$
une flèche de la catégorie d'indices $I$. En notant encore $f :
C^{\sim}_i \to C^{\sim}_j$ le morphisme de topos déterminé par la
structure de site fibré, on a (avec les notations de 7.4.3.4)
un diagramme de morphismes de $𝒰$-topos
\[\vcenter{
\xymatrix@C=2.25cm@R=1.25cm{
C^{\sim}_i \ar[d]_{f} \ar[r]^{\alpha_i} & \Top(C)\\
C_j  \ar[ru]_{\alpha_j}&
}}\leqno{(7.4.5.1)}
\]\etiquette[7.4.5.1]{VI.eq7.4.5.1}
où $\Top(C)$\pageoriginale est le topos total (7.4.3.3).
Le diagramme \eqhyperref[VI.eq7.4.5.1]{7.4.5.1}
\emph{n'est pas commutatif en général}. Nous allons définir un
morphisme canonique de morphismes de topos:
\begin{equation*}
\rho_f : \alpha_i \to \alpha_j \circ f{\quoi}.\tag{7.4.5.2}
\end{equation*}\etiquette[7.4.5.2]{VI.eq7.4.5.2}
Il suffit de définir un tel morphisme au niveau des images inverses,
\ie de définir un morphisme de foncteurs:
\begin{equation*}
\beta_f^{\ast} : f^{\ast}\circ \alpha_j^{\ast} \to
\alpha_i^{\ast}{\quoi}, \tag{7.4.5.3}
\end{equation*}\etiquette[7.4.5.3]{VI.eq7.4.5.3}
où encore, en utilisant l'adjonction entre les foncteurs $f^{\ast}$ et
$f_{\ast}$, il suffit de définir un morphisme de foncteurs:
\begin{equation*}
\gamma_f : \alpha^{\ast}_j \to f_{\ast} \alpha^{\ast}_i{\quoi}.\tag{7.4.5.4}
\end{equation*}\etiquette[7.4.5.4]{VI.eq7.4.5.4}
Soient donc $Y$ un objet de $C_j$ et $F$ un faisceau sur le site
total $C$. On a, par définition, $\alpha^{\ast}_j F(Y) =
F(\alpha_{j!} (Y))$ et  $f_{\ast} \alpha^{\ast}_i F(Y) =
F(\alpha_{i!} (f^{\ast} (Y)))$, où dans le deuxième membre
$f^{\ast} : C_j \to C_i$ désigne le foncteur image inverse pour la
structure fibrée. Mais, par définition de ce foncteur image inverse,
on a un morphisme cartésien canonique:
\begin{equation*}
m_Y : \alpha_{i!} f^{\ast} \to \alpha_{j!}{\quoi};\tag{7.4.5.5}
\end{equation*}\etiquette[7.4.5.5]{VI.eq7.4.5.5}
d'où, en appliquant le foncteur $F$, une application fonctorielle en $Y$;
\begin{equation*}
\gamma_f (F) (Y) = F(m_Y) : \alpha_{j}^{\ast} F(Y) \to f_{\ast}
\alpha^{\ast}_i F(Y){\quoi},\tag{7.4.5.6}
\end{equation*}\etiquette[7.4.5.6]{VI.eq7.4.5.6}
application définissant un morphisme de faisceaux sur $C_j$:
\begin{equation*}
\gamma_f (F) : \alpha_{j}^{\ast} F \to f_{\ast} \alpha^{\ast}_i
F{\quoi};\tag{7.4.5.7}
\end{equation*}\etiquette[7.4.5.7]{VI.eq7.4.5.7}
d'où le morphisme de foncteurs $\gamma_f$ de \eqhyperref[VI.eq7.4.5.4]{7.4.5.4}.

Soient alors\pageoriginale $f : i \to j$ et $g : j \to k$ deux morphismes
composables de $I$. On a un isomorphisme canonique:
$$
c'_{g, f} : g_{\ast} f_{\ast} \to (gf)_{\ast}
$$
et on vérifie immédiatement la formule de compatibilité:
\begin{equation*}
c'_{g, f} \circ g_{\ast} (\gamma_f) \circ \gamma_g =
\gamma_{gf}{\quoi}.\tag{7.4.5.8}
\end{equation*}\etiquette[7.4.5.8]{VI.eq7.4.5.8}

\subsubsection{}\etiquette[7.4.6]{VI.subsubsec7.4.6}

Introduisons alors le $𝒰$-topos fibré $p^{\sim/I} : C^{\sim/I} \to I$
associé à $p : C \to I$ (\hyperref[VI.subsubsec7.2.6]{7.2.6}) et la catégorie fibrée sur $I^{\circ}$
correspondante $(p^{\sim/I})' :(C^{\sim/I})'\to I^{\circ}$ \eqhyperref[VI.eq7.1.3.2]{7.1.3.2}. Les
considérations précédentes permettent d'associer à tout objet $F$ de
$\Top(C)$ une famille $F_i = (\alpha^{\ast}_i F)_{i \in \ob (I)}$
d'objets des $C^{\sim}_i$ et une famille de morphismes
$$
\gamma_f F : F_j \to f_{\ast} F_i,\quad f\in F1(I),\quad f : i \to j{\quoi},
$$
cette famille de morphismes étant soumise aux conditions de
compatibilités de \eqhyperref[VI.eq7.4.5.8]{7.4.5.8}. On a donc associé à $F$ un
$I^{\circ}$-foncteur de la catégorie $I^{\circ}$ dans la catégorie
$(C^{\sim/I})'$ \ie un objet $\Theta_C (F)$ de la catégorie
$\Hhom_{I^{\circ}} (I^{\circ}, (C^{\sim/I})')$. Comme cette
construction est fonctorielle en $F$, on a en définitive un foncteur:
\begin{equation*}
\Theta_C : \Top(C) \to \Hhom_{I^{\circ}} (I^{\circ},
(C^{\sim/I})'){\quoi}.\tag{7.4.6.1}
\end{equation*}\etiquette[7.4.6.1]{VI.eq7.4.6.1}

\begin{enonce*}{Proposition 7.4.7}\etiquette[7.4.7]{VI.prop7.4.7}
Soit $C \to I$ un $\mathcal{U}$-site fibré avec $I$ équivalente à une
petite catégorie. Le foncteur $\Theta_C$ \eqhyperref[VI.eq7.4.6.1]{7.4.6.1} est une équivalence
de catégories.
\end{enonce*}

Nous nous contenterons de décrire un foncteur quasi-inverse à
$\Theta_C$. Soit \\ $\sigma \in \Hom_{I^{\circ}} (I^{\circ},
(C^{\sim/I})')$ \ie un $I^{\circ}$-foncteur $i \mapsto \sigma(i)$ de
$I^{\circ}$ dans $(C^{\sim/I})'$. Pour tout objet $X$ de $C$ au-dessus
de $p(X) = i \in \ob (I)$, on dispose d'un faisceau $\sigma(p(X)) \in
\ob C^{\sim}_i$; d'où un ensemble $\sigma (p(X)) (X)$. Soit $m : X \to
Y$ un morphisme de $C$ au-dessus du morphisme $f : i \to j$ de $I$. Le
morphisme $m$ se factorise d'une manière unique en un morphisme $m': X
\to f^{\ast} Y$ de $C_i$ et le morphisme cartésien canonique $f(Y) \to
Y$. De même le morphisme $\sigma(f)$ se factorise de manière unique en
un morphisme $\gamma_f(\sigma) : \sigma (p(Y)) \to f_{\ast} \sigma(p(X))$
et le morphisme $I^{\circ}$-cartésien canonique $f_{\ast} \sigma(p(X))
\to \sigma(p(X))$. On a donc une application de $\sigma(p(Y))(Y)$ dans
$\sigma(p(X))(X)$,\pageoriginale obtenue en composant les applications:
$$
\sigma(p(Y))(Y) \xrightarrow{\gamma_f(\sigma)(Y)} f_{\ast} (p(X))(Y) =
\sigma(p(X))(f^{\ast}Y) \xrightarrow{\sigma(p(X))(m')} \sigma(p(X))(X){\quoi}.
$$
On a donc associé à tout objet $X$ de $C$ un ensemble
$\sigma(p(X))(X)$ et à tout morphisme $m : X \to Y$ une application
$\sigma(p(Y))(Y) \to \sigma(p(X))(X)$. On vérifie qu'on a bien
déterminé ainsi un préfaisceau sur $C$, préfaisceau qui ne dépend pas
du choix des foncteurs images inverses utilisés pour le construire. Il
résulte de \hyperref[VI.prop7.4.4]{7.4.4} que le préfaisceau $X \mapsto \sigma(p(X))(X)$ est un
faisceau sur le site total $C$, et  il est clair que ce faisceau
dépend fonctoriellement de l'objet $\sigma$ de $\Hhom_{I^{\circ}}
(I^{\circ}, (C^{\sim/I})')$. Il est aussi clair, en revenant à la
définition du foncteur $\Theta_{C}$, que le foncteur de $\Hhom_{I^{\circ}}
(I^{\circ}, (C^{\sim/I})')$ dans $\Top(C)$ qu'on vient de construire
est un foncteur quasi-inverse de $\Theta_C$.

\begin{enonce*}[remark]{Remarque 7.4.8}\etiquette[7.4.8]{VI.rem7.4.8}
Il résulte en particulier de \hyperref[VI.prop7.4.7]{7.4.7} que la catégorie
$\Hhom_{I^{\circ}} (I^{\circ}, (C^{\sim/I})')$ est un
$𝒰$-topos. Ce dernier fait peut se voir directement. On voit
immédiatement, en utilisant I 9.21.10, que les conditions a), b) et
c) de IV 1.1.2 sont satisfaites, et l'existence d'une petite
catégorie génératrice résulte de I 9.25.
\end{enonce*}

\setcounter{subsubsection}{8}
\subsubsection{}\etiquette[7.4.9]{VI.subsubsec7.4.9}

Soient $p : C \to I$ et $q : D \to I$ deux $𝒰$-sites fibrés sur une
petite catégorie et soit $m$ un morphisme du site fibré $C$ dans le
site fibré $D$ (\hyperref[VI.subsubsec7.2.2]{7.2.2}). Notons $m^{\sim/I} : C^{\sim/I} \to
D^{\sim/I}$ le morphisme correspondant entre les $𝒰$-topos fibrés
associés. Le morphisme de topos fibrés $m^{\sim/I}$ est défini à
isomorphisme canonique près (\hyperref[VI.prop7.2.7]{7.2.7}). Il lui correspond un foncteur
image directe $m^{\sim/I}_{\ast} : (C^{\sim /I})'\to (D^{\sim/I})'$ (\hyperref[VI.subsubsec7.1.7]{7.1.7})
qui fournit, en passant aux sections sur $I^{\circ}$, un foncteur:
\begin{equation*}
\Hom_{I^{\circ}} (I^{\circ}, m^{\sim/I}_{\ast}) : \Hhom_{I^{\circ}}
(I^{\circ}, (C^{\sim/I})') \to \Hhom_{I^{\circ}} (I^{\circ},
(D^{\sim/I})'){\quoi}.\tag{7.4.9.1}
\end{equation*}\etiquette[7.4.9.1]{VI.eq7.4.9.1}

Par ailleurs, le foncteur $m$ est un foncteur continu entre les
sites totaux (\hyperref[VI.prop7.4.4]{7.4.4}) et par suite fournit, par
composition, un foncteur $m_{\ast} : \Top(C) \to \Top(D)$.

En résumé, on a un diagramme de foncteurs:
\[\vcenter{
\xymatrix@C=2.75cm@R=1.5cm{
\Top(C) \ar[r]^{\Theta_C} \ar[d]_{m_{\ast}}& \Hhom_{I^{\circ}}(I^{\circ},(C^{\sim/I})') \ar[d]_{\Hom_{I^\circ}(I^{\circ},m_{\ast}^{\sim/I})}\\
\Top(D)  \ar[r]^{\theta_D} & \Hhom_I^{\circ}(I^{\circ},(D^{\sim/I})')
}}\leqno{(7.4.9.2)}
\]\etiquette[7.4.9.2]{VI.eq7.4.9.2}

\begin{enonce*}{Proposition 7.4.10}\etiquette[7.4.10]{VI.prop7.4.10}
Le diagramme\pageoriginale \eqhyperref[VI.eq7.4.9.2]{7.4.9.2} est commutatif à isomorphisme
près. Lorsque $m$
est cartésien, $m$ est un morphisme du site total $C$ dans le site
total $D$.
\end{enonce*}

Nous laissons au lecteur le soin de vérifier la commutativité à
isomorphisme près du diagramme \eqhyperref[VI.eq7.4.9.2]{7.4.9.2}. Pour montrer que
$m$ est un morphisme entre les sites totaux, il suffit, compte tenu
de \hyperref[VI.prop7.4.7]{7.4.7}, de montrer que le foncteur adjoint à gauche
au foncteur $\Hom_{I^{\circ}}(I^{\circ}, m^{\sim/I}_{\ast})$ est
exact à gauche. Nous allons tout d'abord décrire le
foncteur adjoint à gauche à $\Hom_{I^{\circ}}(I^{\circ}, m^{\sim/I}_{\ast})$. Notons $(m^{\sim/I})^{\ast} : D^{\sim/I} \to C^{\sim/I}$ le foncteur image
inverse du morphisme $m^{\sim/I}$ (\hyperref[VI.prop7.2.7]{7.2.7}), et pour tout
objet $i$ de $I$, notons $m^{\ast}_i : D^{\sim}_i \to C^{\sim}_i$ le
foncteur induit par $(m^{\sim/I})^{\ast}$ sur les fibres en $i$.
Choisissons pour tout morphisme $f : i \to j$ de $I$ des morphismes
de topos, notés dans les deux cas $(f^{\ast}, f_{\ast}) : C^{\sim}_i
\to C^{\sim}_j$ et $(f^{\ast}, f_{\ast}) : D^{\sim}_i \to
D^{\sim}_j$. Comme $(m^{\sim/I})^{\ast}$ est un $I^{\circ}$-foncteur
cartésien, on a, pour tout $f : i \to j$, un isomorphisme canonique
\begin{equation*}
m^{\ast}_i f^{\ast} \to f^{\ast} m^{\ast}_j{\quoi},\tag{7.4.10.1}
\end{equation*}\etiquette[7.4.10.1]{VI.eq7.4.10.1}
et comme les foncteurs $f^{\ast}$ et $f_{\ast}$ sont adjoints, on a
des morphismes canoniques (morphismes d'adjonction)
\begin{align*}
& \phi : \id \to f_{\ast} f^{\ast}\\
& \psi : f^{\ast} f_{\ast} \to \id{\quoi}. \tag{7.4.10.2}
\end{align*}\etiquette[7.4.10.2]{VI.eq7.4.10.2}
Considérons alors la suite de morphismes fonctoriels:
\begin{equation*}
m^{\ast}_j f^{\ast} \xrightarrow{(1)} f_{\ast} f^{\ast} m^{\ast}_j
f_{\ast} \xrightarrow{(2)} f_{\ast} m^{\ast}_i f^{\ast} f_{\ast}
\xrightarrow{(3)} f_{\ast} m^{\ast}_i{\quoi},\tag{7.4.10.3}
\end{equation*}\etiquette[7.4.10.3]{VI.eq7.4.10.3}
où le morphisme (1) est déduit de $\phi$, le morphisme (2) est déduit
de l'isomorphisme \eqhyperref[VI.eq7.4.10.1]{7.4.10.1} et le morphisme (3) est déduit de
$\psi$. En composant les différents morphismes\pageoriginale de la
suite \eqhyperref[VI.eq7.4.10.3]{7.4.10.3},
on obtient un morphisme fonctoriel:
\begin{equation*}
\can_f : m^{\ast}_j f_{\ast} \to f_{\ast} m^{\ast}_i{\quoi}.\tag{7.4.10.4}
\end{equation*}\etiquette[7.4.10.4]{VI.eq7.4.10.4}
Soit alors $\tau$ un objet de $\Hhom_{I^{\circ}} (I^{\circ},
(D^{\sim/I})')$. On a donc, pour tout objet $i$ de $I$ un objet
$\tau(i)$ de $D^{\sim}_i$ et pour tout morphisme $f : i \to j$, un
morphisme
\begin{equation*}
\gamma_{f} (\tau) : \tau(j) \to f_{\ast} \tau(i){\quoi},\tag{7.4.10.5}
\end{equation*}\etiquette[7.4.10.5]{VI.eq7.4.10.5}
où la famille des $\gamma_f(\tau)$ possède une propriété de
compatibilité analogue à \eqhyperref[VI.eq7.4.5.8]{7.4.5.8}. Posons alors $\sigma(i) =
m^{\ast}_i (\tau(i))$ et notons
\begin{equation*}
\gamma_{f} (\sigma) : \sigma(j) \to f_{\ast} \sigma(i)\tag{7.4.10.6}
\end{equation*}\etiquette[7.4.10.6]{VI.eq7.4.10.6}
le morphisme composé $m^{\ast}_j \tau(j) \xrightarrow{m_j^{\ast}
  \gamma_f(\tau)} m_j^{\ast} f_{\ast} \tau(i) \xrightarrow{\can_f(\tau(i))}
f_{\ast} m^{\ast}_i \tau(i)$. On vérifie
  que les $\gamma_f(\sigma)$ possèdent la propriété de compatibilité
  de \eqhyperref[VI.eq7.4.5.8]{7.4.5.8} et par suite qu'on a défini ainsi un objet de
  $\Hhom_{I^{\circ}}(I^{\circ}, (C^{\sim/I})')$. Un « adjoint functor
  chasing » montre alors que le foncteur de
  $\Hhom_{I^{\circ}}(I^{\circ}, (D^{\sim/I})')$ dans
  $\Hhom_{I^{\circ}}(I^{\circ}, (C^{\sim/I})')$ qu'on vient de
  construire est adjoint à gauche au foncteur
  $\Hom_{I^{\circ}}(I^{\circ}, m^{\sim/I}_{\ast})$. Reste à montrer
  que cet adjoint à gauche commute aux limites projectives finies, ce
  qui résulte immédiatement du fait que dans les catégories de
  sections les limites projectives se calculent fibre par fibre et que
  tous les foncteurs utilisés pour construire cet adjoint à gauche
  commutent aux limites projectives finies.

\begin{enonce*}{Corollaire 7.4.11}\etiquette[7.4.11]{VI.cor7.4.11}
Soient $p : C \to I$ un $𝒰$-site fibré sur une petite catégorie,
$p^{\sim/I} : C^{\sim/I} \to I$ le $𝒰$-topos fibré associé (\hyperref[VI.subsubsec7.2.6]{7.2.6}),
$\epsilon_I : C \to C^{\sim/I}$ le morphisme canonique du $𝒰$-site
fibré $C^{\sim/I}$ dans le $𝒰$-site fibré $C$. Le morphisme $\epsilon_{I}$ induit une équivalence
$\Top(\epsilon_I) : \Top(C) \to \Top(C^{\sim/I})$ entre les topos
totaux correspondants
\end{enonce*}

Résulte de la commutativité du diagramme \eqhyperref[VI.eq7.4.9.2]{7.4.9.2}
et du fait que le morphisme $\epsilon_{I}$ fournit une équivalence
entre les $𝒰$-topos fibrés associés.

\begin{enonce*}[remark]{Remarque 7.4.11.1}\etiquette[7.4.11.1]{VI.rem7.4.11.1}
Il résulte de \hyperref[VI.cor7.4.11]{7.4.11} que dans les questions concernant
le topos total on peut remplacer les sites fibrés par les topos
fibrés correspondants.
\end{enonce*}

\begin{enonce*}{Lemme 7.4.12}\etiquette[7.4.12]{VI.lem7.4.12}
Soient\pageoriginale $p : C \to I$ un site fibré sur une petite
catégorie $I$ et $i$ un objet final de $I$ \emph{(I 10)}. Le
foncteur $\alpha_{i!} : C^{\sim}_i \to \Top(C)$ (7.4.3.4) est exact
et pleinement fidèle. Les couples de foncteurs $\beta_i =
(\alpha_{i!}, \alpha^{\ast}_i) : \Top(C) \to C^{\sim}_i$ et
$\alpha_i = (\alpha_{i}^{\ast}, \alpha_{i
  \ast}) : C^{\sim}_i \to \Top(C)$ sont des morphismes de topos. Le
morphisme composé $\beta_i \alpha_i : C^{\sim}_i \to C^{\sim}_i$ est
isomorphe au morphisme identique.
\end{enonce*}

Pour tout objet $j$ de $I$, notons $f_j : j \to i$ l'unique
morphisme de $j$ dans $i$. En utilisant l'équivalence $\Theta_C :
\Top(C) \xrightarrow{\sim} \Hhom_{I^{\circ}} (I^{\circ},
(C^{\sim/I})')$ (\hyperref[VI.prop7.4.7]{7.4.7}) on voit que pour tout objet $X$
de $C_i$, $\alpha_{i!}(X)$ est la section $j \mapsto f^{\ast}_j(X)$ et
que pour toute section $j \mapsto \sigma(j) \in \Hom_{I^{\circ}}
(I^{\circ}, (C^{\sim/I})')$, $\alpha^{\ast}_i (j \mapsto \sigma(j))$
est  l'objet $\sigma(i)$. Par suite $\alpha^{\ast}_{i} \alpha_{i!}$
est isomorphe à l'identité. Donc $\alpha_{i!}$ est pleinement
fidèle. Comme les foncteurs $f^{\ast}_{j}$ sont exacts à gauche, le
foncteur $\alpha_{i!}$ est exact à gauche. Donc $\beta_i$ est un
morphisme de topos, et comme $\alpha^{\ast}_{i} \alpha_{i!}$ est
isomorphe à l'identité, le morphisme $\beta_i \alpha_i$ est
isomorphe au morphisme identique.


\begin{enonce*}{Théorème 7.4.13.1}\etiquette[7.4.13.1]{VI.thm7.4.13.1}
Soient $p : C \to I$ et $q : D \to I$ deux $𝒰$-sites fibrés sur une
petite catégorie $I$ et $m : D \to C$ un $I$-foncteur. Pour que
$m$ soit un morphisme du site fibré $C$ dans le site fibré $D$
\emph{(\hyperref[VI.subsubsec7.2.2]{7.2.2})}, il suffit que $m$ soit un morphisme du site total $C$
dans le site total $D$.
\end{enonce*}

\setcounter{subsubsection}{13}
\setcounter{paragraph}{1}
\paragraph{}\etiquette[7.4.13.2]{VI.par7.4.13.2}

Il faut montrer que si un $I$-foncteur $m : D \to C$ est un morphisme
du site total $C$ dans le site total $D$, le foncteur $m$ est un
morphisme du site fibré $C$ dans le site fibré $D$, \ie pour tout
objet $i$ de $I$, le foncteur $m_i : D_i \to C_i$ induit par $m$ sur
les fibres en $i$ est un morphisme du site $C_i$ dans le site
$D_i$. La démonstration de cette dernière assertion occupe les alinéas
\hyperref[VI.par7.4.13.3]{7.4.13.3} à \hyperref[VI.par7.4.13.7]{7.4.13.7}.

\paragraph{}\etiquette[7.4.13.3]{VI.par7.4.13.3}
Montrons que pour tout objet $i$ de $I$, le foncteur $m_i : D_i \to
C_i$ est continu. En effet on a un diagramme commutatif :\etiquette[133]{page133}
\[
\xymatrix@C=2cm@R=1.5cm{
D_i \ar[r]^{m_i} \ar[d]_{\alpha_{i!}}& C_i \ar[d]^{\alpha_{i!}}\\
D   \ar[r]^{m}& C
}
\]
et il\pageoriginale suffit de montrer que pour tout crible couvrant $R
\hookrightarrow X$ de $C_i / X$, le morphisme $m_{i!} (R) \to m_i(X)$
est bicouvrant (III 1.2 et 1.5). Mais, en vertu de la commutativité du
diagramme de la page \pagehyperref[page133]{133} et du fait que les foncteurs $\alpha_{i!}$ et
$m$ sont continus, le morphisme $\alpha_{i!} m_{i!} (R) \to
\alpha_{i!} m_{i} (X)$ est bicouvrant. L'assertion résulte donc de
\hyperref[VI.lem7.4.2.2]{7.4.2.2}.

\paragraph{}\etiquette[7.4.13.4]{VI.par7.4.13.4} \emph{Réduction au cas où $C$ et $D$ sont des
  $𝒰$-topos fibrés.} Comme les foncteurs $m_i$ sont continus, ils admettent des
prolongements naturels aux catégories de faisceaux que nous noterons
$m^{\ast}_{i}: D^{\sim}_{i}\to C^{\sim}_{i}$ (III 1.2 iv)). Soient
$D^{\sim/I}$ et $C^{\sim/I}$ les $𝒰$-topos fibrés associés aux sites fibrés $D$ et $C$
respectivement, et pour tout morphisme $f$ de $I$, notons $f^{\ast}$
les foncteurs images inverses pour les quatre catégories fibrées $C$,
$D$, $C^{\sim /I}$ et $D^{\sim/I}$. Comme le foncteur $m$ est un
$I$-foncteur, on a pour tout $f : i \to j$ un morphisme canonique
$$
f^{\ast} mj \leftarrow m_i f^{\ast}{\quoi}.
$$
Comme les foncteurs $f^{\ast}$, $m_i$ et $m_j$ sont continus, on en
déduit, en passant aux catégories de faisceaux, des isomorphismes
canoniques
$$
f^{\ast} m_j^{\ast} \longleftarrow m_i^{\ast} f^{\ast}{\quoi}.
$$
Ces isomorphismes possèdent des propriétés de compatibilité,
permettant de construire un foncteur cartésien $m^{\ast} : D^{\sim/I}
\to C^{\sim /I}$ qui induit sur les topos fibres les foncteurs
$m^{\ast}_i$. De plus, le diagramme
\[
\xymatrix@C=2.5cm@R=1.5cm{
D \ar[r]^{m} \ar[d]_{\varepsilon_I}& C \ar[d]^{\varepsilon_I}\\
D^{\sim/I}   \ar[r]^{m^{\ast}}& C^{\sim/I}
}
\]
est commutatif à isomorphisme près. Comme les foncteurs $\epsilon_I$
induisent des équivalences sur les topos totaux (\hyperref[VI.cor7.4.11]{7.4.11}) et comme $m$
est un morphisme entre les sites totaux, le foncteur $m^{\ast} :
D^{\sim/I} \to C^{\sim/I}$ est un morphisme entre les sites totaux. On
est donc ramené à démontrer l'assertion de \hyperref[VI.par7.4.13.2]{7.4.13.2} lorsque $D$ et $C$
sont des $𝒰$-topos fibrés, ce que nous supposerons désormais.

\paragraph{}\etiquette[7.4.13.5]{VI.par7.4.13.5}
\emph{Pour tout\pageoriginale objet $i$ de $I$, le foncteur $m_i : D_i
  \to C_i$ transforme
l'objet final de $D_i$ en l'objet final de $C_{i}$}. Notons $m :
  D^{\sim} \to C^{\sim}$ le
prolongement naturel de $m$ aux topos totaux (III 1.2). Utilisant les
équivalences $\Theta_D :D^{\sim} \to \Hhom_{I^{\circ}} (I^{\circ}, D')$
(\hyperref[VI.prop7.4.7]{7.4.7}), on vérifie immédiatement que $m$ associe à toute section
$\sigma \in \Hom_{I^{\circ}}({I^{\circ}}, D')$ la section $(i \mapsto
m_i \sigma(i)) \in \Hom_{I^{\circ}} ({I^{\circ}}, C')$. Comme $m$ est
un morphisme de sites, le foncteur $m$ est exact à gauche et en
particulier transforme l'objet final de $\Hhom_{I^{\circ}}
({I^{\circ}}, D')$ en l'objet final de $\Hhom_{I^{\circ}}({I^{\circ}},
C')$. Pour tout objet $i$ de $I$, notons $e_i$ un objet final de
$D_i$. Un objet final de $\Hhom_{I^{\circ}}({I^{\circ}}, D')$ est la
section $i \mapsto e_i$. Donc la section $i \mapsto m_i (e_i)$ est un
objet final de $\Hhom_{I^{\circ}}({I^{\circ}}, C')$ et par suite, pour
tout objet $i$ de $I$, $m_i(e_i)$ est un objet final de $C_i$.

\paragraph{}\etiquette[7.4.13.6]{VI.par7.4.13.6}
\emph{Réduction au cas où $i$ est un objet final de} $I$. Soit $i$ un objet
de $I$. Le foncteur $\alpha_{i!} : D_i \to D$ se factorise en
$$
D_i \xrightarrow{\bar{\alpha}_{i!}} D_{/e_i} \xrightarrow{\loc(e_i)} D{\quoi},
$$
où $\bar{\alpha}_{i}$ est le foncteur évident et $\loc(e_i)$ le
  foncteur d'oubli. Notons $m/i : D_{/e_i} \to C_{/m(e_i)}$ le foncteur
  qui associe à tout objet $u : X \to e_i$ de $D_{/e_i}$, l'objet
  $m(u) : m(X) \to m(e_i)$ de $C_{/m(e_i)}$. On a un diagramme
  commutatif:
\[
\xymatrix@C=2.25cm@R=1.5cm{
D_i \ar[r]^{m_i} \ar[d]_{\bar{\alpha}_{i!}}& C_i\ar[d]^{\bar{\alpha}_{i!}}\\
D_{/e_i}  \ar[r]^{m/i} & C_{/m(e_i)}
}
\]
Comme $e_i$ et $m(e_i)$ sont des objets finaux de $D_i$ et $C_i$
respectivement (\hyperref[VI.par7.4.13.5]{7.4.13.5}), les catégories $D_{/e_{i}}$
et $C_{/m(e_i)}$ sont des $𝒰$-topos fibrés sur $I_{/i}$ et
le foncteur $m_{/i}$ est un $I_{/i}$-foncteur. Comme la propriété
pour un foncteur d'être un morphisme de sites se localise (IV 4.9),
le foncteur $m_{/i}$ est un morphisme du site total $C_{/m(e_i)}$
dans le site total $D_{/(e_i)}$. On est donc ramené à démontrer que
$m_i$ est un morphisme de topos lorsque $i$ est un objet final de
$I$, ce que nous supposerons désormais.\pageoriginale

\paragraph{}\etiquette[7.4.13.7]{VI.par7.4.13.7}
\emph{Fin de la démonstration}. Notons $\alpha_{i!} : D_i \to
D^{\sim}$, $m^{\sim} : D^{\sim} \to C^{\sim}$, $\alpha_{i!} : C_i \to
C^{\sim}$, les prolongements naturels aux catégories de faisceaux des
  foncteurs $\alpha_{i!}$ et $m$. On a un diagramme commutatif à
  isomorphisme près:
\[
\xymatrix@C=2.5cm@R=1.5cm{
D_i \ar[r]^{m_i} \ar[d]_{\alpha_{i!}}& C_i\ar[d]^{\alpha_{i!}}\\
D^{\sim}  \ar[r]^{m^{\sim}} & C .
}
\]
Comme $i$ est un objet final, le foncteur $\alpha^{\ast}_i \alpha_{i!} :
C_i \to C_i$ est isomorphe à l'identité (\hyperref[VI.lem7.4.12]{7.4.12}) et par suite $m_i$
est isomorphe à $\alpha^{\ast}_i m^{\sim} \alpha_{i!}$. De plus les
foncteurs $\alpha^{\ast}_i$, $m^{\sim}$ et $\alpha^{\ast}_i$ sont exacts à
gauche (\hyperref[VI.lem7.4.12]{7.4.12}). Par suite $m_i$ est exact à gauche et est donc un
morphisme du site $C_i$ dans le site $D_i$.

\begin{enonce*}[remark]{Exercice 7.4.14}\etiquette[7.4.14]{VI.exe7.4.14}
Soient $X = (X_i)_{i \in I}$ un topos fibré sur $I$ et $T$ un
topos. Montrer que les catégories $\Hhomtop(X_i, T)$ sont les fibres
d'une catégorie fibrée qu'on notera $\Hhomtop_I(X, T\times
I)$. Montrer que la catégorie des sections sur $I$ de $\Hhomtop_I(X,
T\times I)$ est canoniquement équivalente à $\Hhomtop(\Top(X), T)$
où $\Top(X)$ est le topos total de $X$.
\end{enonce*}

\begin{enonce*}[remark]{Exercice 7.4.15}\etiquette[7.4.15]{VI.exe7.4.15}
Soit $X = (X_i)_{i \in I}$ un topos fibré sur $I$. Définir un
morphisme cartésien $m$ de $X$ dans le topos fibré constant de fibre
$(\Ens)$. Montrer que le topos total du topos fibré constant de fibre
$(\Ens)$ est canoniquement équivalent à
$I^{\widehat{\phantom{,}}}$. En déduire un morphisme de topos
$$
\Top(m) :\Top(X) \to I^{\widehat{\phantom{,}}}{\quoi}.
$$
Soit $F = (F_i)_{i \in I}$ (où $F_i = F | X_i$) un objet de $\Top(X)$
(\hyperref[VI.prop7.4.7]{7.4.7}). Montrer que
$$
\Top(m)(F) = (i \mapsto \Gamma (X_i, F_i)){\quoi}.
$$
Déduire de \hyperref[VI.lem7.4.12]{7.4.12} que si $F$ est abélien injectif, $F_i$ est abélien
injectif pour tout $i \in \ob I$. En déduire que
$$
R^q \Top(m)(F) = (i \mapsto H^{q} (X_i, F_i)){\quoi}.
$$
Montrer\pageoriginale que la suite spectrale de Cartan-Leray du
morphisme $m$ (\hyperrefext[V][V.sec5]{5}) est:
$$
H^{p+q} (\Top(X), F) \Leftarrow {\varprojlim_{I}}^{(p)} H^{q}
(X_i, F|X_i){\quoi}.
$$
\end{enonce*}

\section{Limites projectives de topos fibrés}\etiquette[8]{VI.sec8}

\subsection{Généralités}\etiquette[8.1]{VI.subsec8.1}

\begin{enonce*}{Définition 8.1.1}\etiquette[8.1.1]{VI.def8.1.1}
Soient $𝒰$ un univers et $p:F \to I$ un $𝒰$-topos
fibré. Un couple $(C, m)$ constitué par un $𝒰$-topos $C$ et
un morphisme cartésien $m : C \times I \to F$ de $𝒰$-topos
fibrés sur $I$ est appelé une limite projective du topos fibré $F$
si pour tout $𝒰$-topos $D$, le foncteur
\begin{equation*}
\Hhomtop(D, C) \to \Hhomtop_{\cart/I}(D\times I, F)\tag{8.1.1.1}
\end{equation*}\etiquette[8.1.1.1]{VI.eq8.1.1.1}
obtenu en composant le foncteur de changement de base
$$
\Hhomtop(D, C) \to \Hhomtop_{I}(D\times I, C \times I)
$$
et le foncteur de composition avec $m$ (\hyperref[VI.subsubsec7.1.8]{7.1.8}), est une équivalence de
catégories. \emph{(\cf \hyperref[VI.par8.1.3.3]{8.1.3.3} pour une formulation plus «   géométrique »)}.
\end{enonce*}

\setcounter{subsubsection}{1}
\subsubsection{}\etiquette[8.1.2]{VI.subsubsec8.1.2}

On utilise les notations de \hyperref[VI.def8.1.1]{8.1.1}. Soient $D$ un $𝒰$-topos et $m_1 : D
\times I \to F$ un morphisme de $𝒰$-topos fibrés. En traduisant la
définition \hyperref[VI.def8.1.1]{8.1.1}, on constate aussitôt qu'il existe un morphisme de
topos $n : D \to C$, et un isomorphisme de morphismes de topos fibrés
sur $I$, $\gamma : m_{1} \xrightarrow{\sim} m_{\circ} (n \times \id_I)$,
rendant commutatif le diagramme:
\[\vcenter{
\xymatrix@C=.75cm@R=.5cm{
D \times I  \ar[rrrr]^{m_1} \ar[dddd]_{m\times \id_I} &&{\phantom{.}}\ar@/_.5pc/@{-->}[dr]_{\gamma}  &&F\\
                         & &{\phantom{a}}& &\\
                         & && & \\
                         & && & \\
C \times I \ar[rrrruuuu]_{m_1}&& &&{\phantom{a}}\\
}}\leqno{(8.1.2.1)}
\]\etiquette[8.1.2.1]{VI.eq8.1.2.1}
De plus, si $(n', \gamma')$ est un autre couple possédant la
propriété ci-dessus, il existe un \emph{unique} isomorphisme $\eta : n
\xrightarrow{\sim} n'$ tel que
\begin{equation*}
\gamma \circ m (\eta \times \id) = \gamma'{\quoi}.\tag{8.1.2.2}
\end{equation*}\etiquette[8.1.2.2]{VI.eq8.1.2.2}
De\pageoriginale là résulte que, si $(D, m_1)$ est une limite
projective de $F$, $n$ est une équivalence de $𝒰$-topos
(IV 3.4) et que la donnée
supplémentaire de l'isomorphisme $\gamma$ faisant commuter le
diagramme \eqhyperref[VI.eq8.1.2.1]{8.1.2.1} détermine le couple $(n, \gamma)$
isomorphisme
canonique près. L'existence de la limite projective sera
démontrée en
\hyperref[VI.thm8.2.3]{8.2.3} lorsque $I$ est cofiltrante.

\subsubsection{}\etiquette[8.1.3]{VI.subsubsec8.1.3}
On note
$$
\limtop{I} F \mbox{ ou } \limtop{I}
F_i \mbox{ ou encore } \underleftarrow{F}
$$
une limite projective du topos fibré $F$, $\mu :
(\limtop{I} F) \times I \to F$ le
morphisme canonique, et pour tout objet $i$ de $I$ on note
$$
\mu_i : \underleftarrow{F} \simeq \underleftarrow{F} \times \{i\} \to F_{i}
$$
le morphisme de topos induit par $\mu$ sur les fibres en $i$.

Choisissons un biscindage de $F$ (\hyperref[VI.subsubsec7.1.4]{7.1.4}) \ie pour tout $f : i \to j$
un morphisme de topos $f\bigdot = (f^{\ast}, f_{\ast}) : F_i \to F_j$ tel
que $f^{\ast}$ soit un foncteur image inverse pour la structure
fibrée. On a alors, pour tout couple $i \xrightarrow{f} j
\xrightarrow{g} k$ de morphismes composables de $I$, un isomorphisme de
morphismes de topos (\hyperref[VI.subsubsec7.1.3]{7.1.3})
$$
c_{f, g} : g.f. \simeq (gf).\quad {\quoi}.
$$
Comme $\mu$ est un morphisme cartésien de topos fibrés, on a pour tout
morphisme $f : i \to j$ de $I$, un isomorphisme de morphismes de topos
(\hyperref[VI.subsubsec7.1.6]{7.1.6})
\begin{equation*}
b_f  : f. \mu_i \xrightarrow{\sim} \mu_j{\quoi},\tag{8.1.3.1}
\end{equation*}\etiquette[8.1.3.1.]{VI.eq8.1.3.1}

et la famille des $b_f$ satisfait à des relations analogues aux
relations \eqhyperref[VI.eq7.1.6.2]{7.1.6.2}

\setcounter{paragraph}{1}
\paragraph{}\etiquette[8.1.3.2]{VI.par8.1.3.2}

Soient $D$ un $𝒰$-topos, $m : D \times I \to F$ un morphisme de topos
fibrés, $m_i : D = D \times \{i\} \to F_i$ les morphismes induits par
$m$ sur les fibres et pour tout $f : i \to j$:
$$
b'_f : f . m_i \xrightarrow{\sim} m_j{\quoi},
$$
l'isomorphisme de transition. Il résulte de \hyperref[VI.subsubsec8.1.2]{8.1.2} qu'il existe un
morphisme de topos $n : D \to
\underleftarrow{F}$\pageoriginale et une famille d'isomorphismes:
$$
\gamma_i : m_i \xrightarrow{\sim} \mu_i n \quad i \in \ob I {\quoi},
$$
tels que pour tout $f : i \to j$ on ait :
$$
b_f \circ f. (\gamma_i) = \gamma_j \circ b'_f{\quoi}.
$$
De plus, le morphisme $n$ et la famille des $\gamma_i$, $i \in ob (I)$
sont déterminés à isomorphisme unique près (\hyperref[VI.subsubsec8.1.2]{8.1.2}).

Réciproquement lorsqu'on se donne un topos $\underleftarrow{F}$, des
morphismes $\mu_i : \underleftarrow{F} \to F_i$ et des isomorphismes $b_f$
\eqhyperref[VI.eq8.1.3.1]{8.1.3.1.}, satisfaisant aux conditions de compatibilités usuelles, et
lorsque toutes ces données possèdent la propriété universelle décrite
ci-dessus, il existe un unique morphisme $\mu : \underleftarrow{F} \times
I \to F$ de topos fibrés qui donnent naissance aux $\mu_{i}$ et aux
$b_f$ et $(\underleftarrow{F}, \mu)$ est une limite projective du topos
fibré $F$.

\paragraph{}\etiquette[8.1.3.3]{VI.par8.1.3.3}

Soit $𝒱$ un univers tel que $\sheaf{\Uu} \in
𝒱$. Les considérations
ci-dessus permettent donc d'interpréter les limites projectives de
topos fibrés comme des « limites projectives », au sens des
2-catégories, des pseudo-foncteurs à valeurs dans la 2-catégorie des
$𝒰$-topos fibrés éléments de $𝒱$ déduits de la
structure fibrée en choisissant les morphismes $f.$. Toutefois le lecteur prendra garde de
ne pas confondre cette notion de « limite projective » au sens des
2-catégories (avec isomorphisme de commutation) avec la notion
stricte (ou naïve) de limite projective (commutation stricte de
diagrammes). Même lorsque la notion stricte a un sens (lorsque le
pseudo-foncteur est un vrai foncteur) les notions de « limite
projective » généralisée et de limite projective stricte ne
coïncident pas en général.

\subsubsection{}\etiquette[8.1.4]{VI.subsubsec8.1.4}

Soient $p : F \to I$ et $q : G \to I$ deux $𝒰$-topos fibrés et $m : F
\to G$ un morphisme de topos fibrés. Supposons que les topos fibrés
$F$ et $G$ possèdent des limites projectives (\hyperref[VI.def8.1.1]{8.1.1}). Il résulte alors
immédiatement de la définition \hyperref[VI.def8.1.1]{8.1.1} qu'il existe un morphisme de
topos:
\begin{equation*}
\underleftarrow{m} : \underleftarrow{F} \to
\underleftarrow{G}{\quoi},\tag{8.1.4.1}
\end{equation*}\etiquette[8.1.4.1]{VI.eq8.1.4.1}
et un isomorphisme de morphismes de topos fibrés $\gamma : m \circ \mu
\xrightarrow{\sim} \mu \circ \underleftarrow{m} \times \id_{I}$, rendant
commutatif\pageoriginale le diagramme
\[\vcenter{
  \xymatrix@C=.6cm@R=.45cm{
\underleftarrow{F}\times I \ar[rrr]^{\mu}
\ar[ddd]_{\underleftarrow{m}\times \id_I}   && & F\ar[ddd]^{m}\\
                          &            & &\\
                          &            & &\ar@/_.55pc/@{-->}[ld]_{\gamma} \\
\underleftarrow{G}\times I \ar[rrr]^{\mu} &       & &  G&.
}}\leqno{(8.1.4.2)}
\]\etiquette[8.1.4.2]{VI.eq8.1.4.2}
De plus, si $(\underleftarrow{m}', \gamma')$ est un autre couple possédant
la propriété ci-dessus il existe un \emph{unique} isomorphisme $\eta:
\underleftarrow{m} \xrightarrow{\sim} \underleftarrow{m}'$ de morphismes de
topos tel que
$$
\gamma \circ \mu (\eta \times \id) = \gamma'{\quoi}.
$$

\subsubsection{}\etiquette[8.1.5]{VI.subsubsec8.1.5}

Soient $p : F \to I$ un $𝒰$-topos fibrés élément de
$𝒱$, $\phi : I'\to I$ un foncteur élément de $𝒱$, $p_{\phi} : F_{\phi} \to I'$ le
$𝒰$-topos fibré déduit de $(F, p)$ par le changement de base $\phi$
(\hyperref[VI.subsubsec7.1.9]{7.1.9}). Soit $(\underleftarrow{F}, \mu)$ une limite projective de $F$. On
obtient par changement de base un morphisme de topos fibrés sur $I'$:
$$
\mu' : \underleftarrow{F} \times I' \to F{\quoi}.
$$
Soit $(\underleftarrow{F}_{\phi}, \mu_{\phi})$ une limite projective de
$F_{\phi}$. On a alors, par la propriété universelle de la limite
projective (\hyperref[VI.subsubsec8.1.2]{8.1.2}) un morphisme de topos
$$
m_{\phi} : \underleftarrow{F} \to \underleftarrow{F}_{\phi}{\quoi},
$$
et un isomorphisme de morphismes de topos fibrés sur $I'$:
$$
\gamma : \mu'\xrightarrow{\sim} \mu_{\phi} \circ (m_{\phi} \times \id_{I'}){\quoi}.
$$
Le couple $(m_{\phi}, \gamma)$ est déterminé à isomorphisme unique
près.

\subsubsection{}\etiquette[8.1.6]{VI.subsubsec8.1.6}

Soient $p : F \to I$ un $𝒰$-site fibré sur une petite catégorie $I$,
$F^{\sim/I} \to I$ le $𝒰$-topos fibré associé (\hyperref[VI.subsubsec7.2.6]{7.2.6}) et $D$ un
$𝒰$-topos. On a tout d'abord un foncteur
\begin{equation*}
\Hhomtop_{\cart / I} (D \times I, F^{\sim/I}) \to
\Mmorsite_{\cart/I}(D \times I, F){\quoi},\tag{8.1.6.1}
\end{equation*}\etiquette[8.1.6.1]{VI.eq8.1.6.1}
qui associe à tout $m$ le foncteur composé de $m^{\ast}$ avec
$\epsilon_{F /I}$ (\hyperref[VI.prop7.2.7]{7.2.7}). De plus, par définition de la catégorie
$\Mmorsite_{\cart /I} (D \times I, F)$ \eqhyperref[VI.eq7.2.2.2]{7.2.2.2}, on a un foncteur
\begin{equation*}
\Mmorsite_{\cart/I}(D \times I, F) \to \Hhom_{\cart / I} (F, D \times
I)^{\circ}{\quoi}.\tag{8.1.6.2}
\end{equation*}\etiquette[8.1.6.2]{VI.eq8.1.6.2}
Enfin,\pageoriginale par définition de la limite inductive (\hyperref[VI.def6.3]{6.3}), on a
un foncteur
\begin{equation*}
\Hhom_{\cart / I} (F, D \times I)^{\circ} \to
\Hhom(\varinjlim_{I^{\circ}}F, D)^{\circ}{\quoi}.\tag{8.1.6.3}
\end{equation*}\etiquette[8.1.6.3]{VI.eq8.1.6.3}
Notons
\begin{equation*}
\Theta : \Hhomtop_{\cart / I}(D \times I, F^{\sim / I}) \to
\Hom(\varinjlim_{I^{\circ}}F, D)^{\circ}\tag{8.1.6.4}
\end{equation*}\etiquette[8.1.6.4]{VI.eq8.1.6.4}
le foncteur composé de ces trois derniers foncteurs et
\begin{equation*}
\Pi : F \to \varinjlim_{I^{\circ}}F{\quoi},\tag{8.1.6.5}
\end{equation*}\etiquette[8.1.6.5]{VI.eq8.1.6.5}
le foncteur canonique.

\begin{enonce*}{Proposition 8.1.7}\etiquette[8.1.7]{VI.prop8.1.7}
Le foncteur $\Theta$ est pleinement fidèle. Son image essentielle est
l'ensemble des foncteurs $m : \varinjlim_{I^{\circ}}F \to D$ tel que
le foncteur
$$
(m \circ \Pi, p) : F \to D \times I
$$
soit un morphisme du site total $D \times I$ dans le site total $F$.
\end{enonce*}

Il résulte de \hyperref[VI.prop7.2.7]{7.2.7} que le foncteur \eqhyperref[VI.eq8.1.6.1]{8.1.6.1} est une équivalence de
catégories, de \hyperref[VI.subsubsec7.2.2]{7.2.2} que le foncteur (8.1.6.2) est pleinement fidèle
et de \hyperref[VI.def6.3]{6.3} que le foncteur \eqhyperref[VI.eq8.1.6.4]{8.1.6.4} est une équivalence de
catégories. Par suite $\Theta$ est pleinement fidèle. Un foncteur $m :
\varinjlim_{I^{\circ}}F \to D$ est isomorphe à l'image
par $\Hhom_{\cart / I} (F, D \times I)^{\circ} \to \Hhom(
\varinjlim_{I^{\circ}}F, D)^{\circ}$ du foncteur $(m \circ
\Pi, p)$, et un tel foncteur est dans l'image essentielle de
\eqhyperref[VI.eq8.1.6.2]{8.1.6.2} si et seulement s'il est un morphisme du site total $D
\times I$ dans le site total $F$ (\hyperref[VI.thm7.4.13.1]{7.4.13.1}); d'où la proposition.

\subsection{Construction de la limite projective lorsque la catégorie
  d'indice est cofiltrante.}\etiquette[8.2]{VI.subsec8.2}

\begin{enonce*}{Lemme 8.2.1}\etiquette[8.2.1]{VI.lem8.2.1}
Soient $I$ et $I'$ deux catégories cofiltrantes \emph{(I 8.1)} et $\phi :
I'\to I$ un foncteur tel que $\phi^{\circ} : {I'}^{\circ} \to
I^{\circ}$ soit cofinal \emph{(I 8.1)}. Soient de plus $p : F \to I$ un $𝒰$-topos fibré, $D$ un
$𝒰$-topos, $m : D \times I \to F$ un morphisme cartésien de topos
fibré, $m': D \times I'\to F$ le morphisme de topos fibrés déduit de
$m$ par le changement de base $\phi$ (\hyperref[VI.subsubsec7.1.9]{7.1.9}). Le couple $(D, m)$ est
une limite projective de $F$ si et seulement si $(D, m')$ est une
limite projective de $F'$.
\end{enonce*}

La démonstration\pageoriginale n'est qu'une vérification de routine assez
longue. Elle est laissée au lecteur patient, qui pourra utiliser le
point de vue des limites projectives de pseudo-foncteurs (\hyperref[VI.par8.1.3.3]{8.1.3.3}).

\begin{enonce*}{Lemme 8.2.2}\etiquette[8.2.2]{VI.lem8.2.2}
Soient $p : F \to I$ un $𝒰$-site fibré sur une petite catégorie
cofiltrante, $\underleftarrow{F}$ la Limite inductive de la catégorie
fibrée $F$ (6.3.9), $\pi : F \to \underleftarrow{F}$ le foncteur
canonique. Munissons $\underleftarrow{F}$ de la topologie la moins fine
rendant continu le foncteur $\pi$ \emph{(III 1)}. Soit $D$ un
$𝒰$-site et $n : \underleftarrow{F} \to D$ un foncteur. Les propriétés suivantes sont
équivalentes:
\begin{enumerate}
\renewcommand{\theenumi}{\roman{enumi}}
\renewcommand{\labelenumi}{\theenumi)}
\item $n$ est un morphisme de sites $D \to \underleftarrow{F}$.

\item $n \circ \pi$ est un morphisme de sites (où $F$ est muni de la
  topologie totale).

\item $(n \circ \pi, p) :  F \to D \times I$ est un morphisme du site
  total $D \times I$ dans le site total $F$.

\item Pour tout objet $i$ de $I$, le foncteur composé $F_i
  \xrightarrow{\alpha_{i!}} F \xrightarrow{n \circ \pi} D$ est un
  morphisme de sites.
\end{enumerate}
\end{enonce*}

On a i) = ii) d'après (\hyperref[VI.exe6.6]{6.6}) et iii) = iv) d'après (7.4.13). Montrons
que iii) = ii). Il suffit pour cela de montrer que le foncteur
première projection $\pr_1 : D \times I \to D$ est un morphisme de
sites. Notons $D^{\simeq}$ le topos des faisceaux sur $D$. Il
résulte de (\hyperref[VI.prop7.4.7]{7.4.7}) que $(D \times I)^{\sim}$ est équivalent au topos
$\Hhom_{I^{\circ}} (I^{\circ}, D^{\sim} \times I^{\circ}) = \Hhom(I^{\circ}, D^{\sim})$. De plus, on
constate immédiatement que le prolongement naturel de $\pr_1 : D
\times I \to D$ aux faisceaux (III 1.2) est le foncteur de $\Hhom(I^{\circ},
D^{\sim})$ dans $D^{\sim}$ qui associe au foncteur $i \to \sigma(i)$
$\Hom(I^{\circ}, D^{\sim})$ l'objet
$\varinjlim_{I^{\circ}}\sigma(i)$. Comme $I^{\circ}$ est
filtrante, les limites inductives suivant $I^{\circ}$ sont exactes (II 4) et
par suite $\pr_1$ est un morphisme de sites (IV 4.9). Il reste à
démontrer que ii = iv). Traitons tout d'abord le cas où $D$ est un
$𝒰$-topos et $F$ un $𝒰$-topos fibré sur $I$.
Soit $i$ un objet de $I$. En raisonnant comme dans \hyperref[VI.par7.4.13.6]{7.4.13.6}, on voit que le foncteur
$\alpha_{i!} : F_i \to F$ se factorise en un morphisme
$\bar{\alpha}_{i!} : F_i \to F/\alpha_{i!} (e_i)$ et le foncteur de
localisation $F/\alpha_{i!}(e_i) \to F$, où $e_i$ est un objet final
de $F_i$. Il suffit donc de montrer que le foncteur composé
$f/\alpha_{i!} (e_i) \xrightarrow{n \circ \pi} D$ est un morphisme de
sites. Mais on a un diagramme commutatif à isomorphisme près:
\[
\xymatrix@C=2.5cm@R=1.25cm{
F/\alpha_{i!}(e_i) \ar[d]\ar[r]^{n\circ
  \pi/\alpha_{i!}(e_i)}& D/n\circ \pi \alpha_{i!}(e_i) \ar[d]\\
F \ar[r]^{n\circ \pi}& D
}
\]
Comme\pageoriginale la propriété d'être un morphisme se localise (IV 5.10), il
suffit de montrer que $n \circ \pi\circ \alpha_{i!}(e_{i})$ est un objet
final de $D$. Pour tout morphisme $f : j \to k$ de $I$, il existe un
et un seul morphisme $e_f : \alpha_{j!}(e_j) \to \alpha_{k!}(e_k)$
au-dessus de $f$ et ce morphisme est cartésien. On a donc une section
cartésienne $j \mapsto \alpha_{j!}(e_j)$ de $F$ sur $I$. Par suite le
morphisme canonique $n \circ \pi \circ \alpha_{i!} (e_{i}) \to \varinjlim_{I} n
\circ \pi \circ \alpha_{j!}(e_j)$
est un isomorphisme, car $\pi$ transforme les morphismes cartésiens en
isomorphismes. Par ailleurs $\varinjlim_I
\alpha_{j!}(e_j)$, où la limite inductive est prise dans
$F^{\widehat{\phantom{,}}}$, est un objet final de
$F^{\widehat{\phantom{,}}}$. Comme $n \circ \pi$ est un morphisme, il
transforme l'objet final de $F^{\sim}$ en un objet final de $D$. Par
suite $\varinjlim_I n \circ \pi \circ
\alpha_{j!}(e_j)$ est un objet final de $D$, ce qui achève la
démonstration dans ce cas. Dans le cas général, le foncteur $(n \pi,
p) : F \to D \times I$ est cartésien et continu (\hyperref[VI.prop7.4.4]{7.4.4}). De plus, le
foncteur $n \circ \pi$ est le foncteur composé
$$
F \xrightarrow{(n \pi, p)} D \times I \xrightarrow{\pr_1} D{\quoi}.
$$
En passant aux $𝒰$-topos fibrés associés, on obtient un diagramme
commutatif à isomorphisme près:
\[
\xymatrix@C=2cm@R=1.75cm{
F       \ar@{-}[r]^{(n\pi,p)} \ar[d]_{\varepsilon_{F/I}}& D\times I
\ar[r]^{\pr_1}\ar[d]^{\varepsilon_D\times \id} & D \ar[d]^{\varepsilon_D}\\
F^{\sim/I}  \ar[r]^{(n\pi,p)^{\sim/I}} & D^{\sim}\times I \ar[r]^{\pr_1} & D^{\sim}&.
}
\]
Comme $(n\pi, p)^{\sim/I}$ est cartésien, il est du type $(n'\pi',
p')$ où $p': F^{\sim/I} \to I$ et $\pi': F^{\sim/I} \to
\varinjlim_{I^{\circ}} F^{\sim}_i$ sont les foncteurs canoniques. On
a donc $\pr_{1^{\circ}}(n\pi, p)^{\sim/I} = n'$. Comme les foncteurs
$\epsilon_{F/I}, \epsilon_D$ induisent une équivalence sur les topos
de faisceaux (\hyperref[VI.cor7.4.11]{7.4.11}), les foncteurs $n$ et $n'$ donnent
des foncteurs isomorphes en passant aux faisceaux associés (III 1.2)
et par suite $n'$ est un morphisme de sites. D'après ce qui
précède,\pageoriginale le foncteur $(n\pi, p)^{\sim/I}$ induit sur
les fibres des morphismes de topos. C'est donc un morphisme du site
total $F$ dans le site total $D^{\sim} \times I$ (7.4.13). Comme les
foncteurs $\epsilon_{F /I}, \epsilon_{D} \times \id$ induisent des
équivalences sur les topos totaux (\hyperref[VI.cor7.4.11]{7.4.11}), le foncteur
$(n \pi, p)$ est un morphisme de sites fibrés (7.4.13), d'où
l'assertion.

\begin{enonce*}{Théorème 8.2.3}\etiquette[8.2.3]{VI.thm8.2.3}
Soient $p : F \to I$ un $𝒰$-site fibré sur une catégorie cofiltrante
essentiellement petite \emph{(I 8.1)}.
\begin{enumerate}
\renewcommand{\labelenumi}{\theenumi)}
\item Notons $\underrightarrow{F}$ le site dont la catégorie sous-jacente est
  la limite inductive de la catégorie fibrée $F \to I$ \emph{(\hyperref[VI.def6.3]{6.3})} et dont
  la topologie est la topologie la moins fine rendant continu le
  foncteur canonique $\pi : F \to \underrightarrow{F}$. Alors
  $\underrightarrow{F}$ est un $𝒰$-site et
$$
(\pi, p) : F \to \underrightarrow{F} \times I
$$
est un morphisme du site fibré $\underrightarrow{F} \times I$ dans le site
fibré $F$.

\item Soit $\mu : \underrightarrow{F}^{\sim} \times I \to F^{\sim/I}$ le
  morphisme de $𝒰$-topos fibrés déduit de $(\pi, p)$ en passant aux
  $𝒰$-topos fibrés associés. Le couple $(\underrightarrow{F}^{\sim}, \mu)$
  est une limite projective de $F^{\sim/I}$.
\end{enumerate}
\end{enonce*}

\begin{enonce*}[remark]{Remarque 8.2.4}\etiquette[8.2.4]{VI.rem8.2.4}
\begin{enumerate}
\renewcommand{\labelenumi}{\theenumi)}
\item Vu les notations introduites en \hyperref[VI.subsubsec8.1.3]{8.1.3}, on peut poser
$$
\limtop{I} F^{\sim} = \underrightarrow{F}^{\sim}
$$

\item Il résulte de \hyperref[VI.prop7.4.4]{7.4.4} que la topologie sur $\underrightarrow{F}$ est
  aussi la topologie la moins fine rendant continue la famille des
  foncteurs $\pi_j : F_j \to F \xrightarrow{\pi} \underrightarrow{F}$, $j
  \in \ob I$.
\end{enumerate}
\end{enonce*}

\begin{enonce*}{Définition 8.2.5}\etiquette[8.2.5]{VI.def8.2.5}
Le site $\underrightarrow{F}$ introduit dans \hyperref[VI.thm8.2.3]{8.2.3} est appelé le site
limite projective du site fibré $F$.
\end{enonce*}

A titre d'exercice, le lecteur pourra comme en \hyperref[VI.lem8.2.1]{8.2.1}
définir la notion de limite projective d'un site fibré et vérifier
que le site $\underrightarrow{F}$ est bien une limite projective du
site $F$. La deuxième assertion du théorème peut donc se paraphraser
ainsi: Le topos des faisceaux sur le site limite projective est
équivalent à la limite projective du topos fibré associé.

\setcounter{subsubsection}{5}
\subsubsection{Démonstration du théorème : Réduction au cas d'une
  petite catégorie d'indices}\etiquette[8.2.6]{VI.subsubsec8.2.6}

Nous\pageoriginale nous bornerons à donner des indications. Soit
$I'$ une sous-catégorie pleine de $I$ telle que ${I'}^{\circ}$ soit
cofinale dans $I^{\circ}$. Notons $F'\to I'$ la catégorie fibrée
déduite de $F$ par le changement de base $I'\to I$, $\pi': F'\to
\underrightarrow{F}'$ le foncteur canonique. On constate tout
d'abord que les catégories $\underrightarrow{F}$ et
$\underrightarrow{F}'$ sont équivalentes et que la topologie sur
$\underrightarrow{F}'$ déduite de la topologie de
$\underrightarrow{F}$ par cette équivalence est la topologie la
moins fine rendant continue le foncteur $\pi' : F'\to
\underrightarrow{F}'$. L'assertion 1) du théorème en résulte alors
lorsqu'elle est démontrée dans le cas où $I$ est petite.
L'assertion 2) dans le cas général se déduit alors de l'assertion 2)
dans le cas où $I$ est petite par le lemme \hyperref[VI.lem8.2.1]{8.2.1}.

\subsubsection{Fin de la démonstration du
  théorème}\etiquette[8.2.7]{VI.subsubsec8.2.7}
La première
assertion résulte du lemme \hyperref[VI.lem8.2.2]{8.2.2} appliqué au cas
où $D =
\underrightarrow{F}$ et $n = \id$. Soit $D$ un $𝒰$-topos. On a un
diagramme commutatif à isomorphisme près:
{\def\objectstyle{\scriptstyle}
\def\labelstyle{\scriptstyle}
\[
\xymatrix@C=.8cm@R=1.15cm{
\Hhomtop(D,{(\underrightarrow{F})}^{\sim}) \ar[r]^{(1)}\ar[d]_{(3)}&
\Mmorsite(D,\underrightarrow{F}) \ar@{^{(}->}[r]^{(2)} \ar[d]_{(3)}&
\Hhom(\underrightarrow{F},D)^{\circ}\ar[dd]_{(5)}\\
\Hhomtop_{\Cart/I}(D\times I,(\underrightarrow{F})^{\sim}\times
I)  \ar[r]^{(1)} \ar[d]_{(4)}& \Mmorsite_{\Cart/I}(D\times I,
(\underrightarrow{F})\times I) \ar[d]&\\
\Hhomtop_{\Cart/I}(D\times I,\underrightarrow{F}^{\sim/I})
\ar[r]^{(1)} & \Mmorsite_{\Cart/I}(D\times I, F)  \ar@{^{(}->}[r]^{(2)} &
\Hhom_{\Cart/I}(F,D\times I)^{\circ},
}
\]}

\noindent où les foncteurs du type (1) sont des équivalences
({IV 4.9.4  et 7.2.7), les foncteurs du type (2) sont
pleinement fidèles ( IV 4.9.1} et 7.2.2), les foncteurs
du type (3) sont des foncteurs de
changement de base, les foncteurs du type (4) sont des foncteurs de
composition (avec $\mu$, \resp $(\pi, p))$ et le foncteur (5) est
l'équivalence déduite de la propriété universelle de la limite
inductive (\hyperref[VI.def6.3]{6.3}). Pour montrer que le foncteur
$\Hhomtop(D, (\underrightarrow{F})^{\sim}) \xrightarrow{(4) \circ
(3)} \Hhomtop_{/I} (D \times I, F^{\sim/I})$ est une équivalence, il
suffit de montrer que le foncteur
$$
\Mmorsite(D, \underrightarrow{F})
\xrightarrow{(4) \circ (3)} \Mmorsite_{/I}(D \times I,F)
$$
est une\pageoriginale équivalence de catégories ou encore il suffit
de montrer que les deux foncteurs pleinement fidèles
\begin{align*}
& \Mmorsite(D, \underrightarrow{F}) \xrightarrow{(5) \circ (2)} \Hhomcart_I(F, D \times I)^{\circ}\\
& \Mmorsite_{/I}(D \times I, F) \xrightarrow{(2)} \Hhomcart_I(F, D \times
I)^{\circ}
\end{align*}
ont mêmes images essentielles, ce qui résulte immédiatement du lemme
\hyperref[VI.lem8.2.2]{8.2.2}.

\subsubsection{}\etiquette[8.2.8]{VI.subsubsec8.2.8}

Soient $p : F \to I$ un $𝒰$-topos fibré sur une catégorie $I$,
$(\limtop{I} F_i, \mu)$ une limite projective de $F$. On a donc un morphisme cartésien
de $𝒰$-topos fibrés sur $I$:
$$
\mu : \limtop{I} F_i \times I \to F{\quoi},
$$
d'où, en passant aux catégories fibrées sur $I^{\circ}$ par les
foncteurs images directes (\hyperref[VI.subsubsec7.1.3]{7.1.3}), un foncteur
cartésien :
$$
\mu_{\ast}: \limtop{I} F_i \times I^{\circ} \to F'
$$
qui est le foncteur image directe par $\mu$ (\hyperref[VI.subsubsec7.1.7]{7.1.7}). On a donc, par
définition de la limite projective d'une catégorie fibrée (\hyperref[VI.subsec6.10]{6.10}), un
foncteur canonique
\begin{equation*}
\Theta : \limtop{I} F_i  \to
\varprojlim_{I, f_{\ast}} F_i = \Hhom_{\cart/I^{\circ}} (I^{\circ},
F'){\quoi}.\tag{8.2.8.1}
\end{equation*}\etiquette[8.2.8.1]{VI.eq8.2.8.1}

\begin{enonce*}{Théorème 8.2.9}\etiquette[8.2.9]{VI.thm8.2.9}
On utilise les notations de \hyperref[VI.subsubsec8.2.8]{8.2.8}, et on suppose
que $I$ est une catégorie cofiltrante essentiellement petite. Alors
le foncteur $\Theta$ est une équivalence de catégories. Supposons
$I$ petite. Posons $\limtop{I} F_i = (\underrightarrow{F})^{\sim}$
(ce qui est licite d'après \hyperref[VI.thm8.2.3]{8.2.3} et \hyperref[VI.rem8.2.4]{8.2.4})
et interprétons la catégorie $\Hhom_{I^{\circ}} (I^{\circ}, F')$
comme la catégorie des faisceaux sur le site total $F$
(\hyperref[VI.prop7.4.7]{7.4.7}). Alors le foncteur composé
$$
(\underrightarrow{F})^{\sim} = \limtop{I} F_i
\xrightarrow{\Theta} \Hhom_{\cart/I^{\circ}} (I^{\circ}, F') \hookrightarrow \Hhom_{I^{\circ}}
(I^{\circ}, F') \approxeq F^{\sim}{\quoi},
$$
n'est autre que le foncteur image directe par le morphisme de sites
défini par le foncteur
$$
\pi : F \to \underrightarrow{F}{\quoi}.
$$
\end{enonce*}

\setcounter{subsubsection}{9}
\subsubsection{}\etiquette[8.2.10]{VI.subsubsec8.2.10}

Quelques\pageoriginale commentaires avant la démonstration du
théorème \hyperref[VI.thm8.2.9]{8.2.9}. À part la solution du problème
d'existence des limites projectives dans le cas où la catégorie
d'indices est cofiltrante, solution qui curieusement n'est pas
triviale dans cette théorie, les théorèmes \hyperref[VI.thm8.2.3]{8.2.3} et
\hyperref[VI.thm8.2.9]{8.2.9} apportent deux informations supplémentaires
essentielles pour le maniement dans la pratique des topos limites
projectives. Tout d'abord le théorème \hyperref[VI.thm8.2.3]{8.2.3} interprète
le topos limite projective comme le topos des faisceaux sur un site
\sisi{donc}{dont} la catégorie sous-jacente est la catégorie limite inductive des
catégories $F_i$ suivant le système des foncteurs images inverses.
Dans les applications, on saura interpréter concrètement cette
catégorie et sa topologie. D'autre part, le théorème
\hyperref[VI.thm8.2.9]{8.2.9} affirme qu'un faisceau de la limite projective est
connu lorsqu'on connaît le système de ses images directes dans les
topos $F_i$ (qui fournit une section cartésienne sur $I^{\circ}$ de
la catégorie $(F^{\sim/I})'$). Ce théorème affirme de plus que
réciproquement lorsqu'on se donne pour tout objet $i$ de $I$ un
faisceau $X_i$ sur $F_i$ et pour tout morphisme $f : i \to j$ de $I$
un isomorphisme $X_j \simeq f_{\ast} X_i$, ces isomorphismes étant
soumis à des conditions de compatibilité, il existe essentiellement
un seul faisceau de la limite projective qui donne naissance par
images directes au système des $X_i$.

\subsubsection{}\etiquette[8.2.11]{VI.subsubsec8.2.11}

L'assertion 2) de \hyperref[VI.thm8.2.3]{8.2.3} et la première assertion de \hyperref[VI.thm8.2.9]{8.2.9}
peuvent être résumées dans la formule frappante :
\begin{equation*}
\varprojlim_{I, f_{\ast}}F_i^{\sim} \approxeq
\limtop{I, f.} F_i^{\sim} \approxeq
(\varinjlim_{I^{\circ}, f^{\ast}} F_i)^{\sim}{\quoi}. \tag{8.2.11.1}
\end{equation*}\etiquette[8.2.11.1]{VI.eq8.2.11.1}

\subsubsection{}\etiquette[8.2.12]{VI.subsubsec8.2.12}
\emph{Démonstration du théorème} \hyperref[VI.thm8.2.9]{8.2.9}
Pour démontrer la première assertion, on se ramène au cas où la
catégorie $I$ est petite en utilisant \hyperref[VI.lem8.2.1]{8.2.1} et un lemme analogue sur
les limites projectives de catégories fibrées. Nous laissons les
détails au lecteur. Supposons désormais que $I$ est une petite
catégorie. On sait que le foncteur
$$
\pi : F \to \varinjlim_{I^{\circ}} F_i
$$
est un morphisme de sites (\hyperref[VI.lem8.2.2]{8.2.2}). Le foncteur image directe
$$
\pi_{\ast} : (\varinjlim_{I^{\circ}} F_i)^{\sim} \to F^{\sim}
$$
qu'il\pageoriginale définit sur les topos est alors pleinement
fidèle, et son image essentielle est l'ensemble des faisceaux sur
$F$ qui transforment les morphismes cartésiens de $F$ en
isomorphismes (\hyperref[VI.exe6.6]{6.6}). Or il est immédiat que ces
faisceaux correspondent dans l'équivalence $F^{\sim} \simeq
\Hhom_{I^{\circ}}(I^{\circ}, F')$ aux sections cartésiennes de $F'$
sur $I^{\circ}$ (\hyperref[VI.prop7.4.7]{7.4.7}). De plus, pour tout $j \in \ob
(I)$, notons $\alpha_{j!} : F_j \to F$ le foncteur d'inclusion qui
donne naissance aux trois foncteurs $\xymatrix{ (\alpha_{j!},
\alpha^{\ast}_j, \alpha_{j\ast}) : F_j  \ar@<4pt>[r] \ar@<-4pt>[r]&
F \ar[l] }$ (\hyperref[VI.rem7.4.3]{7.4.3}). Il résulte de la définition de
l'équivalence $F' \simeq \Hhom_{I^{\circ}}(I^{\circ},F)$
(\hyperref[VI.prop7.4.7]{7.4.7}) qu'à un faisceau $X$  sur
$\varinjlim_{I^{\circ}} F_i$ correspond par le foncteur
$(\varinjlim_{I^{\circ}} F_i)^{\sim} \xrightarrow{\pi_{\ast}} F
\simeq \Hhom_{I^{\circ}} (I^{\circ}, F')$ la section cartésienne $j
\mapsto \alpha^{\ast}_{j} \pi_{\ast}(X)$. Or il résulte de la
description du foncteur $\mu : (\varinjlim_{I^{\circ}} F_i)^{\sim}
\times I \to F$ donnée dans \hyperref[VI.thm8.2.3]{8.2.3} que pour tout objet
$j$ de $I$, le foncteur $\mu_{j_\ast} : (\varinjlim_{I^{\circ}}
F_i)^{\sim} \to F_j$ image directe par le morphisme $\mu_j :
(\varinjlim_{I^{\circ}} F_i)^{\sim} \to F_j$ déduit de $\mu$ par
passage aux fibres en $j$, n'est autre que le foncteur
$\alpha^{\ast}_j \pi_{\ast}$. Par suite le foncteur
$(\varinjlim_{I^{\circ}} F_i)^{\sim} \xrightarrow{\pi_{\ast}} F
\simeq \Hhom_{I^{\circ}} (I^{\circ}, F')$ n'est autre que le
foncteur $(\varinjlim_{I^{\circ}} F_i)^{\sim} \xrightarrow{\Theta}
\Hhom_{\cart/I^{\circ}} (I^{\circ}, F') \hookrightarrow
\Hhom_{I^{\circ}}(I^{\circ}, F')$, d'où le théorème.

\subsection{Topologie du site limite projective : Cas des topos
  cohérents}\etiquette[8.3]{VI.subsec8.3}

\subsubsection{}\etiquette[8.3.1]{VI.subsubsec8.3.1}

Dans ce numéro, $p : F \to I$ est un $𝒰$-site fibré sur une catégorie
cofiltrante essentiellement petite, dont les foncteurs images inverses
$f^{\ast} :F_i \to F_j$ commutent aux produits fibrés. On se donne de
plus, pour tout objet $X$ d'une catégorie fibre $F_i$, un ensemble
$\Cov(X)$ de familles couvrantes pour la topologie de $F_i$ qui possèdent les
propriétés (PTO) et (PT1) de II 1.3 et qui engendrent la topologie de
$F_i$. Enfin on suppose que pour tout $f : i \to j$, tout $Y \in \ob
F_j$, et toute famille $(Y_{\alpha} \to Y)_{\alpha \in A}$ de
$\Cov(Y)$, la famille $(f^{\ast}(Y_{\alpha}) \to f^{\ast} (Y))_{\alpha
  \in A}$ appartient à $\Cov(f^{\ast}(Y))$. Ces conditions sur $\Cov(X)$
sont vérifiées par exemple si on prend $\Cov(X) = $ \{toutes les
  familles couvrantes de $X$ dans $F_i$\}.
On se\pageoriginale propose d'étudier dans ce numéro le site limite
projective de $F$ sous des hypothèses de finitude convenables (\hyperref[VI.thm8.3.13]{8.3.13}).

\subsubsection{}\etiquette[8.3.2]{VI.subsubsec8.3.2}

Soient $\pi : F \to \underrightarrow{F}$ le foncteur canonique et $Y$ un
objet de $\underrightarrow{F}$. Désignons par $\Cov(Y)$ l'ensemble des
familles $(Y_{\alpha} \xrightarrow{n_{\alpha}} Y)_{\alpha \in A}$ du
type suivant:

Il existe un $i \in I$, un objet $X$ de $F_i$, une famille
$(X_{\alpha} \xrightarrow{{n'}_\alpha} X)_{\alpha \in A} \in \Cov(X)$,
un isomorphisme $\phi : \pi(X) \simeq Y$ et une famille
d'isomorphismes $\phi_{\alpha} : \pi(X_{\alpha}) \simeq Y_{\alpha}$,
$\alpha \in A$, tels que pour tout $\alpha \in A$ on ait un diagramme
commutatif:
\[
\xymatrix@C=2.5cm@R=2cm{
\pi(X_{\alpha})
\ar@{-}[r]^{\displaystyle{\mathop{\sim}^{\phi_{\alpha}}}} \ar[d]^{\pi(n'_{\alpha})}& Y_{\alpha}\ar[d]^{n_{\alpha}} \\
\pi(X)  \ar@{-}[r]^{\displaystyle{\mathop{\sim}^{\phi}}} & Y .
}
\]

\begin{enonce*}{Proposition 8.3.3}\etiquette[8.3.3]{VI.prop8.3.3}
\begin{enumerate}
\renewcommand{\labelenumi}{\theenumi)}
\item L'ensemble des familles $\in \Cov (Y)$, $Y \in \ob
  \underrightarrow{F}$, engendre la topologie du site limite projective
  (\hyperref[VI.def8.2.5]{8.2.5}).

\item La famille $Y \mapsto \Cov (Y)$, $Y \in \ob \underrightarrow{F}$,
  possède les propriétés \emph{(PTO)} et \emph{(PT1)} de \emph{II 1.3}.
\end{enumerate}
\end{enonce*}

Nous aurons besoin du lemme suivant:

\begin{enonce*}{Lemme 8.3.4}\etiquette[8.3.4]{VI.lem8.3.4}
Soient $i$ un objet de $I$, $n : X'\to X$ un morphisme quarrable de
$F_i$ tel que pour tout $f : j \to i$, $f^{\ast}(n)$ soit
quarrable. Alors $\pi(n)$ est quarrable. Les foncteurs :
$$
\pi_{j} : F_j \hookrightarrow F \xrightarrow{\pi} \underrightarrow{F},\quad j \in
\ob I{\quoi},
$$
commutent aux produits fibrés.
\end{enonce*}

Rappelons que tout objet $Y$ de $\underrightarrow{F}$ est égal à un objet
$\pi(Z)$, $Z \in \ob F$, et que tout morphisme $m : Y \to \pi (X)$ est
de la forme:
$$
\pi(Z) \xrightarrow{\pi(s)^{-1}} \pi(Z') \xrightarrow{\pi(m')} \pi(X){\quoi},
$$
où $s$ est un morphisme cartésien de $F$. Par suite, pour montrer que
$\pi(n)$ est quarrable,\pageoriginale il suffit de montrer que le
produit fibré de tout diagramme
\[
\xymatrix@C=2.5cm@R=1.25cm{
                           &  \pi(X') \ar[d]^{\pi(u)}\\
 \pi(Z')  \ar[r]^{\pi(m')} & \pi(X)
}
\]
est représentable. Mais le morphisme $m': Z' \to X$ se factorise en
$Z'\xrightarrow{m''} X'' \xrightarrow{s'} X$, où $p(m'')$ est
l'identité de $p(Z') = j$ et où $s'$ est cartésien au-dessus de $p(s')
= f$. Posons $X'''= f^{\ast} (X')$, $n' = f^{\ast}(n)$. On a alors
un diagramme commutatif
\[
\xymatrix@C=2.5cm@R=1.5cm{
X'''         \ar[r]^{s''} \ar[d]^{n'}& X'\ar[d]^{n}        \\
X''         \ar[r]^{s'}  & X,
}
\]
d'où un diagramme commutatif dans $\underrightarrow{F}$:
\[
\xymatrix@C=2.5cm@R=1.5cm{
        & \pi(X''\;')   \ar[r]^{\sim} \ar[d]_{\pi(u')}& \pi(X')\ar[d]^{\pi(n)}\\
\pi(Z')  \ar@{-}[r]^{\pi(m')} & \pi(X'') \ar[r]^{\sim} & \pi(X)
}
\]
où $n'$ est un morphisme quarrable de $F_j$. Pour montrer que $\pi(n)$
est quarrable, il suffit donc de montrer que le produit $\pi(Z')
\times_{\pi(X'')} \pi(X''')$ est représentable, et par suite il
suffit de montrer que $\pi_j : F_j \hookrightarrow F \to
\underrightarrow{F}$ commute aux produits fibrés.

Soient
\[
\xymatrix@C=2.5cm@R=1.25cm{
Y'  \ar[r] \ar[d]& X' \ar[d]\\
Y  \ar[r]  & X
}
\]
un diagramme cartésien de $F_j$ et $W$ un objet de $F$ au-dessus de $k
\in \ob I$. On a (\hyperref[VI.prop6.5]{6.5}):
\def\labelstyle{\scriptstyle}
$$\displaystyle{\Hom(\pi(W),\pi(Y'))\simeq\mathop{\varinjlim_{
\xymatrix@R=.3cm@C=.3cm@M=.05cm{
            & j\\
        l \ar[ru]^{f} \ar[rd]_{g} &    \\
            & k
}}}\Hom_{F_{\ell}}(g^{\ast}(W),f^{\ast}(Y'))}$$

Mais comme\pageoriginale $f^{\ast} : F_j \to F_{\ell}$ commute aux
produits fibrés, on a
$$
\Hom(g^{\ast}(W), f^{\ast}(Y')) \simeq \Hom(g^{\ast}(W), f^{\ast}(Y))
\times_{\Hom(g^{\ast}(W), f^{\ast}(X))} \Hom(g^{\ast}(W){\quoi},
f^{\ast}(X')).
$$
Utilisant alors la commutation des limites filtrantes aux produits
fibrés (I 2) et la formule (\hyperref[VI.prop6.5]{6.5}) on obtient un
isomorphisme:
$$
\Hom(\pi(W), \pi(Y')) \simeq \Hom(\pi(W), \pi(Y)) \times_{\Hom(\pi(W),
  \pi(X))} \Hom(\pi(W), \pi(X')){\quoi},
$$
ce qui montre que le diagramme:
\[
\xymatrix@C=2.5cm@R=1.5cm{
\pi(Y')  \ar[r] \ar[d]& \pi(X')\ar[d]\\
\pi(Y)    \ar[r] & \pi(X)
}
\]
est cartésien dans $\underrightarrow{F}$.

\setcounter{subsubsection}{4}
\subsubsection{}\etiquette[8.3.5]{VI.subsubsec8.3.5}
\emph{Démonstration de} \hyperref[VI.prop8.3.3]{8.3.3}. Démontrons d'abord la
deuxième assertion. Il résulte de \hyperref[VI.lem8.3.4]{8.3.4} que les familles
de $\Cov(Y)$, $Y \in \underrightarrow{F}$, sont composées de
morphismes quarrables, d'où la propriété (PTO). Pour montrer que
ces familles possèdent la propriété (PT1) (stabilité par changement
de base), on est ramené, en procédant à une suite de réductions
comme dans la démonstration de \hyperref[VI.prop8.3.3]{8.3.3}, à démontrer
l'assertion suivante:

Pour tout $i \in \ob I$, tout $Y \in \ob F_i$, tout $(Y_{\alpha} \to
Y)_{\alpha \in A} \in \Cov(Y)$, et tout morphisme $u : Z \to Y$ de
$F_i$, la famille $(\pi(Z) \times_{\pi(Y)} \pi(Y) \to \pi(Z)_{\alpha
  \in A}$ appartient à $\Cov((Z))$.

Cette dernière assertion résulte du fait que les $\pi_i : F_i
\hookrightarrow F \xrightarrow{\pi} \underrightarrow{F}$ commutent aux
produits fibrés (\hyperref[VI.prop8.3.3]{8.3.3}) et que les familles $X \mapsto \Cov(X)$ de
$F_i$ sont stables par changement de base. Démontrons la première
assertion. Il est clair que les familles qui appartiennent aux
$\Cov(Y)$, $Y \in \underrightarrow{F}$, sont couvrantes pour la topologie
$T$ la moins fine rendant continu le foncteur $\pi$ (III 1). Donc la
topologie $T'$ engendrée par ces familles est moins fine que $T$. Pour
montrer qu'elle est plus fine que\pageoriginale $T$ (et par suite
égale à $T$), il suffit de montrer que tout faisceau $M$ pour $T'$ est tel que $M \circ
\pi$ est un faisceau sur $F$, ou encore (\hyperref[VI.rem8.2.4]{8.2.4}) que $M \circ \pi_i$
est un faisceau sur $F_i$ pour tout $i$ ob $I$. Or, comme les foncteurs
$\pi_i : F_i \to \underrightarrow{F}$ commutent aux
produits fibrés (\hyperref[VI.lem8.3.4]{8.3.4}), $M$ est un faisceau pour $T'$ si et seulement
si (II 2.3 et I 2.12) pour tout $i \in \ob I$, pour tout $Y \in \ob
F_{i}$, pour tout $(Y_{\alpha} \to Y)_{\alpha \in A} \in
\Cov(Y)$, la suite d'ensembles
$$
M(\pi_i (Y)) \to \prod_{\alpha \in A} M (\pi_i (Y_{\alpha}))
\rightrightarrows \prod_{(\alpha, \beta) \in A \times A} M
(\pi_i (\fprod{Y_{\alpha}}{Y_{\beta}}{Y}))
$$
est exacte, \ie (\loccit) si et seulement si pour $i \in \ob I$, $M
\circ \pi_i$ est un faisceau sur $F_i$.

\begin{enonce*}{Proposition 8.3.6}\etiquette[8.3.6]{VI.prop8.3.6}
On utilise les notations et hypothèses de \hyperref[VI.subsubsec8.3.1]{8.3.1} et \hyperref[VI.subsubsec8.3.2]{8.3.2}. Si pour tout
$i \in \ob I$, $X \mapsto \Cov(X)$ est une prétopologie sur $F_i$ \emph{(II
1.3)} et si pour tout $X \in \ob F_i$, les familles couvrantes de
$\Cov(X)$ sont finies, alors pour tout $Y \in \underrightarrow{F}$, les
familles couvrantes de $\Cov(Y)$ sont finies et $Y \mapsto \Cov(Y)$ est
une prétopologie sur $\underrightarrow{F}$.
\end{enonce*}

\setcounter{subsubsection}{6}
\subsubsection{}\etiquette[8.3.7]{VI.subsubsec8.3.7}
La seule chose à démontrer est que les familles $Y \mapsto \Cov(Y)$
possèdent la propriété (PT2) de (II 1.3) (stabilité par
composition). Soit $i \in \ob I$, $Y$ un objet de $F_i$, $(Y_\alpha \to
Y)_{\alpha \in A} \in \Cov(Y)$. Soient de plus $i_\alpha$, $\alpha \in
A$, une famille d'objets de $I$, pour tout $\alpha$, $Z_{\alpha}$ un
objet de $F_i$ et $(Z_{\alpha \beta} \xrightarrow{n_{\alpha \beta}}
Z_{\alpha})_{\alpha \beta \in B_{\alpha}}$ une famille de
$\Cov(Z_{\alpha})$. Enfin donnons nous pour tout $\alpha \in A$, un
isomorphisme $\phi_{\alpha} : \pi (Z_{\alpha}) \xrightarrow{\sim} \pi
(Y_{\alpha})$. En revenant à la définition des familles de $\Cov(X)$,
$X \in \ob \underrightarrow{F}$ (\hyperref[VI.subsubsec8.3.2]{8.3.2}), on voit qu'il s'agit de démontrer
que la famille
$$
(\pi (Z_{\alpha \beta}) \xrightarrow{\pi(n_{\alpha}) \circ \phi_{\alpha} \circ
  \pi(n_{\alpha \beta})} \pi(Y))_{\alpha \beta \in \mathop{\coprod}_A
  B_{\alpha}}
$$
appartient à $\Cov(\pi(Y))$ ce que nous ferons après cinq
réductions.

\setcounter{subsubsection}{7}
\subsubsection{}\etiquette[8.3.8]{VI.subsubsec8.3.8}
\emph{Réduction au cas où  pour tout $\alpha \in A$, on a}
$\phi_{\alpha} = \pi(m_{\alpha})$, $m_{\alpha} : Z_{\alpha} \to
Y_{\alpha}$. \emph{Pour tout $\alpha \in A$, on a $\phi_{\alpha} =
\pi(m_{\alpha}) \pi(s_{\alpha})^{-1}$ $(\hyperref[VI.prop6.5]{6.5})$ où $s_{\alpha}$ est un
morphisme cartésien au-dessus de $f_{\alpha} : i'_{\alpha} \to
i_{\alpha}$. On a donc des diagrammes commutatifs}
\[
\xymatrix@C=2.5cm@R=1.25cm{
Z_{\alpha\beta} \ar[d]_{n_{\alpha\beta}} & f^{\ast}(Z_{\alpha\beta})
\ar[l]_{s_{\alpha\beta}} \ar[d]^{f^{\ast}_{\alpha}(n_{\alpha\beta})}\\
Z_{\alpha} &   f^{\ast}_{\alpha}(Z_{\alpha}) \ar[l]_{s_{\alpha}},  &
**[l]\alpha\beta\in \displaystyle{\coprod_A} B_{\alpha},
}
\]
qui\pageoriginale fournissent des diagrammes commutatifs
\[
\xymatrix@C=2.5cm@R=1.25cm@H=.2cm{
\pi(Z_{\alpha\beta})  \ar[r]^{\sim} \ar[d]^{\pi(n_{\alpha\beta})}&
\pi(f^{\ast}_{\alpha}(Z_{\alpha\beta})) \ar[d]^{\pi(f^{\ast}_{\alpha}(m_{\alpha\beta}))}\\
\pi(Z_{\alpha})  \ar[r]^{\displaystyle\mathop{\sim}^{\pi(s_\alpha)^{-1}}} & \pi(f^{\ast}_{\alpha}(Z_{\alpha}))
\ar[r]^{\pi(m_{\alpha})} & \pi(Y_\alpha),\quad
\alpha\beta\in \displaystyle{\coprod_A} B_{\alpha}.}
\]
Comme les familles
$(f^{\ast}_{\alpha}(Z_{\alpha\beta})\xrightarrow{f^{\ast}(n_{\alpha})}
f^{\ast}(Z_{\alpha}))$ appartiennent à
$\Cov(f^{\ast}_{\alpha}(Z_{\alpha}))$, on peut se ramener au cas où
pour tout $\alpha$, on a $\phi_{\alpha} = \pi(m_{\alpha})$.

\subsubsection{}\etiquette[8.3.9]{VI.subsubsec8.3.9}

\emph{Réduction au cas où, pour tout $\alpha \in A$, $i_{\alpha} = j$ et}
  $p(m_{\alpha}) = k : j \to i$. Posons $g_{\alpha} =
  p(m_{\alpha})$. Comme $I$ est cofiltrante et comme $A$ est fini, il
  existe un objet $j$ de $I$ et des morphismes $h_{\alpha} : j \to
  i_{\alpha}$  tels que pour tout couple $(\alpha, \alpha')$ on ait
  $g_{\alpha} h_{\alpha} = g_{\alpha'} h_{\alpha'} = k$. En utilisant
  les foncteurs images inverses $h^{\ast}_{\alpha}$ comme en \hyperref[VI.subsubsec8.3.8]{8.3.8}, on
  se ramène au cas décrit.

\subsubsection{}\etiquette[8.3.10]{VI.subsubsec8.3.10}

\emph{Réduction au cas où} $i = j$ et $k : j \to i$ \emph{est
l'identité}. On a, pour tout $\alpha$, un diagramme commutatif:
\[
\xymatrix@C=2.5cm@R=1.5cm{
Z_{\alpha}  \ar[r]^{m'_\alpha} &  k^{\ast}(Y_{\alpha}) \ar[d]_{k^{\ast}(n_{\alpha})}
\ar[r]^{t_{\alpha}} & Y_{\alpha} \ar[d]^{n_{\alpha}}\\
                    & k^{\ast}(Y)  \ar[r]^{t} & Y,
}
\]
où les morphismes $t$ et $t_{\alpha}$, $\alpha \in A$, sont
cartésiens, où $m'_{\alpha}$ est au-dessus de l'identité de $j$ et
où $t_{\alpha} \circ  m'_{\alpha} = m_{\alpha}$. En remplaçant
la famille $(Y_{\alpha} \to Y) \alpha \in A$ par la\pageoriginale famille
$(k^{\ast} (Y_{\alpha}) \to k^{\ast} (Y))_{\alpha \in A} \in
\Cov (k^{\ast} (Y))$, et les morphismes $m_{\alpha}$ par les
morphismes $m'_{\alpha}$, on est ramené au cas décrit.

\begin{enonce*}{Lemme 8.3.11}\etiquette[8.3.11]{VI.lem8.3.11}
Soient $j$ un objet de $I$ et $m : X' \to X$ un morphisme de $F_j$ tel
que $\pi(m)$ soit un isomorphisme. Il existe un morphisme $\ell : j'
\to j$ tel que $\ell^{\ast}(m)$ soit un isomorphisme.
\end{enonce*}

Montrons qu'il existe un morphisme $f : i \to j$ et un morphisme $q
: f^{\ast} (X) \to f^{\ast}(X')$ tels que $f^{\ast}(m) q =
\id_{f^{\ast}(X)}$ et tels que $\pi(q)$ soit un isomorphisme. En
effet, il existe un isomorphisme $n : \pi (X) \to \pi (X')$ tel que
$\pi(m) \circ n = \id_{\pi(X)}$. Il existe donc (\hyperref[VI.prop6.5]{6.5})
deux morphismes $X \xleftarrow{s_1} Y_1 \xrightarrow{q_1} X'$,
où $s_1$ est cartésien, tels que $n = \pi(q_1) \pi(s_1)^{-1}$.
On a donc $\pi(m) \circ \pi(q_1) = \pi(s_1)$. Les morphismes $mq_1$
et $s_1$ de $Y_{1}$ dans $X$ ont même image dans
$\underrightarrow{F}$. Par suite (\cf la description des morphismes
dans la limite inductive (\hyperref[VI.prop6.5]{6.5})), il existe un morphisme
cartésien $s_2 : Y_2 \to Y_1$ tel que $mq_1s_2 = s_1 s_2$. Posons
$q_1 s_2 = q_2$ et notons $s_3$ le morphisme cartésien $s_1 s_2$. On
a donc $mq_2 = s_3$. Le morphisme $q_2 : Y_2 \to X'$ se factorise de
manière essentiellement unique en $Y_2 \xrightarrow{q} Y_3
\xrightarrow{s_4} X'$ où $s_4$ est cartésien et où $q$ est
au-dessus de l'identité de $p(Y_2)$. On a donc un diagramme
commutatif
\[\vcenter{
\xymatrix@C=2.25cm@R=1.5cm{
Y_3 \ar[d]_{s_4} & Y_2 \ar[l]_{q}\ar[d]^{s_3}\\
X'  \ar[r]^{m} & X.          .
}}\leqno{(8.3.11.1)}
\]\etiquette[8.3.11.1]{VI.eq8.3.11.1}
Posons $p(s_4) = p(s_3) = f : i \to j$. On a $f^{\ast} (X') \simeq
Y_3$, $f^{\ast} (X) \simeq Y_2$, et $f^{\ast}(m) : Y_2 \to Y_3$ rend
commutatif le diagramme:
\[
\xymatrix@C=2.25cm@R=1.5cm{
Y_3  \ar[r]^{f^{\ast}(m)} \ar[d]_{s_4} & Y_2 \ar[d]^{s_3}\\
X'  \ar[r]^{m} & X.
}
\]
On a donc\pageoriginale $s_3 = s_3 (f^{\ast}(m) \circ q)$; comme
$s_3$ est cartésien et comme $f^{\ast} (m) \circ q$ est au-dessus de
l'identité de $p(Y_2)$, on a $f^{\ast}(m) \circ q = \id_{Y_2}$. De
plus, il résulte de la commutativité du diagramme \eqhyperref[VI.eq8.3.11.1]{8.3.11.1}
que $\pi (q)$ est un isomorphisme.

Appliquons alors le résultat précédent au morphisme $q : f^{\ast}
(X) \to f^{\ast}(X')$: il existe un morphisme $f' : j'\to i$ et un
morphisme $q': f'^{\ast}(X') \to f'^{\ast}(X)$, tels que
$f'^{\ast}(q) q'= \id_{f'^{\ast}(f^{\ast}(X'))}$. Mais en posant
$\ell^{\ast} = ff'$; on a un isomorphisme $f'^{\ast}f^{\ast} \simeq
\ell^{\ast}$. On a donc un morphisme $f'^{\ast}(q) : \ell^{\ast}(X)
\to \ell^{\ast}(X')$ et un morphisme $q':\ell^{\ast}(X')\to
\ell^{\ast}(X)$, tels que
$\ell^{\ast}(m){f'}^{\ast}(q)=\id_{\ell^{\ast}(X)}$ et
${f'}^{\ast}(q)q' = \id_{\ell^{\ast}(X')}$. Par suite
$\ell^{\ast}(m)$ est un isomorphisme.

\setcounter{subsubsection}{11}
\subsubsection{}\etiquette[8.3.12]{VI.subsubsec8.3.12}
\emph{Fin de la démonstration}. Comme $I$ est cofiltrante et $A$
fini, il résulte de \hyperref[VI.lem8.3.11]{8.3.11} qu'il existe un foncteur
image inverse $\ell^{\ast}$  qui transforme tous les morphismes
$m_{\alpha}$, $\alpha \in A$, en isomorphismes. On se ramène donc au
cas où les $m_{\alpha}$ sont des isomorphismes. Mais alors les
familles couvrantes $(Z_{\alpha \beta} \to Z)_{\alpha \beta \in
B_{\alpha}}$ sont déduites par le changement de base $m_{\alpha} :
Z_\alpha \to Y_{\alpha}$ de familles couvrantes $(Y_{\alpha \beta}
\to Y_{\alpha})_{\alpha \beta \in B_{\alpha}} \in \Cov(Y_{\alpha})$.
On est donc ramené au cas où $Z_\alpha = Y_{\alpha}$ et
$\Phi_{\alpha} = \id_{\pi(Y_{\alpha})}$. Mais dans ce cas la famille
$(Y_{\alpha\beta} \xrightarrow{n_{\alpha} \circ n_{\alpha
    \beta}} Y)_{\alpha \beta \in \mathop{\coprod}_A B_{\alpha}}$
appartient à $\Cov(Y)$ (propriété (PT2)). Donc son image par $\pi$
appartient à $\Cov(\pi(Y))$ (\hyperref[VI.subsubsec8.3.2]{8.3.2}).

\begin{enonce*}{Théorème 8.3.13}\etiquette[8.3.13]{VI.thm8.3.13}
Soient $p : F \to I$ un $𝒰$-topos fibré sur une catégorie cofiltrante
essentiellement petite. Si pour tout objet $i$ de $I$ le topos fibre
$F_i$ est cohérent $(\hyperref[VI.def2.3]{2.3})$ et si pour tout morphisme $f : i \to j$, le
morphisme de topos $f\bigdot = (f^{\ast}, f_{\ast}) : F_i \to F_j$ est
cohérent (\hyperref[VI.def3.1]{3.1}), alors le topos $\limtop{I} F_i$
est cohérent et pour tout objet $i$ de $I$, le morphisme canonique de
topos
$$
\mu_i : \limtop{I} F_i \to F_i
$$
est cohérent.

De plus, en notant $F_{\coh}$ la sous-catégorie pleine de $F$ définie
par les objets de $F$ qui sont cohérents dans leur fibre (\hyperref[VI.def1.13]{1.13}), la
catégorie $F_{\coh}$ est fibrée sur $I$ et $\underrightarrow{F}_{\coh}$ est
canoniquement équivalente à la catégorie des objets cohérents de
$\limtop{I} F_i$
\end{enonce*}

Pour tout\pageoriginale $i \in \ob I$, notons $(F_i)_{\coh}$ le
$𝒰$-site (muni de la topologie induite) des objets cohérents de
$F_i$ (\hyperref[VI.def1.13]{1.13}). Comme $F_i$
est cohérent, $F_i$ est le topos des faisceaux sur $(F_i)_{\coh}$. Pour
tout morphisme $f : i \to j$, les morphismes $f\bigdot = (f^{\ast},
f_{\ast})$ sont cohérents et par suite $f^{\ast}(F_j)_{\coh} \subset
(F_i)_{\coh}$. On a donc un $𝒰$-site fibré $p_{\coh} : F_{\coh} \to I$,
dont les sites fibres sont les sites $(F_i)_{\coh}$ et dont le topos
fibré associé est équivalent à $p : F \to I$. Par suite on a (\hyperref[VI.thm8.2.3]{8.2.3}):
\begin{equation*}
\limtop{I, f.} F_i \approxeq
(\underrightarrow{F}_{\coh})^{\sim}\tag{8.3.13.1}
\end{equation*}\etiquette[8.3.13.1]{VI.eq8.3.13.1}
Comme les topos $F_i$ sont cohérents, les produits finis et les
produits fibrés sont représentables dans $(F_i)_{\coh}$
(\hyperref[VI.prop2.2]{2.2}). De plus les foncteurs $f^{\ast} : (F_i)_{\coh}
\to (F_j)_{\coh}$ sont exacts à gauche. Enfin, comme les objets de
$(F_i)_{\coh}$ sont quasi-compacts, les familles couvrantes finies
forment une prétopologie. Il résulte alors de (\hyperref[VI.lem8.3.4]{8.3.4})
que les produits finis et les produits fibrés sont représentables
dans $\varprojlim_{I^{\circ},F} (F_i)_{\coh}$ (remords au lemme
\hyperref[VI.lem8.3.4]{8.3.4}. : Montrer que les foncteurs canoniques
$(F_i)_{\coh} \to \underrightarrow{F}_{\coh}$ transforment l'objet
final en l'objet final) et de (\hyperref[VI.prop8.3.6]{8.3.6}) que les objets de
$\underrightarrow{F}_{\coh}$ sont quasi-compacts. Par suite
(\hyperref[VI.subsubsec2.4.5]{2.4.5}), $\limtop{I} (F_i)$  est cohérent. Enfin
les morphismes canoniques $\mu_{i} : \limtop{I}(F_i) \to F_i$ se
déduisent, par passage aux topos des faisceaux, des morphismes de
sites (\hyperref[VI.thm8.2.3]{8.2.3}):
$$
(F_i)_{\coh} \to \underrightarrow{F}_{\coh}
$$
donc (\hyperref[VI.cor3.3]{3.3}) les morphismes $\mu_i$ sont cohérents.

Démontrons la dernière assertion. On a \eqhyperref[VI.eq8.3.13.1]{8.3.13.1}
$$
\limtop{ℐ} (F_i) \simeq
(\underrightarrow{F}_{\coh})^{\sim}{\quoi}.
$$
Soit $X$ un objet cohérent de $(F_{\coh})^{\sim}$. Comme $X$ est
quasi-compact, il existe une famille couvrante finie $(Y_{\alpha} \to
X)$ où $Y_{\alpha} \in \ob \underrightarrow{F}_{\coh}$ (\hyperref[VI.def1.1]{1.1}). Quitte à
prendre un indice $i \in \ob I$ assez petit, on peut supposer que les
$Y_{\alpha}$ sont les images par
$\mu_{i}:F_{i\coh}\to\underrightarrow{F}_{\coh}$ d'une famille $Y'_{\alpha}$
d'où, en prenant la somme directe des $Y'_{\alpha}$
(\hyperref[VI.cor1.15]{1.15}), on voit qu'il existe un indice $i \in \ob I$, un $Y' \in F_{i
\coh}$ et un morphisme surjectif $\mu_i(Y') \to X$. La relation
d'équivalence $R = \fprod{\mu_i (Y')}{\mu_{i} (Y')}{X}$ est alors un
objet\pageoriginale cohérent car $(\underrightarrow{F}_{\coh})^{\sim}$
est un topos cohérent (\hyperref[VI.prop2.2]{2.2}). Montrons tout d'abord que $R$ est un objet de
$\underrightarrow{F}_{\coh}$. Par définition $R$ est un sous-objet de
$\mu_i(Y'\times Y')$ et comme $R$ est cohérent, on peut, quitte à
changer l'indice $i$, trouver une flèche $m : Z \to (Y'\times Y')$ de
$F_{i \coh}$ telle que $\text{Im}(\mu_i(m)) = R$. Il en résulte que, $F_i$
étant cohérent, $\text{Im}(m)$ est cohérent dans $F_{i}$ (\hyperref[VI.cor1.17.1]{1.17.1}) et par suite
appartient à $F_{i \coh}$. Il existe donc un indice $i$ et un diagramme
$\def\labelstyle{\scriptstyle} \xymatrix@C=1cm{
R \ar@<3pt>[r]^{p_1}\ar@<-3pt>[r]_{p_2} & Y'
}
$
de $F_{i \coh}$ tel que
$\def\labelstyle{\scriptstyle}
\xymatrix@C=1cm{
  \mu_i(R)\ar@<3pt>[r]^{\mu_i(p_1)} \ar@<-3pt>[r]_{\mu_i(p_2)}& \mu_i(Y')
}
$
soit une relation d'équivalence et tel que $X =
\mu_i(Y')/\mu_i(R)$. Posons $X_i = \coker(p_1, p_2)$ et $R'=
\fprod{Y'}{Y'}{X'_i}$. On a un morphisme canonique dans $F_i : u :
R \to R'$. De plus $\mu^{\ast}_i(u)$ est un isomorphisme. Donc, quitte
à changer l'indice $i$, on peut supposer que $R = R'$, \ie que $R$ est
une relation d'équivalence. Mais alors $X_i$ appartient à $F_{i \coh}$
(\hyperref[VI.cor1.17.1]{1.17.1}) et par suite $X = \mu_{i}(X_i)$ appartient à
$\underrightarrow{F}_{\coh}$.

\begin{enonce*}{Corollaire 8.3.14}\etiquette[8.3.14]{VI.cor8.3.14}
Soient $I$ une catégorie cofiltrante essentiellement petite, $p : F
\to I$ et $q : G \to I$ deux $𝒰$-topos fibrés sur $I$ tels que pour
tout objet $i$ de $I$ les topos $F_i$ et $G_i$ soient cohérents et
tels que pour tout morphisme $f : i \to j$ les morphismes $f. : F_i
\to F_j$ et $G_i \to G_j$ soient cohérents. Soit de plus $m : F \to G$
un morphisme cartésien de $𝒰$-topos fibrés tels que pour tout $i \in
\ob (I)$, $m_i : F_i \to G_i$ soit cohérent. Alors le morphisme
$\underleftarrow{m}$ déduit de $m$ par passage à la limite projective est
cohérent.
\end{enonce*}

Pour tout objet $i$ de $I$, le morphisme $m_i$ induit un foncteur
$m^{\ast}_{i, \coh} : G_{i, \coh} \to F_{i, \coh}$, d'où un
foncteur cartésien $m^{\ast}_{\coh} = G_{\coh} \to F_{\coh}$ qui est
un morphisme du site fibré $F_{\coh}$ dans le site fibré $G_{\coh}$
(7.4.13). Le foncteur $\varinjlim_{I^{\circ}} m^{\ast}_{\coh} :
\varinjlim_{I^{\circ}} G_{\coh} \to \varinjlim_{I^{\circ}} F_{\coh}$
est un morphisme de sites qui fournit en passant aux topos
correspondants le morphisme $\underleftarrow{m}$. L'assertion
résulte alors de \hyperref[VI.thm8.3.13]{8.3.13} et de \hyperref[VI.cor3.3]{3.3}.

\begin{enonce*}[remark]{Exercice 8.3.15}\etiquette[8.3.15]{VI.exe8.3.15}
Avec les notations de \hyperref[VI.thm8.3.13]{8.3.13}, montrer que si les $F_i$, $i \in I$,
sont algébriques (\hyperref[VI.def2.3]{2.3}) et si les morphismes $f. : F_i \to F_j$, $f \in
\text{\rm Fl}(I)$, sont  cohérents,\pageoriginale $\limtop{I} F_i$ est algébrique
et que $(\limtop{I} F_i)_{\coh} \xrightarrow{\sim} \underrightarrow{F}_{\coh}$.
\end{enonce*}

\subsection{Exemple: Topos locaux}\etiquette[8.4]{VI.subsec8.4}

\subsubsection{}\etiquette[8.4.1]{VI.subsubsec8.4.1}

Soit $X = (X_i)_{i \in I}$ un topos fibré sur une catégorie $I$. On
en déduit un topos fibré sur la catégorie $\Pro(I)$ (I 8.10) par la
formule
\begin{equation*}
X_{\sheaf{\alpha}} = \limtop{i_{\alpha}} X_{i \alpha}\tag{8.4.1.1}
\end{equation*}\etiquette[8.4.1.1]{VI.eq8.4.1.1}
pour tout pro-objet $\sheaf{\alpha} = (i_{\alpha})$ de $I$ (pour voir ceci
on pourra généraliser aux pseudo foncteurs l'exercice I \hyperref[VI.subsubsec8.2.8]{8.2.8}).

\subsubsection{}\etiquette[8.4.2]{VI.subsubsec8.4.2}

Soient $E$ un topos et $p : \text{Fl}(E) \to E$ le topos fibré
considéré dans \hyperref[VI.exem7.3.1]{7.3.1}. En prolongeant comme en
\hyperref[VI.subsubsec8.4.1]{8.4.1}, on obtient un topos fibré
$\bar{\text{Fl}(E)} \to \Pro(E)$. La catégorie $\Point(E)$ s'envoie
par un foncteur pleinement fidèle dans $\Pro(E)$ (IV 6.8.5),
d'où, par changement de base (\hyperref[VI.subsubsec7.1.9]{7.1.9}), un topos
fibré $\Loc(E) \to \Point(E)$. La fibre $\Loc_p(E)$ en un point $p$
de $E$ est appelée le \emph{topos
  localisé de $E$ en le point} $p$. Ce topos dépend, d'après ce qui
précède, de façon covariante du point $p$. On a, par définition,
\begin{equation*}
\Loc_p(E) = \Limtop_{X \in \Vois(p)} E_{/X}{\quoi},\tag{8.4.2.1}
\end{equation*}\etiquette[8.4.2.1]{VI.eq8.4.2.1}
où $\Vois(p)$ est la catégorie des voisinages de $p$ (IV 6.8).

\subsubsection{}\etiquette[8.4.3]{VI.subsubsec8.4.3}

Soit $m : p'\to p$ un morphisme de $\Point(E)$ (IV 6). On en déduit,
pour tout $X \in \Vois(p)$ un point $p'_X$ de $E_{/X}$, d'où, par la
propriété universelle de $\limtop{}$ \hyperref[VI.subsec8.1]{8.1} un point
$\theta_p(m)$ de $\Loc_p(E)$. On définit ainsi un foncteur
\begin{equation*}
\theta_p : \Point(E)_{/p} \to \Point(\Loc_p (E)){\quoi},\tag{8.4.3.1}
\end{equation*}\etiquette[8.4.3.1]{VI.eq8.4.3.1}
dont on constate immédiatement que c'est une équivalence de
catégories. Par suite $\Point(\Loc_p(E))$ est canoniquement équivalent
à la catégorie des générisations de $p$ (IV 7.1.8). On a de plus une
équivalence canonique de topos:
$$
\can \bigdot : \Loc_{\theta(m)} (\Loc_p (E)) \to \Loc_{p'} (E){\quoi},
$$
qui s'insère dans un diagramme commutatif à isomorphisme près
\[
\xymatrix@R=1.5cm{
\Loc_{\theta(m)} (\Loc_p(E))  \ar[rr]^{\can\bigdot} \ar[rd]_{\phi}  &&\Loc_{p'}(E)\ar[ld]^{m}\\
                                  & \Loc_p (E)   &       \\
}
\]
où $\phi$\pageoriginale est le morphisme canonique \eqhyperref[VI.eq8.4.2.1]{8.4.2.1}.

\subsubsection{}\etiquette[8.4.4]{VI.subsubsec8.4.4}

Les topos localisés ne semblent présenter un intérêt que dans le cas
des topos provenant de la géométrie algébrique ou tout au moins que
dans le cas des topos possédant des propriétés de finitude
convenables. Ainsi, lorsque $E$ est le topos des faisceaux sur un
espace topologique séparé, on constate immédiatement (à l'aide par
exemple de \hyperref[VI.thm8.2.9]{8.2.9}) que les topos localisés sont des topos
ponctuels. Par ailleurs dans ce cas, la catégorie $\Point(E)$ est
discrète et la situation décrite en \hyperref[VI.subsubsec8.4.2]{8.4.2} est
triviale. En revanche, soient $E = \Top(X)$ le topos des faisceaux
pour la topologie de Zariski sur un schéma $X$, $x$ un point de $X$,
$O_x$ l'anneau local de $X$ en $x$, $Y = \spec(O_x)$. On constate
que
$$
\Loc_x(E) = \Top(Y){\quoi}.
$$
Pour d'autres exemples provenant de la topologie étale, nous renvoyons
à l'exposé VIII du présent séminaire.

\subsubsection{}\etiquette[8.4.5]{VI.subsubsec8.4.5}

Lorsque $E$ est localement cohérent (\hyperref[VI.def2.3]{2.3}), les topos localisés
$\Loc_p(E)$ sont cohérents (on peut dans \eqhyperref[VI.eq8.4.2.1]{8.4.2.1} se borner aux $X \in
\Vois(p)$ qui sont algébriques et cohérents et appliquer alors
\hyperref[VI.thm8.3.13]{8.3.13}). Pour tout morphisme de points $m : p'\to p$, le morphisme de
topos $m\bigdot : \Loc_{p'}(E) \to \Loc_p(E)$ est cohérent (\hyperref[VI.cor8.3.14]{8.3.14}).

\subsubsection{}\etiquette[8.4.6]{VI.subsubsec8.4.6}

On appelle \emph{topos local} un topos $X$ tel que le foncteur
$\Gamma(X, -)$ = « section sur $X$ » (IV 4.3) soit un foncteur
fibre (IV 6). Le point correspondant à ce foncteur fibre est appelé le
\emph{centre} du topos local. Lorsque $E$ est localement cohérent, les
topos localisés sont des topos locaux (1.2.3, \hyperref[VI.prop8.5.2]{8.5.2} et \hyperref[VI.cor8.5.7]{8.5.7}). Le
centre de $\Loc_p(E)$ est canoniquement isomorphe à $\theta_p(\id_p)$
\eqhyperref[VI.eq8.4.3.1]{8.4.3.1}). L'image du centre de $\Loc_p(E)$ dans $E$ est isomorphe à
$p$. Le topos localisé $\Loc_p(E)$ muni du morphisme canonique\pageoriginale
$\Loc_p(E) \to E$ est la solution du problème universel (2-universel !)
qui consiste à envoyer des topos locaux dans $E$ de façon à
envoyer le centre « sur » $p$.

\subsubsection{}\etiquette[8.4.7]{VI.subsubsec8.4.7}

A propos des topos locaux, il se pose un certain nombre de problèmes
que les rédacteurs n'ont pas abordés. Ainsi, si $X$ est un topos
local, l'objet final $e$ de $X$ possède un ouvert maximal $U \neq
e$. Le complémentaire $Y$ de $U$ est un topos local qui ne possède
pas d'ouvert non trivial. Un tel topos est-il ponctuel? Soit $U$ un
topos et $X = (U, \Phi : U \to (\Ens))$ un topos obtenu par
recollement. Quelles sont les conditions sur le foncteur de
recollement $\Phi$ pour que $X$ soit un topos local ?

\subsection{Structure des faisceaux d'une limite projective filtrante
  de topos}\etiquette[8.5]{VI.subsec8.5}

\subsubsection{}\etiquette[8.5.1]{VI.subsubsec8.5.1}

Dans ce numéro, $F = (F_i)_{i \in I}$ est un $𝒰$-topos fibré sur une
petite catégorie cofiltrante $I$. On choisit un biscindage (\hyperref[VI.subsubsec7.1.4]{7.1.4}) de
$F$, \ie pour tout $f \in \text{Fl}(I) : i \to j$ on choisit des morphismes
de topos $f. : F_i \to F_j$ tels que $f^{\ast} : F_j \to F_i$ soit le
foncteur image inverse pour la structure fibrée. On a alors des
isomorphismes canoniques $c_{f, g}$ possédant une propriété de
cocycles (\hyperref[VI.subsubsec7.1.3]{7.1.3}). On note $\underleftarrow{F}$ la limite projective du
topos fibré $F$ (\hyperref[VI.thm8.2.3]{8.2.3}) et pour tout $i \in \ob I$, on note
\begin{equation*}
\mu_i : \underleftarrow{F} \to F_i{\quoi},\tag{8.5.1.1}
\end{equation*}\etiquette[8.5.1.1]{VI.eq8.5.1.1}
le morphisme canonique (\hyperref[VI.subsubsec8.1.3]{8.1.3}). On note $\Top(F)$ le topos total de
$F$ (7.4.3,3). D'après \hyperref[VI.thm8.2.9]{8.2.9}, le foncteur canonique $F \to
\underrightarrow{F}$ définit un morphisme de topos
\begin{equation*}
Q : \underleftarrow{F} \to \Top(F){\quoi}.\tag{8.5.1.2}
\end{equation*}\etiquette[8.5.1.2]{VI.eq8.5.1.2}
Le topos $\underleftarrow{F}$ s'identifie à $\Hhom_{\Cart /I^{\circ}} (I^{\circ}, F')$
(\hyperref[VI.thm8.2.9]{8.2.9}) et le
topos $\Top(F)$ s'identifie à $\Hhom_{I^{\circ}}(I^{\circ}, F')$ (\hyperref[VI.prop7.4.7]{7.4.7}). Ces
identifications faites, le morphisme $Q$ n'est autre que le morphisme
de plongement de $\Hhom_{\cart/I^{\circ}}(I^{\circ}, F')$ dans
$\Hhom_{I^{\circ}}(I^{\circ}, F')$
(\hyperref[VI.thm8.2.9]{8.2.9}) et pour tout objet $M$ de $\underleftarrow{F}$, on a
\begin{equation*}
Q_{\ast} (M) = (i \to \mu_{i\ast}(M)){\quoi}.\tag{8.5.1.3}
\end{equation*}\etiquette[8.5.1.3]{VI.eq8.5.1.3}
D'après 7.4.3.4, on a pour tout $i \in \ob I$, un morphisme de topos
\begin{equation*}
\alpha_i : F_i \to \Top(F){\quoi}.\tag{8.5.1.4}
\end{equation*}\etiquette[8.5.1.4]{VI.eq8.5.1.4}
Le diagramme\pageoriginale
\[\vcenter{
\xymatrix@R=1.25cm{
{\underleftarrow{F}}  \ar[rr]^{\mu_i}\ar[rd]_{Q}&                 &F_i\ar[ld]^{\alpha_i}\\
                                   &  \Top(F)      &
}}\leqno{(8.5.1.5)}
\]\etiquette[8.5.1.5]{VI.eq8.5.1.5}
n'est pas commutatif en général (même à isomorphisme près). Mais on a,
par définition des morphismes en présence, un isomorphisme canonique
\begin{equation*}
\alpha^{\ast}_i Q_{\ast} \simeq \mu_{i\ast}\tag{8.5.1.6}{\quoi}.
\end{equation*}\etiquette[8.5.1.6]{VI.eq8.5.1.6}
Lorsque $i$ est un \emph{objet final} de $I$, $\alpha_i$ est un
plongement admettant une rétraction $\beta_i : \Top(F) \to F_i$
(\hyperref[VI.lem7.4.12]{7.4.12}) et le diagramme
\sisi{\def\SGAVIEightFiveOneSevenTarget{\Top(X)}}{\def\SGAVIEightFiveOneSevenTarget{\Top(F)}}
\[\vcenter{
\xymatrix@R=1.25cm{
{\underleftarrow{F}}  \ar[rr]^{\mu_i}\ar[rd]_{Q}&                 &F_i\ar[ld]^{\beta_i}\\
                                   &  \SGAVIEightFiveOneSevenTarget      &       ,
}}\leqno{(8.5.1.7)}
\]\etiquette[8.5.1.7]{VI.eq8.5.1.7}
est commutatif à isomorphisme canonique près.

\begin{enonce*}{Proposition 8.5.2}\etiquette[8.5.2]{VI.prop8.5.2}
Soit $j \mapsto M_j$ un objet de $\Top(F)$. Il existe un
isomorphisme fonctoriel
\begin{equation*}
Q^{\ast} (j \mapsto M_j)\simeq \mathop{\varinjlim}_{I^{\circ}}
\mu_j^{\ast} M_{j}{\quoi}.\tag{8.5.2.1}
\end{equation*}\etiquette[8.5.2.1]{VI.eq8.5.2.1}
\end{enonce*}
\noindent
Un objet de $\Top(F)$ (\hyperref[VI.subsubsec7.1.3]{7.1.3}) consiste en la donnée d'une application
$j \mapsto M_j$, $M_j \in \ob F_j$, et en la donnée, pour tout
morphisme $f : i \to j$, d'un morphisme
\begin{equation*}
\beta_f : M_j \to f_{\ast} M_i\tag{8.5.2.2}
\end{equation*}\etiquette[8.5.2.2]{VI.eq8.5.2.2}
ou de manière équivalence par adjonction, d'un morphisme
\begin{equation*}
\beta'_f:  f^{\ast} M_j \to M_i{\quoi},\tag{8.5.2.3}
\end{equation*}\etiquette[8.5.2.3]{VI.eq8.5.2.3}
les morphismes $\beta'_f$ étant soumis à la condition que pour tout
couple de morphismes composables $i \xrightarrow{f} j \xrightarrow{g}
k$, le diagramme ci-après soit commutatif:
\[\vcenter{
\xymatrix@C=2.25cm@R=1.5cm{
f^{\ast}g^{\ast}M_k    \ar[r]^{f^{\ast}\beta'_g} \ar@{-}[d]_{c^{\ast}_{g,f}}&
f^{\ast}M_j \ar@{->}[d]^{\beta'_f}\\
(gf)^{\ast}M_k         \ar[r]^{\beta'_{gf}}      & M_i
}}\leqno{(8.5.2.4)}
\]\etiquette[8.5.2.4]{VI.eq8.5.2.4}
Rappelons\pageoriginale de plus \eqhyperref[VI.eq8.1.3.1]{8.1.3.1.} que pour tout morphisme $f :
i \to j$, on a un isomorphisme $b_f : f. \mu_i \xrightarrow{\sim} \mu_j$, et que
pour tout couple de morphismes composables $i \xrightarrow{f} j
\xrightarrow{g} k$, on a un diagramme commutatif:
\[
\xymatrix@C=2cm@R=1.25cm{
g\bigdot f\bigdot \mu_i  \ar[r]^{c_{g,f}} \ar[d]^{g\bigdot (b_f)}&
(gf) \bigdot \mu_i \ar[d]^{b_{gf}}\\
g \bigdot \mu_j  \ar[r]^{b_g} & \mu_k
}\leqno{(8.5.2.5)}
\]\etiquette[8.5.2.5]{VI.eq8.5.2.5}
Soit alors $f : i \to j$ un morphisme de $I$. Notons
\begin{equation*}
t_f : \mu^{\ast}_j (M_j) \to \mu_i^{\ast} (M_i)\tag{8.5.2.6}
\end{equation*}\etiquette[8.5.2.6]{VI.eq8.5.2.6}
le morphisme composé $\mu^{\ast}_j (M_j)\xrightarrow{b^{\ast}_f}
\mu^{\ast}_i(f^{\ast}(M_j)) \xrightarrow{\mu^{\ast}_i(\beta'_f)}
\mu^{\ast}_i(M_i)$.
Nous laissons au lecteur le soin de vérifier que la commutativité des
diagrammes \eqhyperref[VI.eq8.5.2.4]{8.5.2.4} et \eqhyperref[VI.eq8.5.2.5]{8.5.2.5}
entraîne que pour tout couple de
morphismes composables $i \to j \to k$, on a
\begin{equation*}
t_f t_g = t_{gf}{\quoi}.\tag{8.5.2.7}
\end{equation*}\etiquette[8.5.2.7]{VI.eq8.5.2.7}
Par suite on a défini un foncteur $I^{\circ} \to \underleftarrow{F}$
dont on peut
considérer la limite inductive $\varinjlim_{I^{\circ}}
\mu^{\ast}_i (M_i)$. Soit alors $N$ un objet de $\underleftarrow{F}$. On a
\begin{equation*}
\Hom(\varinjlim_{I^{\circ}} \mu^{\ast}_i (M_i), N) \simeq
\varprojlim_{I^{\circ}} \Hom (\mu^{\ast}_i (M_i), N){\quoi},\tag{8.5.2.8}
\end{equation*}\etiquette[8.5.2.8]{VI.eq8.5.2.8}
d'où, par adjonction, un isomorphisme
\begin{equation*}
\Hom(\varinjlim_{I^{\circ}} \mu^{\ast}_i (M_i), N) \simeq
\varprojlim_{I^{\circ}} \Hom_{F_i} (M_i, \mu_{i \ast} (N)){\quoi}.\tag{8.5.2.9}
\end{equation*}\etiquette[8.5.2.9]{VI.eq8.5.2.9}
On\pageoriginale se propose d'interpréter le deuxième membre de
\eqhyperref[VI.eq8.5.2.9]{8.5.2.9}. Pour cela, notons, pour tout morphisme $f : i \to
j$, par
$$
d_f : \Hom_{F_i}(M_i, \mu_{i\ast}(N)) \to \Hom_{F_j} (M_j,
\mu_{j\ast} (N)){\quoi}.
$$
le morphisme de transition du système projectif qui figure dans
\eqhyperref[VI.eq8.5.2.9]{8.5.2.9}. Il résulte de la définition des morphismes $t_f$ \eqhyperref[VI.eq8.5.2.6]{8.5.2.6}
que l'application $d_f$ associe à un morphisme
$$
u_i : M_i \to \mu_{i\ast}(N)
$$
le morphisme
$$
d_f(u_i) : M_{j} \to \mu_{j\ast}(N){\quoi},
$$
obtenu en composant les morphismes
$$
M_j \xrightarrow{\beta_f} f_{\ast}(M_i) \xrightarrow{f_{\ast} (u_i)}
f_{\ast} \mu_{i\ast} (N) \xrightarrow{b_f} \mu_{j\ast} (N){\quoi}.
$$
Par suite un élément du deuxième membre de \eqhyperref[VI.eq8.5.2.9]{8.5.2.9} s'interprète
comme une famille de morphismes $u_i : M_{i} \to \mu_{i\ast} (N)$, $i
\in \ob I$, telle que pour tout $f : i \to j$, le diagramme
\[
\xymatrix@C=2cm@R=1.25cm{
M_j  \ar[r] \ar[d]_{u_j}& f_{\ast}(M_i) \ar[d]^{f_{\ast}(u_i)}\\
\mu_{j\ast}(N)  \ar[r]^{\sim} & f_{\ast}\mu_{i\ast}(N)
}
\]
soit commutatif, c'est-à dire comme un morphisme de la section $(i
\mapsto M_i) \in \Hhom_{I^{\circ}} (I^{\circ}, F')$ dans la
section $Q_{\ast}(N)$. On a donc un isomorphisme fonctoriel
$$
\Hom(\varinjlim_{I^{\circ}} \mu_i^{\ast} (M_i), N)
\simeq \Hom((i \mapsto M_i), Q_{\ast}(N)){\quoi};
$$
d'où la proposition par adjonction.

\begin{enonce*}{Proposition 8.5.3}\etiquette[8.5.3]{VI.prop8.5.3}
Si pour tout $f : i \to j \in \text{\rm Fl}(I)$ les foncteurs
$f_{\ast} : F_i \to F_j$ commutent aux petites limites inductives filtrantes, on a, pour toute
section $(i \mapsto M_i) \in \Hom_{I^{\circ}}(I^{\circ}, F')$:
\begin{equation*}
Q_{\ast} Q^{\ast} (i \mapsto M_i) \simeq i \mapsto
\varinjlim_{f : j \to i} f_{\ast} M_j{\quoi}.\tag{8.5.3.1}
\end{equation*}\etiquette[8.5.3.1]{VI.eq8.5.3.1}
\end{enonce*}

\noindent Nous\pageoriginale laissons au lecteur le soin
d'expliciter en termes des $\beta_f$ \eqhyperref[VI.eq8.5.2.2]{8.5.2.2} et des $c_{f,
g}$ (\hyperref[VI.subsubsec8.5.1]{8.5.1}) les morphismes de transition des
systèmes inductifs $(f : j \to i) \mapsto f_{\ast}(M_j)$ et les
morphismes de transition de la section $i \mapsto \varinjlim_{f : j
\to i} f_{\ast} M_{j}$. Comme les foncteurs $f_{\ast}$ commutent aux
limites inductives filtrantes, la section $i\mapsto
\varinjlim_{f:j\to i} f_{\ast}M_{j}$ est cartésienne. De plus on a
un morphisme naturel $u$ de $(i \mapsto M_i)$ dans la section $(i
\mapsto \varinjlim_{f
  : j \to i}f_{\ast} M_j)$, et il est clair que tout morphisme de $(i
\mapsto M_i)$ dans une section cartésienne se factorise d'une manière
unique par $u$, d'où l'isomorphisme \eqhyperref[VI.eq8.5.3.1]{8.5.3.1}.

\begin{enonce*}{Corollaire 8.5.4}\etiquette[8.5.4]{VI.cor8.5.4}
Sous les hypothèses de \hyperref[VI.prop8.5.3]{8.5.3} le foncteur $Q_{\ast}$ commute aux
petites limites inductives filtrantes.
\end{enonce*}

Soit $\alpha \mapsto N^{\alpha}$ un petit système inductif filtrant de
$\underleftarrow{F}$. Posons $M^{\alpha} = Q_{\ast} (N^{\alpha})$, de
sorte que $M^{\alpha}$ est une section $i \mapsto M^{\alpha}_i$. On a
$N^{\alpha} \simeq Q^{\ast} Q_{\ast} (N^{\alpha})$ et comme $Q^{\ast}$
commute aux limites inductives, on a
$\varinjlim_{\alpha}{N^{\alpha}}\simeq Q^{\ast}
(\varinjlim_{\alpha}{M^{\alpha}})$. Les limites
inductives dans $\Top(F) = \Hom_{I^{\circ}}(I^{\circ}, F')$ se
calculent fibre par
fibre. On a donc $\varinjlim_{\alpha} {M^{\alpha}} \simeq
(i \mapsto \varinjlim_{\alpha}
{M^{\alpha}_{i}})$. En vertu de \eqhyperref[VI.eq8.5.3.1]{8.5.3.1}, on a
$$
Q_{\ast} (\varinjlim_{\alpha} {N}^{\alpha}) \simeq Q_{\ast}
Q^{\ast} (\varinjlim_{\alpha} {M}^{\alpha}) \simeq i \mapsto
\varinjlim_{f : j \to i} f_{\ast}
(\varinjlim_{\alpha} {M}^{\alpha}_j){\quoi}.
$$
Comme les foncteurs $f_{\ast}$ commutent aux limites inductives
filtrantes, il vient:
$$
Q_{\ast} (\varinjlim_{\alpha} {N}^{\alpha}) \simeq i \mapsto
\varinjlim_{f: j\to i} \varinjlim_{\alpha} f_{\ast}
(M^{\alpha}_j) \simeq i \mapsto \varinjlim_{\alpha} \varinjlim_{f:
  j\to i} f_{\ast}
(M^{\alpha}_j) \simeq \varinjlim_{\alpha} Q_{\ast} (N^{\alpha}){\quoi}.
$$

\begin{enonce*}{Corollaire 8.5.5}\etiquette[8.5.5]{VI.cor8.5.5}
Si les foncteurs $f_{\ast} : F_{i} \to F_{j}$ commutent aux petites
limites inductives filtrantes, on a, pour tout objet $i$ de $I$ et
pour tout $M_i \in \ob F_i$
\begin{equation*}
\mu_{i\ast} \mu_{i}^{\ast} (M_i) \simeq \varinjlim_{f : j \to
      i} f_{\ast} f^{\ast}(M_i){\quoi}.\tag{8.5.5.1}
\end{equation*}\etiquette[8.5.5.1]{VI.eq8.5.5.1}
\end{enonce*}
\noindent
Quitte à faire le changement de base $I/i \to I$, on peut supposer que
$i$ est un objet final de $I$ (\hyperref[VI.lem8.2.1]{8.2.1}). Le morphisme $\mu_{i} :
\underleftarrow{F} \to F_i$ est alors le morphisme composé
des\pageoriginale morphismes (\hyperref[VI.thm8.2.9]{8.2.9}) et (\hyperref[VI.lem7.4.12]{7.4.12}):
\begin{align*}
Q &: \underleftarrow{F} \to F^{\sim}{\quoi},\\
\beta_i &: (\alpha_{i!}, \alpha^{\ast}_i) : F^{\sim} \to F_i{\quoi}.
\end{align*}
On a donc $\mu_i^{\ast} (M_i) \simeq Q^{\ast} \beta_i^{\ast} (M_i)$,
où $\beta_i^{\ast} (M_i)$ est la section $(f : i \to j) \mapsto
f^{\ast}(M_i)$ dont les morphismes de transition sont déduits des
$c_{f,g}$. L'assertion résulte alors de \hyperref[VI.prop8.5.3]{8.5.3}.

\begin{enonce*}{Corollaire 8.5.6}\etiquette[8.5.6]{VI.cor8.5.6}
Sous les hypothèses de \hyperref[VI.cor8.5.5]{8.5.5}, les foncteurs
$$
\mu_{i\ast} : \underleftarrow{F} \to F_i
$$
commutent aux limites inductives filtrantes.
\end{enonce*}

Le foncteur $\mu_{i\ast}$ est composé du foncteur $Q_{\ast}$ qui
commute aux limites inductives filtrantes, et du foncteur « restriction à la fibre en $i$ », qui commute aux limites inductives.

\begin{enonce*}{Corollaire 8.5.7}\etiquette[8.5.7]{VI.cor8.5.7}
Sous les hypothèses de \hyperref[VI.prop8.5.3]{8.5.3}, soient $i$ un objet de $I$ et $X$ un
objet de $F_i$ tel que le foncteur $\Hom(X, -)$ sur $F_i$ commute aux
limites inductives filtrantes. Alors le foncteur $\Hom(\mu^{\ast}_i
(X), -)$ commute aux limites inductives filtrantes et pour tout objet
$M_i$ de $F_i$ on a
\begin{equation*}
\Hom(\mu^{\ast}_i(X), \mu^{\ast}_i(M_i)) \simeq \varinjlim_{f
: j \to i} \Hom(f^{\ast}(X), f^{\ast}(M_i))\tag{8.5.7.1}{\quoi}.
\end{equation*}\etiquette[8.5.7.1]{VI.eq8.5.7.1}
\end{enonce*}
\noindent
Pour tout objet $(j \mapsto M_j)$ de $\Top(F)$ on a
\begin{equation*}
\Hom(\mu^{\ast}_i(X), Q^{\ast}(j \mapsto M_j)) \simeq \varinjlim_{f
: j \to i} \Hom(f^{\ast}(X), M_{j}){\quoi}.\tag{8.5.7.2}
\end{equation*}\etiquette[8.5.7.2]{VI.eq8.5.7.2}

Le foncteur $\Hom(\mu^{\ast}_i (X), -)$ est isomorphe au foncteur
$\Hom(X, \mu_{i \ast} (-))$ et ce dernier commute aux limites
inductives filtrantes (\hyperref[VI.cor8.5.4]{8.5.4}). D'après \eqhyperref[VI.eq8.5.5.1]{8.5.5.1}, on a
$$
\Hom(\mu^{\ast}_i(X), \mu^{\ast}_i (M_i)) \simeq \Hom (X, \mu_{i\ast}
  \mu^{\ast}_i (M_i)) \simeq  \Hom(X, \varinjlim_{f : j \to
  i} f_{\ast} f^{\ast} (M_i)){\quoi},
$$
d'où la formule \eqhyperref[VI.eq8.5.7.1]{8.5.7.1}. De même, d'après \eqhyperref[VI.eq8.5.3.1]{8.5.3.1}, on a
$$
\Hom(\mu^{\ast}_i (X), Q^{\ast}(j \mapsto M_j)) \simeq \Hom(X,
\mu_{i\ast} Q^{\ast} (j \mapsto M_j)) \simeq \Hom(X,
\varinjlim_{f : j \to i} f_{\ast} (M_j)){\quoi},
$$
d'où\pageoriginale la formule \eqhyperref[VI.eq8.5.7.2]{8.5.7.2}.

\begin{enonce*}{Proposition 8.5.8}\etiquette[8.5.8]{VI.prop8.5.8}
Soient $G \to I$ un $𝒰$-topos fibré et $m : F \to G$  un morphisme
cartésien de topos fibrés (\sisi{7.1.15}{7.1.5}). On
utilise pour $G$ les notations
introduites en \hyperref[VI.def8.1.1]{8.1.1}. Soit $\underleftarrow{m} : \underleftarrow{F} \to
\underleftarrow{G}$ le morphisme déduit de $m$ par passage à la limite
projective (\hyperref[VI.subsubsec8.1.4]{8.1.4}). Le diagramme de topos et de morphismes de topos:
\[\vcenter{
\xymatrix@C=2.25cm@R=1.25cm{
\underset{\leftarrow}{F}  \ar[r]^{Q} \ar[d]^{\underset{\leftarrow}{m}}& \Top(F)\ar[d]^{m^{\sim}}\\
\underset{\leftarrow}{G}  \ar[r]^{Q} & \Top(G)
}}\leqno{(8.5.8.1)}
\]\etiquette[8.5.8.1]{VI.eq8.5.8.1}
est commutatif à isomorphisme près. Pour tout objet $N$ de
$\underleftarrow{F}$, on a
\begin{equation*}
\underleftarrow{m}_{\ast} (N) \simeq \varinjlim_{I^{\circ}} \mu_j^{\ast}
m_{j\ast} \mu_{j\ast} (N){\quoi}.\tag{8.5.8.2}
\end{equation*}\etiquette[8.5.8.2]{VI.eq8.5.8.2}
\end{enonce*}

La commutativité du diagramme \eqhyperref[VI.eq8.5.8.1]{8.5.8.1} résulte
immédiatement des définitions (\hyperref[VI.subsubsec8.2.8]{8.2.8} et
\hyperref[VI.subsubsec8.5.1]{8.5.1}). On a $\underleftarrow{m}_{\ast} (N) \simeq
Q^{\ast} Q_{\ast} \underleftarrow{m}_{\ast} (N)$ et, en vertu de la
commutativité de \eqhyperref[VI.eq8.5.8.1]{8.5.8.1}, $\underleftarrow{m}_{\ast} (N)
\simeq Q^{\ast}m^{\sim}_{\ast} Q_{\ast} (N)$. Par suite, on a un
isomorphisme $\underleftarrow{m}_{\ast} (N) \simeq Q^{\ast}
m^{\sim}_{\ast} (j \mapsto \mu_{j\ast} (N))$. Le foncteur
$m_{\ast}^{\sim}$ n'est autre que le foncteur $\Hom_{I^{\circ}}
(I^{\circ}, m_{\ast})$ (\hyperref[VI.prop7.4.10]{7.4.10}), et par suite on a un
isomorphisme canonique $\underleftarrow{m}_{\ast} (N) \simeq
Q^{\ast} (j \mapsto m_{j\ast} \mu_{j\ast} (N))$. La formule
\eqhyperref[VI.eq8.5.8.2]{8.5.8.2} résulte alors de \hyperref[VI.prop8.5.2]{8.5.2}.

\begin{enonce*}{Proposition 8.5.9}\etiquette[8.5.9]{VI.prop8.5.9}
Avec les hypothèses et notations de \hyperref[VI.prop8.5.8]{8.5.8}, on suppose de plus que pour
tout objet $i$ de $I$, le foncteur $m_{i\ast} : F_{i} \to G_i$ commute aux
limites inductives filtrantes et que pour tout morphisme $f : i \to j$
de $I$ le foncteur $f_{\ast} : F_i \to F_j$ commute aux limites
inductives filtrantes. Alors le foncteur $\underleftarrow{m}_{\ast} :
\underleftarrow{F} \to \underleftarrow{G}$ commute aux limites inductives
filtrantes, et pour tout objet $(i \mapsto M_i)$ de $\Top(F)$ on a un
isomorphisme canonique:
\begin{equation*}
\underleftarrow{m}_{\ast} Q^{\ast} (i \mapsto M_i) \simeq
\varinjlim_{I^{\circ}} \mu^{\ast}_j m_{j\ast} (M_j){\quoi}.\tag{8.5.9.1}
\end{equation*}\etiquette[8.5.9.1]{VI.eq8.5.9.1}
\end{enonce*}

La première\pageoriginale assertion résulte de \eqhyperref[VI.eq8.5.8.2]{8.5.8.2} et de (\hyperref[VI.cor8.5.6]{8.5.6}). D'après
\eqhyperref[VI.eq8.5.8.2]{8.5.8.2} et \eqhyperref[VI.eq8.5.3.1]{8.5.3.1}, on a un isomorphisme fonctoriel:
$$
\underleftarrow{m}_{\ast} Q^{\ast} (i \mapsto M_i) \simeq
\varinjlim_i \mu^{\ast}_{i} m_{i\ast} (\varinjlim_{f
  : j \to i} f_{\ast} (M_{j})){\quoi}.
$$
En utilisant la commutation des $m_{i\ast}$ aux limites inductives
filtrantes et les isomorphismes $m_{i\ast} f_{\ast} \simeq f_{\ast}
m_{j\ast}$ (\hyperref[VI.subsubsec7.1.6]{7.1.6}), on obtient:
$$
\underleftarrow{m}_{\ast} Q^{\ast} (i \mapsto M_i) \simeq
\varinjlim_i \mu^{\ast}_{i} (\varinjlim_{f : j \to
  i} f_{\ast} m_{j\ast} (M_j)){\quoi}.
$$
Comme les foncteurs $\mu^{\ast}_{i}$ commutent aux limites inductives,
il vient:
$$
\underleftarrow{m}_{\ast}Q^{\ast}(i\mapsto M_{i})\simeq \varinjlim_{i}
\varinjlim_{f:j\to i} \mu^{\ast}_{i}f_{\ast}m_{j\ast}(M_{j}){\quoi}.
$$
Ce dernier objet peut être interprété comme une limite inductive sur
la catégorie $\text{Fl}(I)^{\circ}$ où $\text{Fl}(I)$ est la catégorie des
morphismes de $I$. Soit alors $\phi : I \to \text{Fl}(I)$ le foncteur qui associe à
tout objet $i$ de $I$ le morphisme identique de $\sisi{I}{i}$. Le foncteur
$\phi^{\circ}
: I^{\circ} \to \text{Fl}(I)^{\circ}$ est cofinal et par suite
(I 8.1) on a un isomorphisme canonique
$$
\underleftarrow{m}_{\ast} Q^{\ast} (i \mapsto M_i) \simeq
\varinjlim_{I^{\circ}} \mu^{\ast}_{j} m_{j\ast}(M_j){\quoi}.
$$

\begin{enonce*}{Corollaire 8.5.10}\etiquette[8.5.10]{VI.cor8.5.10}
Sous les conditions de \hyperref[VI.prop8.5.9]{8.5.9}, pour tout objet $i$ de $I$ et tout objet
$M_i$ de $F_i$, on a un isomorphisme canonique
\begin{equation*}
\underleftarrow{m}_{\ast} \mu^{\ast}_i (M_i) \simeq
\varinjlim_{f:j\to i} \mu^{\ast}_{j} m_{j\ast}f^{\ast}
(M_{i}){\quoi}.\tag{8.5.10.1}
\end{equation*}\etiquette[8.5.10.1]{VI.eq8.5.10.1}
\end{enonce*}

La démonstration est analogue à celle du corollaire \hyperref[VI.cor8.5.5]{8.5.5}.

\begin{enonce*}[remark]{Remarque 8.5.11}\etiquette[8.5.11]{VI.rem8.5.11}
Rappelons que le foncteur image directe par un morphisme cohérent
entre topos cohérents commute aux limites inductives filtrantes
(\hyperref[VI.thm5.1]{5.1}). Par suite les propositions \hyperref[VI.prop8.5.3]{8.5.3} à \hyperref[VI.cor8.5.7]{8.5.7} et
\hyperref[VI.prop8.5.9]{8.5.9}, \hyperref[VI.cor8.5.10]{8.5.10}
s'appliquent lorsque les topos fibrés envisagés sont cohérents et les
morphismes de topos fibrés envisagés sont cohérents.
\end{enonce*}

\subsection{$𝒰$-topos fibrés annelés}\etiquette[8.6]{VI.subsec8.6}

\subsubsection{}\etiquette[8.6.1]{VI.subsubsec8.6.1}

On dit qu'un $𝒰$-topos fibré est \emph{annelé} s'il est muni
d'un faisceau d'anneaux sur le site total (\hyperref[VI.def7.4.1]{7.4.1}). Soit
$p : F \to I$ un $𝒰$-topos fibré. Choisissons un
biscindage\pageoriginale (\hyperref[VI.subsubsec7.1.4]{7.1.4}). Il résulte de
\hyperref[VI.prop7.4.7]{7.4.7} que se donner un faisceau d'anneaux sur $F$
revient à se donner, pour tout objet $i$ de $I$, un anneau $A_i$ de
$F_i$, et pour tout morphisme $f : i \to j$, un morphisme $\phi_f :
A_j \to f_{\ast} A_i$, la famille des $\phi_f$ étant soumise à des
conditions de compatibilité explicitées dans \loccit. Les morphismes
de topos $f. : F_i \to F_j$ sont donc des morphismes de topos
annelés (respectivement par $A_i$ et $A_j$) (IV 11), et la structure
annelée sur $F$ et le choix des morphismes $f$. fournit un
pseudo-foncteur de $I$ dans la 2-catégorie des topos annelés.
Réciproquement, lorsqu'on se donne un tel pseudo-foncteur, on peut
reconstruire un $𝒰$-topos fibré annelé qui lui donne
naissance.

\subsubsection{}\etiquette[8.6.2]{VI.subsubsec8.6.2}

On peut comme en \hyperref[VI.subsec8.1]{8.1}, définir la notion de limite
projective d'un $𝒰$-topos annelé $(F_i, A_i, i \in I)$. Nous
supposerons que cette généralisation immédiate a été faite et nous
appliquerons librement à cette situation le langage introduit dans
\hyperref[VI.subsec8.1]{8.1}. Lorsqu'elle existe, cette limite projective est
un $𝒰$-topos annelé $(H, B)$, et on a pour tout objet $i$ de
$I$ des morphismes de topos annelés $\mu_i : (H, B) \to (F_i, A_i)$.
Si le topos fibré $F$ (non annelé) admet une limite projective
$\limtop{I} F_i$, et si le système inductif d'anneaux $i \mapsto
\mu_i^{\ast}(A_i)$ admet une limite inductive
$\varinjlim_{I^{\circ}}\mu_i^{\ast} A_i$ dans la catégorie des
anneaux de $\limtop{I} F_i$ (ce qui est toujours le cas si $I$ est
petite) alors le topos annelé $(\limtop{I} F_i$,
$\varinjlim_{I^{\circ}} \mu^{\ast}_i A_i$) est de façon évidente
une limite projective du topos fibré annelé considéré.

\subsubsection{}\etiquette[8.6.3]{VI.subsubsec8.6.3}

Les formules du numéro \hyperref[VI.subsec8.5]{8.5}. établies pour les
faisceaux d'ensembles sont valables pour les faisceaux abéliens.
Elles sont aussi valables pour les faisceaux de modules, à condition
d'interpréter les foncteurs images inverses qui y figurent comme des
foncteurs images inverses au sens des faisceaux de modules (IV 1.1).
Nous en laissons la vérification au lecteur.

\subsubsection{}\etiquette[8.6.4]{VI.subsubsec8.6.4}

Soit $(p : F \to I, A)$ un $𝒰$-topos fibré annelé sur une catégorie
cofiltrante. On dit que $(p : F \to I, A)$ est $𝒰$-topos fibré annelé
plat à droite (\resp à gauche) si pour tout morphisme $f : i \to j$,
le morphisme $f. : (F_i, A_i) \to (F_j, A_j)$ est un morphisme plat à
droite (\resp à gauche) de topos annelés (\hyperrefext[V][V.def1.8]{1.8}). Supposons $I$
essentiellement petite. Alors le morphisme de topos annelé
 $Q : (\limtop{I} F_i,
\varinjlim_{I^{\circ}} \mu^{\ast}_i (A_i)) \to (\Hom_{I^{\circ}}
(I^{\circ}, F'), A)$\pageoriginale (\hyperref[VI.subsubsec8.5.1]{8.5.1}) est
plat à droite (\resp à gauche), comme il résulte de la définition et
de \hyperref[VI.prop8.5.2]{8.5.2}, et ceci est valable sans hypothèses de
platitude sur les morphismes de transition $f.$. Si on suppose de
plus que $(p : F \to I, A)$ est un $𝒰$-topos fibré annelé
plat à droite (\resp à gauche), les morphismes $\mu_i : (\limtop{I}
F_i, \varinjlim_{I^{\circ}} \mu^{\ast}_i A_i) \to (F_i, A_i)$ sont
plats à droite (\resp à gauche). Ceci résulte de ce que l'image
inverse d'un module plat est plat (\hyperrefext[V][V.cor1.7.1]{1.7.1}), et de ce
qu'une limite inductive filtrante de Modules plats est un Module
plat. On démontre par les mêmes arguments le fait suivant : Si $m :
(F_i, A_i, i \in I) \to (G_i, A_i, i \in I)$ est un morphisme
cartésien de $𝒰$-topos fibrés annelés sur une catégorie
cofiltrante essentiellement petite et si pour tout objet $i \in \ob
(I)$, le morphisme $m_i$ est plat à droite (\resp à gauche), alors
le morphisme $\underleftarrow{m}$ déduit de $m$ par passage à la
limite projective est plat à droite (\resp à gauche).

\subsection{Cohomologie des faisceaux d'une limite projective de
  topos}\etiquette[8.7]{VI.subsec8.7}

\subsubsection{}\etiquette[8.7.1]{VI.subsubsec8.7.1}

Dans ce numéro, $I$ désigne une petite catégorie cofiltrante, $(p : F
\to I, A)$ et $(q : G \to I, B)$ deux $𝒰$-topos fibrés annelés, $m :
(p : F \to I, A) \to (q : G \to I, B)$ un morphisme de $𝒰$-topos
fibrés annelés. On suppose que pour tout morphisme $f : i \to j$ de
$I$, les foncteurs dérivés $R^n f_{\ast}$, $n \in \NN$, du
foncteur $f_{\ast} : \Mod(F_i, A_i)\to \Mod(F_j, A_j)$ pour les modules,
commutent aux limites inductives filtrantes, et que pour tout objet $i
\in I$, les foncteurs dérivés $R^{n} m_{i\ast}$ du foncteur
$m_{i\ast} : (F_i, A_i) \to \sisi{(G_i, B_j)}{(G_i, B_i)}$ pour les modules\footnote{Pour
fixer les idées nous prendrons les modules à gauche.} commutent aux
limites inductives filtrantes. Rappelons que ces hypothèses sont
satisfaites lorsque les catégories $F_i$, $G_i$ $(i \in \ob (I))$ sont
\emph{cohérentes} et lorsque les morphismes $f. : F_i \to F_j$, $(f :
i \to j \in \operatorname{Fl}(I))$ et $m_i : F_i \to G_i$, $(i \in I)$ sont
\emph{cohérents} (\hyperref[VI.cor5.2]{5.2}). On utilise les notations de \hyperref[VI.subsubsec8.5.1]{8.5.1}. De plus on
note $\underrightarrow{A}$ (\resp $\underrightarrow{B}$) le faisceau
$\varinjlim\mu^{\ast}_i (A_i)$ (\resp $\varinjlim\mu^{\ast}_i
(B_i)$). Enfin pour les Modules, on utilise la notation $\mu^{-1}_i$
pour désigner l'image inverse au sens des faisceaux abéliens en
réservant la notation $\mu_i^{\ast}$ pour l'image inverse au sens des
Modules (IV 11).

\begin{enonce*}{Lemme 8.7.2}\etiquette[8.7.2]{VI.lem8.7.2}
Soit\pageoriginale $j \mapsto M_j$ un $A$-Module injectif de $\Top(F)$. Alors
$Q^{\ast} (j \mapsto M_j)$ est un $A$-Module acyclique pour
$\underleftarrow{m}_{\ast}$ et pour tout $i \in \ob(I)$, $M_i$ est flasque.
\end{enonce*}

Soient $i \in \ob (I)$ et $e_i$ un objet final de $F_i$. Le
$U$-topos fibré $F_{/e_{i}} \to I_{/i}$ est déduit de $F \to I$ par
le changement de base $I_{/i} \to I$. Le faisceau $j \to M_j$ étant
injectif est flasque (\hyperrefext[V][V.subsec4.6]{4.6}). Sa restriction au topos
localisé $F_{/e_i}$ est flasque (\hyperrefext[V][V.subsec4.12]{4.12}). Le foncteur de
restriction de $F_{/e_i}$ à $F_i$ est un foncteur image directe par
un morphisme de topos (\hyperref[VI.lem7.4.12]{7.4.12}). Par suite, il transforme
les Modules flasques en Modules flasques (\hyperrefext[V][V.prop5.2]{5.2}). Donc
$M_i$ est flasque.

Démontrons la première assertion du lemme. Posons $N' = Q^{\ast} (j
\mapsto M_j)$. D'après \eqhyperref[VI.eq8.5.3.1]{8.5.3.1}, on a, pour tout $i \in \ob(I)$
$$
\mu_{i\ast}(N') = \mu_{i\ast} Q^{\ast} (j \mapsto M_j) \simeq
\varinjlim_{f : j \to i} f_{\ast} (M_j){\quoi}.
$$
En utilisant l'hypothèse (\hyperref[VI.subsubsec8.7.1]{8.7.1}), on voit que le
faisceau $N'$ possède
la propriété suivante:

(P)~ Pour tout objet $i \in \ob (I)$, $\mu_{i\ast}(N')$ est
$m_{i\ast}$-acyclique, et pour tout morphisme $f : i \to j$ de $I$,
$\mu_{i \ast} (N')$ est $f_{\ast}$-acyclique.

Il suffit de montrer que tout objet $N'$ qui possède la propriété
(P) est $\underleftarrow{m}_{\ast}$-acyclique. Comme les injectifs
possèdent la propriété (P), il suffit de vérifier les deux
propriétés suivantes (\hyperrefext[V][V.prop0.4]{0.4}): Pour toute suite exacte $0
\to N' \to \sisi{}{N \to }N'' \to 0$, où $N'$ et $N$ possèdent la propriété
(P), (a) $N''$ possède la propriété (P) et (b) $\underleftarrow{m}_{\ast} (N) \to
\underleftarrow{m}_{\ast} (N'')$ est un épimorphisme. La vérification de (a) est
triviale. Vérifions (b). Posons $K = \coker (Q_{\ast}(N') \to
Q_{\ast} (N))$ de sorte que $K$ est la section $(i \mapsto K_i =
\coker (\mu_{i\ast} (N') \to \mu_{i\ast} (N)))$. Pour tout $f : i
\to j$, on a un diagramme commutatif
\[
\xymatrix{
0  \ar[r] & f_{\ast}\mu_{i\ast}(N') \ar[r] & f_{\ast}\mu_{i\ast}(N)\ar[r]
 & f_{\ast}(K_i) \ar[r]& 0\\
0  \ar[r] & \mu_{j\ast}(N')\ar[u] \ar[r] & \mu_{j\ast}(N) \ar[u] \ar[r] & K_j
\ar[r] \ar[u]&0.
}
\]
Comme\pageoriginale $\mu_{i \ast}(N')$ est $f_{\ast}$-acyclique, la
suite horizontale du haut est exacte. Comme $i \mapsto \mu_{i\ast}
(N')$ et $i \mapsto \mu_{i\ast} (N)$ sont des sections cartésiennes,
les deux premiers morphismes verticaux sont des isomorphismes. Par
suite le morphisme canonique $K_j \to f_{\ast} (K_i)$ est un
isomorphisme et $i \mapsto K_i$ est une section cartésienne. Donc
$K$ est de la forme $Q_{\ast} (K')$, et les morphismes canoniques
$Q_{\ast}(N) \to K$ et $K \to Q_{\ast}(N'')$ sont de la forme
$Q_{\ast}(u)$ et $Q_{\ast}(v)$ respectivement, car $Q_{\ast}$ est
pleinement fidèle. Comme $Q^{\ast}$ est exact et comme
$Q^{\ast}Q_{\ast}$ est isomorphe à l'identité, $v$ est un
isomorphisme et par suite $Q_{\ast}(v)$ est un isomorphisme. La
suite $0 \to Q_{\ast} (N') \to Q_{\ast}(N) \to Q_{\ast}(N'') \to 0$
est donc exacte et par suite, pour tout objet $i$ de $I$, la suite
$0 \to \mu_{i\ast}(N') \to \mu_{i\ast}(N) \to \mu_{i\ast}(N'') \to
0$ est exacte. Comme $\mu_{i\ast}(N')$ est $m_{i\ast}$-acyclique, la
suite $0 \to m_{i\ast} \mu_{i\ast} (N') \to m_{i\ast} \mu_{i\ast}
(N) \to m_{i\ast} \mu_{i\ast} (N'') \to 0$ est exacte. De plus on a
un isomorphisme $\underleftarrow{m}_{\ast} \simeq
\varinjlim_{I^{\circ}} \mu^{-1}_j m_{j\ast} \mu_{j\ast}$
\eqhyperref[VI.eq8.5.8.2]{8.5.8.2}. Le foncteur $\mu^{-1}_j$ est exact et les limites
inductives filtrantes sont exactes. Donc la suite $0 \to
\underleftarrow{m}_{\ast}(N') \to \underleftarrow{m}_{\ast}(N) \to
\underleftarrow{m}_{\ast}(N'') \to 0$ est exacte.

\begin{enonce*}{Théorème 8.7.3}\etiquette[8.7.3]{VI.thm8.7.3}
Les notations et hypothèses sont celles de \hyperref[VI.subsubsec8.7.1]{8.7.1}. Pour tout entier $n
\in \NN$, le foncteur dérivé $R^{n}\underleftarrow{m}_{\ast}$ commute
aux limites inductives filtrantes, et on a un isomorphisme fonctoriel
pour tout $A$-Module $j \mapsto M_j$:
\begin{equation*}
R^{n} \underleftarrow{m}_{\ast} Q^{\ast} (j \to M_j) \sisi{\sim}{\simeq} \varinjlim_{I^{\circ}}
\mu^{\ast}_j R^n m_{j\ast} (M_j)\tag{8.7.3.1}{\quoi}.
\end{equation*}\etiquette[8.7.3.1]{VI.eq8.7.3.1}
\end{enonce*}

On a $Q^{-1} (A) = \underrightarrow{A}$, et par suite l'image réciproque
pour les Modules est isomorphe à l'image réciproque pour les faisceaux
abéliens. On en déduit par \hyperref[VI.eq8.5.8.2]{8.5.8.2} que pour toute section $j
\mapsto M_j$, on a un
isomorphisme canonique $\varinjlim_{I^{\circ}} \mu^{-1}_j (M_j) \simeq
\varinjlim_{I^{\circ}} \mu^{\ast}_j (M_j)$. De plus,
sous les hypothèses de \hyperref[VI.subsubsec8.7.1]{8.7.1}, on a un isomorphisme canonique \eqhyperref[VI.eq8.5.9.1]{8.5.9.1}
$\underleftarrow{m}_{\ast} Q^{\ast} (j \mapsto M_j) \simeq
\varinjlim_{I^{\circ}} \mu^{-1}_j
m_{j\ast}(M_j)$. Comme pour tout injectif $j \mapsto M_j$, les $M_j$
sont flasques et $Q^{\ast} (j \to M_j)$ est acyclique pour
$\underleftarrow{m}_{\ast}$ (\hyperref[VI.lem8.7.2]{8.7.2}), on a un isomorphisme $R^{n}
\underleftarrow{m}_{\ast} (Q^{\ast} (j \mapsto M_j)) \simeq
\varinjlim_{I^{\circ}} \mu^{-1}_j R^n m_{j\ast}(M_j)$,
d'où la formule \eqhyperref[VI.eq8.7.3.1]{8.7.3.1} en utilisant ce qui
précède. La première
assertion résulte immédiatement de la formule \eqhyperref[VI.eq8.7.3.1]{8.7.3.1}.

\begin{enonce*}{Corollaire 8.7.4}\etiquette[8.7.4]{VI.cor8.7.4}
On\pageoriginale a des isomorphismes fonctoriels pour $n \in \ZZ$:
\begin{equation*}
R^n \underleftarrow{m}_{\ast} (N) \simeq \varinjlim_{I^{\circ}}
\mu^{\ast}_j R^n m_{j\ast} \mu_{j\ast} (N){\quoi}.\tag{8.7.4.1}
\end{equation*}\etiquette[8.7.4.1]{VI.eq8.7.4.1}
\end{enonce*}

En effet $N$ est isomorphe à $Q^{\ast}(j \mapsto \mu_{j \ast}(N))$.

\begin{enonce*}{Corollaire 8.7.5}\etiquette[8.7.5]{VI.cor8.7.5}
Pour tout objet $i$ de $I$, \sisi{pout entier}{pour tout entier} $n$ et tout $A_i$-Module
$M_i$, on a un isomorphisme fonctoriel
\begin{equation*}
R^n \underleftarrow{m}_{\ast} \mu^{\ast}_i (M_i) \simeq
\varinjlim_{f : j \to i} \mu^{\ast}_j R^n m_{j\ast} f^{\ast}
(M_i){\quoi}.\tag{8.7.5.1}
\end{equation*}\etiquette[8.7.5.1]{VI.eq8.7.5.1}
\end{enonce*}

Quitte à faire le changement de base $I_{/i} \to I$, on peut supposer
que $i$ est un objet final de $I$ (\hyperref[VI.lem8.2.1]{8.2.1}). L'objet $\mu_{i}^{\ast}
(M_i)$ est alors isomorphe à $Q^{\ast} (f : j \to i) \mapsto
f^{\ast}(M_i))$ (\cf \hyperref[VI.cor8.5.5]{8.5.5}).

\begin{enonce*}{Corollaire 8.7.6}\etiquette[8.7.6]{VI.cor8.7.6}
Soit $i$ un objet de $I$ et $n$ un entier. On a un isomorphisme
fonctoriel en la section $j \mapsto M_j$:
\begin{equation*}
R^n \mu_{i\ast} Q^{\ast} (j \mapsto M_j) \simeq \varinjlim_{f :
  j \to i} R^n f_{\ast} (M_j){\quoi}.\tag{8.7.6.1}
\end{equation*}\etiquette[8.7.6.1]{VI.eq8.7.6.1}
On a un isomorphisme en le $\underrightarrow{A}$-Module $N$:
\begin{equation*}
R^n \mu_{i\ast} (N) \simeq \varinjlim_{f:j\to i}
R^n f_{\ast} \mu_{j\ast} (N){\quoi}.\tag{8.7.6.2}
\end{equation*}\etiquette[8.7.6.2]{VI.eq8.7.6.2}
On a un isomorphisme fonctoriel en le $A_i$-Module $M_i$:
\begin{equation*}
R^n \mu_{i\ast} \mu^{\ast}_i (M_i) \simeq \varinjlim_{f:j\to i}
R^n f_{\ast} f^{\ast} (M_i){\quoi}.\tag{8.7.6.3}
\end{equation*}\etiquette[8.7.6.3]{VI.eq8.7.6.3}
Les foncteurs $R^n \mu_{i \ast}$ commutent aux limites inductives filtrantes.
\end{enonce*}

En faisant le changement de base $I_{/i} \to I$ on se ramène au cas où
$i$ est un objet final de $I$. Les formules \eqhyperref[VI.eq8.7.6.1]{8.7.6.1} et \eqhyperref[VI.eq8.7.6.2]{8.7.6.2}
sont alors des cas particuliers des formules \eqhyperref[VI.eq8.7.3.1]{8.7.3.1}, \eqhyperref[VI.eq8.7.4.1]{8.7.4.1} et
\eqhyperref[VI.eq8.7.5.1]{8.7.5.1}, respectivement obtenus en prenant pour $G$ le topos fibré
constant de fibre $F_i$ et pour morphisme $m$ le morphisme $(f., f \in
\ob (I_{/i}))$. La dernière assertion résulte de \eqhyperref[VI.eq8.7.6.2]{8.7.6.2} compte tenu de
(\hyperref[VI.cor8.5.4]{8.5.4}).

\begin{enonce*}{Corollaire 8.7.7}\etiquette[8.7.7]{VI.cor8.7.7}
Soient\pageoriginale $i$ un objet de $I$, $X$ un objet de $F_i$ tel que les
foncteurs $H^{n} (X, -)$, $n \in \NN$, commutent aux limites
inductives filtrantes. On a pour tout $n$ un isomorphisme canonique
fonctoriel en le $A_i$-Module $M_i$:
\begin{equation*}
H^n (\mu_{i}^{\ast}(X), \mu^{\ast}_i (M_i)) \simeq \varinjlim_{f:j\to i}
H^n (f^{\ast} (X), f^{\ast} (M_i))\tag{8.7.7.1}
\end{equation*}\etiquette[8.7.7.1]{VI.eq8.7.7.1}
et les foncteurs $H^{n} (\mu^{\ast}_i (X), -)$ commutent aux limites
inductives filtrantes. En particulier si les foncteurs $H^{n} (F_i,
-)$ commutent aux limites inductives filtrantes, les foncteurs $H^n
(\underleftarrow{F}, -)$ commutent aux limites inductives filtrantes, et
on a des isomorphismes canoniques, fonctoriels en les $A_i$-Modules
$M_i$
\begin{equation*}
H^n (\underleftarrow{F}, \mu_{i}^{\ast} (M_i)) \simeq \varinjlim_{f:j\to i}
H^n (F_j, f^{\ast} (M_i)){\quoi}.\tag{8.7.7.2}
\end{equation*}\etiquette[8.7.7.2]{VI.eq8.7.7.2}
\end{enonce*}

\noindent
Il résulte de (\hyperrefext[V][V.prop5.3]{5.3}) qu'on a deux suites spectrales
\begin{align*}
'E_{2}^{p, q} &= \varinjlim_{f:j\to i} H^p(X, R^q f_{\ast}
f^{\ast} (M_i)) \Rightarrow \varinjlim_{f:j\to i}
H^{p+q}(f^{\ast} (X), f^{\ast} (M_i))\\
''E_{2}^{p, q} &=  H^p(X, R^q \mu_{i\ast} \mu^{\ast}_i (M_i))
\Rightarrow H^{p+q} (\mu^{\ast}_i (X), \mu^{\ast}_i (M_i)){\quoi},
\end{align*}
et de \eqhyperref[VI.eq8.5.7.1]{8.5.7.1} qu'on a un morphisme entre ces deux suites
spectrales. Comme les $H^{p+q} (X, -)$ commutent aux limites
inductives filtrantes, il résulte de \eqhyperref[VI.eq8.7.6.2]{8.7.6.2} que ce morphisme de
suites spectrales est un isomorphisme au niveau des $E^{p, q}_2$. Par
suite il induit un isomorphisme sur les aboutissements, d'où
\eqhyperref[VI.eq8.7.7.1]{8.7.7.1}. Pour tout $\underrightarrow{A}$-Module $N$, on a une suite
spectrale (\hyperrefext[V][V.prop5.3]{5.3}).
$$
E^{p, q}_2 = H^{p} (X, R^q \mu_{i\ast} (N)) \Rightarrow H^{p+q}
(\mu^{\ast}_i (X), N){\quoi}.
$$
La commutation aux limites inductives filtrantes des foncteurs $H^{n}
(\mu^{\ast}_i (X), -)$ résulte alors des propriétés analogues des
foncteurs $H^{n} (X, -)$ et $R^n \mu_{i\ast}$ (\hyperref[VI.cor8.7.6]{8.7.6}).

\begin{enonce*}[remark]{Remarque 8.7.8}\etiquette[8.7.8]{VI.rem8.7.8}
Lorsque les topos $F_j$, $j \in \ob (I)$ sont cohérents et lorsque les
morphismes de transition sont cohérents, l'hypothèse faite sur $X$
(\resp $F_i$) dans \sisi{8.6.7}{8.7.7} est satisfaite
lorsque $X$ (\resp $F_i$) est
cohérent (\hyperref[VI.cor5.2]{5.2}). On sait d'ailleurs que dans ce cas $\mu^{\ast}_i (X)$
(\resp $\underleftarrow{F}$) est cohérent (\hyperref[VI.thm8.3.13]{8.3.13}).
\end{enonce*}

\begin{enonce*}{Corollaire 8.7.9}\etiquette[8.7.9]{VI.cor8.7.9}
Soient\pageoriginale $i$ \sisi{un objet $I$}{un objet de $I$}, $M_i$ un $A_i$-Module à gauche, $L_i$ un
$A_i$-Module à droite possédant la résolution du type:
$$
P_{i, k} \to P_{i, k-1} \to \cdots \to P_{i, 0} \to L_i \to 0
$$
où, pour tout entier $k$, $P_{i, k}$ est isomorphe à une somme directe
finie d'objets de la forme $A_{i|X}$ \emph{(IV 11.3.3)}, où $X$ vérifie les
hypothèses de \hyperref[VI.cor8.7.7]{8.7.7}. Si, outre les hypothèses de \hyperref[VI.subsubsec8.7.1]{8.7.1}, le topos fibré
$F$ est plat à droite (\hyperref[VI.subsec8.6]{8.6}), on a des isomorphismes canoniques, pour
$n \leq k$-$1$:
\begin{equation*}
\Ext^n_{\underrightarrow{A}}(\underleftarrow{F}, \mu^{\ast}_i(L_i),
\mu^{\ast}_i (M_i)) \simeq \varinjlim_{f : j \to i}
\Ext^{n}_{A_j} (F_j, f^{\ast}(L_i), f^{\ast} (M_i)){\quoi}.\tag{8.7.9.1}
\end{equation*}\etiquette[8.7.9.1]{VI.eq8.7.9.1}
\end{enonce*}

\noindent
Notons $P_{i,.}$ la résolution de $L_i$. Comme $F$ est plat à
droite, pour tout $f : j \to i$, $f^{\ast}(P_{i,.})$ est une
résolution de $f^{\ast}(L_i)$ et $\mu^{\ast}_i(P_{i,.})$ est une
  résolution de $\mu^{\ast}_i(L_i)$ (\hyperref[VI.subsec8.6]{8.6}). Ces résolutions permettent
  de construire deux suites spectrales qui convergent respectivement
  vers les deux membres de \eqhyperref[VI.eq8.7.9.1]{8.7.9.1}, et un morphisme entre ces deux
  suites spectrales. Au niveau des $E^{p, q}_1$ ce morphisme est un
  isomorphisme \eqhyperref[VI.eq8.7.7.1]{8.7.7.1}. Il induit donc un isomorphisme sur les
  aboutissements.

\section{Appendice. Critère d'existence de points}\etiquette[9]{VI.sec9}

\centerline{\textit{par} P. \textsc{Deligne}}

\medskip
\begin{enonce*}{Proposition 9.0}\etiquette[9.0]{VI.prop9.0}
Tout topos localement cohérent $S$ a assez de points.
\end{enonce*}

La question étant locale sur $S$, on peut supposer le topos $S$
défini par un site $\Ss$, dans lequel \sisi{est}{les} limites projectives finies
sont représentables et tel que tout recouvrement $f_i : U_i \to U$
admette un sous-recouvrement fini. Il suffit de prouver que si $f :
\Ff \to \Gg$ n'est pas un monomorphisme, alors, il existe un point
$x$ de $\Ss$ tel que $f_x$ ne soit pas injectif. Par hypothèse, il
existe $U \in \ob \Ss$ et $\ss$, $\ss'\in \Ff(U)$ tels que $\ss \neq
\ss'$ et $f(\ss) = f(\ss')$. Remplaçant $\Ss$ par $\Ss/U$, on se
ramène au

\begin{enonce*}{Lemme 9.1}\etiquette[9.1]{VI.lem9.1}
Si\pageoriginale $\ss$ et $\ss'$ sont deux sections globales distinctes d'un faisceau
$\mathfrak{F}$ sur un site $\Ss$ vérifiant les hypothèses précédentes,
alors, il existe un point $x$ de $\Ss$ tel que $\ss_x
\neq \ss'_x$.
\end{enonce*}

Soit $P = (U_i)_{i \in I}$ un système projectif dans $\Ss$, indexé par
un ensemble ordonné filtrant $I$. Pour tout faisceau $\Hh$ sur
$\Ss$, on pose $P(\Hh) = \varinjlim \Hh(U_i)$, et pour
tout $V \in \Ob \Ss$, on pose $P(V) = \varinjlim \Hom (U_i, V)$. Les
foncteurs $P(\Hh)$ et $P(V)$ commutent aux produits fibrés. Si
le foncteur $P(V)$ transforme les recouvrements en familles
surjectives de fonctions, c'est un morphisme de sites du topos
ponctuel dans $\Ss$, et $P(\Hh)$ est le foncteur fibre
correspondant.

Si $P = (U_i)_{i \in I}$ et $Q = (V_j)_{j \in J}$ sont deux systèmes
projectifs dans $\Ss$, on dit que $Q$ raffine $P$ si $I$ est une partie
de $J$ munie de l'ordre induit et si $P$ est la restriction de $Q$ à
$J$. On dispose alors de morphismes de foncteurs de $P(\Hh)$
dans $Q(\Hh)$ et de $P(V)$ dans $Q(V)$.

Quel que soit $P$ comme plus haut, on notera $\ss_p$ et
$\ss'_p$ les images dans $P(\Ff)$ de $\ss$ et $\ss'$. Le lemme \hyperref[VI.lem9.1]{9.1} résulte du
lemme suivant, dans lequel on prend pour $P$ le système projectif
indexé par $\{0\}$ et réduit à l'objet final de $\Ss$.

\begin{enonce*}{Lemme 9.2}\etiquette[9.2]{VI.lem9.2}
Quel que soit $P$ comme plus haut vérifiant $\ss_p \neq \ss'_p$, il
existe $Q$ raffinant $P$, vérifiant $\ss_Q \neq \ss'_Q$, et tel que,
quels que soient le recouvrement $f_i : V_i \to V$ et $f \in Q(V)$,
$f$ soit dans l'image d'un des $Q(V_i)$.
\end{enonce*}

On prouvera tout d'abord :

\begin{enonce*}{Lemme 9.3}\etiquette[9.3]{VI.lem9.3}
Soient $P = (U_i)_{i \in I}$ vérifiant les hypothèses du lemme \hyperref[VI.lem9.2]{9.2},
$f_i : V_i \to V$ un recouvrement fini dans $\Ss$, et $f \in P(V)$. Il
existe $Q$ raffinant $P$, vérifiant $\ss_Q \neq
\ss'_Q$ et tel que l'image de $f$ dans $Q(V)$ soit dans l'image de l'un des $\sisi{(V_i)}{Q(V_i)}$.
\end{enonce*}

Il existe $i_{\circ} \in I$ et $f'\in \Hom(U_{i_{\circ}}, V)$ qui
définissent $f$. Pour $i \geq i_\circ$, posons $U_{i, k} =
\fprod{U_i}{V_k}{V}$, ce produit fibré étant défini par $f'$. On sait
que $(V_k)$ est un recouvrement de $V$, donc que $U_{i, k}$ est un
recouvrement de $U_i$; la flèche suivante sera injective (pour $i >
i_{\circ}$):
$$
\Ff (U_i) \to \prod_k \Ff (U_{i, k})
$$
Passant à\pageoriginale la limite, compte tenu de ce que les produits
\emph{finis} commutent aux limites inductives filtrantes, on voit que
$$
P(\Ff) \to \prod_k \varinjlim \Ff (U_{i, k})
$$
est injectif, donc qu'il existe $k$ tel que $\ss_p$ et
$\ss'_p$ aient des images distinctes dans $\varinjlim_{i > i_\circ} \Ff
(U_{i, k})$. Soit alors $I_{\circ} = \{i | i \geq i_{\circ} \mbox{ et }
  i \in I\}$, et $J = I \coprod I_{\circ}$. On ordonne $J$ en disant
que $j' \leq j''$ si les images de $j'$ et $j''$ dans $I$ satisfont à
$j' \leq j''$ et si on n'a pas $j' \in I_{\circ}$, $j''\in I$. Les
$U_i$ et $U_{i, k}$ sont indexés par $J$, et forment un raffinement $Q$ de $P$, tel
que $\ss_Q \neq \ss'_Q$ , et que l'image de $f$ dans
$Q(V)$ soit dans l'image de $Q(V_k)$.

\begin{enonce*}{Lemme 9.4}\etiquette[9.4]{VI.lem9.4}
Soit $P = (U_i)_{i \in I}$ vérifiant les hypothèses du lemme \hyperref[VI.lem9.2]{9.2}. Il
existe $Q$ raffinant $P$, tel que $\ss_Q \neq \ss'_Q$,
et tel que pour tout recouvrement fini $(V_k)$ d'un ouvert $V$ de $S$, et tout $f \in
P(V)$, l'image de $f$ dans $Q(V)$ soit dans l'image de l'un des $Q(V_k)$.
\end{enonce*}

Soit $E$ l'ensemble des triples formés d'un ouvert $V$ de $S$, d'un
recouvrement fini $(V_k)$ de $V$ et de $f \in P(V)$. Soient $\leq$
un bon ordre sur $E$ et $\bar{E}$ l'ensemble déduit de $E$ par
adjonction d'un plus grand élément, noté $\infty$. On va définir par
récurrence transfinie sur $e \in \bar{E}$ des raffinements $Q_e$ de
$P$, vérifiant $\ss_{Q_e} \neq \ss'_{Q_e}$ et tels que
\begin{enumerate}
\renewcommand{\theenumi}{\roman{enumi}}
\renewcommand{\labelenumi}{(\theenumi)}
\item si $e'\leq e$, alors $Q_e$ raffine $Q_{e'}$.

\item si $e = (V, (V_k), f)$, alors l'image de $f$ \sisi{sans}{dans} $Q_e(V)$ se
  trouve dans l'image de l'un des $Q_{e} (V_k)$.
\end{enumerate}

Supposons les $Q_{e'}$ déjà définis pour $e'< e$. Si $e$ est le
premier élément de $E$, posons $Q'_e = P$. Si $e$ a un prédécesseur
$e-1$, posons $Q'_e = Q_{e-1}$. Sinon, soit $Q'_e$ le système
projectif d'ensemble d'indices $\cup_{e' < e}
I_{e'}$ qui raffine les $Q_{e'}$ pour $e' < e$. Dans ce cas, on a
$$
Q'_e (\Ff) = \varinjlim_{e' < e} Q_{e'} (\Ff)
$$
de \pageoriginale sorte que, dans tous les cas, $Q'_e$ raffine les
$Q_{e'}$ pour $e'< e$ et vérifie $\ss_{Q'_e} \neq \ss'_{Q'_e}$.

On obtient le système projectif $Q_e$ requis en appliquant le lemme
\hyperref[VI.lem9.3]{9.3} à $Q'_e$ et à $e = (V, (V_k), f)$ (\resp en prenant $Q_e =
Q'_e$ si $e = \infty$).

Le système projectif $Q_{\infty}$ vérifie le lemme \hyperref[VI.lem9.4]{9.4}.

Le lemme \hyperref[VI.lem9.4]{9.4}\sisi{.}{} permet de définir, par récurrence sur $n$, une suite
$Q_n = (U_i)_{i \in I_n}$ de systèmes projectifs dans $\Ss$ telle que
\begin{enumerate}
\renewcommand{\theenumi}{\roman{enumi}}
\renewcommand{\labelenumi}{(\theenumi)}
\item $Q_0 = P$

\item $Q_{n+1}$ raffine $Q_n$

\item $\ss_{Q_n} \neq \ss'_{Q_n}$

\item Quels que soient le recouvrement $f_k : V_k \to V$ et $f \in Q_n
  (V)$, l'image de $f$ dans $Q_{n+1} (V)$ se trouve dans l'image de
  l'un des $Q_{n+1} (V_k)$.
\end{enumerate}

Le système projectif $Q$, d'ensemble d'indices la réunion des $I_n$,
qui prolonge les différents $Q_n$, vérifie alors le lemme \hyperref[VI.lem9.2]{9.2}. La
démonstration montre de plus que:

\begin{enonce*}{Corollaire 9.5}\etiquette[9.5]{VI.cor9.5}
Soit $\Ss$ un site, dans lequel les produits fibrés sont représentables
et tel que tout recouvrement dans $\Ss$ admette un sous-recouvrement
fini. Soit \sisi{$e$ le cardinal $\sup(\quad)_{\circ}$, $\card(\text{\rm Fl }
\Ss))$}{$c=\sup(\aleph_0,\operatorname{card}(\operatorname{Fl}\Ss))$}. Alors il existe un ensemble très dense de points de
$\Ss$, tel que
\begin{enumerate}
\renewcommand{\theenumi}{\roman{enumi}}
\renewcommand{\labelenumi}{(\theenumi)}
\item $\card(X) \leq 2^c$

\item si $x \in X$ et $U \in \Ob \Ss$, alors $\card(U_x) \leq c$.
\end{enumerate}
\end{enonce*}


\begin{thebibliography}{99}\pageoriginale
\bibitem{SGA4-06:1} P.\textsc{Gabriel} et M.\textsc{Zisman}: Homotopy Theory
  and Calculus   of \sisi{Fraction}{Fractions}, Ergebnisse der Mathematik, Bd 35.

\bibitem{SGA4-06:2} J. \textsc{Giraud} : Méthode de la descente. Mémoire de la S.M.F.

\bibitem{SGA4-06:3} A. \textsc{Grothendieck} et J. \textsc{Dieudonne} : Eléments de
  Géométrie Algébrique I, 2ème édition.

\bibitem{SGA4-06:4} A. \textsc{Grothendieck} et J. \textsc{Dieudonne} : Eléments de
  Géométrie Algébrique IV 1, I.H.E.S.~${\rm n^{\circ}}$

\bibitem{SGA4-06:5} \SGA. 1 VI, par A. \textsc{Grothendieck}, I.H.E.S.

\bibitem{SGA4-06:6} \SGA. 3 IV 6.3 par M. \textsc{Demazure}, Lecture Notes
  ${\rm n^{\circ}}$, Springer Verlag.

\bibitem{SGA4-06:7} A. \textsc{Grothendieck} : Sur quelques points d'algèbres
  homologiques, Tohoku Math. J.
\end{thebibliography}

\ifx\danslelivre\undefined
\end{document}
\fi
